% ======================================================================
%  Theory appendix: proofs and derivations for Section 2 (sec-theory.tex)
%  Input after sec-appendix-integrated.tex inside \appendix\onecolumn.
%  Eight pieces, one \subsection each; labels are those of the verified
%  theory pieces (suffix -p1, fourbulk, buffer, linear, bridge, pred, gen, scope).
% ======================================================================
\section{Theory: proofs and derivations}
\label{app:theory}

This appendix collects the full statements and proofs behind \cref{sec:theory}. Throughout, $\Delta_t=1-e^{-2t}$, the random-feature score model is $s_A(x_t)=\pwidth^{-1/2}A\tanh(Wx_t/\sqrt{\dlat})=A\phi(x_t)$ with $W\in\R^{\pwidth\times\dlat}$ i.i.d.\ $\mathcal N(0,1)$ frozen and $A\in\R^{\dlat\times\pwidth}$ trained from $A(0)=0$, and $\Umat=\tfrac1{\nsamp}\sum_\mu\E_\eta[\phi(x_t^\mu)\phi(x_t^\mu)^{\top}]$ is the feature second-moment operator. Every statement is labelled \emph{proved}, \emph{sketched} or \emph{conjectured}; every displayed edge and timescale carries its $\dlat$-dependence explicitly. Sections are ordered as in the main text: setup and spectral timescales (\cref{app:theory-setup}), the four spectral blocks of $\Umat$ (\cref{app:theory-fourbulk}), the buffer ratio (\cref{app:theory-buffer}), the exactly solvable linear score model (\cref{app:theory-linear}), the bridge from residual sample-bulk error to near-duplicate generation (\cref{app:theory-bridge}), the fraction-level predictor (\cref{app:theory-predictor}), general latent spectra and posterior collapse (\cref{app:theory-general}), and regime boundaries (\cref{app:theory-scope}, with a closing list of configurations outside the hypotheses).

% ----------------------------------------------------------------------
\subsection{Setup, standing assumptions, and spectral definitions of \texorpdfstring{$\taugen$}{tau\_gen} and \texorpdfstring{$\taumem$}{tau\_mem}}
\label{app:theory-setup}

\paragraph{Setting.}
Data: $x^\mu=Q(m_{c(\mu)}+\xi^\mu)$, $\mu=1,\dots,\nsamp$, with $k$ centres $m_c$ of norm $s$ in $\operatorname{span}(e_1,\dots,e_{\dint})$, $\xi^\mu\sim\mathcal N(0,\diag(\sigsig^2I_{\dint},\signoise^2I_{\dlat-\dint}))$, $Q\in O(\dlat)$ Haar. Since $W$ is Gaussian and the loss is rotation-equivariant we set $Q=I$. Forward process $x_t=e^{-t}x+\sqrt{\Delta_t}\eta$; $M_t=\tfrac1{\nsamp}\sum_\mu\E_\eta[x_t^\mu x_t^{\mu\top}]=e^{-2t}(\hat C+\hat S)+\Delta_tI$ with $\hat C$ the empirical centre second moment and $\hat S$ the empirical within-cluster covariance; population block eigenvalues $\alpha_t^2=e^{-2t}(s^2/\dint+\sigsig^2)+\Delta_t$, $\beta_t^2=e^{-2t}\signoise^2+\Delta_t$. Features $\phi(x)=\pwidth^{-1/2}\tanh(Wx/\sqrt{\dlat})$; the code stores $W/\sqrt{\dlat}$ as $\mathcal N(0,1/\dlat)$ entries and omits $\pwidth^{-1/2}$ in its logged $\Umat$, so \emph{code eigenvalues are $\pwidth$ times ours}. Two feature operators are distinguished throughout: the single-draw $\Umat^{(1)}_k$ (one noise draw per point; rank $\le\min(\pwidth,\nsamp)$) and its expectation $\Umat=\E_\eta\Umat^{(1)}_k$ (full rank; the operator governing the mean training dynamics). The released eigenvalue files average $50$ draws and are of the second kind: they are full rank and show only a soft step at index $\nsamp$ ($\lambda_{\nsamp}/\lambda_{\nsamp+1}=1.00,1.7,1.7$ at $\dlat=20,60,200$). Empirical objective $L_{\nsamp}$, population objective $L(A)=\tfrac12\E\|A\phi(x_t)-y\|^2$, $\bar\Umat=\E[\phi\phi^{\top}]$, $\bar V=\E[y\phi^{\top}]=\E[s^\star(x_t)\phi(x_t)^{\top}]$ (tower property), $s^\star=\nabla\log p_t$. Gaussian-equivalence coefficients $a_1(q)=\E[\tanh'(\sqrt qz)]$, $a_*(q)^2=\E[\tanh(\sqrt qz)^2]-a_1(q)^2q$, $z\sim\mathcal N(0,1)$ \citep{pennington2017nonlinear,benigni2021eigenvalue,goldt2022gaussian,hu2022universality,mei2022generalization}.

\paragraph{Related work and scope.}
\citet{bonnaire2025} prove the spectral characterisation of $\Umat$ for an arbitrary input spectral measure in the proportional regime; \citet{george2025dsm} give a linear pencil with block dimensions $(\pwidth,\dlat-\dint,\dint,\nsamp)$ for data on a $\dint$-dimensional subspace with $\signoise=0$. Everything in this subsection is exact finite-size algebra for the expected gradient flow given the spectrum of $\Umat$; the only asymptotic inputs are the block scales of \cref{app:theory-fourbulk}, which are those theorems evaluated at the two-atom spectral measure $\{\alpha_t^2,\beta_t^2\}$ with $\signoise>0$.

\subsubsection{Standing assumptions}

\begin{assumption}[Standing hypotheses]\label{ass:standing-p1}
\begin{enumerate}
\item[(A1)] \emph{Fixed diffusion time.} $t>0$ is fixed, $\Delta_t=1-e^{-2t}$; the readout is trained at this single $t$ (repository: $t=0.01$ for the RFNN, $t=0.1$ for the real-data spectral clock).
\item[(A2)] \emph{Two feature operators; over-parametrisation.} Write $\Umat^{(1)}_k=\tfrac1{\nsamp}\sum_\mu\phi(x^{\mu,k}_t)\phi(x^{\mu,k}_t)^{\top}$ for the \emph{single-draw} operator built from one noise draw $\eta^{\mu,k}$ per training point, and $\Umat=\E_\eta\Umat^{(1)}_k$ for its expectation. Then $\operatorname{rank}\Umat^{(1)}_k\le\min(\pwidth,\nsamp)$ exactly (this is the machine-precision cliff at index $\nsamp$ in single-draw four-bulk spectra), whereas $\Umat$ is generically of full rank $\pwidth$: its $\pwidth-\nsamp$ lowest eigenvalues form a \emph{tail block} $\mathcal Z$ that is \emph{contiguous} with the bottom of the sample block (repository spectra, $50$ draws, code units: $\lambda_{\nsamp}/\lambda_{\nsamp+1}=1.00$ at $\dlat=20$, $1.7$ at $\dlat=200$; a $(\dlat,\pwidth,\nsamp)=(20,1280,500)$ simulation gives tail max $1.79\times10^{-6}$ against sample min $1.81\times10^{-6}$, stable from $50$ to $400$ draws). Since the training loop draws fresh noise at every step, the mean training trajectory follows the flow generated by $\Umat$, not $\Umat^{(1)}$ (\cref{prop:sgdmean-p1}); all dynamical statements below refer to $\Umat$. We require $\pwidth\ge\nsamp$ for the tail block to exist and for the four-block picture; the repository's rule $\pwidth=64\dlat$ \emph{violates this at $\dlat=5$} ($\pwidth=320<\nsamp=500$), where $\mathcal M=(\dlat,\pwidth]$ and $\mathcal Z=\emptyset$; \cref{lem:gf-p1} is unchanged there.
\item[(A3)] \emph{Data model and scaling regime.} $x^\mu=Q(m_{c(\mu)}+\xi^\mu)$ with $k$ centres of norm $s$ in a fixed $\dint$-dimensional subspace, $\xi^\mu\sim\mathcal N(0,\diag(\sigsig^2I_{\dint},\signoise^2I_{\dlat-\dint}))$, $Q$ Haar. $(k,s,\dint,\sigsig,\signoise,\nsamp)$ are fixed and $\dlat\to\infty$ at fixed $\psi_p=\pwidth/\dlat$ (repository: $64$). Hence $\psi_n=\nsamp/\dlat$ is \emph{not} fixed; we require $\psi_n>1$ so that the sample block is nonempty, and every statement that separates sample from population modes degrades as $\psi_n\downarrow1$ (repository: $\psi_n=2.5$ at $\dlat=200$). Because $\dint$ is fixed, the signal block is a finite-rank (BBP-type) spike set, not a bulk.
\item[(A4)] \emph{Operating point through $q$.} $q=q(\dlat)=\operatorname{tr}M_t/\dlat\in[\beta_t^2,q(\dint)]$ (\cref{lem:q-p1}); $a_1(q)=\E[\tanh'(\sqrt qz)]\in(0,1]$ and $a_*(q)^2=\E[\tanh(\sqrt qz)^2]-a_1(q)^2q>0$ are always evaluated at $q(\dlat)$ and never treated as $\dlat$-independent. No condition $s/\sqrt{\dlat}\ll1$ is used.
\item[(A5)] \emph{Block ordering.} The sorted spectrum of $\Umat$ is partitioned by index into $\mathcal S=\{1,\dots,\dint\}$, $\mathcal N=\{\dint+1,\dots,\dlat\}$, $\mathcal M=\{\dlat+1,\dots,\nsamp\}$, $\mathcal Z=\{\nsamp+1,\dots,\pwidth\}$, with $\lambda^{\mathcal S}_{\min}>\lambda^{\mathcal N}_{\max}$ and $\lambda^{\mathcal N}_{\min}>\lambda^{\mathcal M}_{\max}$ (repository: $\lambda^{\mathcal N}_{\min}/\lambda^{\mathcal M}_{\max}=7.7$ at $\dlat=20$, $12$ at $\dlat=200$). No gap is assumed between $\mathcal M$ and $\mathcal Z$. In the isotropic case $\signoise=\sigsig$ the first inequality is dropped and $\mathcal S\cup\mathcal N$ is one population block of $\dlat$ modes.
\item[(A6$'$)] \emph{Weighted spectral alignment (used only for the gap).} Let $\bar\Umat=\E[\phi\phi^{\top}]$, $\bar V=\E[y\phi^{\top}]$, $\bar u_{ij}=v_i^{\top}\bar\Umat v_j$, $\bar u_i=\bar u_{ii}$, and define the \emph{alignment coefficient} $\gamma_i:=\langle a_i^\star,\bar Vv_i\rangle/(\lambda_i\pi_i)$ for $\pi_i>0$ ($\gamma_i=1$ for a mode whose empirical optimum equals its population optimum, $\gamma_i=0$ for a pure sample mode). Let $G_{\rm thr}:=\min\{G_\infty,\,0.01\dlat/\Delta_t\}$ be the smallest gap level that matters (\cref{rem:proxy-p1}). We assume, uniformly in $T\ge0$,
\[
\Big|\tfrac12\sum_{i\ne j}(1-e_i)(1-e_j)\bar u_{ij}\langle a_i^\star,a_j^\star\rangle\Big|=o(G_{\rm thr}),\qquad
\sum_{i\in\mathcal S\cup\mathcal N}\lambda_i\pi_i\Big(\big|\tfrac{\bar u_i}{\lambda_i}-1\big|+2|\gamma_i-1|\Big)=o(G_{\rm thr}).
\]
This is a hypothesis stated in the \emph{mass-weighted} form that the exact gap formula requires: the target masses are large ($\Pi_{\mathcal N}\approx1.1\times10^4$, $\Pi_{\mathcal M}\approx2.4\times10^7$, $\Pi_{\mathcal Z}\approx6\times10^8$ at $(\dlat,\pwidth,\nsamp)=(20,1280,500)$, since $a_i^\star=Vv_i/\lambda_i$), so an unweighted $o(1)$ in $\bar u_{ij}$ would not control the gap; at $\dlat=20$, $\lambda^{\mathcal N}\Pi_{\mathcal N}\approx60$ against $0.01\dlat/\Delta_t=10$, so the population blocks must be aligned to relative accuracy $\ll10\%$. Nothing is assumed about $\gamma_i$ on $\mathcal M\cup\mathcal Z$: those coefficients are measured (\cref{def:taumem-p1}).
\end{enumerate}
\end{assumption}

\noindent\emph{Status.} Not a theorem. (A2): the rank bound of the single-draw operator is exact linear algebra ($\Umat^{(1)}_k$ is a sum of $\nsamp$ rank-one terms); contiguity of the tail is read from the $\signoise=0.5$ RFNN spectra ($50$ draws): $(\lambda_{\nsamp},\lambda_{\nsamp+1})$ in code units $=(0.00239,0.00238)$ at $\dlat=20$, $(0.0122,0.0072)$ at $60$, $(0.0557,0.0323)$ at $200$. (A5) is the repository's four-bulk observation. (A6$'$) is a hypothesis to be tested by computing $\bar\Umat,\bar V$ on held-out data (Prediction~2 below).

\paragraph{Numerical status of (A2).}
At $(\dlat,\pwidth,\nsamp,t)=(20,1280,500,0.01)$ with the repository's \texttt{compute\_U}: one draw gives tail eigenvalues $\le10^{-17}$; $50$ and $400$ draws give an identical tail (max $1.79\times10^{-6}$, mean $5.0\times10^{-7}$, $0.11\%$ of $\operatorname{tr}\Umat$), so this is the true $\E_\eta$ tail and not Monte-Carlo noise. The tail top equals the sample-block bottom ($1.81\times10^{-6}$; the Marchenko--Pastur lower edge $a_*^2(1-\sqrt{\nsamp/\pwidth})^2/\nsamp=4.6\times10^{-6}$ lies above the realised minimum). Under the clock $T=0.02\dlat\times$steps, $\lambda_{\mathcal Z}T\in[0.06,0.2]$ at $3\times10^5$ steps, and because $a_i^\star=Vv_i/\lambda_i$ with $\|VP_{\mathcal Z}\|_F^2/\|V\|_F^2=4.3\times10^{-4}$, the tail target mass is $\Pi_{\mathcal Z}\approx6\times10^8$ against $\Pi_{\mathcal M}\approx2.4\times10^7$. The tail is therefore slow but not inert in Frobenius norm; it is inert in \emph{loss} only to the extent of \cref{prop:gap-p1}'s bound (weights $\bar u_i,\lambda_i\lesssim2\times10^{-6}$).

\subsubsection{Mean dynamics, gradient flow and residual spectral error}

\begin{proposition}[Mean of fresh-noise full-batch SGD is the $\E_\eta$ gradient flow; proved]\label{prop:sgdmean-p1}
Let the training step be $A_{k+1}=A_k-2\Delta_t\eta_{\rm lr}\big(A_k\Umat^{(1)}_k-V^{(1)}_k\big)$ with $(\Umat^{(1)}_k,V^{(1)}_k)$ built from the $k$-th noise draw, $V^{(1)}_k=\tfrac1{\nsamp}\sum_\mu y^{\mu,k}\phi(x^{\mu,k}_t)^{\top}$, draws i.i.d.\ across $k$, $A_0=0$, momentum $0$. Then $\bar A_k:=\E[A_k]$ satisfies exactly $\bar A_{k+1}=\bar A_k-2\Delta_t\eta_{\rm lr}(\bar A_k\Umat-V)$ with $\Umat=\E\Umat^{(1)}_k$, $V=\E V^{(1)}_k$: the mean trajectory is the Euler discretisation of the flow of \cref{lem:gf-p1} with $\delta T=2\Delta_t\eta_{\rm lr}$.
\end{proposition}
\begin{proof}
$A_k$ is a function of draws $1,\dots,k-1$ only, hence independent of $(\Umat^{(1)}_k,V^{(1)}_k)$. Taking expectations, $\E[A_k\Umat^{(1)}_k]=\E[A_k]\,\E[\Umat^{(1)}_k]=\bar A_k\Umat$ and $\E[V^{(1)}_k]=V$. The recursion is linear in $A_k$, so no further approximation is involved. (Numerically, the mean over $400$ replicas of fresh-noise SGD at $(\dlat,\pwidth,\nsamp)=(4,30,12)$ matches the Euler scheme of the $\E_\eta$ flow to $1\%$, the Monte-Carlo error.)
\end{proof}

\begin{lemma}[Closed-form gradient flow; proved]\label{lem:gf-p1}
Let $L_{\nsamp}(A)=\tfrac1{2\nsamp}\sum_\mu\E_\eta\|A\phi(x^\mu_t)-y^\mu\|^2$ with $y^\mu=-\eta^\mu/\sqrt{\Delta_t}$, $\Umat=\tfrac1{\nsamp}\sum_\mu\E_\eta[\phi(x^\mu_t)\phi(x^\mu_t)^{\top}]=\sum_{i=1}^{\pwidth}\lambda_iv_iv_i^{\top}$ (eigenvalues sorted, $\lambda_i\ge0$), $V=\tfrac1{\nsamp}\sum_\mu\E_\eta[y^\mu\phi(x^\mu_t)^{\top}]$. Then (i) every row of $V$ lies in $\operatorname{range}(\Umat)$, so $A^\star:=V\Umat^{+}$ is the minimum-norm minimiser of $L_{\nsamp}$ and $A^\star\Umat=V$; (ii) the gradient flow $\dot A=-\nabla_AL_{\nsamp}(A)=-(A\Umat-V)$ from $A(0)=0$ has the unique solution $A(T)=A^\star(I-e^{-\Umat T})$; (iii) in the eigenbasis, $a_i(T):=A(T)v_i=a_i^\star(1-e^{-\lambda_iT})$ with $a_i^\star=Vv_i/\lambda_i$ for $\lambda_i>0$, and $a_i(T)\equiv0$ for $\lambda_i=0$. When $\Umat$ has full rank (fresh-noise expectation, (A2)) every one of the $\pwidth$ modes is active, including the tail block $\mathcal Z$, whose rates $\lambda_i\lesssim2\times10^{-6}$ (theory units, $\dlat=20$) give $\lambda_iT\approx0.06$--$0.2$ at the repository's $3\times10^5$ steps: slow, but not inert. (iv) Full-batch gradient descent with step $\eta_{\rm lr}$ on the logged loss $\|\sqrt{\Delta_t}s_A(x_t)+\eta\|^2$ (summed over coordinates, averaged over the batch) $=2\Delta_tL_{\nsamp}+{\rm const}$ is the Euler discretisation of this flow with $T=2\Delta_t\eta_{\rm lr}\times(\text{steps})$, faithful while $2\Delta_t\eta_{\rm lr}\lambda_1\ll1$; the repository's rule $\eta_{\rm lr}=0.01\dlat/\Delta_t$ gives $T=0.02\,\dlat\times(\text{steps})$, and the logged column $\tau_{\rm log}=\eta_{\rm lr}\times(\text{steps})$ satisfies $T=2\Delta_t\tau_{\rm log}$.
\end{lemma}
\begin{proof}
(i) If $u\perp\operatorname{range}(\Umat)$ then $0=u^{\top}\Umat u=\tfrac1{\nsamp}\sum_\mu\E_\eta(u^{\top}\phi(x^\mu_t))^2$, hence $u^{\top}\phi(x^\mu_t)=0$ a.s.\ for every $\mu$, hence $Vu=0$. Thus $V=V\Umat\Umat^{+}$ and $A^\star\Umat=V$. Since $L_{\nsamp}(A)=\tfrac12\operatorname{tr}(A\Umat A^{\top})-\operatorname{tr}(AV^{\top})+c_{\nsamp}=\tfrac12\operatorname{tr}\big((A-A^\star)\Umat(A-A^\star)^{\top}\big)-\tfrac12\operatorname{tr}(A^\star\Umat A^{\star\top})+c_{\nsamp}$, $A^\star$ minimises $L_{\nsamp}$ and is the unique minimiser with rows in $\operatorname{range}(\Umat)$. (ii) $\tfrac{d}{dT}(A-A^\star)=-(A-A^\star)\Umat$ is linear with constant coefficients, so $A(T)-A^\star=-A^\star e^{-\Umat T}$. (iii) Multiply on the right by $v_i$. The numerical value of $\lambda_{\mathcal Z}T$ uses $T=0.02\dlat\times3\times10^5=1.2\times10^5$ at $\dlat=20$ and tail eigenvalues $5\times10^{-7}$ (mean) to $1.8\times10^{-6}$ (max). (iv) $\nabla_A\|\sqrt{\Delta_t}A\phi+\eta\|^2=2\Delta_t(A\phi-y)\phi^{\top}$, so one step is $A\leftarrow A-2\Delta_t\eta_{\rm lr}(A\Umat-V)$ (in expectation, \cref{prop:sgdmean-p1}); the mode-$i$ multiplier $1-2\Delta_t\eta_{\rm lr}\lambda_i$ approximates $e^{-\lambda_i\delta T}$ when $2\Delta_t\eta_{\rm lr}\lambda_1\ll1$. Brute-force check at $(\dlat,\pwidth,\nsamp)=(4,30,12)$: closed form versus Euler with $h=10^{-3}$ agree to $10^{-4}$--$10^{-6}$.
\end{proof}

\begin{definition}[Residual spectral error and target masses]\label{def:rse-p1}
Let $\pi_i:=\|A^\star v_i\|^2=\|A^\star P_i\|_F^2$, $P_i=v_iv_i^{\top}$, be the target mass in mode $i$, $i=1,\dots,\pwidth$ ($\pi_i=0$ only if $\lambda_i=0$; for the full-rank $\Umat$ of (A2) the tail masses are the largest of all, $\pi_i=\|Vv_i\|^2/\lambda_i^2$). For a block $\mathcal B\subseteq\{1,\dots,\pwidth\}$ define
\[
E_{\mathcal B}(T):=\sum_{i\in\mathcal B}\pi_ie^{-2\lambda_iT},\qquad \Pi_{\mathcal B}:=E_{\mathcal B}(0)=\sum_{i\in\mathcal B}\pi_i,\qquad
\mathcal L_{\mathcal B}(T):=\tfrac12\sum_{i\in\mathcal B}\lambda_i\pi_ie^{-2\lambda_iT},
\]
and $E(T)=E_{\{1..\pwidth\}}(T)$. Because $\Pi_{\mathcal Z}\gg\Pi_{\mathcal M}\gg\Pi_{\mathcal N}\gg\Pi_{\mathcal S}$ (ratios $\sim10^1$--$10^3$ at $\dlat=20$: $\Pi_{\mathcal S}=44$, $\Pi_{\mathcal N}=1.1\times10^4$, $\Pi_{\mathcal M}=2.4\times10^7$, $\Pi_{\mathcal Z}=6\times10^8$ in a $(20,1280,500)$ simulation), the \emph{global} $E(T)$ is dominated by the tail and is not a useful observable; every timescale below is defined through a block-restricted $E_{\mathcal B}$ or through loss-weighted quantities, in which mode $i$ enters with weight $\lambda_i\pi_i=\|Vv_i\|^2/\lambda_i$ or $\lambda_i^2\pi_i=\|Vv_i\|^2$.
\end{definition}

\begin{proposition}[Exactness of the residual spectral error; proved]\label{prop:rse-p1}
Under \cref{lem:gf-p1}, for all $T\ge0$: $\|A(T)-A^\star\|_F^2=E(T)$ and $L_{\nsamp}(A(T))-L_{\nsamp}(A^\star)=\mathcal L_{\{1..\pwidth\}}(T)=\tfrac12\sum_{i\le\pwidth}\lambda_i\pi_ie^{-2\lambda_iT}$. For any block $\mathcal B$, $\|A(T)P_{\mathcal B}\|_F^2=\sum_{i\in\mathcal B}\pi_i(1-e^{-\lambda_iT})^2$, so the observable absorption deficit $D_{\mathcal B}(T):=\Pi_{\mathcal B}-\|A(T)P_{\mathcal B}\|_F^2$ satisfies $E_{\mathcal B}(T)\le D_{\mathcal B}(T)\le2\sqrt{\Pi_{\mathcal B}E_{\mathcal B}(T)}$. In particular the tail block satisfies $\|A(T)P_{\mathcal Z}\|_F^2\le T^2\|VP_{\mathcal Z}\|_F^2$ and $\mathcal L_{\mathcal Z}(0)-\mathcal L_{\mathcal Z}(T)\le T\|VP_{\mathcal Z}\|_F^2$ for all $T$.
\end{proposition}
\begin{proof}
By \cref{lem:gf-p1}(ii), $A(T)-A^\star=-\sum_ie^{-\lambda_iT}(A^\star v_i)v_i^{\top}$; the $v_i$ are orthonormal so $\|A(T)-A^\star\|_F^2=\sum_ie^{-2\lambda_iT}\pi_i$. The loss identity follows from $L_{\nsamp}(A)-L_{\nsamp}(A^\star)=\tfrac12\operatorname{tr}((A-A^\star)\Umat(A-A^\star)^{\top})$. For the absorption, $A(T)P_{\mathcal B}=\sum_{i\in\mathcal B}(1-e^{-\lambda_iT})(A^\star v_i)v_i^{\top}$, and $D_{\mathcal B}(T)=\sum_{i\in\mathcal B}\pi_i(2e_i-e_i^2)$ with $e_i=e^{-\lambda_iT}$; since $e_i^2\le2e_i-e_i^2\le2e_i$, $E_{\mathcal B}\le D_{\mathcal B}\le2\sum_{\mathcal B}\pi_ie_i\le2\sqrt{\Pi_{\mathcal B}E_{\mathcal B}}$ by Cauchy--Schwarz. Tail bounds: $(1-e^{-x})^2\le x^2$ and $1-e^{-2x}\le2x$ with $\pi_i\lambda_i^2=\|Vv_i\|^2$. Brute-force check at $(4,30,12)$: all three identities hold to six decimals at $T=0.3,2,10$.
\end{proof}

\subsubsection{Operating point}

\begin{lemma}[Operating point along the sweep; proved]\label{lem:q-p1}
Under (A3), with $M_t=\tfrac1{\nsamp}\sum_\mu\E[x_t^\mu x_t^{\mu\top}]$, $\alpha_t^2=e^{-2t}(s^2/\dint+\sigsig^2)+\Delta_t$ and $\beta_t^2=e^{-2t}\signoise^2+\Delta_t$,
\begin{equation}\label{eq:q-p1}
\E\,q(\dlat)=\frac{\E\operatorname{tr}M_t}{\dlat}=\beta_t^2+\frac{e^{-2t}\big(s^2+\dint(\sigsig^2-\signoise^2)\big)}{\dlat}=\beta_t^2+\frac{\dint(\alpha_t^2-\beta_t^2)}{\dlat},
\end{equation}
so $q(\dlat)$ decreases monotonically to $\beta_t^2\ge\Delta_t>0$ as $\dlat\to\infty$ (never to $0$). Along it $a_1(q(\dlat))$ is increasing in $\dlat$ (proved) and $a_*(q(\dlat))^2$ is decreasing in $\dlat$ (verified numerically for $q\in[0.01,20]$; see proof). At $t=0.01$, $\signoise=0.5$ (Gauss--Hermite, $200$ nodes): $\beta_t^2=0.2649$, $\alpha_t^2=2.7644$, and
\[
\begin{array}{c|ccccc}
\dlat&5&20&60&200&\infty\\\hline
q&2.764&0.890&0.473&0.327&0.265\\
a_1(q)^2&0.180&0.393&0.541&0.625&0.674\\
a_*(q)^2&0.0785&0.0235&0.0090&0.0047&0.0030
\end{array}
\]
At $\signoise=0.01$ the sweep reaches only $q(200)=0.0885$ ($a_1^2=0.854$, $a_*^2=2.6\times10^{-4}$, $18\times$ below its $\signoise=0.5$ value); the limit is $q(\infty)=\beta_t^2=0.0199$. As $q\to0$, $a_1(q)=1-q+2q^2+O(q^3)$ and $a_*(q)^2=\tfrac23q^3+O(q^4)$; the $O(q^4)$ term is large ($\tfrac23q^3$ overshoots $a_*^2$ by $16\%$ at $q=0.02$ and by $77\%$ at $q=0.0885$), so the expansion is not quantitatively usable at any point of either sweep and the tabulated values should be used instead.
\end{lemma}
\begin{proof}
$\E\operatorname{tr}M_t=e^{-2t}\E\|x^\mu\|^2+\dlat\Delta_t$ and $\E\|x^\mu\|^2=s^2+\dint\sigsig^2+(\dlat-\dint)\signoise^2$ ($Q$ orthogonal). Divide by $\dlat$ and regroup; the second form uses $\dint\alpha_t^2=e^{-2t}(s^2+\dint\sigsig^2)+\dint\Delta_t$. $a_1(q)=\E[1-\tanh^2(\sqrt qz)]$ is decreasing in $q$ because $\tanh^2(\sqrt q z)$ is increasing in $q$ pointwise; hence $a_1(q(\dlat))$ increases with $\dlat$. For $a_*^2$: it equals $\sum_{k\ge3}c_k(q)^2$ with $c_k(q)=\E[\tanh(\sqrt qz)\mathrm{He}_k(z)]/\sqrt{k!}$ (Hermite coefficients; $c_1^2=a_1^2q$ by Stein's lemma), and each $c_k(q)^2$ is not obviously monotone, so we do not claim a proof; on a $200$-node Gauss--Hermite grid $a_*^2$ is strictly increasing in $q$ on $[0.01,20]$ and bounded by $1$. Expansion: $\tanh u=u-u^3/3+2u^5/15+O(u^7)$ gives $\E\tanh^2(\sqrt qz)=q-2q^2+\tfrac{17}3q^3+O(q^4)$, $a_1=1-q+2q^2+O(q^3)$, $a_1^2q=q-2q^2+5q^3+O(q^4)$, $a_*^2=\tfrac23q^3+O(q^4)$. Numerical check: $a_*^2$ vs $\tfrac23q^3$ $=6.2\times10^{-7}$ vs $6.7\times10^{-7}$ ($q=0.01$), $4.6\times10^{-6}$ vs $5.3\times10^{-6}$ ($0.02$), $3.5\times10^{-4}$ vs $6.7\times10^{-4}$ ($0.1$). The identity $a_1^2q+a_*^2=\E\tanh^2(\sqrt qz)$ holds exactly by definition and is the trace check of \cref{app:theory-fourbulk}.
\end{proof}

\noindent\emph{Trace check against the released spectra.} $\operatorname{tr}\Umat/\pwidth=0.545\to0.209$ over $\dlat=5\to200$ versus $\E\tanh^2(\sqrt qz)=0.576\to0.209$; the sample-block mean is $0.72$--$0.89\times a_*^2/\nsamp$ (Jensen dispersion of the per-sample $q_\mu$); the noise-dimension medians $5.22\to9.15$ (code units) versus the leading-order $a_1^2\beta_t^2\psi_p=4.76\to10.6$ overshoot by $10$--$15\%$ at large $\dlat$ and saturate from $\dlat=100$. This is why no edge and no timescale below is $\dlat$-independent.

\subsubsection{Clocks}

\begin{remark}[Three clocks and their conversion]\label{rem:clock-p1}
The paper reports timescales in three units which differ by $\dlat$-dependent factors; mixing them produces contradictory statements about whether $\taugen$ grows or falls with $\dlat$. With $\lambda_{\rm code}=\pwidth\lambda$ the eigenvalue in the repository's normalisation (which omits $\pwidth^{-1/2}$ in $\phi$) and $\psi_p=\pwidth/\dlat=64$:
\[
\begin{array}{l|l|l}
\text{unit} & \text{relation} & \text{e-fold time of mode }i\\\hline
\text{flow time }T & \text{\cref{lem:gf-p1}} & 1/\lambda_i=\pwidth/\lambda_{i,\rm code}\\
\text{logged }\tau_{\rm log}=\eta_{\rm lr}\times\text{steps} & T=2\Delta_t\tau_{\rm log} & 1/(2\Delta_t\lambda_i)\\
\text{repository steps} & T=0.02\,\dlat\times\text{steps} & \psi_p/(0.02\lambda_{i,\rm code})=3200/\lambda_{i,\rm code}
\end{array}
\]
Spectral times in flow units are $\pwidth/\lambda_{\rm code}$ (\cref{app:theory-fourbulk}): on the released $\signoise=0.5$ spectra $\taugen=1/\lambda^{\mathcal S}_{\min}$ grows $11.0\times$ and the index-based onset $1/\lambda^{\mathcal M}_{\max}$ grows $221\times$ over $\dlat=5\to200$; these are \emph{not} gradient-descent steps. In repository steps a per-step decay rate is $0.02\lambda_{\rm code}/\psi_p=3.1\times10^{-4}\lambda_{\rm code}$, independent of $\dlat$ at fixed $\psi_p$, so step counts are read off code eigenvalues directly: signal e-fold $158\to43$ steps, noise-dimension e-fold $3200/\lambda^{\mathcal N}_{\min,\rm code}=620$ ($\dlat=20$) to $1509$ ($\dlat=200$) steps, index-based sample-top e-fold $3200/\lambda^{\mathcal M}_{\max,\rm code}=3.3\times10^3\to1.8\times10^4$ steps. Ratios of timescales are identical in all three clocks and are the only quantities compared across $\dlat$ below. (Sources: learning rate $0.01\dlat/\Delta_t$, loss summed over coordinates and averaged over the batch, \texttt{compute\_U} omitting $\pwidth^{-1/2}$; code eigenvalues $\lambda^{\mathcal S}_{\min,\rm code}=20.30,41.98,73.88$ at $\dlat=5,20,200$; $\lambda^{\mathcal N}_{\min,\rm code}=5.16,2.12$ at $\dlat=20,200$; $\lambda^{\mathcal M}_{\max,\rm code}=0.957,0.667,0.173$ at $\dlat=5,20,200$.)
\end{remark}

\subsubsection{Generalisation time}

\begin{definition}[Generalisation time]\label{def:taugen-p1}
For a tolerance $\varepsilon\in(0,1)$,
\[
\taugen(\varepsilon):=\inf\{T\ge0:\ E_{\mathcal S}(T)\le\varepsilon\,\Pi_{\mathcal S}\},\qquad
\tau_{\rm pop}(\varepsilon):=\inf\{T\ge0:\ E_{\mathcal S\cup\mathcal N}(T)\le\varepsilon\,\Pi_{\mathcal S\cup\mathcal N}\}.
\]
$\taugen$ is the time at which the signal-subspace score (the part that determines sample quality) is absorbed to relative tolerance $\varepsilon$; $\tau_{\rm pop}$ additionally requires the noise-dimension block (target $s^\star_{\mathcal N}(x)=-x_{\mathcal N}/\beta_t^2$, exactly linear) to be absorbed. In the isotropic case only $\tau_{\rm pop}$ is defined. Both are block functionals of $\mathcal S$ (resp.\ $\mathcal S\cup\mathcal N$) only and are unaffected by the tail block $\mathcal Z$ or by the sample block.
\end{definition}

\begin{proposition}[$\taugen$ is set by the slowest signal mode; proved given the block scale of \cref{app:theory-fourbulk}]\label{prop:taugen-p1}
Let $\mathcal B\in\{\mathcal S,\mathcal S\cup\mathcal N\}$, $\lambda^{\mathcal B}_{\min}=\min\{\lambda_i:i\in\mathcal B,\pi_i>0\}$ attained at $i_\circ$, and $\theta_{\mathcal B}:=\log(\Pi_{\mathcal B}/\pi_{i_\circ})\ge0$. Then for every $\varepsilon\in(0,1)$,
\begin{equation}\label{eq:taugen-bounds-p1}
\frac{\log(1/\varepsilon)-\theta_{\mathcal B}}{2\lambda^{\mathcal B}_{\min}}\ \le\ \tau_{\mathcal B}(\varepsilon)\ \le\ \frac{\log(1/\varepsilon)}{2\lambda^{\mathcal B}_{\min}},
\end{equation}
and for the observable deficit $D_{\mathcal B}$, $\tau^D_{\mathcal B}(\varepsilon):=\inf\{T:D_{\mathcal B}(T)\le\varepsilon\Pi_{\mathcal B}\}$ satisfies
\[
\frac{\log(1/\varepsilon)-\theta_{\mathcal B}}{\lambda^{\mathcal B}_{\min}}\ \le\ \tau^D_{\mathcal B}(\varepsilon)\ \le\ \frac{\log(2/\varepsilon)}{\lambda^{\mathcal B}_{\min}} .
\]
The slowest signal mode is not the block scale $a_1(q)^2\alpha_t^2/\dlat$ but
\begin{equation}\label{eq:lammin-signal-p1}
\lambda^{\mathcal S}_{\min}=\frac{a_1(q(\dlat))^2\,(e^{-2t}\ell_{\min}+\Delta_t)}{\dlat}\,(1+o(1)),
\end{equation}
where $\ell_{\min}$ is the smallest eigenvalue of the signal-subspace second-moment matrix $(\hat C+\hat S)_{\mathcal S}$ of the data (\cref{app:theory-fourbulk}); $\ell_{\min}\ge\sigsig^2(1+o(1))$ and $\ell_{\min}$ is $\dlat$-independent when the centre configuration is held fixed along the sweep (repository: $\lambda^{\mathcal S}_{\min}/\bar\lambda^{\mathcal S}\approx0.6$). Consequently, in gradient-flow time $\taugen\propto\dlat/(a_1(q(\dlat))^2(e^{-2t}\ell_{\min}+\Delta_t))$ \emph{grows} with $\dlat$ (repository spectra: $1/\lambda^{\mathcal S}_{\min}$ grows $10.9\times$ over $\dlat=5\to200$), whereas in repository steps $\taugen^{\rm steps}=\taugen/(0.02\dlat)\propto1/a_1(q(\dlat))^2$ \emph{falls} ($3200/\lambda^{\mathcal S}_{\min,\rm code}=158\to43$ steps for $\log(1/\varepsilon)=2$). Neither is $\dlat$-independent, and the two clocks must not be mixed (\cref{rem:clock-p1}).
\end{proposition}
\begin{proof}
Upper bound: $E_{\mathcal B}(T)\le\Pi_{\mathcal B}e^{-2\lambda^{\mathcal B}_{\min}T}$, so $T=\log(1/\varepsilon)/(2\lambda^{\mathcal B}_{\min})$ is admissible. Lower bound: $E_{\mathcal B}(T)\ge\pi_{i_\circ}e^{-2\lambda^{\mathcal B}_{\min}T}>\varepsilon\Pi_{\mathcal B}$ for $T<(\log(1/\varepsilon)-\theta_{\mathcal B})/(2\lambda^{\mathcal B}_{\min})$. For $D_{\mathcal B}$ use $\pi_{i_\circ}e^{-\lambda_{\min}T}\le D_{\mathcal B}(T)\le2\Pi_{\mathcal B}e^{-\lambda_{\min}T}$ (proof of \cref{prop:rse-p1}); the factor $2$ is genuinely attained (one mode, $\pi=\lambda=1$: $D(\log(1/\varepsilon))=2\varepsilon-\varepsilon^2>\varepsilon$), which is why $\log2$ sits in the upper bound. The formula for $\lambda^{\mathcal S}_{\min}$ is the linear-response (Gaussian-equivalent) block scale $a_1(q)^2\times$(data eigenvalue)$/\dlat$ applied eigenvalue by eigenvalue inside the finite-rank signal block, where the data eigenvalues of $M_t$ are $e^{-2t}\ell_j+\Delta_t$, $\ell_j$ the eigenvalues of $(\hat C+\hat S)_{\mathcal S}$, whose mean is $s^2/\dint+\sigsig^2$ and whose minimum is at least the within-cluster variance; its validity is \cref{prop:fourbulk,lem:fourbulk-bbp}, checked here numerically: $(20,1280,500)$ simulation, signal mean $0.0526$ (scale $0.0545$), min $0.0314$. The step-count numbers use \cref{lem:gf-p1}(iv) with code eigenvalues $\lambda^{\mathcal S}_{\min,\rm code}=20.3$ ($\dlat=5$) and $73.9$ ($\dlat=200$).
\end{proof}

\paragraph{Why the signal block.}
The noise-dimension target is exactly linear, $s^\star_{\mathcal N}(x_t)=-x_{t,\mathcal N}/\beta_t^2$ (cluster means have no $\mathcal N$-component; \cref{app:theory-linear}), with per-coordinate second moment $1/\beta_t^2=3.8$ at $\signoise=0.5$, whereas the signal-subspace target has per-coordinate second moment at most $1/(e^{-2t}\sigsig^2+\Delta_t)=1.0$ (convexity of Fisher information for the Gaussian mixture); hence for $\dlat\gg\dint$ the reducible population loss is dominated by $\mathcal N$ and a test-loss plateau measures $\tau_{\rm pop}$ (rate $\lambda^{\mathcal N}_{\min}\simeq a_1(q)^2\beta_t^2/\dlat$) rather than $\taugen$ (rate $\lambda^{\mathcal S}_{\min}$). Sample quality (score error in the signal subspace, FID after decoding) is a functional of the $\mathcal S$-block score only, so $\taugen$ is defined through $\mathcal S$. The block $\mathcal N$ is the buffer: its absorption costs time and buys no sample quality.

\subsubsection{Generalisation gap and memorisation times}

\begin{proposition}[Exact spectral generalisation gap; exact formula proved, reduction under (A6$'$)]\label{prop:gap-p1}
Let $L(A)=\tfrac12\E\|A\phi(x_t)-y\|^2$ be the population objective (its held-out estimate is the logged test loss), $\bar\Umat=\E[\phi\phi^{\top}]$, $\bar V=\E[y\phi^{\top}]=\E[s^\star(x_t)\phi(x_t)^{\top}]$, and $G(T):=[L(A(T))-L_{\nsamp}(A(T))]-[L(0)-L_{\nsamp}(0)]$. With $e_i=e^{-\lambda_iT}$, $\bar u_{ij}=v_i^{\top}\bar\Umat v_j$, $\gamma_i$ as in (A6$'$), exactly
\begin{equation}\label{eq:gap-exact-p1}
G(T)=\tfrac12\sum_{i,j\le\pwidth}(1-e_i)(1-e_j)\,\bar u_{ij}\,\langle a_i^\star,a_j^\star\rangle-\sum_{i\le\pwidth}(1-e_i)\gamma_i\lambda_i\pi_i+\tfrac12\sum_{i\le\pwidth}\lambda_i\pi_i(1-e_i^2).
\end{equation}
Define the diagonal per-mode contribution $g_i(T):=\tfrac12\pi_i\big[\bar u_i(1-e_i)^2-2\gamma_i\lambda_i(1-e_i)+\lambda_i(1-e_i^2)\big]$, which vanishes identically when $\bar u_i=\lambda_i$ and $\gamma_i=1$. Under (A6$'$),
\[
G(T)=\sum_{i\in\mathcal M\cup\mathcal Z}g_i(T)+o(G_{\rm thr})\quad\text{uniformly in }T .
\]
For the sub-block $\mathcal M_{\rm blk}:=\{i\in\mathcal M:\gamma_i\le\gamma_0\}$ (\cref{def:taumem-p1}) write $G_{\rm blk}(T):=\tfrac12\sum_{i\in\mathcal M_{\rm blk}}\pi_i[\bar u_i(1-e_i)^2+\lambda_i(1-e_i^2)]$; then $\big|\sum_{\mathcal M_{\rm blk}}g_i-G_{\rm blk}\big|\le\gamma_0\sum_{\mathcal M_{\rm blk}}\lambda_i\pi_i$, and $G_{\rm blk}$ is nonnegative, increasing, $G_{\rm blk}(0)=0$, $G_\infty^{\rm blk}=\tfrac12\sum_{\mathcal M_{\rm blk}}\pi_i(\bar u_i+\lambda_i)$, $G_{\rm blk}'(0)=\sum_{\mathcal M_{\rm blk}}\lambda_i^2\pi_i$, and
\begin{equation}\label{eq:gap-bounds-p1}
G_\infty^{\rm blk}\big(1-e^{-\lambda^{\rm blk}_{\min}T}\big)^2\ \le\ G_{\rm blk}(T)\ \le\ G_\infty^{\rm blk}\big(1-e^{-2\lambda^{\rm blk}_{\max}T}\big).
\end{equation}
The tail contribution is bounded for all $T$ by $\big|\sum_{\mathcal Z}g_i(T)\big|\le\tfrac12\|VP_{\mathcal Z}\|_F^2\big(T^2\sup_{\mathcal Z}\bar u_i+2T(1+\sup_{\mathcal Z}|\gamma_i|)\big)$.
\end{proposition}
\begin{proof}
Write $A(T)=\sum_{i\le\pwidth}(1-e_i)a_i^\star v_i^{\top}$ (\cref{lem:gf-p1}). Then $\operatorname{tr}(A\bar\Umat A^{\top})=\sum_{ij}(1-e_i)(1-e_j)\bar u_{ij}\langle a_i^\star,a_j^\star\rangle$, $\operatorname{tr}(A\bar V^{\top})=\sum_i(1-e_i)\langle a_i^\star,\bar Vv_i\rangle=\sum_i(1-e_i)\gamma_i\lambda_i\pi_i$, $\operatorname{tr}(A\Umat A^{\top})=\sum_i\lambda_i\pi_i(1-e_i)^2$ and $\operatorname{tr}(AV^{\top})=\sum_i(1-e_i)\langle a_i^\star,Vv_i\rangle=\sum_i\lambda_i\pi_i(1-e_i)$ since $Vv_i=\lambda_ia_i^\star$. The additive constants cancel in $G$; the empirical part is $\sum_i\lambda_i\pi_i[(1-e_i)-\tfrac12(1-e_i)^2]=\tfrac12\sum_i\lambda_i\pi_i(1-e_i^2)$. (Brute-force check at $(4,30,12)$: six decimals at $T=0.3,2,10$.) Under (A6$'$) the off-diagonal double sum is $o(G_{\rm thr})$ and, for $i\in\mathcal S\cup\mathcal N$, $|g_i(T)|\le\tfrac12\pi_i\lambda_i(|\bar u_i/\lambda_i-1|+2|\gamma_i-1|)$ (use $(1-e_i)^2\le1$, $1-e_i\le1$ and $(1-e_i)^2-2(1-e_i)+(1-e_i^2)=0$), whose sum is $o(G_{\rm thr})$. The $\gamma_0$ bound uses $1-e_i\le1$. Properties of $G_{\rm blk}$: each bracket is $0$ at $T=0$, increasing in $T$, and tends to $\bar u_i+\lambda_i$; $\tfrac{d}{dT}(1-e_i)^2|_0=0$, $\tfrac{d}{dT}(1-e_i^2)|_0=2\lambda_i$. Bounds: $(1-e)^2\le1-e^2$ on $[0,1]$, so the normalised bracket lies in $[(1-e_i)^2,1-e_i^2]$; multiply by $\tfrac12\pi_i(\bar u_i+\lambda_i)$ and sum. Tail: $(1-e^{-x})^2\le x^2$, $1-e^{-x}\le x$, $1-e^{-2x}\le2x$ with $\pi_i\lambda_i^2=\|Vv_i\|^2$.
\end{proof}

\begin{definition}[Spike/bulk split of the sample block; memorisation times]\label{def:taumem-p1}
Fix $\gamma_0\in(0,\tfrac12]$. Split $\mathcal M=\mathcal M_{\rm spk}\cup\mathcal M_{\rm blk}$ with $\mathcal M_{\rm spk}:=\{i\in\mathcal M:\gamma_i>\gamma_0\}$ (population-aligned sample modes: their empirical optimum $a_i^\star$ has a nonnegligible projection on the population target, so absorbing them \emph{lowers} the test loss) and $\mathcal M_{\rm blk}$ the remainder. For clustered data $\mathcal M_{\rm spk}$ contains the cluster-residual spikes above the Marchenko--Pastur edge ($\lambda^{\mathcal M}_{\max}/\bar\lambda^{\mathcal M}=21.5,15.4,5.2,3.0,1.8$ at $\dlat=5,20,60,100,200$ on the repository spectra), and its size is expected to grow as $\psi_n\downarrow1$.
\begin{enumerate}
\item[(i)] \emph{Onset:} $\taumem^{\rm onset}:=1/\lambda^{\rm blk}_{\max}$, the e-folding time of the fastest \emph{bulk} sample mode. The index-based statistic $1/\lambda^{\mathcal M}_{\max}$ (used in the tables) is a lower bound on it, with equality when $\mathcal M_{\rm spk}=\emptyset$.
\item[(ii)] \emph{Gap-based:} for $\delta\in(0,1)$, $\taumem^{\rm gap}(\delta):=\inf\{T:G_{\rm blk}(T)\ge\delta\,G^{\rm blk}_\infty\}$. By \cref{prop:gap-p1},
\[
\frac{\log\frac1{1-\delta}}{2\lambda^{\rm blk}_{\max}}\ \le\ \taumem^{\rm gap}(\delta)\ \le\ \frac{\log\frac1{1-\sqrt\delta}}{\lambda^{\rm blk}_{\min}},
\]
so the onset controls the earliest time the gap can reach a given fraction of its asymptote and the slowest bulk mode its completion; the position of $\taumem^{\rm gap}(\delta)$ within the interval depends on the within-bulk distribution of $(\lambda_i,\pi_i,\bar u_i)_{i\in\mathcal M_{\rm blk}}$. The tail block $\mathcal Z$ enters $G$ through the bound of \cref{prop:gap-p1} and, on the repository's horizon, cannot be dropped from Frobenius-norm quantities: at $(20,1280,500)$ and $3\times10^5$ steps the absorbed tail mass $\sum_{\mathcal Z}\pi_i(1-e_i)^2=2.1\times10^6$ is two thirds of the absorbed sample-block mass $3.2\times10^6$; its loss weight, however, is $\le\sup_{\mathcal Z}(\bar u_i,\lambda_i)\lesssim2\times10^{-6}$ per unit mass.
\item[(iii)] \emph{Fraction-level (placeholder, completed in \cref{app:theory-bridge,app:theory-predictor}):} the paper's primary statistic is the memorised fraction $F(T)=\tfrac1{\nsamp}\sum_\mu\mathbf 1[\rho_\mu(T)<1/3]$ with $\rho_\mu$ the nearest-neighbour ratio of \citet{somepalli2023diffusion}. Its spectral surrogate is $\hat F(T)=\tfrac1{\nsamp}\sum_\mu\mathbf 1[m_\mu(T)\ge\vartheta]$, $m_\mu(T)=\sum_{i\in\mathcal M_{\rm blk}}\pi_{i,\mu}(1-e^{-\lambda_iT})^2/\sum_{i\in\mathcal M_{\rm blk}}\pi_{i,\mu}$, where $\pi_{i,\mu}$ is the share of mode $i$'s target mass attributable to training point $\mu$; a single sample edge makes $\hat F$ a step function, so the map $\mu\mapsto(\pi_{i,\mu})_i$ and the spread of $\lambda_i$ over $\mathcal M_{\rm blk}$ are required and are supplied in \cref{app:theory-predictor}.
\end{enumerate}
The displayed bounds are \cref{prop:gap-p1} inverted. The spike ratios are $\lambda^{\mathcal M}_{\max}/\operatorname{mean}(\lambda_{\mathcal M})$ from the released $\signoise=0.5$ spectra with the sample block cut at index $\nsamp$ (index-based; $\gamma_i$ has not yet been computed, see Prediction~5).
\end{definition}

\paragraph{Consequences.}
(i) Under (A6$'$) the population blocks contribute nothing to the gap at the level $G_{\rm thr}$: the model's coefficient on a population direction converges to the population-optimal one, so train and test loss fall together on $\mathcal S\cup\mathcal N$ (the ``both curves drop'' phase of \cref{sec:rfnn}). (ii) A pure sample mode ($\gamma_i=0$) raises the test loss by $\tfrac12\bar u_i\pi_i(1-e^{-\lambda_iT})^2$ and lowers the train loss by $\tfrac12\lambda_i\pi_i(1-e^{-2\lambda_iT})$; a population-aligned sample mode ($\gamma_i\approx1$, $\bar u_i\approx\lambda_i$) lowers both. The repository logs show the latter: the test loss keeps falling until $35$k ($\dlat=40$), $55$k ($100$), $90$k ($200$) steps, $20$--$60\times$ later than the slowest population e-fold ($3200/\lambda^{\mathcal N}_{\min,\rm code}=1.5$k steps at $\dlat=200$), so the top of the index-based sample block ($\lambda_{\rm code}\sim0.1$--$0.17$) is in $\mathcal M_{\rm spk}$ at large $\dlat$. The squared form $(1-e^{-\lambda_iT})^2$ is the origin of the shape used heuristically in the released predictor (\cref{sec:spectral-predictor}), but there the eigenvalues are those of $M_t$, not of the sample block, a discrepancy addressed in \cref{app:theory-predictor}. (iii) $\taumem^{\rm onset}$ and $\taugen$ are single-eigenvalue functionals, so $\taumem^{\rm onset}/\taugen(\varepsilon)=2\lambda^{\mathcal S}_{\min}/(\lambda^{\rm blk}_{\max}\log(1/\varepsilon))\,(1+O(\theta_{\mathcal S}/\log(1/\varepsilon)))$ is clock-free. On the released spectra with the \emph{index-based} sample top, $\lambda^{\mathcal S}_{\min}/\lambda^{\mathcal M}_{\max}=34,63,143,211,279,340,416,426$ for $\dlat=10,20,40,60,80,100,150,200$ (monotone; $\dlat=5$ gives $21$ but there $\pwidth<\nsamp$); since $\lambda^{\rm blk}_{\max}\le\lambda^{\mathcal M}_{\max}$ these are lower bounds on the bulk-onset ratio, whose monotonicity is to be re-established after the $\gamma$-cut (Prediction~5). No statement of the form ``$\taumem-\taugen\ge(\dlat-\dint)\times(\text{a timescale})$'' follows from a diagonal flow, since modes decay in parallel; the buffer is a statement about the ratio of two eigenvalues (\cref{app:theory-buffer}).

\subsubsection{Measured proxies}

\begin{remark}[Correspondence with the measured proxies; sketched]\label{rem:proxy-p1}
(a) \emph{Gap proxy.} The logged RFNN loss is per-coordinate, $\ell=\|\sqrt{\Delta_t}s_A+\eta\|^2/\dlat$, so $\ell_{\rm test}-\ell_{\rm train}=(2\Delta_t/\dlat)G(T)+{\rm const}$ and ``gen-gap $>0.02$'' is the crossing $G(T)=0.01\dlat/\Delta_t$: an \emph{absolute} threshold, i.e.\ $\taumem^{\rm gap}(\delta_{\rm eff})$ with $\delta_{\rm eff}=0.01\dlat/(\Delta_tG_\infty)$ varying with $\dlat$ through $G_\infty$ (final gaps $0.019\to0.093$ over $\dlat=5\to200$). Moreover the logged train loss uses one fresh noise draw per evaluation and the test loss one fixed draw, so the logged gap carries a noise floor of order $\sqrt{2/(\nsamp\dlat)}$: $0.028$ ($\dlat=5$), $0.014$ ($20$), $0.0045$ ($200$), comparable to the $0.02$ threshold for $\dlat\le20$; this, together with $\delta_{\rm eff}$, accounts for the erratic $\taumem^{\rm gap>0.02}=20,15,25,10,10$k steps at $\dlat=5..60$. Any re-extraction should use a many-draw gap or the exact $G(T)$ of \cref{prop:gap-p1}.
(b) \emph{Test-loss proxy.} ``Test loss within $5\%$ of its minimum'' fires at the first evaluation ($5$k-step grid) for every $\dlat$. The reason is the clock, not a vacuous tolerance: per-step decay rates are $0.02a_1(q)^2\alpha_t^2\approx0.01$--$0.035$ for signal modes and $0.02a_1(q)^2\beta_t^2\approx0.001$--$0.0033$ for noise-dimension modes ($\dlat=5\to200$), so both population blocks are absorbed within $\sim10^2$--$1.5\times10^3$ steps, inside the first evaluation interval. The tolerance is \emph{not} larger than the reducible loss except at $\dlat\le10$: the reducible per-coordinate loss is bounded by $\Delta_t[\dint/(e^{-2t}\sigsig^2+\Delta_t)+(\dlat-\dint)/\beta_t^2]/\dlat=0.020,0.047,0.061,0.068,0.073$ at $\dlat=5,10,20,40,200$ (measured $1-\ell_{\min}=0.052$--$0.063$ for $\dlat\ge40$), against $0.05\,\ell_{\min}=0.049,0.048,0.047,0.047,0.046$. Thus for $\dlat\ge20$ the $5\%$ proxy fires when $65$--$80\%$ of the reducible loss is absorbed, and at $\signoise=0.01$ (reducible $0.5$--$0.97$) it is informative throughout. The tighter $1\%$ and $2\%$ crossings are \emph{not} on the grid floor: $15/30/50$k and $10/20/35$k steps at $\dlat=40/100/200$, and the test-loss argmin is at $35/55/90$k steps, $20$--$60\times$ later than the slowest population e-fold ($1.5$k steps at $\dlat=200$). Those late decreases come from modes with $\lambda_{\rm code}\sim0.1$, i.e.\ the top of the sample block ($\lambda^{\mathcal M}_{\max,\rm code}=0.17$ at $\dlat=200$), which are therefore in $\mathcal M_{\rm spk}$ ($\gamma_i$ not small); they measure neither $\taugen$ nor $\tau_{\rm pop}$ but the absorption of population-aligned sample modes, and they are the observable that determines the $\gamma_0$-cut of \cref{def:taumem-p1}. For the same reason the reducible population loss is dominated by $\mathcal N$ for $\dlat\gg\dint$: per-coordinate target second moment $1/\beta_t^2=3.8$ on $\mathcal N$ against at most $1/(e^{-2t}\sigsig^2+\Delta_t)=1.0$ on $\mathcal S$ (Fisher-information convexity of the mixture score), so a test-loss plateau tracks $\tau_{\rm pop}$, not $\taugen$. The signal-block absorption deficit $D_{\mathcal S}(T)$, logged at fine resolution, is the observable that tests \cref{def:taugen-p1}; the repository's per-bulk absorption routine (\texttt{bulk\_indices}) has correct signal and noise-dimension slices but still cuts sample/tail at $\dlat+\nsamp$, so its sample and rank-null masses are not the blocks $\mathcal M,\mathcal Z$ of (A5) and must be re-cut at index $\nsamp$.
(c) For the MLP and real-data sweeps the memorised fraction $F(T)$ crossing $0.01$ (synthetic) or a fixed level (real) is the proxy for \cref{def:taumem-p1}(iii). Ratios $\taumem/\taugen$ are clock-independent; absolute step counts carry the factor $(0.02\,\dlat)^{-1}$ of \cref{lem:gf-p1}(iv) (\cref{rem:clock-p1}).
\end{remark}

\noindent Concretely, the plotting code defines $\taugen^{\rm meas}$ = first step with test loss $\le1.05\times$ its run minimum; $\taumem^{\rm meas}$ = first step with memorised fraction $>0.01$ when available (MLP, real data), otherwise first step with gen-gap $>0.02$ (RFNN). In the RFNN logs (seed 42, evaluation every $5\,000$ steps) $\taugen^{\rm meas}\le10^4$ steps for every $\dlat$; $\taumem^{\rm meas}$ moves from $2\times10^4$ ($\dlat=5$) to $8\times10^4$ ($\dlat=200$) and the final gap from $0.019$ to $0.093$; the theory attributes the modest growth to the $\dlat$-dependence of $G_\infty$ entering $\delta_{\rm eff}$ and to the single-draw noise floor $\sqrt{2/(\nsamp\dlat)}$, and predicts a much steeper growth for a relative-threshold, many-draw extraction on $\mathcal M_{\rm blk}$.

\paragraph{Fraction-level definition.}
The fraction $F(T)$ is a per-training-point statistic. Writing $\pi_i=\sum_\mu\pi_{i,\mu}$ for a decomposition of each bulk sample mode's target mass over training points, point $\mu$ is memorised once the absorbed share $m_\mu(T)$ of \cref{def:taumem-p1}(iii) exceeds a level $\vartheta$ set by the nearest-neighbour test; since the empirical-score reverse dynamics condense onto a training point once its weight dominates, which at $t=0.01$ holds for all $\dlat$ in the sweep ($\Delta_t=0.02\ll d_{\rm NN}^2/(2\log\nsamp)$, \citealp{biroli2024dynamical}), memorisation at fixed $t$ is a question of how much bulk sample-mode mass has been absorbed. A single edge makes $m_\mu$ identical across $\mu$ and $\hat F$ a step; the required spread of $(\lambda_i,\pi_{i,\mu})$ over $\mathcal M_{\rm blk}$ is derived in \cref{app:theory-bridge,app:theory-predictor}.

\paragraph{Predictions and how to test them.}
\begin{enumerate}
\item \emph{$\taugen$ in flow time is set by $\lambda^{\mathcal S}_{\min}=a_1(q)^2(e^{-2t}\ell_{\min}+\Delta_t)/\dlat$ and grows with $\dlat$ ($10.9\times$ over $\dlat=5..200$), whereas in repository steps it falls ($158\to43$).} Test: log the signal-block absorption deficit $D_{\mathcal S}$ every $\sim50$ steps, fit its late-time exponential rate, and compare with $3.1\times10^{-4}\lambda^{\mathcal S}_{\min,\rm code}$ per step; compare $\lambda^{\mathcal S}_{\min}$ across $\dlat$ with \eqref{eq:lammin-signal-p1} using $\ell_{\min}$ from the realised $(\hat C+\hat S)_{\mathcal S}$. Zero free parameters.
\item \emph{Zero-parameter reconstruction of the logged curves as expected dynamics:} with $\Umat,V$ the $\E_\eta$ operators, the mean train-loss excess is $\tfrac12\sum_i\lambda_i\pi_ie^{-2\lambda_iT}$ exactly and the mean gap is \eqref{eq:gap-exact-p1} with $T=0.02\dlat\times$steps, tail block included. Test: re-run one RFNN configuration saving $A$ every $500$ steps; compute $\Umat,V,\bar\Umat,\bar V$ (held-out data, many draws), $A^\star,\pi_i,\gamma_i,\bar u_{ij}$; overlay the closed forms on the logged curves converted by $2\Delta_t/\dlat$. A train-loss mismatch falsifies the Euler/expected-flow correspondence; a gap mismatch falsifies (A6$'$).
\item \emph{The absolute-threshold proxy gen-gap $>0.02$ conflates onset with asymptote and is contaminated by the single-draw noise floor; a relative threshold ($\delta=0.2$ of the many-draw asymptotic gap) on $\mathcal M_{\rm blk}$ will show a much steeper growth with $\dlat$ than the logged $4\times$.} Test: re-extract with $\ge50$-draw train and test losses at each evaluation and compare growth factors with $3200/\lambda^{\rm blk}_{\max,\rm code}$.
\item \emph{Operating point:} $a_*(q)^2$ falls $0.0785\to0.0047$ over $\dlat=5..200$ at $\signoise=0.5$ and equals $2.6\times10^{-4}$ at $\dlat=200$ for $\signoise=0.01$, so every memorisation timescale in flow units is pushed out by more than an order of magnitude at $\signoise=0.01$ relative to $0.5$ at equal $\dlat$. Test: compare $\operatorname{tr}\Umat/\pwidth$ minus the data-block mass across the two $\signoise$ sweeps with $a_*(q(\dlat))^2$ (expect $0.7$--$0.9\times a_*^2/\nsamp$ for the sample-block mean).
\item \emph{Alignment coefficients $\gamma_i$ on held-out data will be $1+o(1)$ on $\mathcal S\cup\mathcal N$, exceed $\tfrac12$ on a set $\mathcal M_{\rm spk}$ of $O(k)$ top sample modes at small $\dlat$ whose share grows as $\psi_n\to1$, and be $O(1/\sqrt{\nsamp})$ on $\mathcal M_{\rm blk}$; the test-loss argmin at $35/55/90$k steps ($\dlat=40/100/200$) will be reproduced by $3200/\lambda_{\rm code}$ of the slowest mode with $\gamma_i>\tfrac12$.} Test: compute $\bar\Umat,\bar V$ on $10^4$ held-out points, histogram $\gamma_i$ by block at $\dlat\in\{5,20,60,200\}$; sweep $k$ at fixed $\nsamp,\dlat$ to check $|\mathcal M_{\rm spk}|\sim k$.
\end{enumerate}

\paragraph{Open gaps.}
(1) The split $\mathcal M_{\rm spk}/\mathcal M_{\rm blk}$ is a definition plus Prediction~5: $\gamma_i$ has not been computed on the repository's data, so the quoted ratio numbers ($34\to426$, monotone) are index-based lower bounds on the bulk onset, and whether monotone growth of $\taumem^{\rm onset}/\taugen$ survives the $\gamma$-cut is open. The tail bound in \cref{prop:gap-p1} is stated through $\|VP_{\mathcal Z}\|_F^2$ and $\sup_{\mathcal Z}\bar u_i$, which have not been evaluated against the threshold $0.01\dlat/\Delta_t$. (2) The rank-$\nsamp$ cliff is a single-draw fact; the released $50$-draw spectra are full rank with a soft step $\le1.7$ at index $\nsamp$, and the main text must say which operator each figure shows. (3) (A6$'$) is a hypothesis; the required relative accuracy on the population blocks is $\sim10\%$ at $\dlat=20$ and tighter at larger $\dlat$ ($\Pi_{\mathcal N}\sim\dlat^2/(a_1^2\beta_t^4)$), and it may fail as $\psi_n\to1$, consistent with the test loss still falling at $90$k steps at $\dlat=200$. (4) The $(1+o(1))$ in \eqref{eq:lammin-signal-p1} is not quantified here; the finite-rank BBP correction is \cref{lem:fourbulk-bbp}'s. (5) Monotonicity of $a_*^2$ in $q$ is numerical on $[0.01,20]$ only. (6) The $\dlat=5$ point violates $\pwidth\ge\nsamp$; all anchors there are flagged and the monotone ratio is quoted from $\dlat=10$. (7) The regime $\nsamp\le\dlat$ (memorisation returning at large $\dlat$, \cref{app:theory-scope}) is outside (A3): there $\mathcal M$ is empty and $G_{\rm blk}\equiv0$; the growth of $\mathcal M_{\rm spk}$ as $\psi_n\to1$ is the natural bridge but is not derived. (8) Real VAE latents have continuous spectra with posterior-collapsed coordinates at exactly $\Delta_t$; the index-based partition of (A5) needs the data-driven cut of \cref{app:theory-general}. (9) The clock $T=0.02\dlat\times$steps holds for full-batch SGD with the stated learning-rate rule and momentum $0$; the MLP sweeps use Adam and no clock statement is made for them. (10) The $O(\theta_{\mathcal S}/\lambda^{\mathcal S}_{\min})$ term in \cref{prop:taugen-p1} scales like $\taugen$ itself unless $\theta_{\mathcal S}=\log(\Pi_{\mathcal S}/\pi_{i_\circ})$ is small; for $10$ random centres in $\dint=5$ it is not, so $\taugen(\varepsilon)$ is pinned to $\lambda^{\mathcal S}_{\min}$ only up to the factor $1-\theta_{\mathcal S}/\log(1/\varepsilon)$. (11) The sample-block mean is $0.72$--$0.89\times a_*^2/\nsamp$ (Jensen dispersion of per-sample $q_\mu$), so $a_1,a_*$ should be defined with $\mu$-averaged Gaussian expectations or the trace check is $10$--$30\%$ off in the sample block; the noise-dimension median also overshoots by $10$--$15\%$ at large $\dlat$. (12) The $\gamma_0$-cut is not canonical; the bounds in \cref{def:taumem-p1} carry an explicit $\gamma_0\sum\lambda_i\pi_i$ error that has not been evaluated.

% ----------------------------------------------------------------------
\subsection{Four spectral blocks of the feature correlation matrix}
\label{app:theory-fourbulk}

Throughout this subsection $\phi(x)=\pwidth^{-1/2}\tanh(Wx/\sqrt{\dlat})$ with $W_{ij}\sim\mathcal N(0,1)$ i.i.d.; the released code uses $\tanh(\tilde Wx)$ with $\tilde W=W/\sqrt{\dlat}$ and no $\pwidth^{-1/2}$, so $\Umat_{\mathrm{code}}=\pwidth\,\Umat$ and all eigenvalues quoted from the repository are $\pwidth$ times the theory values; timescales in theory units are $\pwidth/\lambda_{\mathrm{code}}$. $\Umat=\tfrac1\nsamp\sum_\mu\E_\xi[\phi(x_t^\mu)\phi(x_t^\mu)^\top]$ is the $\xi$-averaged matrix ($K=50$ draws in \texttt{compute\_U}); $\Umat^{(1)}=\Phi^\top\Phi/\nsamp$ denotes the single-copy matrix with $x_t^\mu=e^{-t}x^\mu+\sqrt{\Delta_t}\,\xi^\mu$, $\Phi\in\R^{\nsamp\times\pwidth}$ the feature matrix of one diffused copy per point, so that $\Umat^{(1)}$ shares its nonzero spectrum with the Gram matrix $G=\Phi\Phi^\top/\nsamp$ and $\operatorname{rank}\Umat^{(1)}\le\nsamp$. Write $q_\mu=\|x_t^\mu\|^2/\dlat$ for the pre-activation variance of point $\mu$ and $q=\tfrac1\nsamp\sum_\mu q_\mu\simeq\operatorname{tr}M_t/\dlat$, which for the synthetic family is explicit (\cref{lem:q-p1}),
\begin{equation}\label{eq:fourbulk-q}
  q(\dlat)=e^{-2t}\Bigl[\tfrac{s^2+\dint\sigsig^2}{\dlat}+\signoise^2\bigl(1-\tfrac{\dint}{\dlat}\bigr)\Bigr]+\Delta_t
  \;\xrightarrow{\;\dlat\to\infty\;}\; e^{-2t}\signoise^2+\Delta_t .
\end{equation}
Let $\sigma_r=a_*(q)^2/\nsamp$ and $\kappa_{\mathrm{MP}}=\pwidth^{-1}(\sqrt{\nsamp-\dlat}+\sqrt{\pwidth-\dlat})^2$. Split $M_t=e^{-2t}(\hat C+\hat S)+\Delta_t I$ into its $\dint\times\dint$ signal block $e^{-2t}(\hat C+\hat S_\parallel)+\Delta_t I$ and its null block $M_\perp=e^{-2t}\hat S_\perp+\Delta_t I$, where $\hat S_\perp$ is $\signoise^2$ times a white Wishart matrix of aspect ratio $\gamma_\perp=(\dlat-\dint)/\nsamp$; set $\gamma_p=\nsamp/\pwidth$.

\subsubsection{Step 1: the feature kernel is exactly a Hermite series}

\begin{lemma}[Kernel of the $\tanh$ features; proved]\label{lem:fourbulk-kernel}
Fix the diffused training points $x_t^1,\dots,x_t^\nsamp$ and let $q_\mu=\|x_t^\mu\|^2/\dlat$, $\rho_{\mu\nu}=\langle x_t^\mu,x_t^\nu\rangle/(\dlat\sqrt{q_\mu q_\nu})$. For $w\sim\mathcal N(0,I_{\dlat})$ the pre-activations $g_\mu=\langle w,x_t^\mu\rangle/\sqrt{\dlat}$ are jointly Gaussian with $\operatorname{Var}g_\mu=q_\mu$ and $\operatorname{Corr}(g_\mu,g_\nu)=\rho_{\mu\nu}$, and
\begin{equation}\label{eq:fourbulk-hermite}
  K_{\mu\nu}:=\E_w\bigl[\tanh(g_\mu)\tanh(g_\nu)\bigr]
  =\sum_{k\ \mathrm{odd}} c_k(q_\mu)c_k(q_\nu)\,\rho_{\mu\nu}^{\,k},\qquad
  c_k(q)=\E\bigl[\tanh(\sqrt q z)\,\mathrm{He}_k(z)\bigr]/\sqrt{k!},
\end{equation}
with only odd $k$ because $\tanh$ is odd. By Stein's identity $c_1(q)=\sqrt q\,\E[\tanh'(\sqrt q z)]=\sqrt q\,a_1(q)$, so the $k=1$ term equals $a_1(q_\mu)a_1(q_\nu)\langle x_t^\mu,x_t^\nu\rangle/\dlat$, and
\begin{equation}\label{eq:fourbulk-Ksplit}
  K = A_1\,\frac{X_tX_t^\top}{\dlat}\,A_1 + K_{\mathrm{res}},\qquad A_1=\diag\bigl(a_1(q_\mu)\bigr),
\end{equation}
where $(K_{\mathrm{res}})_{\mu\mu}=a_*(q_\mu)^2$ and $(K_{\mathrm{res}})_{\mu\nu}=\sum_{k\ge3\ \mathrm{odd}}c_k(q_\mu)c_k(q_\nu)\rho_{\mu\nu}^{\,k}=O(\rho_{\mu\nu}^3)$ for $\mu\ne\nu$.
\end{lemma}
\begin{proof}
Joint Gaussianity is immediate since $g$ is linear in $w$. Expanding $\tanh(\sqrt{q}\,\cdot)$ in normalized Hermite polynomials and using Mehler's identity $\E[\mathrm{He}_k(g_\mu/\sqrt{q_\mu})\mathrm{He}_l(g_\nu/\sqrt{q_\nu})]=\delta_{kl}\,k!\,\rho_{\mu\nu}^k$ gives \eqref{eq:fourbulk-hermite}; oddness kills even $k$ (in particular $c_0=0$: the features are centred, so no rank-one $\mathbf 1\mathbf 1^\top$ term appears, unlike for ReLU). Stein's identity $\E[f(\sqrt q z)z]=\sqrt q\,\E[f'(\sqrt q z)]$ gives $c_1$. The diagonal of $K_{\mathrm{res}}$ is $\E\tanh(\sqrt{q_\mu}z)^2-a_1(q_\mu)^2q_\mu=a_*(q_\mu)^2$ by definition.
\end{proof}

For the synthetic family, $\rho_{\mu\nu}=O(\dlat^{-1/2})+O\bigl(s^2\langle m_{c(\mu)},m_{c(\nu)}\rangle/(\dlat q)\bigr)$ for points in different clusters (the generator uses $\dint$ orthogonal centres plus $k-\dint$ random unit centres, so cross-cluster centre overlaps are $O(s^2)$, not zero) and $\rho_{\mu\nu}\simeq s^2/(\dlat\sqrt{q_\mu q_\nu})$ for the same cluster ($0.65$ at $\dlat=5$, $0.5$ at $\dlat=20$). The off-diagonal part of $K_{\mathrm{res}}$ is $O(\rho^3)$: $K_{\mathrm{res}}$ is diagonal-dominant with a same-cluster component of rank $\lesssim k$ and top eigenvalue $\simeq(\nsamp/k)\,c_3(q)^2\rho_{\mathrm{same}}^3$ ($c_3^2=0.052,0.033,0.019,0.011,0.006,0.004$ at $\dlat=5,10,20,40,100,200$; $\approx10\,a_*^2$ at $\dlat=5$, $\approx5\,a_*^2$ at $\dlat=20$). This is the precise form of the failure of the orthonormality ansatz $v_\mu^\top v_\nu=\delta_{\mu\nu}+o(1)$ for clustered data. The mixture is not Gaussian, but \cref{lem:fourbulk-kernel} needs only Gaussianity \emph{over $W$}, which is exact.

\subsubsection{Step 2: Gaussian equivalence and the spiked-covariance form}

\begin{assumption}[Gaussian equivalence of the feature matrix; (a) imported, (b)--(c) conjectured]\label{ass:fourbulk-ge}
Let $G=\Phi\Phi^\top/\nsamp$. (a) [Bulk] In the limit $\pwidth,\nsamp,\dlat\to\infty$ with $\gamma_p=\nsamp/\pwidth$, $\gamma_\perp=(\dlat-\dint)/\nsamp$ fixed and $\max_\mu q_\mu$ bounded, the empirical spectral distribution of $G$ coincides with that of
\begin{equation}\label{eq:fourbulk-GE}
  G^{\mathrm{GE}}=\frac{1}{\nsamp}\,K^{1/2}\,\frac{ZZ^\top}{\pwidth}\,K^{1/2},\qquad Z\in\R^{\nsamp\times\pwidth}\ \text{i.i.d.}\ \mathcal N(0,1),
\end{equation}
a sample covariance with population $K/\nsamp$ and $\pwidth$ samples. (b) [Outliers] The $\dint$ largest eigenvalues (a finite-rank set) and the $\dlat-\dint$ noise-dimension eigenvalues of $G$ are located by the same replacement to relative accuracy $o(1)$. (c) [$\xi$-average] The same holds for $\Umat$ with the Gram of the $K\nsamp$ noise draws in place of $K$.
\end{assumption}
Part (a) is a theorem: the deterministic equivalent of \citet{louart2018random} for the resolvent of $\Phi^\top\Phi/\nsamp$ with Gaussian $W$, Lipschitz activation and arbitrary bounded-norm data is exactly the Silverstein equivalent of a sample covariance with population $K=\E_w[\sigma(w^\top X)^\top\sigma(w^\top X)]$; \citet{liao2018spectrum} treat Gaussian-mixture inputs with $O(1)$ same-cluster overlaps through the same Hermite coefficients, and \citet{fanwang2020spectra} give the Marchenko--Pastur--Silverstein form for the conjugate kernel; the replica counterpart is Theorem~3.1 of \citet{bonnaire2025} (stated for an arbitrary input covariance), and the Gaussian-equivalence principle for the learning problem is \citet{goldt2022gaussian,hu2022universality,mei2022generalization}. Parts (b) and (c) are open: the signal directions are a finite-rank spike set at fixed $\dint$ (the isolated-eigenvalue results of \citet{liao2018spectrum} require growth conditions on the class means that the fixed $s=3$, $\dint=5$ regime does not meet), and the $K\nsamp$ rows of the averaged feature matrix are strongly dependent (same-point draws have $\rho\simeq1-\Delta_t/q$). This is why \cref{prop:fourbulk} is labelled sketched.

\paragraph{Identity in distribution and the correct perturbation model.}
In the homogeneous case $a_1(q_\mu)\equiv a_1$, let $L=[a_1X_t/\sqrt{\dlat},\;K_{\mathrm{res}}^{1/2}]\in\R^{\nsamp\times(\dlat+\nsamp)}$, so $LL^\top=K$. For $\tilde Z\in\R^{(\dlat+\nsamp)\times\pwidth}$ Gaussian, each column of $L\tilde Z$ is $\mathcal N(0,K)$, hence $L\tilde Z/\sqrt{\pwidth}\overset d=K^{1/2}Z/\sqrt{\pwidth}$. Writing $\tilde Z=[W^\top;\Theta]$, the $\pwidth\times\pwidth$ picture
\begin{equation}\label{eq:fourbulk-Upicture}
  \Umat\simeq\frac{a_1^2}{\pwidth\,\dlat}\,W\hat M_tW^\top+\frac{1}{\pwidth\,\nsamp}\,\Theta^\top K_{\mathrm{res}}\Theta+\text{cross},\qquad \hat M_t=X_t^\top X_t/\nsamp,
\end{equation}
is therefore the same object as \eqref{eq:fourbulk-GE}, and the data directions are spikes of the \emph{population} covariance $K/\nsamp=D+R$ of a Wishart-type matrix (multiplicative model), not additive rank-$\dlat$ perturbations of a Wishart matrix; the cross terms in \eqref{eq:fourbulk-Upicture} are what the multiplicative model resolves. Two consequences that an additive model gets wrong: the data part $a_1^2W\hat M_tW^\top/(\pwidth\dlat)$ has the nonzero spectrum of $\hat M_t^{1/2}(W^\top W/\pwidth)\hat M_t^{1/2}\,a_1^2/\dlat$, so every data block is multiplied by a white Wishart of ratio $\dlat/\pwidth$ (Step~3); and the outlier shift and residual mass follow the spiked-covariance rules of \cref{lem:fourbulk-bbp} (Step~4).

\subsubsection{Step 3: spectrum of $K/\nsamp$, the data blocks and the four counts}

\begin{lemma}[Counts by interlacing; proved]\label{lem:fourbulk-counts}
Let $D=A_1X_tX_t^\top A_1/(\dlat\nsamp)$ (PSD, $\operatorname{rank}D\le\dlat$, $=\dlat$ a.s.\ for $\nsamp>\dlat$) and $R=K_{\mathrm{res}}/\nsamp$. Then at least $\nsamp-\dlat$ eigenvalues of $K/\nsamp=D+R$ lie in $[\lambda_{\min}(R),\lambda_{\max}(R)]$ and at most $\dlat$ lie above $\lambda_{\max}(R)$; exactly $\nsamp-\dlat$ lie inside iff every data direction is an outlier, which holds under the margin condition of \cref{lem:fourbulk-bbp} (measured margin $\ge20\times$ at $\signoise=0.5$; violated at $\signoise=0.01$, \cref{rem:fourbulk-sigperp0}). For the single-copy matrix $\Umat^{(1)}=\Phi^\top\Phi/\nsamp$ with $\pwidth\ge\nsamp$ there are exactly $\pwidth-\nsamp$ zero eigenvalues; for the $K$-draw average $\operatorname{rank}\Umat\le\min(\pwidth,K\nsamp)$. Hence the sample block has $\nsamp-\dlat$ modes and the last non-residual index is $\nsamp$, not $\dlat+\nsamp$.
\end{lemma}
\begin{proof}
Cauchy interlacing for a PSD perturbation of rank $r\le\dlat$: $\lambda_{i}(R)\le\lambda_{i}(D+R)\le\lambda_{i-r}(R)$, so $\lambda_i(D+R)\in[\lambda_{\min}(R),\lambda_{\max}(R)]$ for all $i>r$, while the top $r$ are $\ge\lambda_r(R)$ and may lie inside or above $\lambda_{\max}(R)$, which is the iff. The zeros follow from $\operatorname{rank}\Umat^{(1)}=\operatorname{rank}\Phi\le\nsamp$; the averaged matrix is $\Phi_{\mathrm{all}}^\top\Phi_{\mathrm{all}}/(K\nsamp)$ with $\Phi_{\mathrm{all}}\in\R^{K\nsamp\times\pwidth}$. The $\nsamp$ non-trivial directions of $G$ are shared between the $\dlat$ data directions and the residual, so there is no room for $\dlat+\nsamp$ non-zero modes; the Marchenko--Pastur--Silverstein transform preserves this bookkeeping.
\end{proof}

\paragraph{Data directions and the $W$ factor.}
The nonzero eigenvalues of the data term are those of $\hat M_t^{1/2}(W^\top W/\pwidth)\hat M_t^{1/2}\,a_1^2/\dlat$ up to the spread of $a_1(q_\mu)$ across points. After $\xi$-averaging $\hat M_t\to M_t=e^{-2t}(\hat C+\hat S)+\Delta_t I$, which is block diagonal in the (rotated) signal/null split up to $O(\nsamp^{-1/2})$ empirical cross-covariances: the signal block $e^{-2t}(\hat C+\hat S_\parallel)+\Delta_t I_{\dint}$ has $\dint$ eigenvalues with mean $\alpha_t^2$, and the null block $M_\perp=e^{-2t}\hat S_\perp+\Delta_t I$ has $\dlat-\dint$ eigenvalues, $\hat S_\perp$ being $\signoise^2$ times a white Wishart with $\nsamp$ samples in $\dlat-\dint$ dimensions (Marchenko--Pastur, ratio $\gamma_\perp$, support $\signoise^2(1\pm\sqrt{\gamma_\perp})^2$). The factor $W_\perp^\top W_\perp/\pwidth$ is an independent white Wishart of ratio $(\dlat-\dint)/\pwidth$, asymptotically free of $M_\perp$, so block (ii) is $[e^{-2t}\signoise^2\mathrm{MP}(\gamma_\perp)+\Delta_t]\boxtimes\mathrm{MP}(\dlat/\pwidth)$ scaled by $a_1^2/\dlat$; the trace of a product of free elements is the product of traces, so the mean is $a_1^2\beta_t^2/\dlat$ exactly, and the support is $\tfrac{a_1^2}{\dlat}[e^{-2t}\signoise^2(1\pm\sqrt{\gamma_\perp})^2+\Delta_t]$ at leading order in $\dlat/\pwidth$ with relative edge corrections $O(\sqrt{\dlat/\pwidth})$. Direct check (repository generator, $\dlat=40$, single copy): noise-dim min/median/max $1.05/2.07/4.07$, $4.63/8.11/13.47$, $18.96/33.5/54.9$ at $\pwidth=16\dlat,64\dlat,256\dlat$ versus the leading-order formula $1.20/2.05/3.25$, $4.81/8.19/13.00$, $19.25/32.75/52.0$: the median is $\pwidth$-independent to $1\%$, the edges move by $-13\%/+25\%$ at $\dlat/\pwidth=1/16$ and by $\pm4\%$ at $1/64$. For the finite-rank signal spikes the $W$ factor contributes only $O(\pwidth^{-1/2})$ (signal-min$/\pwidth$: $0.0161,0.0156,0.0158$ across the three widths). The signal block at fixed $\dint=5$, $k=10$ is a finite set of eigenvalues of the realized centre geometry, not a self-averaging bulk: $a_1^2(\pwidth/\dlat)\lambda_j(M_t)$ computed from the seed-42 training set reproduces all five signal eigenvalues to $2$--$5\%$ at every $\dlat$ (e.g.\ $\dlat=200$: measured $167.8,135.4,99.8,87.1,73.9$ vs $168.8,141.4,99.4,88.0,75.1$; $\dlat=5$: $46.6,41.9,29.7,22.0,20.3$ vs $47.2,39.5,27.3,24.3,20.2$), whereas the mean-only value $a_1^2\alpha_t^2\psi_p$ overshoots the median by $10$--$20\%$. For the single-copy matrix $\Delta_t$ enters the Wishart variance instead and the leading-order support is $\tfrac{a_1^2}{\dlat}\beta_t^2(1\pm\sqrt{\gamma_\perp})^2$ (\cref{rem:fourbulk-singlecopy}).

\paragraph{Residual directions.}
$R=K_{\mathrm{res}}/\nsamp$ has $\operatorname{tr}R=\tfrac1\nsamp\sum_\mu a_*(q_\mu)^2$, which equals $a_*(q)^2$ to $<1\%$ for $\dlat\ge20$ ($0.0733$ vs $0.0785$ at $\dlat=5$). If the residual is homogeneous ($R=\sigma_rI$) the sample block is the plain Marchenko--Pastur law of \cref{lem:fourbulk-bbp}(b), with top edge $\kappa_{\mathrm{MP}}\sigma_r$. In general the block is the MP--Silverstein image of $\operatorname{spec}(P_\perp RP_\perp)$ with $\pwidth-\dlat$ effective samples, whose width is set by the spread of $a_*(q_\mu)^2$ (range $0.0036$--$0.153$ at $\dlat=5$, $0.0017$--$0.035$ at $\dlat=40$, $0.0024$--$0.0089$ at $\dlat=200$, because $q_\mu$ inherits the cross term $2\langle m_{c(\mu)},\xi^\mu\rangle/\dlat$ of the mixture) and by the same-cluster component of Step~1; both shrink with $\dlat$. We do not have a closed form for the resulting top edge $\kappa_{\mathrm{top}}\sigma_r$; it is characterised by this transform and measured in \cref{tab:fourbulk-numbers}, where it decreases from $19.0$ to $1.45$ and reaches $\kappa_{\mathrm{MP}}=1.31$ at $\dlat=200$ once the $\approx11\%$ inflation caused by the $K=50$ average is removed. On the $\pwidth=\dlat+\nsamp+300$ spectra ($K=500$) the measured values $14.5,12.1,8.0,3.10,2.00$ at $\dlat=10,\dots,200$ compare with $\kappa_{\mathrm{MP}}=3.14,3.07,2.94,2.59,2.08$; the additive-model value $(1+\sqrt{\gamma_p})^2=2.91$ at $\dlat=200$ would wrongly put the measured edge below the homogeneous one.

\begin{proposition}[Four spectral blocks of $\Umat$; sketched]\label{prop:fourbulk}
Assume $\pwidth\ge\nsamp>\dlat>\dint$, $\dint$ and $t>0$ fixed, and take $\dlat,\nsamp,\pwidth\to\infty$ with $\gamma_\perp,\gamma_p$ fixed, under \cref{ass:fourbulk-ge}. Let $\varepsilon=\bar k\,\Delta_t/q$ with $\bar k=\sum_{k\ge3}k\,c_k(q)^2/\sum_{k\ge3}c_k(q)^2\in[3.2,3.9]$. The ordered spectrum of $\Umat$ consists of four index blocks:
\begin{itemize}
\item[(i)] $\dint$ \emph{signal} outliers, $\lambda^{\mathrm{sig}}_j=\tfrac{a_1(q)^2}{\dlat}\bigl[e^{-2t}\lambda_j(\hat C+\hat S_\parallel)+\Delta_t\bigr]\bigl(1+O(\pwidth^{-1/2})\bigr)+\sigma_r\bigl(1+O(\nsamp/\pwidth)\bigr)$, population mean $a_1(q)^2\alpha_t^2/\dlat$;
\item[(ii)] $\dlat-\dint$ \emph{noise-dimension} eigenvalues, those of $\tfrac{a_1(q)^2}{\dlat}M_\perp^{1/2}(W_\perp^\top W_\perp/\pwidth)M_\perp^{1/2}$, i.e.\ $\bigl[e^{-2t}\signoise^2\mathrm{MP}(\gamma_\perp)+\Delta_t\bigr]\boxtimes\mathrm{MP}(\dlat/\pwidth)$ scaled by $a_1^2/\dlat$: mean $a_1(q)^2\beta_t^2/\dlat$ exactly, support $\tfrac{a_1(q)^2}{\dlat}\bigl[e^{-2t}\signoise^2(1\pm\sqrt{\gamma_\perp})^2+\Delta_t\bigr]$ at leading order in $\dlat/\pwidth$, edges corrected by $O(\sqrt{\dlat/\pwidth})$, median $\pwidth$-independent;
\item[(iii)] $\nsamp-\dlat$ \emph{sample} eigenvalues: the Marchenko--Pastur--Silverstein image, with $\pwidth-\dlat$ effective samples, of the residual kernel $P_\perp K_{\mathrm{res}}P_\perp/\nsamp$ (diagonal $a_*(q_\mu)^2$, $P_\perp$ off the data directions); mean $\sigma_r(1-\dlat/\pwidth)(1-\varepsilon)$, top edge $\kappa_{\mathrm{top}}\sigma_r$, where $\kappa_{\mathrm{top}}=\kappa_{\mathrm{MP}}$ for a homogeneous residual and a single copy, and in general is set by the spread of $a_*(q_\mu)^2$ and the same-cluster part of $K_{\mathrm{res}}$;
\item[(iv)] $\pwidth-\nsamp$ \emph{residual-floor} eigenvalues of total mass $\simeq\varepsilon a_*(q)^2$; for a single diffused copy these are exact zeros (\cref{rem:fourbulk-singlecopy}).
\end{itemize}
The blocks satisfy the exact trace identity
\begin{align}\label{eq:fourbulk-trace}
  \underbrace{\dint\tfrac{a_1^2\alpha_t^2}{\dlat}+(\dlat-\dint)\tfrac{a_1^2\beta_t^2}{\dlat}}_{=\,a_1(q)^2q}
  +\underbrace{\dlat\,\sigma_r\bigl(1+\tfrac{\nsamp-\dlat}{\pwidth}\bigr)}_{\text{outlier shifts}}+\underbrace{(\nsamp-\dlat)\,\sigma_r\bigl(1-\tfrac{\dlat}{\pwidth}\bigr)}_{\text{sample bulk}}
  =a_1(q)^2q+a_*(q)^2=\E\bigl[\tanh(\sqrt q z)^2\bigr],
\end{align}
while $\operatorname{tr}\Umat=\tfrac1\nsamp\sum_\mu\E\tanh(\sqrt{q_\mu}z)^2$ exactly, which differs from $\E\tanh(\sqrt qz)^2$ by a Jensen correction of $2$--$5\%$ over the sweep. All data-side eigenvalues scale as $a_1(q(\dlat))^2/\dlat$ and are $\pwidth$-independent at leading order in $\dlat/\pwidth$ (the code normalization $\Umat_{\mathrm{code}}=\pwidth\,\Umat$ reintroduces $\psi_p=\pwidth/\dlat$); the sample block scales as $a_*(q(\dlat))^2/\nsamp$.
\end{proposition}
\begin{proof}[Proof sketch]
Step 1 (\cref{lem:fourbulk-kernel}): $K=\nsamp D+K_{\mathrm{res}}$ exactly. Step 2 (\cref{ass:fourbulk-ge}): $G$ is a sample covariance with population $K/\nsamp$ and $\pwidth$ samples, so its bulk is the MP--Silverstein transform of $\operatorname{spec}(K/\nsamp)$ at ratio $\gamma_p$. Step 3: the data part of $\Umat$ is $a_1^2W\hat M_tW^\top/(\pwidth\dlat)$ up to the spread of $a_1(q_\mu)$, whose nonzero spectrum is that of $\hat M_t^{1/2}(W^\top W/\pwidth)\hat M_t^{1/2}a_1^2/\dlat$; after $\xi$-averaging $\hat M_t\to M_t$, block diagonal in the signal/null split up to $O(\nsamp^{-1/2})$ cross terms; the null block is $e^{-2t}\signoise^2\times$ a white Wishart of ratio $\gamma_\perp$ plus $\Delta_t I$, and $W_\perp^\top W_\perp/\pwidth$ is an independent white Wishart of ratio $(\dlat-\dint)/\pwidth$, giving (ii) by asymptotic freeness. The signal block is a finite set of $\dint$ eigenvalues of the realized centre geometry (i). Counts by \cref{lem:fourbulk-counts}. Step 4: \cref{lem:fourbulk-bbp} (spiked population covariance) gives the outlier shift and the sample-bulk mass $(\nsamp-\dlat)\sigma_r(1-\dlat/\pwidth)$, and the trace identity closes exactly at the level of expectations. Step 5: the $\xi$-average replaces the residual of each point by a coherent part $a_*^2(1-\varepsilon)$ plus $K-1$ incoherent modes of weight $a_*^2\varepsilon/K$ each (\cref{rem:fourbulk-singlecopy} and the paragraph following it), giving the factor $(1-\varepsilon)$ in (iii) and the floor (iv). The heterogeneous residual is characterised, not solved: its MP--Silverstein image has width set by the spread of $a_*(q_\mu)^2$ (a factor $40$ at $\dlat=5$, $4$ at $\dlat=200$) and by the same-cluster component of size $\simeq(\nsamp/k)c_3(q)^2\rho_{\mathrm{same}}^3$, both of which shrink with $\dlat$; the measured $\kappa_{\mathrm{top}}$ homogenises to $\kappa_{\mathrm{MP}}$ within the $\sim10\%$ inflation caused by the $K=50$ average.
\end{proof}

\subsubsection{Step 4: spiked-covariance outliers and the trace check}

\begin{lemma}[Outliers of the spiked population covariance; proved]\label{lem:fourbulk-bbp}
Let $G=\Sigma^{1/2}(ZZ^\top/\pwidth)\Sigma^{1/2}$, $Z\in\R^{\nsamp\times\pwidth}$ i.i.d.\ Gaussian, $\Sigma=D+\sigma_rI$ with $D\succeq0$ of rank $\dlat$ and eigenvalues $\theta_1\ge\dots\ge\theta_{\dlat}$ (homogeneous residual; for the heterogeneous residual replace the MP law by the MP--Silverstein image of $\operatorname{spec}(P_\perp K_{\mathrm{res}}P_\perp)/\nsamp$ and $\sigma_r$ in (b) by the mean residual eigenvalue). (a) For fixed $\dlat$ and $\gamma_p=\nsamp/\pwidth<1$, the eigenvalue of $G$ associated with $\theta_j$ is
\begin{equation}\label{eq:fourbulk-bbp}
  \lambda(\theta_j)=(\theta_j+\sigma_r)\Bigl(1+\frac{\gamma_p\sigma_r}{\theta_j}\Bigr)=\theta_j+\sigma_r(1+\gamma_p)+\frac{\gamma_p\sigma_r^2}{\theta_j}
  \qquad\text{if}\qquad\theta_j>\theta_c=\sigma_r\sqrt{\gamma_p},
\end{equation}
and sticks to the edge $\sigma_r(1+\sqrt{\gamma_p})^2$ otherwise. (b) For any $\dlat<\min(\nsamp,\pwidth)$, conditionally on the $\dlat$ spiked rows $Z_1$ of $Z$, the restriction of $G$ to the complement of the spiked directions is $(\sigma_r/\pwidth)Z_2P^\perp Z_2^\top$, $P^\perp$ the projector off the row space of $Z_1$: a white Wishart with $\pwidth-\dlat$ samples in $\nsamp-\dlat$ dimensions, i.e.\ Marchenko--Pastur of ratio $(\nsamp-\dlat)/(\pwidth-\dlat)$, scale $\sigma_r(1-\dlat/\pwidth)$, expected mass $(\nsamp-\dlat)\sigma_r(1-\dlat/\pwidth)$ and top edge $\sigma_r\kappa_{\mathrm{MP}}$, with $\kappa_{\mathrm{MP}}=\pwidth^{-1}(\sqrt{\nsamp-\dlat}+\sqrt{\pwidth-\dlat})^2$ (the same formula holds for $\pwidth<\nsamp$, where the block has $\pwidth-\dlat$ modes). Since $\E\operatorname{tr}G=\operatorname{tr}\Sigma$, the $\dlat$ outliers carry expected mass $\sum_j\theta_j+\dlat\sigma_r(1+\tfrac{\nsamp-\dlat}{\pwidth})$: a mean shift of $\sigma_r(1+\tfrac{\nsamp-\dlat}{\pwidth})$ per outlier, consistent with \eqref{eq:fourbulk-bbp} at leading order in $\dlat/\nsamp$.
\end{lemma}
\begin{proof}
(a) is the spiked sample-covariance theorem \citep{baik2005phase,baik2006eigenvalues}, equivalently the multiplicative case of \citet{benaychgeorges2011eigenvalues}: a population eigenvalue $\ell=\theta+\sigma_r$ above $\sigma_r(1+\sqrt{\gamma_p})$ produces the outlier $\ell(1+\gamma_p\sigma_r/(\ell-\sigma_r))$. (b) In the spike eigenbasis $\Sigma^{1/2}Z=[(\Theta+\sigma_rI)^{1/2}Z_1;\ \sqrt{\sigma_r}Z_2]$; by interlacing for the $\pwidth\times\pwidth$ dual $Z^\top\Sigma Z/\pwidth$, the non-spiked eigenvalues of $G$ are those of $\sigma_rZ_2P^\perp Z_2^\top/\pwidth$ up to $\dlat$ interlacing corrections, and $Z_2P^\perp$ has, given $Z_1$, i.i.d.\ Gaussian rows supported on a $(\pwidth-\dlat)$-dimensional subspace, which gives the Wishart statement, its mass $\sigma_r(\nsamp-\dlat)(\pwidth-\dlat)/\pwidth$ and its edge. Subtracting from $\operatorname{tr}\Sigma=\sum_j\theta_j+\nsamp\sigma_r$ gives the outlier mass.
\end{proof}
Numerical check of (b) (homogeneous, $\nsamp=500$, $\dlat=40$, $\pwidth=2560$, $5$ repetitions): bulk mass$/(\sigma_r(\nsamp-\dlat))=0.9840$ vs $1-\dlat/\pwidth=0.9844$; top edge$/\sigma_r=1.987$ vs $\kappa_{\mathrm{MP}}=2.005$, excluding the additive-model value $(1+\sqrt{\gamma_p})^2=2.079$. Individual outliers fluctuate by $\theta_j\sqrt{2/\pwidth}$, so the per-outlier shift ($\sigma_r$ vs $1.18\sigma_r$) is not resolvable in a single run; (b) is a statement in expectation. In code units the threshold is $\theta_c\pwidth\in[0.03,0.1]$ across the sweep while the smallest noise-dimension eigenvalue is $2$--$5$ at $\signoise=0.5$: every data direction is an outlier by a margin $\ge20\times$, and the shift $\sigma_r$ is $0.15$--$0.25\%$ of $\lambda^{\mathrm{sig}}_{\min}$ ($a_*^2\pwidth/\nsamp=0.05$--$0.12$ vs $20$--$74$), which is the operational content of ``the data-side eigenvalues are $\pwidth$-independent'' beyond the $W$ factor of Step~3.

\paragraph{Trace arithmetic.}
Exactly, $\operatorname{tr}\Umat=\tfrac1\nsamp\sum_\mu\E\|\phi(x_t^\mu)\|^2=\tfrac1\nsamp\sum_\mu\E\tanh(\sqrt{q_\mu}z)^2=\tfrac1\nsamp\sum_\mu[a_1(q_\mu)^2q_\mu+a_*(q_\mu)^2]$. On the released spectra this heterogeneous prediction matches $\operatorname{tr}\Umat_{\mathrm{code}}/\pwidth$ to $1.5\%$ at every $\dlat$ ($0.548,0.449,0.361,0.292,0.233,0.209$ predicted vs $0.545,0.456,0.363,0.292,0.233,0.209$ measured, $\dlat=5,\dots,200$), while the homogeneous value $\E\tanh(\sqrt qz)^2=a_1(q)^2q+a_*(q)^2$ is $2$--$5\%$ high at small $\dlat$ because $q\mapsto\E\tanh(\sqrt qz)^2$ is concave (Jensen; $\operatorname{sd}_\mu q_\mu=1.37$ at $\dlat=5$, $0.04$ at $\dlat=200$). On the block side, with $\operatorname{tr}M_t=\dint\alpha_t^2+(\dlat-\dint)\beta_t^2=\dlat q$ and \cref{lem:fourbulk-bbp}(b), the identity \eqref{eq:fourbulk-trace} holds because $\dlat(\nsamp-\dlat)/\pwidth-(\nsamp-\dlat)\dlat/\pwidth=0$. An edge formula that divides every block by $\psi_p=\pwidth/\dlat$ would make the left side $O(\dlat/\pwidth)$ and fails this check at $\dlat=40$ by a factor $64$. Block-level closure is weaker than trace-level closure on the released spectra: the noise-dim block sum at $\dlat=40$ is $0.118$ vs $a_1^2(\dlat-\dint)\beta_t^2/\dlat=0.114$, but the sample block carries only $0.57,0.65,0.72,0.77,0.80,0.81$ of $a_*^2(1-\dlat/\nsamp)$ at $\dlat=5,\dots,200$ (theory units, $0.0088$ vs $0.0115$ at $\dlat=40$). The factor $(1-\dlat/\pwidth)$ accounts for $1.6\%$, the $\xi$-average floor for $0$--$20\%$ (growing with $\dlat$), and heterogeneity of the trace split for $<1\%$; the remaining deficit, largest at small $\dlat$, has the size and $\dlat$-dependence of the same-cluster component of $K_{\mathrm{res}}$, which lives in the span of the cluster-indicator vectors and is therefore absorbed by the signal outliers rather than by the sample block. We have not verified this decomposition and record it as open.

\subsubsection{Step 5: separation conditions and limits}

\begin{corollary}[Separation conditions; sketched]\label{cor:fourbulk-separation}
Under \cref{prop:fourbulk}, at leading order in $\dlat/\pwidth$ the blocks are disjoint iff
\begin{align}
  \text{(signal $\succ$ noise-dim)}\qquad& \lambda_{\min}(\hat C+\hat S_\parallel)\;>\;\signoise^2\bigl(1+\sqrt{\gamma_\perp}\bigr)^2,
  \label{eq:fourbulk-sep1}\\
  \text{(noise-dim $\succ$ sample)}\qquad& \frac{\nsamp}{\dlat}\,\frac{a_1(q)^2}{a_*(q)^2}\Bigl[e^{-2t}\signoise^2\bigl(1-\sqrt{\gamma_\perp}\bigr)^2+\Delta_t\Bigr]\;>\;\kappa_{\mathrm{top}} .
  \label{eq:fourbulk-sep2}
\end{align}
In \eqref{eq:fourbulk-sep1} the factors $a_1^2$, $1/\dlat$ and $\Delta_t$ cancel exactly; with population means as proxies it reads $(\alpha_t^2-\Delta_t)/(\beta_t^2-\Delta_t)=(s^2/\dint+\sigsig^2)/\signoise^2>(1+\sqrt{\gamma_\perp})^2$ ($11.2$ versus $\le2.7$ for $\dlat\le240$ at $\signoise=0.5$): signal/null separation is governed by the anisotropy against a Marchenko--Pastur spread of ratio $\dlat/\nsamp$, and $\pwidth$ never enters. In \eqref{eq:fourbulk-sep2} the left side carries the explicit factor $\nsamp/\dlat$, the de-saturation ratio $a_1^2/a_*^2$ (which grows from $2.3$ to $134$ over the sweep) and a bracket that closes as $\dlat\to\nsamp$ or $\signoise\to0$, with an $O(\sqrt{\dlat/\pwidth})$ correction from the $W$ factor (order one for $\pwidth=\dlat+\nsamp+300$ at $\dlat=200$, where $\dlat/\pwidth=0.2$).
\end{corollary}
\begin{proof}
Compare $\lambda^{\mathrm{sig}}_{\min}$ with the upper edge of block (ii) and the lower edge of block (ii) with $\kappa_{\mathrm{top}}\sigma_r$, using the edge formulas of \cref{prop:fourbulk}; the common factor $a_1^2/\dlat$ and the additive $\Delta_t$ cancel in the first comparison, and $\alpha_t^2-\Delta_t=e^{-2t}(s^2/\dint+\sigsig^2)$, $\beta_t^2-\Delta_t=e^{-2t}\signoise^2$. The realized $\lambda_{\min}(\hat C+\hat S_\parallel)\approx1.7$--$1.9$ exceeds $0.25(1+\sqrt{\gamma_\perp})^2\le0.7$.
\end{proof}
The measured gap $\lambda^{\mathrm{nd}}_{\min}/\lambda^{\mathrm{samp}}_{\max}$ is $4.4,7.7,13,18,12$ ($\pwidth=64\dlat$) and $5.1,5.9,8.2,11.4,6.1$ ($\pwidth=\dlat+\nsamp+300$) at $\dlat=10,20,40,100,200$: non-monotone; it opens because $a_1^2/a_*^2$ grows and $\kappa_{\mathrm{top}}$ falls faster than $\nsamp/\dlat$ shrinks, then closes because $(1-\sqrt{\gamma_\perp})^2\to0$ and $\nsamp-\dlat\to0$ as $\dlat\to\nsamp$; at $\dlat=500=\nsamp$ the released spectra show $\lambda^{\mathrm{nd}}_{\min}=0.77$ (code units) and at $\dlat=1000$ the noise-dim median has collapsed to $1.00$.

\begin{remark}[Isotropic limit; sketched]\label{rem:fourbulk-isotropic}
For $\signoise=\sigsig$ condition \eqref{eq:fourbulk-sep1} requires $(s^2/\dint+1)=2.8>(1+\sqrt{\gamma_\perp})^2$, which fails as soon as $\gamma_\perp\gtrsim0.1$: blocks (i)--(ii) merge into a single population bulk with $\dint$ BBP-type outliers at its top, which is $\rho_2$ of \citet{bonnaire2025}, and (iii) is their $\rho_1$. Block (ii) does not disappear: at $\dlat=100$, $\signoise=1$ ($q=1.088$, $a_1^2=0.348$) its leading-order lower edge $a_1^2\pwidth[e^{-2t}(1-\sqrt{\gamma_\perp})^2+\Delta_t]/\dlat=7.40$ (measured $7.4$--$7.5$ in a direct simulation with the repository generator, $\pwidth=6400$, $K=50$; $W$-factor correction $\approx-4\%$) sits $12\times$ above the sample top $0.61$, with signal median $61.8$ (predicted $61.6$) and noise-dim max $44.8$ (predicted $45.5$). The $\dlat-\dint$ noise-dimension modes therefore \emph{exist} in the isotropic case and remain separated from the sample bulk by \eqref{eq:fourbulk-sep2}; anisotropy adds the internal split of $\rho_2$ into (i) and (ii), not the existence of (ii).
\end{remark}

\begin{remark}[$\signoise\to0$; sketched]\label{rem:fourbulk-sigperp0}
As $\signoise\to0$ the noise-dimension population eigenvalues collapse onto the diffusion floor $\theta_\perp=a_1^2\Delta_t/\dlat$, which is comparable to the top edge $\kappa_{\mathrm{top}}\sigma_r$ of the heterogeneous sample bulk ($\kappa_{\mathrm{top}}\approx10$--$16$ at $\dlat\le40$). By \cref{lem:fourbulk-bbp} such directions are not outliers far above the bulk but are repelled to just above its edge, so blocks (ii) and (iii) merge, which is the smooth data/sample transition observed at $\signoise=0.01$: noise-dim min/median/max $0.71/0.87/1.08$, $0.58/0.68/0.96$, $0.64/0.81/1.00$ (code units, $\dlat=10,20,40$) versus the floor $a_1^2\Delta_t\pwidth/\dlat=0.38,0.57,0.76$ and the sample top $0.66,0.58,0.56$; the gap ratio is $1.08,1.01,1.14$. The $\signoise=0$ pencil of \citet{george2025dsm} is this merged limit: the $\dlat-\dint$ null directions are no longer separated from the sample block.
\end{remark}

\paragraph{Relation to prior spectral results.}
Theorem~3.1 of \citet{bonnaire2025} gives the limiting density of $\Umat$ for an \emph{arbitrary} population spectral measure $\rho_\Sigma$; \cref{prop:fourbulk} is that theorem evaluated at the two-atom measure $\rho_\Sigma=\tfrac{\dint}{\dlat}\delta_{\alpha_t^2}+(1-\tfrac{\dint}{\dlat})\delta_{\beta_t^2}$, with one caveat: at fixed $\dint$ the signal atom has vanishing mass in the proportional limit, so block (i) is a finite-rank spike set governed by spiked-covariance perturbation theory rather than a bulk. The bulk part of the Gaussian-equivalence hypothesis is a theorem for deterministic data \citep{louart2018random,liao2018spectrum,fanwang2020spectra}. \citet{george2025dsm} analyze a linear pencil with block dimensions $(\pwidth,\dlat-\dint,\dint,\nsamp)$ for data supported exactly on a $\dint$-dimensional subspace ($\signoise=0$), which is the merged limit of \cref{rem:fourbulk-sigperp0}. What is new here is strictly positive $\signoise$, which lifts (ii) off the sample bulk by $e^{-2t}\signoise^2$ and gives it a Marchenko--Pastur width of ratio $\gamma_\perp$, together with the gradient-flow reading of the blocks as learning-rate bands (\cref{sec:rfnn}).

\subsubsection{The $\xi$-averaged matrix used in the code}

\begin{remark}[Single copy vs.\ $\xi$-average; proved]\label{rem:fourbulk-singlecopy}
For $\Umat^{(1)}$ the statements of \cref{prop:fourbulk} hold with $\hat M_t$ in place of $M_t$ (so block (ii) has leading-order support $\tfrac{a_1^2}{\dlat}\beta_t^2(1\pm\sqrt{\gamma_\perp})^2$: lower edge $1.50$ vs $2.17$ code units at $\dlat=200$; measured $4.58$ vs $4.52$ at $\dlat=40$, $K=1$), $\varepsilon=0$, and exactly $\pwidth-\nsamp$ zeros ($\lambda_{\nsamp+1}/\lambda_\nsamp\sim3\times10^{-6}$ in the float32 code path, $10^{-16}$ in float64): a machine-precision cliff at index $\nsamp$. On the released $K=50$ spectra there is no visible drop at index $\nsamp$ for $\dlat\le20$ ($\lambda_\nsamp/\lambda_{\nsamp+1}=1.01$; $\gamma_p\ge0.4$ pushes the lower MP edge of the sample block to $\approx0$) and a drop by $1.32$, $1.68$, $1.73$ at $\dlat=40,100,200$; there is no drop at index $\dlat+\nsamp$ at any $\dlat$.
\end{remark}
The released \texttt{compute\_U} averages $\Phi_k^\top\Phi_k/\nsamp$ over $K=50$ independent noise draws, so $\operatorname{rank}\Umat\le\min(\pwidth,K\nsamp)=\pwidth$ and the zeros of block (iv) are replaced by a floor. Its size follows from \cref{lem:fourbulk-kernel}: two draws of the same point have $\rho\simeq1-\Delta_t/q$, so their residual kernel is $\sum_{k\ge3}c_k^2(1-\Delta_t/q)^k\simeq a_*^2(1-\varepsilon)$ with $\varepsilon=\bar k\Delta_t/q$, $\bar k=\sum_{k\ge3}kc_k^2/\sum_{k\ge3}c_k^2=3.9,3.7,3.5,3.3,3.2,3.2$ at $\dlat=5,\dots,200$. The residual Gram block of one point is therefore $\tfrac{a_*^2}{K\nsamp}[(1-\varepsilon)\mathbf 1\mathbf 1^\top+\varepsilon I]$: one coherent mode of weight $a_*^2(1-\varepsilon)/\nsamp$ that belongs to the sample block, and $K-1$ incoherent modes of weight $a_*^2\varepsilon/(K\nsamp)$ each, total $\simeq\varepsilon a_*^2$ over the training set, which populate the $\pwidth-\nsamp$ floor when $\pwidth-\nsamp\gg\nsamp$ and stay inside the sample block otherwise. Measured tail mass $\sum_{i>\nsamp}\lambda_i/(\pwidth a_*^2)=0.001,0.017,0.065,0.139,0.196$ at $\dlat=10,\dots,200$ vs $\varepsilon=0.05,0.08,0.11,0.16,0.19$ (agreement at $\dlat\ge100$, where $\pwidth-\nsamp\ge5900$). For $\pwidth=\dlat+\nsamp+300$ the ratio $\lambda_\nsamp/\lambda_{\nsamp+1}$ is $1.00$--$1.01$ throughout. The average also inflates the sample top by $\approx11\%$ at $\dlat=40$ ($0.371\to0.413$, code units), so the $\kappa_{\mathrm{top}}$ column of \cref{tab:fourbulk-numbers} is $\xi$-averaged data compared with a single-copy formula. The data blocks (i)--(ii) are affected only through $\hat M_t\to M_t$, which is why $\Delta_t$ enters (ii) as an exact shift rather than through the Wishart variance. (\cref{app:theory-buffer} writes the same coherent fraction as $\theta(q)=\kappa_{\rm res}(q-\Delta_t)/\kappa_{\rm res}(q)$, the exact Hermite-series form of $1-\varepsilon$.)

\subsubsection{Numerical checks and timescales}\label{app:fourbulk-numerics}

\Cref{tab:fourbulk-numbers} compares the released $\Umat_{\mathrm{code}}$ spectra (seed 42, $\signoise=0.5$, $\pwidth=64\dlat$, $K=50$) with the parameter-free predictions above; $q(\dlat)$ is \eqref{eq:fourbulk-q} and $a_1,a_*,c_k$ are evaluated by $200$-node Gauss--Hermite quadrature. The noise-dimension block is reproduced to $2$--$8\%$ in median, lower and upper edge (max: $6.44/5.69$, $9.52/8.98$, $13.8/13.0$, $19.75/19.68$, $27.2/26.7$ measured/predicted at $\dlat=10,\dots,200$) once its MP shape at ratio $\gamma_\perp$ is used; mean-only predictions overshoot the median by up to $15\%$ at $\dlat=200$. The signal eigenvalues are reproduced to $2$--$5\%$ by the realized centre geometry. On the $\pwidth=\dlat+\nsamp+300$ spectra ($K=500$) the theory-unit noise-dim medians are $0.0083/0.0074$, $0.0051/0.0052$, $0.0030/0.0033$, $0.0014/0.0016$, $0.0007/0.0008$ (measured/predicted), confirming $\pwidth$-independence of the median. In theory units ($\pwidth/\lambda_{\mathrm{code}}$) the gradient-flow timescales on the Table spectra are $\taugen=1/\lambda^{\mathrm{sig}}_{\min}=15.8,20.1,30.5,47.7,94.8,173$ ($10.96\times$, equal to $40\times(0.180/0.625)\times(1.76/1.85)$: the $\dlat/a_1^2$ factor times the realized $\lambda_5$ ratio, with no residual), $\taumem(\text{onset})=1/\lambda^{\mathrm{samp}}_{\max}=334,688,1918,6843,32249,73857$ ($221\times=16.7$ from $a_*^2$ $\times$ $13.1$ from $\kappa_{\mathrm{top}}$) and $\taumem/\taugen=21.2,34.3,62.9,143,340,426$, monotone ($20.1\times$). A fixed-$\mu_1$ theory that ignores the drift of $q(\dlat)$ predicts flat block positions ($64.9$ and $6.22$ in code units) and is rejected by the same data.

\begin{table}[h]
\centering\small
\caption{Predicted vs.\ measured block statistics of $\Umat_{\mathrm{code}}=\pwidth\,\Umat$ ($\signoise=0.5$, $\pwidth=64\dlat$, $\nsamp=500$, $K=50$, $t=0.01$, seed 42). Predictions use $q(\dlat)$ from \eqref{eq:fourbulk-q} and no fitted parameter. Signal min: fifth signal eigenvalue (sets $\taugen$) vs.\ $a_1^2(\pwidth/\dlat)\lambda_5(M_t)$ of the realized training set; noise-dim median and min: MP median and lower edge at ratio $\gamma_\perp$ (leading order in $\dlat/\pwidth$); $\kappa_{\mathrm{top}}$: sample-bulk top in units of $\sigma_r\pwidth$, vs.\ the homogeneous single-copy value $\kappa_{\mathrm{MP}}$ (the $K=50$ average inflates the measured value by $\approx10\%$); $\operatorname{tr}\Umat$: measured vs.\ $\tfrac1\nsamp\sum_\mu\E\tanh^2(\sqrt{q_\mu}z)$ (homogeneous value in parentheses). $^\dagger$At $\dlat=5$, $\pwidth=320<\nsamp$: block (iv) is empty, the sample block has $\pwidth-\dlat$ modes, and $\kappa_{\mathrm{top}}$ is reported in the same units for comparability only.}
\label{tab:fourbulk-numbers}
\resizebox{\textwidth}{!}{%
\begin{tabular}{r r r r | r r | r r | r r | r r | r r}
\toprule
$\dlat$ & $q$ & $a_1^2$ & $a_*^2$ & \multicolumn{2}{c|}{signal min} & \multicolumn{2}{c|}{noise-dim median} & \multicolumn{2}{c|}{noise-dim min} & \multicolumn{2}{c|}{$\kappa_{\mathrm{top}}$} & \multicolumn{2}{c}{$\operatorname{tr}\Umat$}\\
 & & & & meas. & pred. & meas. & pred. & meas. & pred. & meas. & $\kappa_{\mathrm{MP}}$ & meas. & pred.\\
\midrule
5$^\dagger$ & 2.76 & 0.180 & 0.0785 & 20.3 & 20.2 & --   & --   & --   & --   & 19.0 & 5.00 & 0.545 & 0.548 (0.576)\\
10  & 1.52 & 0.281 & 0.0447 & 31.9 & 31.5 & 5.22 & 4.75 & 4.10 & 3.92 & 16.3 & 3.49 & 0.456 & 0.449 (0.470)\\
20  & 0.89 & 0.393 & 0.0235 & 42.0 & 44.2 & 7.03 & 6.60 & 5.16 & 4.71 & 11.1 & 2.57 & 0.363 & 0.361 (0.373)\\
40  & 0.58 & 0.494 & 0.0125 & 53.6 & 55.9 & 8.36 & 8.19 & 4.90 & 4.81 & 5.85 & 2.01 & 0.292 & 0.292 (0.297)\\
100 & 0.39 & 0.586 & 0.0064 & 67.5 & 67.8 & 9.32 & 9.34 & 3.57 & 3.66 & 2.41 & 1.54 & 0.233 & 0.233 (0.235)\\
200 & 0.33 & 0.625 & 0.0047 & 73.9 & 75.1 & 9.22 & 9.31 & 2.12 & 2.17 & 1.45 & 1.31 & 0.209 & 0.209 (0.209)\\
\bottomrule
\end{tabular}}
\end{table}

\paragraph{Predictions and how to test them.}
\begin{enumerate}
\item \emph{Block (ii) is $\mathrm{MP}(\gamma_\perp)\boxtimes\mathrm{MP}(\dlat/\pwidth)$ scaled by $a_1(q)^2/\dlat$: its median is $\pwidth$-independent (to $\sim1\%$), its edges move by $O(\sqrt{\dlat/\pwidth})$, and its lower edge sinks toward $a_1^2\Delta_t/\dlat$ as $\dlat\to\nsamp$.} Test: compare noise-dim medians across $\pwidth\in\{16,64,256\}\dlat$ at $\dlat=40$ (single copy); run an $\nsamp$-sweep at fixed $\dlat=40$ ($\nsamp=250\ldots2000$) and check the lower edge follows $(1-\sqrt{(\dlat-5)/\nsamp})^2$.
\item \emph{The sample-bulk top edge homogenises to $\kappa_{\mathrm{MP}}\sigma_r$ from above as $\dlat$ grows; the small-$\dlat$ excess is set by the spread of $a_*(q_\mu)^2$ and a same-cluster component of size $\sim(\nsamp/k)c_3(q)^2\rho_{\rm same}^3$; the $K$-draw average inflates the top by $\sim10\%$.} Test: compare the sample-block distribution with the MP--Silverstein transform of $\{a_*(q_\mu)^2\}$; regenerate with $k=1$ at fixed norm; recompute with $K=1$.
\item \emph{$\taumem(\text{onset})=1/(\kappa_{\mathrm{top}}\sigma_r)$ decomposes into the $a_*(q(\dlat))^2$ trajectory ($16.7\times$, closed form) and the homogenisation of $\kappa_{\mathrm{top}}$ ($13.1\times$, characterised); a different $\signoise$ changes the predicted $a_*^2$ curve before any run.} Test: regress $\log\taumem(\text{onset})$ on $\log a_*^2$ and $\log\kappa_{\mathrm{top}}$; predict $a_*(q(\dlat))^2$ for $\signoise\in\{0.1,0.3,0.01\}$ and compare the sample-top eigenvalue on existing spectra.
\item \emph{$\taugen=\dlat/(a_1(q(\dlat))^2[e^{-2t}\lambda_5(\hat C+\hat S_\parallel)+\Delta_t])$ in theory units; $a_1^2$ rising $0.18\to0.625$ makes the growth $\sim11\times$ over $\dlat=5..200$ rather than $40\times$.} Test: plot $\taugen(\dlat)a_1(q)^2/\dlat$ against $1/(e^{-2t}\lambda_5+\Delta_t)$ of the realized training set; repeat over centre-geometry seeds.
\item \emph{Memorisation returns as $\dlat\to\nsamp$ because the noise-dim/sample gap closes and the sample block empties; the gap ratio is non-monotone with a peak near $\dlat=100$.} Test: check the $\dlat=500,1000$ spectra for the noise-dim minimum and the drop at index $\nsamp$; train RFNNs at $\dlat\in\{300,400,480\}$ and check the fraction rises.
\item \emph{With one diffused copy $\Umat$ has exactly $\pwidth-\nsamp$ zeros; with the $K$-draw average the block becomes a floor of mass $\sim\bar k(\Delta_t/q)a_*^2$, visible as a $1.3$--$1.8\times$ drop at index $\nsamp$ only for $\gamma_p\le0.2$.} Test: recompute $\Umat$ with $K\in\{1,50\}$ and $t\in\{0.001,0.01,0.05\}$; the tail mass should scale as $\Delta_t/q$.
\item \emph{Trace-level bookkeeping closes to $1.5\%$; block-level bookkeeping does not: the sample block carries $0.57$--$0.81$ of its predicted mass.} Test: decompose $K_{\mathrm{res}}$ of the realized training set into diagonal, same-cluster and cross-cluster parts and check whether the same-cluster mass accounts for the deficit.
\item \emph{As $\signoise\to0$ blocks (ii) and (iii) merge smoothly.} Test: sweep $\signoise\in\{0.01,0.05,0.1,0.25,0.5\}$ at $\dlat=40$ and track $\lambda^{\mathrm{nd}}_{\min}/\lambda^{\mathrm{samp}}_{\max}$; it should cross $1$ where $a_1^2e^{-2t}\signoise^2(1-\sqrt{\gamma_\perp})^2/\dlat\sim(\kappa_{\mathrm{top}}-1)\sigma_r$.
\end{enumerate}

\paragraph{Open gaps.}
(1) The proposition is stated for the $\xi$-averaged $\Umat$; the single-copy exact-rank statement is \cref{rem:fourbulk-singlecopy}, and on the released $K=50$ spectra there is no drop at index $\nsamp$ for $\dlat\le20$: ``cliff at $\nsamp$'' is a single-copy fact, and any figure colouring the sample block as $[\dlat+1,\dlat+\nsamp]$ must be re-cut to $[\dlat+1,\nsamp]$ and $[\nsamp+1,\pwidth]$. (2) The additive-spike model is the wrong model for $G^{\mathrm{GE}}$; the spiked-population-covariance version is used, but its per-outlier mean shift is proved only in expectation (simulations give $0.97\pm0.2\,\sigma_r$ and $1.41\pm0.15\,\sigma_r$, neither discriminating from $1.18\sigma_r$). (3) The $W$ factor corrects edges by $O(\sqrt{\dlat/\pwidth})$; whether this factor carries the $\psi_p$-dependence of Theorem~3.1 of \citet{bonnaire2025}, and whether their zero-mass atom is $1-\nsamp/\pwidth$ or $1-(\nsamp+\dlat)/\pwidth$, has not been checked against their supplement. (4) Block-level mass: the sample block carries only $0.57$--$0.81$ of $a_*^2(1-\dlat/\nsamp)(1-\dlat/\pwidth)$; the candidate explanation (same-cluster component of $K_{\mathrm{res}}$ absorbed by the signal outliers) is unverified. (6) Gaussian equivalence: the bulk part is imported and the hypotheses of \citet{louart2018random} (Lipschitz activation, bounded data norm) have not been re-verified beyond boundedness of $q_\mu$; the $\dint$ outliers and the $\xi$-averaged matrix remain conjectural. (7) $\kappa_{\mathrm{top}}$ is characterised, not derived; homogenisation to $\kappa_{\mathrm{MP}}$ is an observation. (8) The floor mass $\bar k\Delta_t/q$ is a heuristic that matches at $\dlat\ge100$ and fails at $\dlat\le20$. (9) The proposition is a proportional-limit statement while the sweeps hold $\nsamp=500$ fixed ($\gamma_\perp\in[0,0.47]$, $\gamma_p$ up to $0.6$); the $\dlat=5$ point has $\pwidth<\nsamp$. (10) Some archived configuration files report a stale $\dlat$; they should be regenerated before those directories are cited.

% ----------------------------------------------------------------------
\subsection{The buffer as a ratio of spectral edges}
\label{app:theory-buffer}

\paragraph{Setting.}
We work with the RFNN of \eqref{eq:rfnn} trained by the gradient flow \eqref{eq:rfnn-gradflow} at fixed $t$, with $\Umat$ as in \eqref{eq:U}. Write $\Umat=\sum_{i=1}^{\pwidth}\lambda_iv_iv_i^\top$ with $\lambda_1\ge\dots\ge\lambda_{\pwidth}\ge0$, and use the a-priori index blocks $I_{\rm sig}=\mathcal S=\{1,\dots,\dint\}$, $I_{\rm noise}=\mathcal N=\{\dint+1,\dots,\dlat\}$, $I_{\rm samp}=\mathcal M=\{\dlat+1,\dots,\nsamp\}$, $I_{\rm null}=\mathcal Z=\{\nsamp+1,\dots,\pwidth\}$ of (A5). Data and notation are as in \cref{app:theory-setup}; in addition $\alpha_0^2:=s^2/\dint+\sigsig^2$ (so $\alpha_t^2=e^{-2t}\alpha_0^2+\Delta_t$), $\psi_\perp:=(\dlat-\dint)/\nsamp$ ($=\gamma_\perp$ of \cref{app:theory-fourbulk}) and $\gamma:=\nsamp/\pwidth$. ``$\dlat$ sweep'' means $\dlat$ varies with $(\dint,\nsamp,k,s,\sigsig,\signoise,t)$ fixed and $\pwidth$ either $\propto\dlat$ or fixed $\ge\nsamp$. Eigenvalues in the released code are $\pwidth\times$ the ones here; all ratios are normalisation-free. The nonzero spectrum of the Monte-Carlo estimate of $\Umat$ with $N_{\rm mc}$ draws equals that of the $(\nsamp N_{\rm mc})\times(\nsamp N_{\rm mc})$ Gram $\Phi\Phi^\top/(\nsamp N_{\rm mc})$, which is how the simulations of this subsection (``sim'': $N_{\rm mc}=4$, seed $0$) are computed; ``repo'' refers to the released $\signoise=0.5$ RFNN spectra ($\pwidth=64\dlat$, $N_{\rm mc}=50$, seed $42$) unless another set is named.

\begin{definition}[$\varepsilon$-absorption times]\label{def:abs-times-buf}
Let $a_i(T):=A(T)v_i\in\R^{\dlat}$ and $a_i^\star:=A^\star v_i$, so that $a_i(T)-a_i^\star=(a_i(0)-a_i^\star)e^{-\lambda_iT}$ (\cref{lem:gf-p1}). For $\varepsilon\in(0,1)$ the $\varepsilon$-absorption time of mode $i$ is $\tau_i(\varepsilon):=\lambda_i^{-1}\log(1/\varepsilon)$. Define
\begin{equation}
\taugen(\varepsilon):=\max_{i\in I_{\rm sig}}\tau_i(\varepsilon)=\frac{\log(1/\varepsilon)}{\lambda^{\rm sig}_{\min}},\qquad
\taumem^{\rm on}(\varepsilon):=\min_{i\in I_{\rm samp}}\tau_i(\varepsilon)=\frac{\log(1/\varepsilon)}{\lambda^{\rm samp}_{\max}},\qquad
R:=\frac{\taumem^{\rm on}}{\taugen}=\frac{\lambda^{\rm sig}_{\min}}{\lambda^{\rm samp}_{\max}},\qquad
R_{\rm full}:=\frac{\lambda^{\rm noise}_{\min}}{\lambda^{\rm samp}_{\max}},
\label{eq:def-R-buf}
\end{equation}
with $\lambda^{\rm sig}_{\min}:=\lambda_{\dint}$, $\lambda^{\rm noise}_{\min}:=\lambda_{\dlat}$, $\lambda^{\rm samp}_{\max}:=\lambda_{\dlat+1}$. $\taugen$ is the time to learn the score on the \emph{signal subspace}; the noise-dimension modes, which carry the exactly linear null score (\cref{app:theory-linear}), are slower, so $R_{\rm full}\le R$ measures the onset relative to the whole data block. Both ratios are independent of $\varepsilon$ and of the learning rate. These are the per-mode versions of \cref{def:taugen-p1,def:taumem-p1}: they coincide with the block functionals there up to the factor $2$ in \eqref{eq:taugen-bounds-p1} and the $\theta_{\mathcal B}$ correction, and $\taumem^{\rm on}$ is the index-based statistic that lower-bounds $\taumem^{\rm onset}$ of \cref{def:taumem-p1}(i).
\end{definition}
\begin{proof}
Each mode is a scalar exponential with rate $\lambda_i$; the max (resp.\ min) of $\lambda_i^{-1}\log(1/\varepsilon)$ over a block is attained at its smallest (resp.\ largest) eigenvalue, and $\log(1/\varepsilon)$ cancels in the ratios. $R_{\rm full}\le R$ because $\lambda_{\dlat}\le\lambda_{\dint}$.
\end{proof}

\subsubsection{Standing hypotheses}

\begin{assumption}[Inputs to the buffer corollary; conjectured]\label{assump:inputs-buf}
\begin{itemize}
\item[(H1)] (\emph{Gaussian equivalence at block scale.}) With $a_1(q)$, $a_*(q)^2$ as before, write $\phi(x)=a_1(q)(\pwidth\dlat)^{-1/2}Wx+a_*(q)\pwidth^{-1/2}z(x)$, where $z(x)\in\R^{\pwidth}$ has zero mean, unit-variance coordinates, is uncorrelated with $Wx$, and has Gram $\E_W[z(x)^\top z(x')]/\pwidth=\kappa_{\rm res}(x\!\cdot\!x'/\dlat)/a_*(q)^2$,
\begin{equation}
\kappa_{\rm res}(\rho)=\sum_{\ell\ {\rm odd}\ \ge3}c_\ell(q)^2\,(\rho/q)^\ell,\qquad c_\ell(q):=\frac{\E[\tanh(\sqrt qz)\,\mathrm{He}_\ell(z)]}{\sqrt{\ell!}},\qquad\kappa_{\rm res}(q)=a_*(q)^2
\label{eq:kres-buf}
\end{equation}
(the Hermite series of \cref{lem:fourbulk-kernel}). We assume $\Umat=\Umat_{\rm lin}+\Umat_{\rm res}+E$ with $\Umat_{\rm lin}=a_1(q)^2(\dlat\pwidth)^{-1}WM_tW^\top$, $\Umat_{\rm res}=a_*(q)^2(\nsamp\pwidth)^{-1}\sum_\mu\E_\xi[z(x^\mu_t)z(x^\mu_t)^\top]$, and $E$ negligible \emph{at the scale of each block}: $\|P_{\rm data}EP_{\rm data}\|=o(a_1^2/\dlat)$, $\|P_{\rm samp}EP_{\rm samp}\|=o(\theta a_*^2/\nsamp)$. A plain $o(1)$ operator-norm statement would be vacuous (every block is $o(1)$). This is stronger than the ESD-level Gaussian-equivalence theorems \citep{goldt2022gaussian,hu2022universality,mei2022generalization}, is of the type proved for the extreme eigenvalues of the conjugate kernel with i.i.d.\ inputs by \citet{benigni2021eigenvalue}, and is assumed here for the $k$-centre mixture; its trace consequence is verified below.
\item[(H2)] (\emph{$\xi$-averaging.}) Two independent diffusion draws of the same $x^\mu$ have $\E[x^\mu_t\!\cdot\!x^{\mu\prime}_t]/\dlat=e^{-2t}\|x^\mu\|^2/\dlat=q-\Delta_t+o(1)$; hence the $\xi$-averaged residual feature $\bar z_\mu:=\E_\xi z(x^\mu_t)$ has coherent self-kernel $\kappa_{\rm res}(q-\Delta_t)=\theta(q)a_*(q)^2$, $\theta(q):=\kappa_{\rm res}(q-\Delta_t)/\kappa_{\rm res}(q)$ (\cref{lem:coeffs-buf}(v)). We assume the incoherent residual mass $(1-\theta)a_*^2$, spread over $\pwidth$ modes (exact $\E_\xi$) or $\nsamp(N_{\rm mc}-1)$ modes ($N_{\rm mc}$ draws), lies below the sample block: $(1-\theta)/\min(\pwidth,\nsamp N_{\rm mc})\ll\theta/\nsamp$. (On the repo's $N_{\rm mc}=50$ spectra $\lambda_{\nsamp+1}/\lambda_{\nsamp}=0.48$--$0.76$ for $\dlat\ge40$: the index-$\nsamp$ step is real but not sharp at finite $N_{\rm mc}$.)
\item[(H3)] (\emph{Asymptotic regime.}) $\dlat,\nsamp,\pwidth\to\infty$ with $\psi_n=\nsamp/\dlat$, $\psi_p=\pwidth/\dlat$, $q$ as parameters; along the sweep these move ($\psi_n$ from $100$ to $\approx1$), so no statement is uniform in $\dlat$, and finite-size corrections are $O(\dlat^{-1/2})$ relative ($\|x^\mu\|^2/\dlat$ fluctuates by $\approx25\%$ at $\dlat=5$: trace deficit $5\%$ there, $<1\%$ for $\dlat\ge100$). $\pwidth\ge\nsamp>\dlat$ is needed for the block counts; for $\pwidth=64\dlat$ it fails at $\dlat\le7$.
\item[(H4)] (\emph{Data-side edges; \cref{app:theory-fourbulk}.}) Let $\tilde\nu_1\ge\dots\ge\tilde\nu_{\dint}$ be the top $\dint$ eigenvalues of the full empirical covariance $\hat C+\hat S$ and $\nu_1\ge\dots\ge\nu_{\dint}$ those of its $\dint\times\dint$ signal block (law independent of $\dlat$; mean $\alpha_0^2$; measured $\nu_{\min}/\alpha_0^2=0.39$--$0.46$). Then $\lambda_i=a_1(q)^2(e^{-2t}\tilde\nu_i+\Delta_t)/\dlat\,(1+o(1))$ for $i\in I_{\rm sig}$, with the spiked-covariance law \citep{baik2005phase}
\begin{equation}
\tilde\nu_i=\nu_i\Bigl(1+\frac{\psi_\perp\signoise^2}{\nu_i-\signoise^2}\Bigr)(1+o(1))\qquad\text{if }\nu_i>\signoise^2(1+\sqrt{\psi_\perp}),
\label{eq:bbp-signal-buf}
\end{equation}
and $I_{\rm noise}$ occupies $a_1(q)^2[e^{-2t}\signoise^2(1\pm\sqrt{\psi_\perp})^2+\Delta_t]/\dlat$ (Marchenko--Pastur law of the empirical null covariance centred at $\beta_t^2$; verified to $3$--$6\%$ at both edges at $\pwidth=64\dlat$; the edges move by $\approx10\%$ between $\pwidth=850$ and $3200$ at $\dlat=50$, the medians by $<0.5\%$).
\item[(H5)] (\emph{Block separation.}) $\lambda^{\rm noise}_{\min}>\lambda^{\rm samp}_{\max}$. On the repo's spectra this holds at $\signoise=0.5$ ($R_{\rm full}=4.4$--$18.8$ over $\dlat=10$--$200$) and fails at $\signoise=0.01$ ($R_{\rm full}=1.0$--$1.15$ at every $\dlat=8$--$40$), where the noise-dimension block sits on the diffusion floor $a_1^2\Delta_t/\dlat$ and merges with the sample block (\cref{rem:fourbulk-sigperp0}). All statements about $\taumem^{\rm on}$, $R$, $R_{\rm full}$ are restricted to $\signoise=0.5$ (and the isotropic control); at $\signoise=0.01$ only $\taugen$ is predicted.
\end{itemize}
\end{assumption}

\begin{lemma}[Gaussian-equivalent coefficients along the sweep; proved except as noted]\label{lem:coeffs-buf}
Assume the data construction above with $\alpha_0^2>\signoise^2$, $\nsamp\to\infty$ (so $\operatorname{tr}\hat C\to s^2$, $\operatorname{tr}\hat S\to\dint\sigsig^2+(\dlat-\dint)\signoise^2$) and $\Delta_t<q$.
\begin{enumerate}
\item[(i)] $q(\dlat)=q_\infty+c_q/\dlat$, $q_\infty:=e^{-2t}\signoise^2+\Delta_t$, $c_q:=e^{-2t}\dint(\alpha_0^2-\signoise^2)>0$; $q$ decreases strictly from $q_{\max}:=q(\dint)=\alpha_t^2$ to $q_\infty$ (this is \eqref{eq:q-p1} rewritten).
\item[(ii)] $a_1(q)=\E[\operatorname{sech}^2(\sqrt qz)]$ is strictly decreasing in $q$, $a_1\le1$, $a_1(q)^2=1-2q+O(q^2)$; $q\,a_1(q)^2=c_1(q)^2$ is strictly increasing in $q$; hence $a_1(q(\dlat))^2$ is strictly increasing in $\dlat$ with limit $a_1(q_\infty)^2$.
\item[(iii)] $a_*(q)^2=\sum_{\ell\ {\rm odd}\ge3}c_\ell(q)^2$ and $a_*(q)^2=\tfrac23q^3(1-8q+50q^2+O(q^3))$ as $q\to0$ (within $4\%$ at $q=0.05$; not usable beyond $q\approx0.1$).
\item[(iv)] (\emph{Trace identity}) $a_1(q)^2q+a_*(q)^2=\E[\tanh^2(\sqrt qz)]=\operatorname{tr}\Umat\,(1+O(\dlat^{-1/2}))$; equivalently $\dint\lambda^{\rm sig}+(\dlat-\dint)\lambda^{\rm noise}+\sum_{I_{\rm samp}}\lambda_i+(\text{floor})=\E\tanh^2(\sqrt qz)=O(1)$, which any proposed edge must satisfy.
\item[(v)] (\emph{$\xi$-averaging factor}) $\theta(q):=\kappa_{\rm res}(q-\Delta_t)/\kappa_{\rm res}(q)$ satisfies $0<\theta(q)\le(1-\Delta_t/q)^3$, with equality to leading order as $q\to0$, and $\theta\to1$ as $\Delta_t/q\to0$.
\end{enumerate}
For the repo sweep ($t=0.01$, $\Delta_t=0.0198$, $\dint=5$, $s=3$, $\sigsig=1$, $\nsamp=500$; $120$-node Gauss--Hermite quadrature) at $\signoise=0.5$: $c_q=12.50$, $q_\infty=0.265$, $d^\star:=c_q/q_\infty=47$, $a_1(q_\infty)^2/a_1(q_{\max})^2=3.73$, and
\begin{center}\small
\begin{tabular}{lccccccccc}
$\dlat$ & 5 & 10 & 20 & 50 & 100 & 200 & 300 & 400 & 480\\ \hline
$q$ & 2.764 & 1.515 & 0.890 & 0.515 & 0.390 & 0.327 & 0.307 & 0.296 & 0.291\\
$a_1(q)^2$ & 0.180 & 0.281 & 0.393 & 0.521 & 0.586 & 0.625 & 0.640 & 0.647 & 0.651\\
$a_*(q)^2\ (\times10^3)$ & 78.5 & 44.7 & 23.5 & 10.4 & 6.44 & 4.67 & 4.11 & 3.85 & 3.72\\
$\theta(q)$ & 0.969 & 0.951 & 0.925 & 0.880 & 0.847 & 0.821 & 0.811 & 0.805 & 0.802\\
$\E\tanh^2$ (pred.) & 0.576 & 0.470 & 0.373 & 0.278 & 0.235 & 0.209 & 0.200 & 0.195 & 0.193\\
$\operatorname{tr}\Umat$ (sim.) & 0.551 & 0.458 & 0.365 & 0.271 & 0.233 & 0.208 & 0.200 & 0.196 & 0.193
\end{tabular}
\end{center}
(repo: $\operatorname{tr}\Umat=0.545,0.456,0.363,0.233,0.209$ at $\dlat=5,10,20,100,200$). At $\signoise=0.01$: $c_q=13.72$, $q_\infty=0.0199$, $d^\star=690$, slope ratio $5.35$, $\theta=0.80,0.66,0.47$ at $\dlat=50,100,200$. At $\signoise=1$: $c_q=8.82$, $q_\infty=1$, $d^\star=8.8$, slope ratio $2.04$, $\theta\ge0.93$.
\end{lemma}
\begin{proof}
(i) $\operatorname{tr}M_t=e^{-2t}[\dint\alpha_0^2+(\dlat-\dint)\signoise^2]+\dlat\Delta_t$; divide by $\dlat$. (ii) $\operatorname{sech}^2$ is even and strictly decreasing in $|x|$, so for $q<q'$ and every $z\ne0$, $\operatorname{sech}^2(\sqrt qz)>\operatorname{sech}^2(\sqrt{q'}z)$; integrate. The expansion uses $\operatorname{sech}^2x=1-x^2+\tfrac23x^4-\dots$. By Stein's lemma $c_1(q)=\E[z\tanh(\sqrt qz)]=\sqrt q\,a_1(q)$, so $q\,a_1(q)^2=c_1(q)^2=(\E[|z|\tanh(\sqrt q|z|)])^2$, increasing in $q$ since $\tanh$ is. Compose with (i). (iii) $\tanh$ is odd so even Hermite coefficients vanish; Parseval gives $\E\tanh^2(\sqrt qz)=\sum_\ell c_\ell^2$, hence $a_*^2=\sum_{\ell\ge3}c_\ell^2$. From $\tanh x=x-\tfrac{x^3}3+\tfrac{2x^5}{15}-\tfrac{17x^7}{315}+\dots$, $x^3=\mathrm{He}_3+3\mathrm{He}_1$, $x^5=\mathrm{He}_5+10\mathrm{He}_3+15\mathrm{He}_1$, $x^7=\mathrm{He}_7+21\mathrm{He}_5+105\mathrm{He}_3+105\mathrm{He}_1$: $c_3(q)=-\tfrac{2}{\sqrt6}q^{3/2}(1-4q+17q^2+O(q^3))$, so $c_3^2=\tfrac23q^3(1-8q+50q^2+O(q^3))$, while $c_5^2=O(q^5)$; quadrature: $c_3^2/[\tfrac23q^3(1-8q+50q^2)]=0.998,0.960,0.73$ at $q=0.02,0.05,0.1$. (iv) $\operatorname{tr}\Umat=\nsamp^{-1}\sum_\mu\E\|\phi(x^\mu_t)\|^2=\E\tanh^2(w\!\cdot\!x^\mu_t/\sqrt{\dlat})$ with $w\!\cdot\!x^\mu_t/\sqrt{\dlat}\sim\mathcal N(0,q_\mu)$, $q_\mu=\|x^\mu_t\|^2/\dlat=q(1+O(\dlat^{-1/2}))$; Parseval again. (v) $\theta=\sum_{\ell\ge3}c_\ell^2(1-\Delta_t/q)^\ell/\sum_{\ell\ge3}c_\ell^2$ is a weighted average of $(1-\Delta_t/q)^\ell$, $\ell\ge3$ odd, with $0<1-\Delta_t/q<1$, hence $\le(1-\Delta_t/q)^3$; as $q\to0$ the weight concentrates on $\ell=3$ by (iii).
\end{proof}

\begin{remark}[Numerical monotonicity of $a_*^2$; conjectured]\label{rem:astar-mono-buf}
On $q\in[0.01,4]$, $a_*(q)^2$ and $\theta(q)a_*(q)^2$ are strictly increasing in $q$ (quadrature on a $400$-point grid; proved only for $q<0.05$ via \cref{lem:coeffs-buf}(iii)). Hence $\theta a_*^2$ is strictly decreasing in $\dlat$ along the sweep: $0.0761\to0.00383$ over $\dlat=5\to200$ at $\signoise=0.5$. We use this only as a numerical input.
\end{remark}

\subsubsection{Scaling of $\taugen$}

\begin{proposition}[$\taugen$ scales with $\dlat$; conditional on (H1), (H3), (H4)]\label{prop:taugen-buf}
With $\tilde\nu_{\min}:=\tilde\nu_{\dint}$ and $\nu_{\min}:=\nu_{\dint}$,
\begin{equation}
\taugen(\varepsilon)=\frac{\dlat\,\log(1/\varepsilon)}{a_1(q(\dlat))^2\,(e^{-2t}\tilde\nu_{\min}(\dlat)+\Delta_t)}\,(1+o(1)),\qquad q(\dlat)=q_\infty+\frac{c_q}{\dlat},\qquad \tilde\nu_{\min}=\nu_{\min}\Bigl(1+\frac{\psi_\perp\signoise^2}{\nu_{\min}-\signoise^2}\Bigr).
\label{eq:taugen-buf}
\end{equation}
The prefactor $\dlat/a_1(q(\dlat))^2$ is strictly increasing in $\dlat$, and for $\psi_\perp\signoise^2\ll\nu_{\min}-\signoise^2$ (i.e.\ $\dlat\ll\nsamp(\nu_{\min}-\signoise^2)/\signoise^2\approx3.6\nsamp$ at $\signoise=0.5$)
\begin{equation}
\frac{\dlat\log(1/\varepsilon)}{a_1(q_\infty)^2(e^{-2t}\nu_{\min}+\Delta_t)}\;\lesssim\;\taugen(\varepsilon)\;\lesssim\;\frac{\dlat\log(1/\varepsilon)}{a_1(q_{\max})^2(e^{-2t}\nu_{\min}+\Delta_t)},
\label{eq:taugen-sandwich-buf}
\end{equation}
with slope ratio $a_1(q_\infty)^2/a_1(q_{\max})^2$ ($3.73$ at $\signoise=0.5$, $5.35$ at $\signoise=0.01$). For $\dlat\lesssim d^\star$ the growth is sublinear ($a_1^2$ still rising); for $d^\star\ll\dlat\ll3.6\nsamp$ it is linear with slope $\log(1/\varepsilon)/[a_1(q_\infty)^2(e^{-2t}\nu_{\min}+\Delta_t)]$; beyond, the BBP factor (linear in $\dlat$) makes it sublinear again, with saturation at $\nsamp(\nu_{\min}-\signoise^2)\log(1/\varepsilon)/[a_1(q_\infty)^2e^{-2t}\nu_{\min}\signoise^2]$ as long as the spike stays above threshold ($\dlat\lesssim13\nsamp$).
\end{proposition}
\begin{proof}
By \cref{def:abs-times-buf}, $\taugen=\log(1/\varepsilon)/\lambda_{\dint}$, and by (H4) $\lambda_{\dint}=a_1(q)^2(e^{-2t}\tilde\nu_{\min}+\Delta_t)/\dlat\,(1+o(1))$. The law of $\nu_{\min}$ is $\dlat$-free (signal coordinates are drawn before the embedding; $WQ^\top\overset d=W$ makes $\operatorname{spec}\Umat$ invariant under $Q$); the $\dlat$-dependence of $\tilde\nu_{\min}$ is the BBP factor \eqref{eq:bbp-signal-buf}. Substituting \cref{lem:coeffs-buf}(i) gives \eqref{eq:taugen-buf}. \emph{Monotonicity of $g(d):=d/a_1(q(d))^2$}: $d=c_q/(q-q_\infty)$, so $g=c_q/h(q)$ with $h(q):=(q-q_\infty)a_1(q)^2=c_1(q)^2-q_\infty a_1(q)^2$; $c_1^2$ is increasing and $-q_\infty a_1^2$ is increasing (\cref{lem:coeffs-buf}(ii)), so $h$ is increasing in $q$, $g$ decreasing in $q$, and since $q$ decreases in $d$, $g$ increases in $d$. The sandwich follows from $a_1(q_{\max})^2\le a_1(q(d))^2\le a_1(q_\infty)^2$ and $\tilde\nu_{\min}=\nu_{\min}(1+o(1))$ in the stated regime; the large-$\dlat$ statements substitute $\psi_\perp=(\dlat-\dint)/\nsamp$ into \eqref{eq:bbp-signal-buf}. Numerically, $\tilde\nu_{\min}/\nu_{\min}=1.01,1.03,1.06,1.11,1.19,1.21,1.23$ at $\dlat=10,50,100,200,300,400,480$ ($\signoise=0.5$), and \eqref{eq:taugen-buf} reproduces the simulated $1/\lambda^{\rm sig}_{\min}$ to $1$--$4\%$ (\cref{tab:buffer-sim}); on the repo's $\dlat=500\to1000$ spectra $\taugen$ grows $1.28\times$ versus $1.62\times$ predicted with $\nu_{\min}=0.41\alpha_0^2$ and $2\times$ for linear growth.
\end{proof}

\subsubsection{Scaling of the onset}

\begin{proposition}[Onset of the sample block; sketched, small-$\dlat$ factor $g$ open]\label{prop:onset-buf}
Assume (H1)--(H3), (H5), clusters of sizes $n_c$, and within-cluster overlaps $\rho_c=e^{-2t}x^\mu\!\cdot\!x^\nu/\dlat$ with $\E\rho_c=e^{-2t}s^2/\dlat$ and non-vanishing relative fluctuations $O(\sigsig/s)$. Let $\bar Z=[\bar z_1,\dots,\bar z_{\nsamp}]\in\R^{\pwidth\times\nsamp}$ and $\bar K:=\E_W[\bar Z^\top\bar Z]/(\pwidth\theta a_*^2)$ (unit diagonal; $\bar K_{\mu\nu}=\kappa_{\rm res}(\rho_{\mu\nu})/\kappa_{\rm res}(q-\Delta_t)$). To leading order:
\begin{enumerate}
\item[(a)] (\emph{Mass and mean.}) The sample block is the compression of the $\nsamp$-mode Wishart $\Umat^{\rm coh}_{\rm res}=\theta a_*^2(\nsamp\pwidth)^{-1}\bar Z\bar Z^\top$ onto the orthogonal complement of the $\dlat$ data-block eigenvectors; it has $\nsamp-\dlat$ modes, mean eigenvalue $\theta(q)a_*(q)^2/\nsamp$ and total mass $(1-\dlat/\nsamp)\theta a_*^2$.
\item[(b)] (\emph{Edges by interlacing.}) Let $\mu_1\ge\dots\ge\mu_{\nsamp}$ be the eigenvalues of $\Umat^{\rm coh}_{\rm res}$ in units of $\theta a_*^2/\nsamp$. Then $\mu_{\dlat+1}\le\lambda^{\rm samp}_{\max}\nsamp/(\theta a_*^2)\le\mu_1$. For $\bar K=I$, $\mu_1=(1+\sqrt\gamma)^2$ and $\mu_{\dlat+1}$ is the Marchenko--Pastur quantile of order $1-\dlat/\nsamp$; for $\bar K=I+\sum_cr_c\mathbf 1_c\mathbf 1_c^\top+E$, $r_c=\kappa_{\rm res}(\rho_c)/\kappa_{\rm res}(q-\Delta_t)$, a cluster spike detaches when $\ell_c:=1+n_cr_c>1+\sqrt\gamma$, at $\mu_1=\ell_c(1+\gamma/(\ell_c-1))$ \citep{baik2005phase}.
\item[(c)] (\emph{Onset.}) With $g(\dlat):=\lambda^{\rm samp}_{\max}\nsamp/(\theta a_*^2)$,
\begin{equation}
\lambda^{\rm samp}_{\max}=\frac{\theta(q(\dlat))\,a_*(q(\dlat))^2}{\nsamp}\,g(\dlat)\,(1+o(1)),\qquad \mu_{\dlat+1}\le g(\dlat)\le\max\Bigl\{(1+\sqrt\gamma)^2,\ \max_c\ell_c\bigl(1+\tfrac{\gamma}{\ell_c-1}\bigr)\Bigr\},
\label{eq:lmax-samp-buf}
\end{equation}
\begin{equation}
\taumem^{\rm on}(\varepsilon)=\frac{\nsamp\log(1/\varepsilon)}{\theta(q(\dlat))\,a_*(q(\dlat))^2\,g(\dlat)}\,(1+o(1)).
\label{eq:tauon-buf}
\end{equation}
\end{enumerate}
$\theta a_*^2$ is strictly decreasing in $\dlat$ (\cref{rem:astar-mono-buf}), so $\taumem^{\rm on}$ grows at least like $1/(\theta a_*^2)$ once $g$ has settled. At $\signoise=0.5$, $g$ itself falls from $\approx5$--$6.5$ ($\dlat\le20$) to $\approx1$ ($\dlat\ge200$) and below $1$ for $\dlat/\nsamp\gtrsim0.6$ (interlacing), adding $\approx3$--$4\times$. We have \emph{no closed form for $g$ at small $\dlat$}: the spike bound in (b) with the full kernel gives $2.5,2.9,2.7,1.7,1.1,1.0$ at $\dlat=5,10,20,50,100,200$ against observed $5.2,6.6,5.1,3.2,1.7,1.25$, and the top sample eigenvector has squared overlap only $0.23$--$0.29$ with the $k$ cluster indicators at $\dlat\le20$ (null $\approx k/\nsamp=0.02$), $0.17$ at $\dlat=50$, $0.03$ at $\dlat=200$. The remainder is attributed to the finite tensor rank of the cubic residual features ($\mathrm{He}_3$ features span $\le D_3=\binom{\dlat+2}{3}=35,220,1540$ dimensions at $\dlat=5,10,20$, versus $\nsamp=500$), which makes $\bar K$ rank-deficient for $\dlat\lesssim13$ and far from $I$ well beyond; an MP description needs at least $\nsamp\ll D_3$. In this subsection $g(\dlat)$ is an empirically calibrated factor (it is $\kappa_{\rm top}/\theta$ in the notation of \cref{prop:fourbulk}). The cluster-spike ratio $n_cr_c$ is roughly constant for $\dlat\lesssim d^\star$ ($\rho_c/q\to s^2/[\dint(\alpha_0^2-\signoise^2)]=0.71$ while $q\approx c_q/\dlat$; $n_cr_c=10.1,8.4,5.6,2.0$ at $\dlat=5,10,20,50$) and decays like $(d^\star/\dlat)^3$ only beyond, where it is below threshold.
\end{proposition}
\begin{proof}[Proof sketch]
Under (H1)--(H2), $\Umat=\Umat_{\rm lin}+\Umat^{\rm coh}_{\rm res}+(\text{floor})+E$; $\Umat^{\rm coh}_{\rm res}$ and $\Umat_{\rm lin}$ are functions of the same $\nsamp$ points, so $\operatorname{rank}\Umat\le\nsamp$. Let $P$ project onto the top-$\dlat$ eigenvectors of $\Umat$ ($\approx$ range of $WM_t^{1/2}$ by (H4)), $P^\perp=I-P$. The eigenvalues of $\Umat$ below the top $\dlat$ are exactly those of $P^\perp\Umat P^\perp$ on $\operatorname{ran}P^\perp$, and $P^\perp\Umat_{\rm lin}P^\perp\approx0$, so the sample block is the spectrum of the compression $P^\perp\Umat^{\rm coh}_{\rm res}P^\perp$ of an $\nsamp$-mode Wishart by a projector of corank $\dlat$; Cauchy interlacing gives $\mu_{i+\dlat}\le\mu^{\rm samp}_i\le\mu_i$, whence (b) with $i=1$; the trace, with the compressed-out directions carrying a $\dlat/\nsamp$ fraction of the residual mass on average (they are uncorrelated with $\bar z$ to leading order), gives (a). Conditionally on the data the rows of $\bar Z$ across features are i.i.d.\ with covariance $\theta a_*^2\bar K$, so $\bar Z^\top\bar Z/\pwidth$ is a sample covariance of $\pwidth$ observations with population $\bar K$: Marchenko--Pastur with edges $(1\pm\sqrt\gamma)^2$ and mean $1$ if $\bar K=I$ \citep{couillet2011rmt}, BBP outliers for the spiked $\bar K$ \citep{baik2005phase}. $\bar K$ is a kernel matrix whose entries are a Hadamard-power series ($\ell\ge3$) in the overlaps, the structure analysed by \citet{ELKaroui2010spectrum}; the cross-cluster entries are $O(\dlat^{-3})$ individually but not negligible collectively at small $\dlat$ because of the rank bound $D_3+D_5+\dots$. Neglected: the $\Umat_{\rm lin}$--$\Umat_{\rm res}$ interaction (uncorrelated, not independent), non-Gaussianity of the residuals for the mixture, and the fluctuations of $\rho_c$ (cross terms $m_c\!\cdot\!\xi^\nu/\dlat\sim\mathcal N(0,s^2\sigsig^2/\dlat^2)$ are a fixed $\approx34\%$ of $\E\rho_c$, so $\E r_c$ exceeds $r_c(\E\rho_c)$ by up to $\approx1.7$ for the cubic term). \emph{Checks} (\cref{tab:buffer-sim}): per-mode mean $\times\nsamp/(\theta a_*^2)=0.60,0.67,0.78,0.88,0.97,1.01,1.03,1.05,1.07$ (sim, $\dlat=5,\dots,480$) and $0.92,0.68,0.78,0.86,0.91,0.93,0.95,0.97,0.98$ (repo, $\dlat=5,10,20,40,60,80,100,150,200$; $0.80$ at $\dlat=200$ if $\theta$ is omitted); in the $\pwidth=\nsamp+\dlat+300$ sweep $g/(1+\sqrt\gamma)^2=6.1,4.8,4.1,2.9,1.2,0.83$ at $\dlat=5,10,20,40,100,200$, the value below $1$ at $\dlat=200$ ($\dlat/\nsamp=0.4$, $\gamma=0.5$) being the interlacing shortfall.
\end{proof}

\subsubsection{The buffer corollary}

\begin{corollary}[The buffer is the ratio $\lambda^{\rm sig}_{\min}/\lambda^{\rm samp}_{\max}$; conditional]\label{cor:buffer-buf}
Under the hypotheses of \cref{prop:taugen-buf,prop:onset-buf},
\begin{equation}
R(\dlat)=\frac{\taumem^{\rm on}}{\taugen}=\frac{\lambda^{\rm sig}_{\min}}{\lambda^{\rm samp}_{\max}}
=\frac{\nsamp}{\dlat}\cdot\frac{a_1(q(\dlat))^2}{\theta(q(\dlat))\,a_*(q(\dlat))^2}\cdot\frac{e^{-2t}\tilde\nu_{\min}(\dlat)+\Delta_t}{g(\dlat)}\,(1+o(1)),
\label{eq:R-buf}
\end{equation}
independent of $\varepsilon$ and of the learning rate. $\dlat$ enters explicitly through the $1/\dlat$ normalisation of the data-side edges (the linear mass $a_1^2q$ is shared among $\dlat$ coordinates by the $\dlat^{-1/2}$ input scaling), through $q(\dlat)$, and through $g$ and $\tilde\nu_{\min}$; it never multiplies an absorption time. $R$ opens with $\dlat$ exactly when $a_1^2(e^{-2t}\tilde\nu_{\min}+\Delta_t)/(\theta a_*^2g)$ grows faster than $\nsamp/\dlat$ shrinks.
\end{corollary}
\begin{proof}
Divide \eqref{eq:tauon-buf} by \eqref{eq:taugen-buf}. With $g\equiv1$, \eqref{eq:R-buf} gives $56,102,212,328,391,362,287,288,284$ at $\dlat=5,\dots,480$ against simulated $11,15,41,101,231,290,262,298,312$: accurate to $\le15\%$ for $\dlat\ge200$, over-predicting by the factor $g$ below.
\end{proof}

\begin{proposition}[Two phases of the buffer; conjectural]\label{prop:phases-buf}
Let $d^\star=c_q/q_\infty$ ($=47$ at $\signoise=0.5$), suppose $(1+\sqrt\gamma)^2$ varies by at most a constant along the sweep, $\dlat<\nsamp$, and $\signoise=0.5$.
(i) \emph{Opening phase} $\dlat\lesssim4d^\star$: $a_1^2/(\theta a_*^2)$ grows from $2.4$ to $163$ over $\dlat=5\to200$, outpacing the $40\times$ fall of $\nsamp/\dlat$ and the $3.5\times$ fall of $(1+\sqrt\gamma)^2$, and $g$ collapses from $\approx5$--$6.5$ to $\approx1$; $R$ increases.
(ii) \emph{Saturated phase} $\dlat\gg d^\star$: $a_1,a_*,\theta\to$ their $q_\infty$ values and
\begin{equation}
R(\dlat)\simeq\frac{\nsamp}{\dlat}\cdot\frac{a_1(q_\infty)^2\,(e^{-2t}\tilde\nu_{\min}(\dlat)+\Delta_t)}{\theta(q_\infty)a_*(q_\infty)^2\,g(\dlat)}
\label{eq:R-saturated-buf}
\end{equation}
is predicted to decrease, roughly like $\nsamp/\dlat$ tempered by the BBP growth of $\tilde\nu_{\min}$ and the interlacing decrease of $g$. Hence $R$ is \emph{conjectured} non-monotone with a broad maximum at $\dlat=\Theta(d^\star)$--$\Theta(\nsamp)$. This is not demonstrated: simulated $R=290,262,298,312$ at $\dlat=200,300,400,480$ is flat, and the repo's spectra with a non-empty sample block stop at $\dlat=200$ with $R$ still rising ($341,419,428$ at $\dlat=100,150,200$). What is observed to close is $R_{\rm full}$: as $\dlat\to\nsamp$ the noise-dimension lower edge $a_1^2[e^{-2t}\signoise^2(1-\sqrt{\psi_\perp})^2+\Delta_t]/\dlat\to a_1^2\Delta_t/\dlat$ approaches $\lambda^{\rm samp}_{\max}$ ($R_{\rm full}=3.1,6.3,12.5,17.2,11.3,6.1,3.7,2.4$ at $\dlat=10,\dots,480$, sim; $4.4\to18.8\to12.5$ over $\dlat=10\to80\to200$, repo) while $|I_{\rm samp}|=\nsamp-\dlat\to0$. The empirical return of memorization at large $\dlat$ (an MLP sweep, \cref{app:theory-scope}) is at most a candidate consequence of this; it is not explained by $R$, which stays $\approx300$.
\end{proposition}
\begin{proof}[Proof sketch]
(ii) is substitution of $q=q_\infty+O(1/\dlat)$ into \eqref{eq:R-buf}; with $g\equiv1$ it gives $391,362,287,288,284$ at $\dlat=100,\dots,480$. (i): numbers from \cref{lem:coeffs-buf} and \cref{tab:buffer-sim}. The closing statement uses the MP lower edge of (H4) (verified to $3$--$10\%$ for $\dlat\ge20$) and the mode count of \cref{lem:fourbulk-counts}.
\end{proof}

\begin{corollary}[Fixed-energy sweep; sketched]\label{cor:fixedenergy-buf}
For the headline control ($\pwidth=1800$ fixed, $\signoise^2=E_\perp/(\dlat-\dint)$, $E_\perp=3.75$), $q(\dlat)=\Delta_t+e^{-2t}(\dint\alpha_0^2+E_\perp)/\dlat$: $q_\infty=\Delta_t=0.0198$, $c_q=17.4$, $d^\star=878>\nsamp$. The whole feasible sweep is in the opening phase, $\theta$ falls to $0.54$ at $\dlat=200$, block separation holds (the noise-dimension block sits at $a_1^2[e^{-2t}E_\perp(1\pm\sqrt{\psi_\perp})^2/(\dlat-\dint)+\Delta_t]/\dlat\gg\theta a_*^2g/\nsamp$), and \eqref{eq:R-buf} predicts a monotonically increasing $R$ over $\dlat\le200$, consistent with the monotone growth reported for that sweep ($R=16.7\to889$ over $\dlat=5\to200$, recomputed from the repository's spectra in \texttt{clean figures/p\_fixed\_energy\_fixed}, $\pwidth=1800$, $\nsamp=500$).
\end{corollary}

\begin{table}[h]
\centering\scriptsize
\caption{Own simulation of \eqref{eq:U} via its Gram ($t=0.01$, $\signoise=0.5$, $\nsamp=500$, $k=10$, $s=3$, $\pwidth=64\dlat$, $N_{\rm mc}=4$, seed $0$; theory normalisation, arbitrary time units). ``pred.'' is \eqref{eq:taugen-buf} with $\tilde\nu_{\min}$ from the same data; $1/\lambda^{\rm MP}_+$ uses $\theta a_*^2(1+\sqrt\gamma)^2/\nsamp$; $g$ here is $\lambda^{\rm samp}_{\max}\nsamp/(\theta a_*^2)/(1+\sqrt\gamma)^2$; ``ov.'' is the squared overlap of the top sample eigenvector with the cluster indicators; $R_{g=1}$ is \eqref{eq:R-buf} with $g=1$. Repo values ($N_{\rm mc}=50$, $\dlat=5,10,20,40,60,80,100,150,200$): $1/\lambda^{\rm sig}_{\min}=16,20,31,48,63,81,95,137,173$; $1/\lambda^{\rm samp}_{\max}=333,684,1927,6866,13457,22778,32303,57587,73819$; $g=3.9,4.9,4.5,3.1,2.6,2.0,1.7,1.3,1.23$; $R=21,34,63,144,212,280,341,419,428$; $R_{\rm full}=-,4.4,7.8,13.0,15.6,18.8,18.1,18.0,12.5$.}
\label{tab:buffer-sim}
\resizebox{\textwidth}{!}{%
\begin{tabular}{r|cc|cc|ccc|cc|cc|c}
\toprule
$\dlat$ & $q$ & $\theta a_*^2\,(\times10^3)$ & $1/\lambda^{\rm sig}_{\min}$ & pred. & $1/\lambda^{\rm samp}_{\max}$ & $1/\lambda^{\rm MP}_+$ & $g$ & ov. & mean$\cdot\nsamp/\theta a_*^2$ & $R$ & $R_{g=1}$ & $R_{\rm full}$\\ \midrule
10 & 1.515 & 42.5 & 33.0 & 32.4 & 506 & 3315 & 6.55 & 0.29 & 0.67 & 15 & 102 & 3.1\\
20 & 0.890 & 21.7 & 41.1 & 41.1 & 1700 & 8708 & 5.12 & 0.29 & 0.78 & 41 & 212 & 6.3\\
50 & 0.515 & 9.14 & 86.7 & 85.6 & 8743 & 28096 & 3.21 & 0.17 & 0.88 & 101 & 328 & 12.5\\
100 & 0.390 & 5.45 & 145.6 & 143.2 & 33629 & 56016 & 1.67 & 0.16 & 0.97 & 231 & 391 & 17.2\\
200 & 0.327 & 3.83 & 250.7 & 251.5 & 72657 & 90985 & 1.25 & 0.03 & 1.01 & 290 & 362 & 11.3\\
300 & 0.307 & 3.33 & 384.4 & 387.3 & 100699 & 111143 & 1.10 & 0.01 & 1.03 & 262 & 287 & 6.1\\
400 & 0.296 & 3.10 & 431.1 & 431.7 & 128379 & 124287 & 0.97 & 0.02 & 1.05 & 298 & 288 & 3.7\\
480 & 0.291 & 2.98 & 464.2 & 465.3 & 145021 & 131962 & 0.91 & 0.01 & 1.07 & 312 & 284 & 2.4\\
\bottomrule
\end{tabular}}
\end{table}

\subsubsection{Remarks}

\begin{remark}[Why the count-times-timescale bound is a non-sequitur]\label{rem:nonsequitur-buf}
An earlier draft asserted $\taumem-\taugen\ge(\dlat-\dint)/\lambda^{\rm noise}_{\min}$, as if the readout had to traverse each noise-dimension mode in turn. Under $\dot a_i=-\lambda_i(a_i-a_i^\star)$ all modes decay simultaneously; the time at which sample mode $j$ reaches tolerance $\varepsilon$ is $\lambda_j^{-1}\log(1/\varepsilon)$ whatever the number or rates of the other modes. Adding $\dlat-\dint$ intermediate modes changes $\taumem^{\rm on}$ only insofar as it changes $\lambda^{\rm samp}_{\max}$ (through $q(\dlat)$, $\theta$, $g$) and $\taugen$ through the $1/\dlat$ normalisation and $\tilde\nu_{\min}$. A mode count can never multiply a timescale in a diagonal linear system, and the asserted bound is violated at every $\dlat\ge20$ on the repo's own spectra. The correct statement is the ratio \eqref{eq:R-buf}.
\end{remark}

\begin{remark}[Onset versus median]\label{rem:median-buf}
The sample block has $\nsamp-\dlat$ modes of mean $\theta a_*^2/\nsamp$. At small $\dlat$ the residual Gram is rank-deficient ($D_3<\nsamp$ for $\dlat\le13$) and spiked, so most of the mass sits in a few top modes and a heavy lower tail: $\mathrm{median}\times\nsamp/(\theta a_*^2)=0.037,0.072,0.28,0.57,0.72,0.80,0.84,0.92,0.95$ at $\dlat=5,\dots,200$ on the repo ($\lambda^{\rm samp}_{\max}/\mathrm{median}\approx540$ at $\dlat=5$, $\approx2$ at $\dlat=200$). As $\dlat$ grows the block becomes a proper Marchenko--Pastur bulk and the median rises toward the mean, while $\lambda^{\rm sig}_{\min}\propto a_1^2\tilde\nu_{\min}/\dlat$ falls. Therefore $\lambda^{\rm sig}_{\min}/\mathrm{median}(\lambda^{\rm samp})$ shrinks with $\dlat$ (repo: $11350\to798$, $14\times$; sim: $5750\to540$ over $\dlat=10\to200$) while $R$ grows (repo: $21\to428$). The buffer is a claim about the \emph{onset} of sample-specific learning; the time to absorb the bulk of the sample modes, which governs the memorization fraction, is a different quantity (\cref{app:theory-bridge,app:theory-predictor}).
\end{remark}

\begin{remark}[The isotropic control]\label{rem:isotropic-buf}
Set $\signoise=\sigsig$. \emph{Structure:} the $\dlat-\dint$ noise-dimension modes still exist, at $a_1^2[e^{-2t}\sigsig^2(1\pm\sqrt{\psi_\perp})^2+\Delta_t]/\dlat$ (verified to $3\%$), between the $\dint$ signal spikes and the sample block; what isotropy removes is the factor $\sigsig^2/\signoise^2$ separating the signal and noise-dimension blocks, so the data side becomes a single Marchenko--Pastur bulk with $\dint$ BBP spikes, the $\rho_2$ of \citet{bonnaire2025} with outliers: the \emph{four-bulk structure} is absent, not the noise-dimension modes (\cref{rem:fourbulk-isotropic}). \emph{Mechanism:} de-saturation is essentially absent ($c_q=8.8$, $q_\infty=1$, $d^\star=8.8$; $a_1^2/(\theta a_*^2)$ rises only $2.4\to13$ over $\dlat=5\to200$), and \eqref{eq:R-buf} with $g\equiv1$ gives $R=56,80,95,89,74,58,50$ at $\dlat=5,10,20,50,100,200,300$, a decline beyond $\dlat\approx20$ driven by $\nsamp/\dlat$ and only partly offset by the strong BBP lift of the signal spikes ($\tilde\nu_{\min}/\nu_{\min}=1.15,1.12,1.52,1.77,2.20,2.93$ at $\dlat=10,\dots,300$). The simulated $R$ is $11,21,49,83,71,60,55$: it rises $7\times$ to $\dlat\approx50$ because $g$ collapses ($5.2,4.0,1.9,1.1,1.0,1.0,0.9$), then follows the predicted decline. Under isotropy $R$ is therefore predicted non-monotone with a peak near $\dlat\approx50$; the earlier claim that de-saturation still makes $R$ grow is withdrawn. A direct simulation of \eqref{eq:U} at $\signoise=\sigsig=1$ ($\nsamp=500$, $\pwidth=4096$, $t=0.01$, three seeds; \texttt{\_next\_steps/scripts/e5\_isotropic.py}) gives $\lambda^{\mathcal S}_{\min}/\lambda^{\mathcal M}_{\max}=14,30,58,87,77,60$ at $\dlat=5,10,20,40,80,160$: a rise to a peak near $\dlat\approx40$ followed by the predicted decline. (An earlier ``$10.7\times$ growth of the gap'' was a difference of code-unit eigenvalues, which carry the factor $\pwidth=64\dlat$; in theory units the same four values are $1.50,2.22,2.45,1.61$ at $\dlat=10,20,40,100$, rising then falling.) With no low-dimensional structure at all ($x\sim\mathcal N(0,I_{\dlat})$, same settings) the data-block/sample-block ratio falls monotonically from $77$ to $6$ over $\dlat=5\to160$.
\end{remark}

\begin{remark}[Relation to prior work]\label{rem:prior-buf}
Theorem~3.1 of \citet{bonnaire2025} is stated for an arbitrary spectral measure of $\Sigmadata$; the four-block spectrum used in (H4) is that theorem evaluated at a two-atom measure, and \citet{george2025dsm} already give a linear pencil with block dimensions $(\pwidth,\dlat-\dint,\dint,\nsamp)$ for data on a $\dint$-dimensional subspace with $\signoise=0$. At fixed $\dlat$, \eqref{eq:tauon-buf} gives $\taumem^{\rm on}\propto\nsamp/(\theta a_*^2g)$, consistent with the linear-in-$\nsamp$ memorization time of \citet{bonnaire2025}; the $\dlat$-dependence through $q(\dlat)$, $\theta$, $\tilde\nu_{\min}$ and $g$, the ratio form \eqref{eq:R-buf}, and the training-time reading for $\signoise>0$ are the new content.
\end{remark}

\begin{remark}[What this subsection does not predict]\label{rem:notpredicted-buf}
$R$ concerns two spectral edges under gradient flow at fixed $t$. It does not by itself predict (a) the memorization fraction at a fixed step budget $T$, which depends on how many sample modes have $\lambda_i^{-1}\log(1/\varepsilon)<T$, i.e.\ on the within-block spread of the sample spectrum, on the map from gradient-flow time to optimizer steps (\cref{rem:clock-p1}), and on the drift of the NN-ratio statistic with ambient dimension; (b) whether generation quality is matched across $\dlat$ at the comparison time, since $\taugen$ grows $\approx11\times$ over $\dlat=5\to200$ and concerns the signal subspace only; (c) the MLP and real-VAE settings, whose spectra are continuous rather than two-atom and where $a_1,a_*,\theta$ must be replaced by learned-feature analogues (\cref{app:theory-general}); (d) anything at $\signoise=0.01$ beyond $\taugen$, since block separation fails there (H5).
\end{remark}

\paragraph{Predictions and how to test them.}
\begin{enumerate}
\item \emph{$\taugen\,a_1(q(\dlat))^2(e^{-2t}\tilde\nu_{\min}(\dlat)+\Delta_t)$ is linear in $\dlat$ through the origin with slope $\log(1/\varepsilon)$, independent of $\dlat$, $\signoise$ and $\pwidth$.} Test: from the saved spectra extract $\lambda_{\dint}$, rebuild the data with the saved seed to get $\tilde\nu_{\min}$, and plot $\dlat/(\lambda_{\dint}a_1(q)^2(e^{-2t}\tilde\nu_{\min}+\Delta_t))$ against $\dlat$; both width rules and both $\signoise$ should collapse onto one line. Without the BBP shift the collapse fails at the $10$--$25\%$ level for $\dlat\ge200$.
\item \emph{$\taugen$ grows linearly only for $d^\star\ll\dlat\ll3.6\nsamp$; beyond, the BBP lift makes it sublinear.} Test: $\taugen(1000)/\taugen(500)$ on the existing spectra should be $\approx1.6$, not $2$ (observed $1.28$ at $\pwidth=64\dlat$, $1.59$ at $\pwidth=\nsamp+\dlat+300$; the residual $20\%$ is unexplained).
\item \emph{The sample-block per-mode mean equals $\theta(q)a_*(q)^2/\nsamp$ with $\theta\approx(1-\Delta_t/q)^3$.} Test: $\operatorname{mean}(\lambda_{\dlat+1..\nsamp})\nsamp/a_*(q)^2$ should equal $\theta(q)=0.97,0.95,0.93,0.88,0.85,0.82$ at $\dlat=5,10,20,50,100,200$; at $\signoise=0.01$ the factor is $0.47$ at $\dlat=200$.
\item \emph{$\lambda^{\rm samp}_{\max}$ lies between $\mu_{\dlat+1}$ and $\mu_1$ of an $\nsamp$-mode Wishart of scale $\theta a_*^2/\nsamp$; for $\dlat\gtrsim150$ it is within $\sim25\%$ of the MP edge and falls below it when $\dlat/\nsamp$ and $\nsamp/\pwidth$ are both $O(1)$; at small $\dlat$ it exceeds the edge by $4$--$6.5\times$, of which only $\sim25\%$ is cluster-coherent.} Test: rebuild $\Umat$ with eigenvectors, compute the squared overlap of the $(\dlat+1)$-th eigenvector with the $k$ cluster indicators.
\item \emph{$\taumem^{\rm on}$ grows by more than two orders of magnitude over $\dlat=5..200$: $\sim20\times$ from $\theta a_*^2$, $\sim3.5\times$ from $(1+\sqrt\gamma)^2$ at $\pwidth=64\dlat$, $\sim3$--$4\times$ from the collapse of $g$.} Test: tabulate $1/\lambda_{\dlat+1}$ against $\nsamp/(\theta a_*^2(1+\sqrt{\nsamp/\pwidth})^2)$ and read off $g$.
\item \emph{$R$ is conjectured non-monotone at $\signoise=0.5$ with a broad maximum between $d^\star\sim50$ and $\sim\nsamp$.} Test: extend the sweep to $\dlat\in\{300,400,480\}$ with $N_{\rm mc}=50$ and several seeds, and to $\nsamp\in\{1000,2000\}$; the peak should scale with $\nsamp$.
\item \emph{$R_{\rm full}$ is non-monotone with a peak near $\dlat\sim60$--$100$ and closes to $O(1)$ as $\dlat\to\nsamp$; this, with the vanishing count $\nsamp-\dlat$, is the candidate mechanism for the return of memorization.} Test: $\lambda_{\dlat}/\lambda_{\dlat+1}$ on new $\dlat\in\{300,400,480\}$ runs; then test whether the RFNN memorization fraction returns where $R_{\rm full}\lesssim3$.
\item \emph{$R$ is independent of $\varepsilon$ and of the learning rate.} Test: re-extract at two plateau thresholds and two learning rates; the ratio is unchanged while the individual times scale.
\item \emph{The median-based ratio decreases with $\dlat$ (opposite to the onset ratio).} Verified on the repo: $\mathrm{median}\cdot\nsamp/(\theta a_*^2)=0.037\to0.95$ and $R_{\rm median}$ shrinks $14\times$.
\item \emph{Isotropic control: $R$ rises $\sim6\times$ to a peak near $\dlat\approx40$ and then declines like $\nsamp/\dlat$; the noise-dimension modes exist but are not separated from the signal.} Status: the rise-then-decline is confirmed by the three-seed simulation in \cref{rem:isotropic-buf} ($14\to87\to60$). Remaining test: trained runs at $\signoise=1$, $\dlat\in\{5,\dots,300\}$; check $\tilde\nu_{\min}/\nu_{\min}=1.5,1.8,2.2$ at $\dlat=50,100,200$.
\item \emph{Changing $\pwidth$ at fixed $\dlat$ leaves $\taugen$ unchanged (to $\sim2\%$) and changes $\taumem^{\rm on}$ by at most the MP-edge factor $[(1+\sqrt{\nsamp/\pwidth_2})/(1+\sqrt{\nsamp/\pwidth_1})]^2$.} Repo: ratio $0.95,1.01,1.09,1.26,1.39,1.31$ at $\dlat=5,\dots,200$ vs MP factor $0.63,\dots,2.04$; $\taugen$ differs by $1$--$7\%$.
\item \emph{At $\signoise=0.01$ block separation fails, so $\taumem^{\rm on}$ and $R$ are undefined while $\taugen$ is still predicted.} Repo: $\lambda_{\dlat}/\lambda_{\dlat+1}=1.0$--$1.15$; $1/\lambda_5$ grows $2.4\times$ over $\dlat=5..40$ vs predicted $2.4\times$.
\end{enumerate}

\paragraph{Open gaps.}
(1) $g(\dlat)$ at small $\dlat$ has no derivation: the cluster-spike model with the full residual kernel accounts for $\sim$half ($2.5$--$2.9\times$) and the top eigenvector is only $\sim25\%$ cluster-coherent; the remainder is attributed to the finite tensor rank $D_3$ of the cubic residual features without a controlled derivation, so \cref{prop:onset-buf} is sketched with the small-$\dlat$ formula explicitly open. (2) The turnover of $R$ (\cref{prop:phases-buf}) is conjectural: the $N_{\rm mc}=4$ simulation is flat over $\dlat=200..480$ and the repo has no spectra with a non-empty sample block beyond $\dlat=200$; only $R_{\rm full}$ is observed to close, and the link to the return of memorization is untested in the RFNN. (4) The machine-precision cliff at index $\nsamp$ does not hold for the finite-$N_{\rm mc}$ $\Umat$ ($\lambda_{\nsamp+1}/\lambda_\nsamp=0.48$--$0.76$ for $\dlat\ge40$); why the Monte-Carlo excess modes sit only $\sim2\times$ below the sample-block bottom rather than $\sim100\times$ is not understood. (5) The rise-then-decline of $R$ under isotropy is confirmed by simulation (\cref{rem:isotropic-buf}); no trained $\signoise=1$ run exists in the repository. (6) At $\signoise=0.01$ block separation fails outright; the fixed-energy control's $R$ numbers are recomputed ($16.7\to889$) but its mechanism is not analysed here. (7) The $\pwidth$-dependence of the onset is a one-sided interlacing bound, partially falsified as an equality ($1.39$ vs $1.86$ at $\dlat=100$). (8) All $o(1)$ statements are proportional-asymptotic with $O(\dlat^{-1/2})$ corrections; nothing is uniform in the sweep. (9) Monotonicity of $a_*^2$ and $\theta a_*^2$ in $q$ is numerical. (10) The BBP shift (H4) is verified to $1$--$4\%$ but is an input from spiked-covariance theory, not derived within the RFNN pencil; in the isotropic case the closed-form factor diverges near threshold and the empirical top-$\dint$ eigenvalues of $\hat C+\hat S$ must be used directly. (11) Main-text statements asserting $\dlat$-independent $\taugen$, sequential traversal of the four bulks, ``no buffer in the isotropic case'', or sample indices $[\dlat+1,\dlat+\nsamp]$ are inconsistent with this subsection and must be revised.

% ----------------------------------------------------------------------
\subsection{The exactly solvable linear score model}
\label{app:theory-linear}

This subsection isolates the part of the problem that is exactly solvable: a linear score model on the block-anisotropic mixture. Throughout, $t>0$ is fixed and expectations are over the noised population distribution $p_t$ unless a subscript $n$ indicates the empirical one. Two values of $t$ occur: the MLP is evaluated at $t=0.1$, while the RFNN is trained \emph{and} evaluated at $t_{\rm fixed}=0.01$. Statements are basis-free; proofs are written in the rotated frame $Q^{\top}x$ where the signal subspace is coordinate-aligned. \emph{Convention.} The risks here carry no factor $\tfrac12$, so the flow is $\dot B=-2(BM_t+I)$ and parameters decay as $e^{-2\lambda T}$; in the convention of \cref{lem:gf-p1} (risk with $\tfrac12$) the same trajectory is $e^{-\lambda T'}$ with $T'=2T$. Ratios of clocks are convention-free.

\begin{definition}[Setting for the linear score model]\label{def:lin-setting}
Fix $t>0$ and write $x_t=(x_{t,\mathsf S},x_{t,\mathsf N})$ in the rotated frame $Q^{\top}x_t$, with $x_{t,\mathsf S}\in\R^{\dint}$ (signal block) and $x_{t,\mathsf N}\in\R^{\dlat-\dint}$ (null block). Under the data model $x=m_c+\xi$, $c\sim\mathrm{Unif}[k]$, $m_c\in\mathrm{span}(\mathsf S)$, $\xi\sim\mathcal N(0,\diag(\sigsig^2 I_{\dint},\signoise^2 I_{\dlat-\dint}))$, and $x_t=e^{-t}x+\sqrt{\Delta_t}\,\eta$, the noised marginal is the Gaussian mixture
\begin{equation}\label{eq:lin-pt}
p_t=\tfrac1k\textstyle\sum_{c}\mathcal N\big(e^{-t}m_c,\,\Sigma_t\big),\qquad \Sigma_t=\diag\big(\gamma_t^2 I_{\dint},\,\beta_t^2 I_{\dlat-\dint}\big),
\end{equation}
with $\gamma_t^2=e^{-2t}\sigsig^2+\Delta_t$ and $\beta_t^2=e^{-2t}\signoise^2+\Delta_t$. Its \emph{uncentered} second moment is $M_t=\E[x_tx_t^{\top}]=\diag(M_{t,\mathsf S},\beta_t^2 I_{\dlat-\dint})$, $M_{t,\mathsf S}=e^{-2t}(C+\sigsig^2 I)+\Delta_t I$, $C=\tfrac1k\sum_c m_cm_c^{\top}$, with eigenvalues $\lambda^{\mathsf S}_1\ge\dots\ge\lambda^{\mathsf S}_{\dint}\ge\gamma_t^2$; they all equal $\alpha_t^2=e^{-2t}(s^2/\dint+\sigsig^2)+\Delta_t$ only for isotropic centers $C=(s^2/\dint)I$, and in general $\alpha_t^2$ is their trace average (exact because $\operatorname{tr}C=s^2$). $M_t$ is the second moment, not the covariance, of $x_t$; the two differ on the signal block by $e^{-2t}\bar m\bar m^{\top}$, $\bar m=\frac1k\sum_cm_c$ ($\|\bar m\|=1.15$ for the repo's centers), and it is $M_t$ that enters the dynamics. Given a training set $\{x^\mu\}_{\mu\le\nsamp}$, the \emph{empirical} second moment is $\widehat M_t=e^{-2t}\frac1n\sum_\mu x^\mu x^{\mu\top}+\Delta_tI$. The \emph{linear score model} is $s_B(x)=Bx$, $B\in\R^{\dlat\times\dlat}$, trained by gradient flow $\dot B=-\nabla_BL$ from $B(0)=B_0$ on either the population score-matching risk $L_t(B)=\E_{p_t}\|Bx_t-s^\star(x_t)\|^2$, $s^\star=\nabla\log p_t$, or the empirical denoising risk $L_{t,n}(B)=\frac1n\sum_\mu\E_\eta\|Bx^\mu_t+\eta/\sqrt{\Delta_t}\|^2$. The population denoising objective equals $L_t$ up to a $B$-independent constant \citep{vincent2011connection}, so it generates the same flow; $L_{t,n}$ with $\E_\eta$ is the mean path of full-batch SGD with a fresh noise draw per step (\cref{prop:sgdmean-p1}), which is what the code runs (momentum $0$).
\end{definition}

\begin{table}[h]
\centering\small
\caption{Constants of the linear model at the two evaluation times used by the code ($\sigsig=1$, $s=3$, $\dint=5$, $k=10$, centers of seed 42: $\operatorname{eig}C=\{0.907,1.033,1.203,2.491,3.367\}$). $\mathcal J_{t,\mathsf S}$ and $R_{\mathrm{irr}}$ are $4\times10^5$-sample Monte-Carlo values.}\label{tab:lin-constants}
\begin{tabular}{lcc}
\toprule
 & MLP, $t=0.1$ & RFNN, $t=0.01$\\
\midrule
$\Delta_t$, $e^{-2t}$ & $0.1813$, $0.8187$ & $0.0198$, $0.9802$\\
$\gamma_t^2$ & $1.000$ & $1.000$\\
$\lambda^{\mathsf S}_i$ & $1.742,1.846,1.985,3.039,3.756$ & $1.889,2.013,2.179,3.441,4.300$\\
$\alpha_t^2$ (trace average) & $2.474$ & $2.764$\\
$\operatorname{tr}M_{t,\mathsf S}^{-1}$, $\mathcal J_{t,\mathsf S}$, $R_{\mathrm{irr}}$ & $2.215$, $3.582$, $1.367$ & $2.008$, $3.686$, $1.678$\\
$\beta_t^2$, $1/\beta_t^2$, $1/(4\beta_t^2)$ at $\signoise=0.5$ & $0.386$, $2.59$, $0.648$ & $0.2649$, $3.78$, $0.944$\\
$\beta_t^2$, $1/\beta_t^2$, $1/(4\beta_t^2)$ at $\signoise=0.01$ & $0.1814$, $5.51$, $1.379$ & $0.0199$, $50.3$, $12.56$\\
\bottomrule
\end{tabular}
\end{table}

\begin{lemma}[The null block of the score is exactly linear; the best linear score is $-M_t^{-1}$; proved]\label{lem:null-exact-lin}
Under \cref{def:lin-setting}:
\begin{enumerate}
\item[(a)] (\emph{exact null block}) For every $x_t$ and every $t>0$,
\begin{equation}\label{eq:lin-null-score}
s^\star_{\mathsf N}(x_t)=-\,x_{t,\mathsf N}/\beta_t^2 ,
\end{equation}
while $s^\star_{\mathsf S}(x_t)=-\gamma_t^{-2}\big(x_{t,\mathsf S}-e^{-t}\sum_c w_c(x_{t,\mathsf S})\,m_c\big)$ with softmax weights $w_c\propto\exp(-\|x_{t,\mathsf S}-e^{-t}m_c\|^2/2\gamma_t^2)$; this is affine if all centers coincide and linear only if all $m_c=0$. The weights depend on $x_t$ only through $x_{t,\mathsf S}$.
\item[(b)] (\emph{Stein identity}) $\E_{p_t}[s^\star(x_t)\,x_t^{\top}]=-I_{\dlat}$.
\item[(c)] (\emph{best linear score}) $L_t$ is a strictly convex quadratic with unique minimizer $B^\star_t=-M_t^{-1}$, which is block diagonal; on the null block $B^\star_{t,\mathsf N}=-\beta_t^{-2}I$ reproduces \eqref{eq:lin-null-score} exactly. The irreducible risk is supported on the signal block and equals
\begin{equation}\label{eq:lin-irr}
R_{\mathrm{irr}}(t):=\min_B L_t(B)=\E\|s^\star_{\mathsf S}\|^2-\operatorname{tr}M_{t,\mathsf S}^{-1}\;\in\;\Big[0,\ \frac{\dint}{\gamma_t^2}-\sum_{i=1}^{\dint}\frac1{\lambda^{\mathsf S}_i}\Big],
\end{equation}
with $R_{\mathrm{irr}}=0$ iff $p_t=\mathcal N(0,M_t)$ (a centered Gaussian). For the repo's centers (\cref{tab:lin-constants}) a Monte-Carlo evaluation gives $R_{\mathrm{irr}}=1.37$ at $t=0.1$ and $1.68$ at $t=0.01$, independent of $\dlat$.
\end{enumerate}
\end{lemma}
\begin{proof}
(a) Because $m_{c,\mathsf N}=0$ and $\Sigma_t$ is block diagonal, each mixture component factorizes, $\mathcal N(x;e^{-t}m_c,\Sigma_t)=\mathcal N(x_{\mathsf S};e^{-t}m_c,\gamma_t^2 I)\,\mathcal N(x_{\mathsf N};0,\beta_t^2 I)$, and the second factor is common to all $c$. Hence $p_t(x)=p_{t,\mathsf S}(x_{\mathsf S})\,\mathcal N(x_{\mathsf N};0,\beta_t^2I)$ and $\nabla_{\mathsf N}\log p_t=-x_{\mathsf N}/\beta_t^2$. The signal-block formula is the gradient of $\log$ of the $\dint$-dimensional mixture $p_{t,\mathsf S}$; with a single center $m$ it is $-\gamma_t^{-2}(x_{\mathsf S}-e^{-t}m)$, affine with nonzero intercept unless $m=0$.
(b) $\E[s^\star x^{\top}]=\int\nabla p_t(x)\,x^{\top}dx=-\int p_t(x)\,\nabla(x^{\top})\,dx=-I$, integrating by parts coordinate-wise; the boundary term $p_t(x)x^{\top}$ vanishes at infinity because $p_t$ has Gaussian tails.
(c) Expanding and using (b),
\[
L_t(B)=\operatorname{tr}(BM_tB^{\top})-2\operatorname{tr}\big(B\,\E[x_ts^{\star\top}]\big)+\E\|s^\star\|^2=\operatorname{tr}(BM_tB^{\top})+2\operatorname{tr}B+\mathcal J_t,
\]
where $\mathcal J_t:=\E\|s^\star\|^2=\operatorname{tr}J(p_t)$ is the Fisher information of $p_t$. Since $M_t\succeq\Delta_tI\succ0$, $L_t$ is strictly convex with $\nabla_BL_t=2BM_t+2I$, so $B^\star_t=-M_t^{-1}$ and $L_t(B^\star_t)=\mathcal J_t-\operatorname{tr}M_t^{-1}$. Independence of $x_{t,\mathsf S}$ and $x_{t,\mathsf N}$ together with $\E x_{t,\mathsf N}=0$ makes $M_t$ block diagonal, hence so is $B^\star_t$; on the null block $B^\star_{t,\mathsf N}x_{t,\mathsf N}=-x_{t,\mathsf N}/\beta_t^2=s^\star_{\mathsf N}$ by (a), so the null block contributes $0$ and $\mathcal J_t-\operatorname{tr}M_t^{-1}=\E\|s^\star_{\mathsf S}\|^2-\operatorname{tr}M_{t,\mathsf S}^{-1}$. Nonnegativity is $L_t\ge0$, with equality iff $s^\star(x)=-M_t^{-1}x$ a.e.\ iff $p_t=\mathcal N(0,M_t)$. For the upper bound, write $p_{t,\mathsf S}=\sum_cw_cq_c$ with $q_c=\mathcal N(e^{-t}m_c,\gamma_t^2I)$; by joint convexity of $(a,b)\mapsto\|a\|^2/b$ on $\R^{\dint}\times\R_{>0}$, $\|\nabla p_{t,\mathsf S}\|^2/p_{t,\mathsf S}\le\sum_cw_c\|\nabla q_c\|^2/q_c$; integrating, $\E\|s^\star_{\mathsf S}\|^2\le\sum_cw_c\operatorname{tr}J(q_c)=\dint/\gamma_t^2$. The Monte-Carlo values use $4\times10^5$ samples of $p_{t,\mathsf S}$ with the seed-42 centers and the exact softmax score.
\end{proof}

\begin{proposition}[Closed-form gradient flow and per-mode risk; proved]\label{prop:lin-gf-closed}
Let $M_t=\sum_{i=1}^{\dlat}\lambda_iv_iv_i^{\top}$ be the spectral decomposition. Under \cref{def:lin-setting}, for every $B_0$ and $T\ge0$,
\begin{equation}\label{eq:lin-closed-form}
B(T)=-M_t^{-1}+\big(B_0+M_t^{-1}\big)\,e^{-2M_tT},
\end{equation}
and for every $B$, $L_t(B)=L_t(B^\star_t)+\operatorname{tr}\big((B-B^\star_t)M_t(B-B^\star_t)^{\top}\big)$, so the excess risk $\mathcal E(T):=L_t(B(T))-L_t(B^\star_t)$ decomposes over the eigenmodes of $M_t$ as
\begin{equation}\label{eq:lin-excess}
\mathcal E(T)=\sum_{i=1}^{\dlat}\lambda_i\,e^{-4\lambda_iT}\,\big\|(B_0+M_t^{-1})v_i\big\|^2\quad\xrightarrow{\;B_0=0\;}\quad \mathcal E(T)=\sum_{i=1}^{\dlat}\frac{e^{-4\lambda_iT}}{\lambda_i}.
\end{equation}
Consequently, for $B_0=0$: (i) the risk carried by mode $i$ decays with the \emph{mode clock} $\tau_i=1/(4\lambda_i)$ (the parameter error on that mode decays with $1/(2\lambda_i)$; all ratios of clocks are the same); (ii) the risk initially carried by mode $i$ is its target energy $\E\langle v_i,B^\star_tx_t\rangle^2=1/\lambda_i$, so the slowest modes are also the largest; (iii) if $B_0$ is block diagonal (in particular $B_0=0$) then so is $B(T)$ for all $T$: the signal and null blocks never interact during population training.
\end{proposition}
\begin{proof}
From the proof of \cref{lem:null-exact-lin}(c), $\dot B=-2(BM_t+I)$. Set $E(T)=B(T)+M_t^{-1}$; then $\dot E=-2EM_t$, with unique solution $E(T)=E(0)e^{-2M_tT}$, which is \eqref{eq:lin-closed-form}. Substituting $B=B^\star_t+E$ into $L_t(B)=\operatorname{tr}(BM_tB^{\top})+2\operatorname{tr}B+\mathcal J_t$ gives $L_t(B)=L_t(B^\star_t)+2\operatorname{tr}(EM_tB^{\star\top}_t)+2\operatorname{tr}E+\operatorname{tr}(EM_tE^{\top})=L_t(B^\star_t)+\operatorname{tr}(EM_tE^{\top})$, since $B^\star_t=-M_t^{-1}$ makes $2\operatorname{tr}(EM_tB^{\star\top}_t)=-2\operatorname{tr}E$. Hence $\mathcal E(T)=\operatorname{tr}\big(E(0)e^{-2M_tT}M_te^{-2M_tT}E(0)^{\top}\big)=\sum_i\lambda_ie^{-4\lambda_iT}\|E(0)v_i\|^2$. With $B_0=0$, $E(0)v_i=v_i/\lambda_i$, giving \eqref{eq:lin-excess}. (i) is read off from $e^{-4\lambda_iT}$ versus $E(T)v_i=e^{-2\lambda_iT}E(0)v_i$. (ii): $\E\langle v_i,B^\star_tx_t\rangle^2=\lambda_i^{-2}v_i^{\top}M_tv_i=1/\lambda_i$. (iii): $M_t$ and $M_t^{-1}$ are block diagonal, so \eqref{eq:lin-closed-form} preserves block diagonality. (Checked against an Euler integration at $\dlat=6$: $\max|B_{\rm Euler}-B_{\rm closed}|=3\times10^{-4}$ at $dt=10^{-3}$.)
\end{proof}

\begin{remark}[Discrete steps, loss normalization, and the code's clocks]\label{rem:lin-gd}
Gradient descent with step $\eta$ on $c_{\rm loss}\cdot L$ (the code multiplies the risk by a normalization $c_{\rm loss}$) gives $B_k+M_t^{-1}=(B_0+M_t^{-1})(I-2\eta_{\rm eff}M_t)^k$ with $\eta_{\rm eff}=\eta c_{\rm loss}$, i.e.\ \eqref{eq:lin-excess} with $e^{-4\lambda_iT}\to(1-2\eta_{\rm eff}\lambda_i)^{2k}$; it is stable iff $\eta_{\rm eff}<1/\lambda_{\max}$ and reduces to the flow with $T=\eta_{\rm eff}k$ when $2\eta_{\rm eff}\lambda_{\max}\ll1$, so that $\tau_i=1/(4\eta_{\rm eff}\lambda_i)$ steps. The repo's two RFNN scripts share $\eta=0.01\,\dlat/\Delta_t$ but differ in $c_{\rm loss}$: the \texttt{code/v3} path sums the residual $\|\sqrt{\Delta_t}s+\eta\|^2$ over dimensions and averages over the batch, so $c_{\rm loss}=\Delta_t$, whereas the script that produced the saved sweeps divides additionally by $\dlat$:
\begin{equation}\label{eq:lin-steps}
T=0.01\,\dlat\,k\ \ (\texttt{code/v3}),\qquad T=0.01\,k\ \ (\text{saved sweeps}).
\end{equation}
Stability holds in both: $2\eta_{\rm eff}\lambda_{\max}(\Umat)\approx0.02\,a_1(q)^2\lambda^{\mathsf S}_{\max}\approx0.05$ (v3) and $\le0.01$ (v2). Both use momentum $0$ and full batch with fresh $\eta$ each step, so the realized trajectory is this ODE plus zero-mean fluctuations (\cref{prop:sgdmean-p1}). The draft's \eqref{eq:rfnn-gradflow} writes $\dot A=-(A-A_\star)\Umat$ with per-mode factor $e^{-\lambda_iT}$: its $T$ is twice ours and its factor is the \emph{parameter} clock, so the risk clock $1/(4\lambda_i)$ here is $1/(2\lambda_i)$ in the draft's units (and \cref{rem:clock-p1}'s $T=0.02\dlat\times$steps is the v3 clock in those units). Adaptive optimizers (the MLP's Adam) are not covered.
\end{remark}

\begin{corollary}[Population risk versus $(\dlat,T)$ in the block-anisotropic model; proved]\label{cor:lin-blocks}
With $B_0=0$ and population training, the excess risk of \cref{prop:lin-gf-closed} is
\begin{equation}\label{eq:lin-total}
\mathcal E(T;\dlat)=\underbrace{\sum_{i=1}^{\dint}\frac{e^{-4\lambda^{\mathsf S}_iT}}{\lambda^{\mathsf S}_i}}_{\text{signal block}}\;+\;\underbrace{(\dlat-\dint)\,\frac{e^{-4\beta_t^2T}}{\beta_t^2}}_{\text{null block}},
\end{equation}
and the per-dimension score error is $\mathrm{err}_\infty(T;\dlat)=\big[\mathcal E(T;\dlat)+R_{\mathrm{irr}}\big]/\dlat$. Writing $a(T)=\sum_i e^{-4\lambda_i^{\mathsf S}T}/\lambda^{\mathsf S}_i+R_{\mathrm{irr}}$ and $b(T)=e^{-4\beta_t^2T}/\beta_t^2$, both independent of $\dlat$:
\begin{enumerate}
\item[(a)] (\emph{affine law}) $\dlat\cdot\mathrm{err}_\infty(T;\dlat)=a(T)+(\dlat-\dint)\,b(T)$ is affine in $\dlat$ at every $T$, with slope $b(T)$ decaying at the single rate $4\beta_t^2$.
\item[(b)] (\emph{monotone in $\dlat$}) $\mathrm{err}_\infty(T;\dlat)=b(T)+\big(a(T)-\dint b(T)\big)/\dlat$ is monotone in $\dlat$ at every $T$: increasing toward the plateau $b(T)$ while $a(T)<\dint b(T)$, which holds at $T=0$ whenever $\signoise<\sigsig$ since $a(0)\le\dint/\gamma_t^2<\dint b(0)$, and decreasing like $R_{\mathrm{irr}}/\dlat$ for $T$ beyond the $\dlat$-independent crossover $T_\times$ solving $a(T_\times)=\dint b(T_\times)$. The late regime and $T_\times$ are $\nsamp=\infty$ statements; at the $\dlat/\nsamp$ of the sweeps they are replaced by \cref{prop:lin-emp}(d).
\item[(c)] (\emph{null dominance}) For $\lambda^{\mathsf S}_i\equiv\alpha_t^2$, $\mathcal E_{\mathsf N}/\mathcal E_{\mathsf S}=\frac{\dlat-\dint}{\dint}\cdot\frac{\alpha_t^2}{\beta_t^2}\cdot e^{4(\alpha_t^2-\beta_t^2)T}$, increasing in both $\dlat$ and $T$; with the actual centers replace $\alpha_t^2$ by the slowest signal eigenvalue $\lambda^{\mathsf S}_{\min}$ for the late-time rate.
\item[(d)] (\emph{$\varepsilon$-time}) The null block's per-dimension excess risk falls below $\varepsilon$ at $T_\varepsilon=(4\beta_t^2)^{-1}\log\big(1/(\beta_t^2\varepsilon)\big)$, independent of $\dlat$ at $\nsamp=\infty$.
\end{enumerate}
\end{corollary}
\begin{proof}
Insert the block spectrum of $M_t$ into \eqref{eq:lin-excess} and add the irreducible term of \eqref{eq:lin-irr}; (a)--(d) are algebraic rearrangements. In (b), $\partial_{\dlat}\mathrm{err}_\infty=-(a-\dint b)/\dlat^2$ has the sign of $\dint b-a$, independent of $\dlat$; $a(0)=\mathcal J_{t,\mathsf S}\le\dint/\gamma_t^2$ by \cref{lem:null-exact-lin}(c). In (c) divide the two terms of \eqref{eq:lin-total}.
\end{proof}

\begin{remark}[Isotropic control and the diffusion floor]\label{rem:lin-iso}
If $\signoise=\sigsig$ then $\beta_t^2=\gamma_t^2$ ($=1$ for all $t$ when $\sigsig=1$): the null block still exists, still carries energy $(\dlat-\dint)/\gamma_t^2$ and still has its own clock $1/(4\gamma_t^2)$; what disappears is only the \emph{separation} $\lambda^{\mathsf S}_i/\beta_t^2$ of its clock from the signal clocks. If instead $\signoise\to0$ then $\beta_t^2\to\Delta_t$: the null block sits on the diffusion floor, its target energy per dimension $1/\Delta_t$ is set by the noising process alone and its clock $1/(4\Delta_t)$ is the slowest attainable at that $t$. In this limit the empirical $\widehat M_t$ has null block $\Delta_tI$ up to $e^{-2t}\signoise^2(1+\sqrt c)^2\le2.5\times10^{-4}$ (relative $1\%$ at $\signoise=0.01$), so the finite-$\nsamp$ effects of \cref{prop:lin-emp} are absent: the $\signoise=0.01$ sweep is the control in which this corollary applies verbatim.
\end{remark}

\paragraph{Zero-initialization score error.}
With $B_0=0$ the predicted per-dimension score error at the first evaluation is $[\mathcal J_{t,\mathsf S}+(\dlat-\dint)/\beta_t^2]/\dlat$ with the Monte-Carlo $\mathcal J_{t,\mathsf S}$ of \cref{tab:lin-constants} (a point prediction, no fitted parameters). At $t=0.1$, $\signoise=0.5$: $0.716$ ($\dlat=5$), $1.654$ ($10$), $2.122$ ($20$), $2.357$ ($40$), $2.497$ ($100$), $2.544$ ($200$), limit $2.59$; at $\signoise=0.01$: $0.716$, $3.115$, $4.315$, $4.914$, $5.274$, $5.394$, limit $5.51$. At the RFNN's $t=0.01$, $\signoise=0.5$: $0.737$, $2.256$, $3.016$, $3.396$, $3.624$, $3.700$, limit $3.78$; at $\signoise=0.01$: $0.737$, $25.5$, $37.9$, $44.1$, $47.8$, $49.0$, limit $50.3$. The code's first logged evaluation is step $1$, which moves the top signal mode by a factor $(1-2\eta_{\rm eff}\lambda_{\max})^2\approx0.9$ and the null modes by less than $1\%$.

\begin{proposition}[Empirical-risk flow: finite-$\nsamp$ null band, overshoot, and late-time floor; (a)--(c) proved, (d) leading order]\label{prop:lin-emp}
Let $\widehat B(T)$ be gradient flow on $L_{t,n}$ from $B_0=0$ and let $M_t=\diag(M_{t,\mathsf S},\beta_t^2I)$ be the population second moment.
\begin{enumerate}
\item[(a)] (\emph{exact}) $\nabla_BL_{t,n}=2(B\widehat M_t+I)$, hence $\widehat B(T)=-\widehat M_t^{-1}\big(I-e^{-2\widehat M_tT}\big)\to-\widehat M_t^{-1}$, and the population excess risk of the trained model is exactly
\begin{equation}\label{eq:lin-emp-excess}
\mathcal E_n(T)=\operatorname{tr}\big(E(T)M_tE(T)^{\top}\big),\qquad E(T)=M_t^{-1}-\widehat M_t^{-1}\big(I-e^{-2\widehat M_tT}\big).
\end{equation}
\item[(b)] (\emph{null block}) Let $\widehat M_{t,\mathsf N}$ be the null--null principal block of $\widehat M_t$, with eigenvalues $\widehat\lambda^{\mathsf N}_i$. Dropping the signal--null cross block of $\widehat M_t$, the null contribution to \eqref{eq:lin-emp-excess} is
\begin{equation}\label{eq:lin-emp-null}
\mathcal E_{n,\mathsf N}(T)=\beta_t^2\sum_{i=1}^{\dlat-\dint}\Big(\frac1{\beta_t^2}-\frac{1-e^{-2\widehat\lambda^{\mathsf N}_iT}}{\widehat\lambda^{\mathsf N}_i}\Big)^2 .
\end{equation}
A mode with $\widehat\lambda^{\mathsf N}_i\ge\beta_t^2$ decays monotonically to its floor $\beta_t^2(1/\widehat\lambda^{\mathsf N}_i-1/\beta_t^2)^2$; a mode with $\widehat\lambda^{\mathsf N}_i<\beta_t^2$ first reaches \emph{zero} population error at $T^0_i=(2\widehat\lambda^{\mathsf N}_i)^{-1}\log\big[\beta_t^2/(\beta_t^2-\widehat\lambda^{\mathsf N}_i)\big]$ and then rises to the same floor: the empirical fit overshoots the population score on under-sampled null directions.
\item[(c)] (\emph{Marchenko--Pastur band}) $\widehat M_{t,\mathsf N}=e^{-2t}\signoise^2\,\mathcal W+\Delta_tI$ with $\mathcal W=\frac1n X_{\mathsf N}^{\top}X_{\mathsf N}$ a white Wishart matrix of aspect ratio $c=(\dlat-\dint)/\nsamp$. As $\nsamp\to\infty$ with $c$ fixed its spectrum follows the shifted, scaled Marchenko--Pastur law \citep{marchenko1967distribution,couillet2011rmt}, with edges
\begin{equation}\label{eq:lin-mp-edges}
\widehat\lambda^{\mathsf N}_\pm=e^{-2t}\signoise^2\big(1\pm\sqrt{c}\big)^2+\Delta_t,\qquad c=\frac{\dlat-\dint}{\nsamp}.
\end{equation}
The null clock is therefore a band $\big[1/(4\widehat\lambda^{\mathsf N}_+),\,1/(4\widehat\lambda^{\mathsf N}_-)\big]$ whose slow edge grows with $\dlat/\nsamp$ and diverges as $c\to1$ (the regime where the sample bulk $\nsamp-\dlat$ vanishes). By Cauchy interlacing the $\dlat-\dint$ smallest eigenvalues of the full $\widehat M_t$ obey the same bracket up to an index shift of $\dint$.
\item[(d)] (\emph{late-time floor}) As $T\to\infty$,
\begin{equation}\label{eq:lin-floor}
\mathrm{err}_n(\infty;\dlat)=\frac{R_{\mathrm{irr}}+\mathcal E^{\rm est}_{\mathsf S}}{\dlat}+\frac{\dlat-\dint}{\dlat}\,F_t(c),\qquad F_t(c)=\beta_t^2\!\int\!\Big(\frac1{\widehat\lambda}-\frac1{\beta_t^2}\Big)^2d\mu^{\rm MP}_{t,c}(\widehat\lambda)=c\,\frac{(e^{-2t}\signoise^2)^2}{\beta_t^6}+O(c^2),
\end{equation}
with $\mathcal E^{\rm est}_{\mathsf S}=O(\dint/\nsamp)$ the signal-block estimation error. $F_t$ vanishes at $c=0$ and at $\signoise\to0$ and increases with $c$; hence at late times $\mathrm{err}_n$ is \emph{U-shaped} in $\dlat$ (the $R_{\mathrm{irr}}/\dlat$ decrease at small $\dlat$, the $F_t(c)$ increase at large $\dlat$), and at large $\dlat/\nsamp$ the trained model is worse on the null block than the zero model (per null dimension $F_t>1/\beta_t^2$ at $\dlat\ge200$, $t=0.01$).
\end{enumerate}
\end{proposition}
\begin{proof}
(a) $\nabla_BL_{t,n}=\frac2n\sum_\mu\E_\eta[(Bx_t^\mu+\eta/\sqrt{\Delta_t})x_t^{\mu\top}]=2B\widehat M_t+\frac{2}{\sqrt{\Delta_t}}\frac1n\sum_\mu\E_\eta[\eta x_t^{\mu\top}]=2(B\widehat M_t+I)$, using $\E_\eta[\eta x_t^{\mu\top}]=\sqrt{\Delta_t}I$ and $\frac1n\sum_\mu\E_\eta[x_t^\mu x_t^{\mu\top}]=\widehat M_t$. The ODE $\dot B=-2(B\widehat M_t+I)$ is solved as in \cref{prop:lin-gf-closed} with $M_t\to\widehat M_t$. The population excess of any $B$ is $\operatorname{tr}((B+M_t^{-1})M_t(B+M_t^{-1})^{\top})$ by the exact expansion in that proposition, and $\widehat B(T)+M_t^{-1}=E(T)$.
(b) With the cross block dropped, $\widehat M_t$ and $E(T)$ are block diagonal; on the null block $M_t=\beta_t^2I$ and $E(T)$ is diagonal in the eigenbasis of $\widehat M_{t,\mathsf N}$ with entries $1/\beta_t^2-g_i(T)$, $g_i(T)=(1-e^{-2\widehat\lambda_iT})/\widehat\lambda_i$, giving \eqref{eq:lin-emp-null}. $g_i$ increases from $0$ to $1/\widehat\lambda_i$; if $\widehat\lambda_i<\beta_t^2$ then $1/\widehat\lambda_i>1/\beta_t^2$ and $g_i(T^0_i)=1/\beta_t^2$ at the stated $T^0_i$, after which $(1/\beta_t^2-g_i)^2$ increases to $(1/\beta_t^2-1/\widehat\lambda_i)^2$. Numerically (seed 42, $\nsamp=500$, $\signoise=0.5$, $t=0.01$, $\dlat=200$) the full-$\widehat M_t$ evaluation of \eqref{eq:lin-emp-excess} gives a late floor of $4.81$ per dimension versus $4.50$ with the cross block dropped, and the null eigenvalues move from $[0.0554,0.661]$ to $[0.0544,0.658]$.
(c) In the rotated frame the null coordinates of the training points are i.i.d.\ $\mathcal N(0,\signoise^2I_{\dlat-\dint})$, independent of the signal coordinates and of the labels, so $\frac1nX_{\mathsf N}^{\top}X_{\mathsf N}=\signoise^2\mathcal W$ with $\mathcal W$ white Wishart; the MP law with variance parameter $1$ has support $[(1-\sqrt c)^2,(1+\sqrt c)^2]$, mean $1$ and variance $c$. Interlacing: for a principal submatrix, $\lambda_{k+\dint}(\widehat M_t)\le\lambda_k(\widehat M_{t,\mathsf N})\le\lambda_k(\widehat M_t)$. At $\nsamp=500$ the measured null band and the MP edges \eqref{eq:lin-mp-edges} agree to $1$--$3\%$ (\cref{tab:lin-band}).
(d) Let $T\to\infty$ in (b) and average over the MP law. Expanding $1/\widehat\lambda-1/\beta_t^2=-(\widehat\lambda-\beta_t^2)/\beta_t^4+O((\widehat\lambda-\beta_t^2)^2)$ and using $\E(\widehat\lambda-\beta_t^2)^2=(e^{-2t}\signoise^2)^2\operatorname{Var}_{\rm MP}=(e^{-2t}\signoise^2)^2c$ gives the leading term (at $t=0.1$, $\dlat=40$: $0.051$ versus the exact $0.060$). On the signal block $\widehat B_{\mathsf S}(\infty)=-\widehat M_{t,\mathsf S}^{-1}$ with $\widehat M_{t,\mathsf S}-M_{t,\mathsf S}=O_P(n^{-1/2})$ of fixed rank $\dint$, so its excess is $O(\dint/\nsamp)$. Monotonicity of $F_t$ in $c$ and the U-shape are verified numerically across the sweep (\cref{tab:lin-band}): the finite-$\nsamp$ term overtakes $R_{\mathrm{irr}}/\dlat$ at $\dlat\approx40$ ($t=0.1$) and $\dlat\approx20$ ($t=0.01$).
\end{proof}

\begin{table}[h]
\centering\scriptsize
\caption{Finite-$\nsamp$ null block of $\widehat M_t$ ($\nsamp=500$, $\signoise=0.5$, seed 42) and the resulting late-time population excess per dimension of the trained linear model (full $\widehat M_t$, including the signal-block estimation error). MP edges from \eqref{eq:lin-mp-edges}.}\label{tab:lin-band}
\resizebox{\textwidth}{!}{%
\begin{tabular}{lccccc}
\toprule
$\dlat$ & 20 & 40 & 100 & 200 & 240\\
\midrule
$c=(\dlat-5)/500$ & $0.03$ & $0.07$ & $0.19$ & $0.39$ & $0.47$\\
$t=0.01$: null band $[\min,\mathrm{med},\max]$ & $[0.195,0.263,0.361]$ & $[0.156,0.261,0.410]$ & $[0.098,0.248,0.516]$ & $[0.054,0.233,0.658]$ & $[0.045,0.225,0.707]$\\
\quad MP edges & $[0.187,0.357]$ & $[0.152,0.412]$ & $[0.098,0.525]$ & $[0.054,0.666]$ & $[0.044,0.716]$\\
\quad slow-edge clock / $1/(4\beta_t^2)$ & $1.4$ & $1.7$ & $2.6$ & $4.8$ & $5.9$\\
\quad late excess per dim & $0.11$ & $0.32$ & $1.21$ & $4.81$ & $7.4$\\
$t=0.1$: null band $[\min,\mathrm{med},\max]$ & $[0.327,0.385,0.466]$ & $[0.295,0.380,0.507]$ & $[0.247,0.369,0.596]$ & $[0.210,0.358,0.714]$ & $[0.203,0.353,0.755]$\\
\quad late excess per dim & $0.02$ & $0.066$ & $0.167$ & $0.365$ & $0.45$\\
$R_{\mathrm{irr}}/\dlat$ ($t=0.1$ / $0.01$) & $0.068/0.084$ & $0.034/0.042$ & $0.014/0.017$ & $0.007/0.008$ & $0.006/0.007$\\
\bottomrule
\end{tabular}}
\end{table}

\begin{corollary}[Random-feature readout: substitute $M_t\to\Umat$; (a) proved, (b) sketched]\label{cor:lin-rfnn}
Let $L_n(A)=\frac1n\sum_{\mu}\E_\eta\|A\phi(x_t^\mu)+\eta/\sqrt{\Delta_t}\|^2$ and $V:=-\frac1n\sum_\mu\E_\eta[\eta\,\phi(x_t^\mu)^{\top}]/\sqrt{\Delta_t}\in\R^{\dlat\times\pwidth}$. Let $\Umat=\sum_{i}\lambda_i(\Umat)u_iu_i^{\top}$ and $A_\star:=V\Umat^{+}$.
\begin{enumerate}
\item[(a)] (\emph{exact}) Gradient flow on $L_n$ from $A_0=0$ gives $A(T)=A_\star\big(I-e^{-2\Umat T}\big)$ and
\begin{equation}\label{eq:rfnn-excess}
L_n(A(T))-\min_AL_n=\sum_{i:\lambda_i(\Umat)>0}\lambda_i(\Umat)\,e^{-4\lambda_i(\Umat)T}\,\|A_\star u_i\|^2 ,
\end{equation}
so every mode of $\Umat$ has clock $\tau_i=1/(4\lambda_i(\Umat))$ and modes in $\ker\Umat$ never move. For a single fixed noise draw $\operatorname{rank}\Umat\le\nsamp$, so at least $\pwidth-\nsamp$ such modes exist (the cliff at index $\nsamp$); for the $\E_\eta$-averaged $\Umat$ the eigenvalues beyond index $\nsamp$ are not exactly zero but $O(10^{-6})$, with clocks far beyond any training horizon. This is \cref{prop:lin-gf-closed} with $M_t\to\Umat$ and $M_t^{-1}\to A_\star$ (and \cref{lem:gf-p1} in the present convention).
\item[(b)] (\emph{data-side identification, sketched}) Under the Gaussian-equivalence decomposition $\phi(x)=Gx+a_*(q)z(x)$, $G=a_1(q)W/\sqrt{\pwidth\dlat}$, restricted to the linear part $\phi_{\rm lin}=Gx$ and with $W^{\top}W=\pwidth I$, the $\dlat$ nonzero eigenvalues of $\Umat_{\rm lin}=G\widehat M_tG^{\top}$ are
\begin{equation}\label{eq:rfnn-eig-id}
\lambda_i(\Umat)=\frac{a_1(q)^2}{\dlat}\,\widehat\lambda_i(\widehat M_t),\qquad q=\frac{\operatorname{tr}\widehat M_t}{\dlat}\simeq\beta_t^2+\frac{\dint(\alpha_t^2-\beta_t^2)}{\dlat},
\end{equation}
and the population excess risk carried by these modes is that of the empirical-risk linear model, \cref{prop:lin-emp}, with every clock multiplied by $\dlat/a_1(q)^2$:
\begin{equation}\label{eq:rfnn-risk-id}
\mathcal E^{\rm RF}(T;\dlat)=\mathcal E_n\big(a_1(q)^2T/\dlat\big),
\end{equation}
so that the null clock is the band
\begin{equation}\label{eq:rfnn-null-clock}
\tau_{\mathsf N}^{\rm RF}(\dlat)\in\frac{\dlat}{4\,a_1(q)^2}\Big[\frac1{\widehat\lambda^{\mathsf N}_+},\frac1{\widehat\lambda^{\mathsf N}_-}\Big],\qquad \text{median}\approx\frac{\dlat}{4\,a_1(q)^2\,\beta_t^2}\ (\text{MP median}/\beta_t^2=0.88\text{ at }\dlat=200),
\end{equation}
with $\widehat\lambda^{\mathsf N}_\pm$ from \eqref{eq:lin-mp-edges}. In gradient-flow units of $L_n$ the median clock grows essentially linearly in $\dlat$, partly offset by the de-saturation of $\tanh$ ($q\downarrow$, $a_1(q)^2\uparrow$): at $t=0.01$, $\signoise=0.5$ it is $26.2\to302$ over $\dlat=5\to200$ ($11.5\times$), and its slow edge is a further $4.8\times$ longer at $\dlat=200$. In \emph{training steps} the code's learning rate $0.01\,\dlat/\Delta_t$ removes the factor $\dlat$ (\eqref{eq:lin-steps}): $\tau^{\rm RF}_{\mathsf N}=25/(a_1(q)^2\beta_t^2)$ steps in \texttt{code/v3}, \emph{decreasing} $524\to151$ across the sweep, and $25\dlat/(a_1(q)^2\beta_t^2)=2620\to30\,190$ steps for the saved sweeps. The trace of \eqref{eq:rfnn-eig-id} over the data-side modes is $a_1(q)^2q$, the linear part of $\E\tanh^2(\sqrt q\,z)=a_1(q)^2q+a_*(q)^2$; the remaining mass $a_*(q)^2$ lives on the $\nsamp-\dlat$ sample modes, which are invisible to this corollary.
\end{enumerate}
\end{corollary}
\begin{proof}
(a) $L_n$ is a quadratic in $A$ with $\nabla_AL_n=2(A\Umat-V)$, so $\dot A=-2(A\Umat-V)$. The rows of $V$ lie in the row space of $\Umat$: if $u\in\ker\Umat$ then $\frac1n\sum_\mu\E_\eta(\phi(x_t^\mu)^{\top}u)^2=0$, so $\phi(x_t^\mu)^{\top}u=0$ a.s.\ and $Vu=0$. Hence $V=A_\star\Umat$ with $A_\star=V\Umat^{+}$; $E:=A-A_\star$ obeys $\dot E=-2E\Umat$ and $E(T)=-A_\star e^{-2\Umat T}$. Expanding $L_n(A_\star+E)=L_n(A_\star)+\operatorname{tr}(E\Umat E^{\top})$ (the cross term vanishes because $A_\star\Umat=V$) gives \eqref{eq:rfnn-excess}; modes with $\lambda_i(\Umat)=0$ contribute nothing and $A(T)u_i=0$ for them. The rank statement for a single noise draw is $\Umat=\frac1n\Phi^{\top}\Phi$ with $\Phi\in\R^{\nsamp\times\pwidth}$; for the $\E_\eta$ average, \texttt{compute\_U} at $\dlat=20$, $\pwidth=1280$, $t=0.01$ gives eigenvalues $2.0\times10^{-6}$ and $1.8\times10^{-6}$ at indices $\nsamp-1$ and $\nsamp$: no exact cliff, but clocks of order $10^5$ flow units.
(b) Empirical version, linear part. The empirical fixed point solves $A_\star G\widehat M_tG^{\top}=-G^{\top}$ (the empirical Stein identity $\frac1n\sum_\mu\E_\eta[(\eta/\sqrt{\Delta_t})x_t^{\mu\top}]=I$ replaces \cref{lem:null-exact-lin}(b)). With $G^{\top}G=cI_{\dlat}$, $c=a_1(q)^2W^{\top}W/(\pwidth\dlat)=a_1(q)^2/\dlat$, right-multiplying by $G$ gives $A_\star G=-\widehat M_t^{-1}$. Also $\Umat_{\rm lin}G=G\widehat M_tG^{\top}G=cG\widehat M_t$, hence $e^{-2\Umat_{\rm lin}T}G=Ge^{-2c\widehat M_tT}$ and the nonzero spectrum of $\Umat_{\rm lin}$ is $c\,\widehat\lambda_i(\widehat M_t)$ with eigenvectors $G\widehat v_i/\|G\widehat v_i\|$, which is \eqref{eq:rfnn-eig-id}. The population excess risk of $s_{A(T)}$ on these modes is $\operatorname{tr}((A(T)G+M_t^{-1})M_t(A(T)G+M_t^{-1})^{\top})$ with $A(T)G+M_t^{-1}=M_t^{-1}-\widehat M_t^{-1}(I-e^{-2c\widehat M_tT})=E(cT)$ in the notation of \eqref{eq:lin-emp-excess}, which is \eqref{eq:rfnn-risk-id}; the band \eqref{eq:rfnn-null-clock} is \cref{prop:lin-emp}(c) rescaled. The step conversions follow from \cref{rem:lin-gd}. Consistency with the measured spectra: at $t=0.01$, $\signoise=0.5$, the idealized $q=\operatorname{tr}M_t/\dlat$ runs $2.764\to0.327$ and $a_1(q)^2=0.180\to0.625$ over $\dlat=5\to200$, exactly the values of \cref{lem:q-p1}; the code-normalized noise-dim medians predicted by \eqref{eq:rfnn-eig-id} with the MP median of $\widehat M_{t,\mathsf N}$ are $64\,a_1(q)^2\,\mathrm{med}(\widehat\lambda^{\mathsf N})=4.77$ ($\dlat=10$) and $9.33$ ($\dlat=200$) against measured $5.22$ and $9.22$; using the population $\beta_t^2$ instead gives $4.76$ and $10.60$, so the $13\%$ shortfall at $\dlat=200$ is the MP median ($0.233/0.265=0.88$), not a finite-$\pwidth$ effect. What is only sketched: (i) $W^{\top}W=\pwidth I$ holds in expectation with relative fluctuations $O(\sqrt{\dlat/\pwidth})$, which perturb the data-side eigenvalues by the same relative amount (the $W$ factor of \cref{app:theory-fourbulk}); (ii) the residual features $a_*(q)z$ are uncorrelated with $x$ under Gaussian equivalence \citep{pennington2017nonlinear,benigni2021eigenvalue,goldt2022gaussian,hu2022universality,mei2022generalization} but not orthogonal to $Gx$ sample by sample, so the data-side eigenvectors of the full $\Umat$ are those of $\Umat_{\rm lin}$ only to leading order; the data-side eigenvalues of the full $\Umat$ measured with \texttt{compute\_U} ($\dlat=20$, $\pwidth=1280$, $t=0.01$) exceed \eqref{eq:rfnn-eig-id} by $4$--$7\%$ (null block $[3.96,5.43,7.63]\times10^{-3}$ versus $[3.82,5.13,7.10]\times10^{-3}$) and $\operatorname{tr}\Umat=0.357$ versus $a_1^2q+a_*^2=0.372$; this few-percent correction is consistent with row-wise fluctuations of the pre-activation variance $q_j=w_j^{\top}\widehat M_tw_j/\dlat$ (Jensen) and with the non-Gaussian mixture pre-activations $w_j^{\top}m_c/\sqrt{\dlat}$, both decreasing with $\dlat$, and is much larger than the sample-bulk mean $a_*(q)^2/\nsamp\approx4\times10^{-5}$; (iii) Gaussian equivalence itself is proved for Gaussian or subspace data with isotropic-plus-spike covariance; a rigorous version of (b) would follow from a Gaussian-equivalence theorem for the conjugate kernel with anisotropic input covariance \citep{hu2022universality,bonnaire2025} evaluated at the two-atom spectral measure, and for data exactly on a subspace ($\signoise=0$) the corresponding block structure appears in \citet{george2025dsm}.
\end{proof}

\begin{remark}[What the linear model cannot and can see]\label{rem:lin-nomem}
Trained on the empirical denoising risk, the linear model converges to $-\widehat M_t^{-1}$, the score of the centered Gaussian with the empirical second moment; it never reproduces training points, because it has no sample bulk. Memorization in the NN-ratio sense requires the residual features $a_*(q)z$, whose second moment is the $\nsamp-\dlat$-mode sample bulk of $\Umat$; \cref{cor:lin-rfnn}(a) gives those modes the same clock law $1/(4\lambda_i(\Umat))$ but with weights $\lambda_i\|A_\star u_i\|^2$ that are not $1/\lambda_i$, and they are treated in \cref{app:theory-buffer,app:theory-predictor}. What the linear model \emph{does} see is the finite-sample overfitting of the null block (\cref{prop:lin-emp}): the null directions are estimated from $\nsamp$ points, the empirical eigenvalues spread over the MP band, and the fitted inverse overshoots on the under-sampled directions with a late-time population penalty $F_t(\dlat/\nsamp)$ that grows with $\dlat$. This is the linear-model form of the empirical observation that memorization returns as $\dlat\to\nsamp$ (\cref{app:theory-scope}), and it lives entirely in the $\dlat$ data-side modes.
\end{remark}

\begin{remark}[Which block sets the generalization clock]\label{rem:lin-taugen}
Since $\beta_t^2<\lambda^{\mathsf S}_{\min}$ whenever $\signoise<\sigsig$ ($0.265$ versus $1.89$ at $t=0.01$), the slowest \emph{data-side} modes are the null band, not the signal block: the population score error is learned on the clock band $[1/(4\widehat\lambda^{\mathsf N}_+),1/(4\widehat\lambda^{\mathsf N}_-)]$ (linear model) or $\dlat/(4a_1(q)^2)$ times that (RFNN), ordered between the signal clocks and the sample-bulk onset. If $\taugen$ is defined from the signal edge (\cref{def:taugen-p1,def:abs-times-buf}), the ordering $\taugen<\tau_{\mathsf N}<\taumem$ holds; if it is defined from the score error, it is the slow edge of $\tau_{\mathsf N}$ itself. At $t=0.01$ the factor $\dlat/a_1(q)^2$ common to all data-side clocks grows $11.5\times$ over $\dlat=5\to200$ in flow units, while the signal-edge $\taugen$ on the released spectra grows $11\times$ with the actual (non-isotropic, empirical) signal spectrum (\cref{app:theory-fourbulk}); the two are the same quantity only up to the spread of $\widehat\lambda^{\mathsf S}$ and the step conversion \eqref{eq:lin-steps}, which must be applied before any absolute comparison of step counts.
\end{remark}

\paragraph{Predictions and how to test them (all for the corrected score-error metric; no saved value constrains them).}
\begin{enumerate}
\item \emph{Zero-initialization score error per latent dimension is $[\mathcal J_{\mathsf S}+(\dlat-\dint)/\beta_t^2]/\dlat$} (values listed above). Test: log the score error at the first evaluation for the RFNN ($A_0=0$) and for an MLP whose last layer is zeroed, at $t=0.01$ and $t=0.1$ separately.
\item \emph{Affine law:} $\dlat\cdot\mathrm{err}(T)$ is affine in $\dlat$ with a single slope $b(T)$ only where the null band is narrow ($\dlat\le40$ at $\signoise=0.5$; the whole sweep at $\signoise=0.01$); beyond, the slope is the MP average of \eqref{eq:lin-emp-null}, larger than $b(T)$ and $\dlat/\nsamp$-dependent. Test: fit lines at several fixed steps; check the fitted slope decays at the rate of \eqref{eq:lin-steps}; check linearity breaks upward at $\dlat\ge100$, $\signoise=0.5$, late steps.
\item \emph{Shape:} monotone increasing in $\dlat$ at early times, U-shaped at late times with minimum near $\dlat\sim20$--$40$ and the large-$\dlat$ branch rising with $F_t(c)$; never n-shaped; non-monotone in steps at $\dlat\ge100$ (early-stopping minimum $\approx0.36$ per dim at $\dlat=200$, $t=0.01$, floor $4.8$). Test: plot $\mathrm{err}$ vs $\dlat$ at a grid of steps and vs steps at $\dlat\ge100$; compare late floors with \cref{tab:lin-band}.
\item \emph{$\signoise=0.01$ control:} the null block is $2.13\times$ ($t=0.1$) or $13.3\times$ ($t=0.01$) larger and slower than at $\signoise=0.5$; slope ratio $b_{0.01}/b_{0.5}=(\beta^2_{0.5}/\beta^2_{0.01})e^{+4(\beta^2_{0.5}-\beta^2_{0.01})T}$; no finite-$\nsamp$ floor, so the late $\dlat$-dependence has opposite sign in the two sweeps.
\item \emph{RFNN null band:} $\dlat\lambda^{\rm nd}(\Umat)/a_1(q)^2$ fills $e^{-2t}\signoise^2(1\pm\sqrt{(\dlat-5)/500})^2+\Delta_t$ with medians $0.263,0.261,0.248,0.233,0.225$ at $\dlat=20,\dots,240$, and per-mode risk decays as $e^{-4\lambda_iT}$. Test: rescale the saved noise-dim block by $1/(\psi_pa_1(q)^2)$ and overlay the MP band; fit per-mode risk decay in the eigenbasis of $\Umat$ to check the factor $4$.
\item \emph{Isotropic control:} the null block persists with energy $\dlat-5$ and clock $1/4$ exactly; only its separation from the signal clocks survives; the finite-$\nsamp$ band is four times wider in absolute terms and the late floor larger. Test: early slope decays at rate $4$; late floor exceeds the $\signoise=0.5$ floor at matched $\dlat$.
\end{enumerate}

\paragraph{Open gaps.}
(1) All RFNN-facing constants are at $t_{\rm fixed}=0.01$ and the null block is carried as the empirical $\widehat M_t$ with its MP band; the late-time shape is U-shaped in $\dlat$ (not monotone at every $T$, and not ``decreasing like $R_{\mathrm{irr}}/\dlat$, never n-shaped'' as an earlier draft said). (2) \cref{prop:lin-emp}(b) drops the signal--null cross block of $\widehat M_t$ ($1\%$ on eigenvalues, $7$--$8\%$ on the late floor at $\dlat=200$); the exact \eqref{eq:lin-emp-excess} needs no such approximation, but a clean analytic bound on the cross-block contribution is not written out; monotonicity of $F_t(c)$ beyond leading order and the location of the U minimum are numerical. (3) \cref{cor:lin-rfnn}(b) needs (i) concentration of $W^{\top}W/\pwidth$, (ii) control of the $a_*(q)z$ cross terms on the data-side eigenvectors of the full $\Umat$ (the $+4$--$7\%$ shift is attributed, without proof, to row-wise $q$ fluctuations and mixture pre-activations), (iii) a Gaussian-equivalence theorem with anisotropic covariance for the two-atom $\rho_\Sigma$. (4) The linear model has no sample bulk and never reproduces training points; it explains the generalizing data-side modes, the score-error curve and the finite-$\nsamp$ null overfitting, not the NN-ratio fraction curve. (5) The MLP is trained with Adam, minibatches, $t\sim U[0.01,3]$ and non-zero default init; every MLP-facing prediction is conjectural for it. (6) The clock is $1/(4\lambda_i)$ for the risk and $1/(2\lambda_i)$ for the parameters in this convention; absolute step counts must be converted with \eqref{eq:lin-steps} before comparison.

% ----------------------------------------------------------------------
\subsection{Bridge from residual sample-bulk error to near-duplicate generation}
\label{app:theory-bridge}

This subsection connects the spectral quantities above to the memorisation \emph{fraction} that the paper reports. It proves two elementary lemmas, gives the sketch of \cref{prop:bridge}, derives its corollaries, and records the empirical anchors on the paper's own synthetic data. Two referee reports on this subsection could not be incorporated by revision; their substantive corrections are applied inline where they are unambiguous (marked ``[R]'') and listed under Open gaps otherwise. Throughout, the sampler is the reverse VP-SDE $dx=(x+2s_T)\,d\tau+\sqrt2\,dW$ ($\tau=-t$) run from $t_{\max}=3$ to $t_{\min}=0.01$ by Euler--Maruyama with $500$ steps and no terminal denoising step, as in the released code; $\Delta_{t_{\min}}=1-e^{-0.02}=0.0198$.

\begin{definition}[Empirical score, residual error, NN-ratio]\label{def:bridge-objects}
Let $\{x^\mu\}_{\mu=1}^{\nsamp}\subset\R^{\dlat}$ be the training set and $\Delta_t=1-e^{-2t}$. The \emph{exact empirical score} is the score of the noised empirical measure $\rho^{\rm emp}_t=\frac1\nsamp\sum_\mu\mathcal N(e^{-t}x^\mu,\Delta_t I)$,
\begin{equation}\label{eq:semp}
 s_{\rm emp}(x,t)=\sum_{\mu=1}^{\nsamp} w_\mu(x,t)\,\frac{e^{-t}x^\mu-x}{\Delta_t},\qquad w_\mu(x,t)=\frac{\exp\!\big(-\|x-e^{-t}x^\mu\|^2/2\Delta_t\big)}{\sum_\nu \exp\!\big(-\|x-e^{-t}x^\nu\|^2/2\Delta_t\big)}
\end{equation}
\citep{scarvelis2023closedform,gu2023memorization}. For a score model $s_T$ at training time $T$ write $\zeta_T(x,t):=s_T(x,t)-s_{\rm emp}(x,t)$ and define the \emph{residual error at diffusion time $t$}
\begin{equation}\label{eq:ET}
 \mathcal E_t(T):=\frac1\nsamp\sum_{\mu}\E_{\eta}\big\|\zeta_T(e^{-t}x^\mu+\sqrt{\Delta_t}\eta,\,t)\big\|^2,\qquad \eta\sim\mathcal N(0,I_{\dlat}),
\end{equation}
and its per-sample version $\mathcal E_{t,\mu}(T)$. For the RFNN at its training time $t$, $\mathcal E_t(T)=\sum_i\lambda_i\|a_i(0)-a_i^\star\|^2e^{-2\lambda_iT}+\mathcal E_t(\infty)$ (\cref{prop:rse-p1}), where $\mathcal E_t(\infty)$ is the $L^2(\rho^{\rm emp}_t)$ projection error of $s_{\rm emp}$ onto the span of the $\pwidth$ features (the gradient-flow fixed point is that projection; [R] this is unrelated to the $\lambda=0$ modes); we split $\mathcal E_t=\mathcal E^{\rm sig}_t+\mathcal E^{\rm null}_t+\mathcal E^{\rm samp}_t+\mathcal E_t(\infty)$ by block. Let $d_{\rm NN}(\mu)=\min_{\nu\neq\mu}\|x^\mu-x^\nu\|$, $G_{\rm NN}$ its empirical CDF over $\mu$, and $d_{\rm NN}^{(r)}$ its $r$-quantile. For a generated point $\hat x$ with nearest and second-nearest training points $\mathrm{NN}_1,\mathrm{NN}_2$, the repository's NN-ratio statistic (both \texttt{code/v3/lib/metrics.py} and the older \texttt{experiment\_v2.py}) is $\varrho(\hat x)=\|\hat x-\mathrm{NN}_1\|/\|\mathrm{NN}_1-\mathrm{NN}_2\|$, the form of \citet{somepalli2023diffusion}. $\hat x$ is memorized iff $\varrho<1/3$, and $F_{\rm mem}(T)=\Pr[\varrho(\hat x)<1/3]$.
\end{definition}

\begin{lemma}[Condensation of the empirical softmax; after \citealp{biroli2024dynamical}; proved]\label{lem:condensation}
Fix $t$, a training index $\mu$, and $\delta\in(0,\tfrac12)$. Let $x=e^{-t}x^\mu+\sqrt{\Delta_t}\,\eta$ with $\eta\sim\mathcal N(0,I_{\dlat})$. If
\begin{equation}\label{eq:cond-rigorous}
 e^{2t}\Delta_t\;\le\;\frac{d_{\rm NN}(\mu)^2}{32\,\log(\nsamp/\delta)},
\end{equation}
then with probability at least $1-2\delta$ over $\eta$ [R], $1-w_\mu(x,t)\le \delta^{8}\nsamp^{-7}$, and on that event $s_{\rm emp}(x,t)=(e^{-t}x^\mu-x)/\Delta_t+r$ with $\|r\|\le \delta^{8}\nsamp^{-7}\max_\nu\|e^{-t}(x^\mu-x^\nu)\|/\Delta_t$. Evaluating the log-weight ratio at $\eta=0$ gives the sufficient condition $e^{2t}\Delta_t<d_{\rm NN}(\mu)^2/(2\log\nsamp)$ for $w_\mu(e^{-t}x^\mu,t)>\tfrac12$, which for $t\ll1$ is the condensation threshold $\Delta_t<d_{\rm NN}^2/(2\log \nsamp)$ of \citet{biroli2024dynamical}. Since $e^{2t}\Delta_t=e^{2t}-1$ is increasing in $t$, condensation at $t_c$ implies condensation at every $t\le t_c$.
\end{lemma}
\begin{proof}
For $\nu\neq\mu$ write $D_\nu=\|x^\mu-x^\nu\|\ge d_{\rm NN}(\mu)$. From \eqref{eq:semp},
\[\log\frac{w_\nu}{w_\mu}=-\frac{\|x-e^{-t}x^\nu\|^2-\|x-e^{-t}x^\mu\|^2}{2\Delta_t}=-\frac{e^{-2t}D_\nu^2}{2\Delta_t}-\frac{e^{-t}\langle\eta,x^\mu-x^\nu\rangle}{\sqrt{\Delta_t}},\]
using $x-e^{-t}x^\nu=e^{-t}(x^\mu-x^\nu)+\sqrt{\Delta_t}\eta$ and $x-e^{-t}x^\mu=\sqrt{\Delta_t}\eta$. The second term is Gaussian with standard deviation $e^{-t}D_\nu/\sqrt{\Delta_t}$; by a two-sided union bound, with probability $\ge1-2\delta$ simultaneously $|\langle\eta,x^\mu-x^\nu\rangle|\le D_\nu\sqrt{2\log(\nsamp/\delta)}$ for all $\nu$. On this event
\[\log\frac{w_\nu}{w_\mu}\le-\frac{e^{-t}D_\nu}{2\sqrt{\Delta_t}}\Big(\frac{e^{-t}D_\nu}{\sqrt{\Delta_t}}-2\sqrt{2\log(\nsamp/\delta)}\Big).\]
Condition \eqref{eq:cond-rigorous} says $e^{-t}d_{\rm NN}/\sqrt{\Delta_t}\ge4\sqrt{2\log(\nsamp/\delta)}$, so the bracket is at least $\tfrac12e^{-t}D_\nu/\sqrt{\Delta_t}$ and $\log(w_\nu/w_\mu)\le-e^{-2t}D_\nu^2/(4\Delta_t)\le-8\log(\nsamp/\delta)$. Summing, $\sum_{\nu\ne\mu}w_\nu/w_\mu\le\nsamp(\delta/\nsamp)^8$, hence $1-w_\mu\le\delta^8\nsamp^{-7}$. The bound on $r=\sum_{\nu\neq\mu}w_\nu\,e^{-t}(x^\nu-x^\mu)/\Delta_t$ follows from $\sum_{\nu\ne\mu}w_\nu=1-w_\mu$. At $\eta=0$ the cross term vanishes and $\sum_{\nu\neq\mu}w_\nu/w_\mu\le(\nsamp-1)e^{-e^{-2t}d_{\rm NN}^2/2\Delta_t}<1$ if $e^{2t}\Delta_t<d_{\rm NN}^2/(2\log\nsamp)$.
\end{proof}

\begin{lemma}[Endpoint law of the exact empirical sampler; proved for the continuous-time sampler]\label{lem:endpoint}
Let the reverse VP-SDE with the exact score $s_{\rm emp}$ be run from $x_{t_c}\sim\rho^{\rm emp}_{t_c}$ to $t_{\min}$. Then $x_{t_{\min}}\sim\rho^{\rm emp}_{t_{\min}}=\frac1\nsamp\sum_\mu\mathcal N(e^{-t_{\min}}x^\mu,\Delta_{t_{\min}}I)$ [R: no condensation hypothesis is needed for this; only the $\mathcal N(0,I)$ initialisation at $t_{\max}=3$ is approximate]. If in addition condensation \eqref{eq:cond-rigorous} holds for $\mu$ on $[t_{\min},t_c]$, the trajectory started near $e^{-t_c}x^\mu$ stays attached to $\mu$, and for the probability-flow ODE with the same score the offset $y_t=x_t-e^{-t}x^\mu$ satisfies $y_{t_{\min}}=\sqrt{\Delta_{t_{\min}}/\Delta_{t_c}}\,y_{t_c}$. In either case the endpoint is $x^\mu$ up to a displacement of per-coordinate variance $\Delta_{t_{\min}}$ (plus the mean shrinkage $e^{-t_{\min}}x^\mu-x^\mu$). Consequently $\varrho(x_{t_{\min}})\approx\sqrt{\dlat\Delta_{t_{\min}}}/d_{\rm NN}(\mu)$ and the reported fraction cannot exceed the \emph{metric ceiling}
\begin{equation}\label{eq:ceiling}
 F^{\max}_{\rm mem}=1-G_{\rm NN}\big(3\sqrt{\dlat\,\Delta_{\rm term}}\big)\qquad(\text{SDE sampler stopped at }t_{\min},\ \dlat\gg1),
\end{equation}
with $\Delta_{\rm term}=\Delta_{t_{\min}}$ for the continuous-time sampler. [R] For the released Euler--Maruyama loop ($500$ steps on $[0.01,3]$, $dt=0.00598$, last-step contraction $0.63$) the terminal per-coordinate variance of the condensed linear score is $\Delta_{\rm term}=0.0269=1.36\,\Delta_{t_{\min}}$ (exact variance recursion and a $4\times10^5$-sample simulation; $1.17\times$ at $1000$ steps, $1.03\times$ at $5000$): discretisation \emph{raises} the terminal variance and lowers the ceiling. At $\dlat\lesssim40$ the $\chi^2_{\dlat}$ spread of $\|\delta\|^2$ (relative s.d.\ $\sqrt{2/\dlat}$) is not negligible and \eqref{eq:ceiling} must be evaluated by Monte Carlo over the endpoint law rather than in the concentrated form.
\end{lemma}
\begin{proof}
The reverse SDE with the exact score transports the time-$t_c$ marginal of the forward VP process started at $\rho^{\rm emp}_0$ to its time-$t_{\min}$ marginal \citep{song2021scoresde}, which is the stated mixture; under condensation $s_{\rm emp}(x,t)=-(x-e^{-t}x^\mu)/\Delta_t$ is the score of the single component $\mathcal N(e^{-t}x^\mu,\Delta_tI)$. For the ODE, $\frac{dx}{d\tau}=x+s_{\rm emp}$ and $\frac{d}{d\tau}(e^{-t}x^\mu)=e^{-t}x^\mu$ give $\frac{dy}{d\tau}=-\frac{e^{-2t}}{\Delta_t}y$, and $\int_{t_{\min}}^{t_c}\frac{e^{-2t}}{1-e^{-2t}}dt=\tfrac12\log(\Delta_{t_c}/\Delta_{t_{\min}})$. For \eqref{eq:ceiling}: $\|x_{t_{\min}}-x^\mu\|^2=\Delta_{\rm term}\|\eta\|^2+O(t_{\min}^2\|x^\mu\|^2)+$ a cross term of s.d.\ $2t_{\min}\sqrt{\Delta_{\rm term}}\|x^\mu\|$, concentrating at $\dlat\Delta_{\rm term}(1+O(\dlat^{-1/2}))$; when this is below $d_{\rm NN}(\mu)^2/9$ we have $\mathrm{NN}_1=x^\mu$ and $\|\mathrm{NN}_1-\mathrm{NN}_2\|=d_{\rm NN}(\mu)$, so $\varrho<1/3$ iff $d_{\rm NN}(\mu)>3\sqrt{\dlat\Delta_{\rm term}}$. Averaging over $\mu$ (the sampler lands on training points with frequency $1/\nsamp$ on average) gives \eqref{eq:ceiling}. The E--M variance is the fixed point of the offset recursion $v_{k+1}=a_k^2v_k+2\,dt$, $a_k=1-(1+e^{-2t_k})dt/\Delta_{t_k}$, evaluated along the code's grid.
\end{proof}

\begin{proposition}[Bridge: residual error acts as an effective terminal noise level; sketched]\label{prop:bridge}
Fix a diffusion time $t\in[t_{\min},t_c]$ at which the sampler makes its final attachment to a training point, and a training time $T$. Assume
\begin{enumerate}
\item[(B1)] (\emph{condensation}) $e^{2t}\Delta_t<d_{\rm NN}(\mu)^2/(2\log\nsamp)$ for the training points $\mu$ under consideration;
\item[(B2)] (\emph{concentrated residual}) [R] conditionally on $\mu$, $\|\zeta_T(x_t,t)\|^2$ with $x_t=e^{-t}x^\mu+\sqrt{\Delta_t}\eta$ concentrates at $\mathcal E_{t,\mu}(T)$ over $\eta$, and its cross term with the terminal noise of \cref{lem:endpoint} is negligible (neither Gaussianity nor isotropy of $\zeta_T$ is used for \eqref{eq:mem-criterion}; the isotropic form is used only for the distribution of $\|\delta\|^2$);
\item[(B3)] (\emph{one-attachment reduction}) the endpoint of the learned-score sampler equals the Tweedie denoiser of the learned score applied at the attachment time $t$, plus the sampler's terminal noise of \cref{lem:endpoint}.
\end{enumerate}
Then the generated point is $\hat x=x^\mu+\delta$ with $\E\|\delta\|^2=\dlat\Delta_{\rm eff,\mu}(T)$,
\begin{equation}\label{eq:Deff}
 \Delta_{\rm eff,\mu}(T)=\Delta_{\rm term}+c_t\,\mathcal E_{t,\mu}(T),\qquad c_t=\frac{e^{2t}\Delta_t^{2}}{\dlat},
\end{equation}
so the endpoint has the law of the \emph{exact} empirical sampler stopped at noise level $\Delta_{\rm eff,\mu}(T)$. Since $\|\delta\|^2=\dlat\Delta_{\rm eff,\mu}(T)(1+O(\dlat^{-1/2}))$, the sample attached to $\mu$ is memorized iff
\begin{equation}\label{eq:mem-criterion}
 \dlat\,\Delta_{\rm eff,\mu}(T)<\tfrac19\,d_{\rm NN}(\mu)^2\quad\Longleftrightarrow\quad \mathcal E_{t,\mu}(T)<\frac{d_{\rm NN}(\mu)^2-9\dlat\Delta_{\rm term}}{9\,e^{2t}\Delta_t^{2}} ,
\end{equation}
and
\begin{equation}\label{eq:Fmem}
 F_{\rm mem}(T)=\frac1\nsamp\sum_{\mu=1}^{\nsamp}\mathbf 1\Big[\,d_{\rm NN}(\mu)>3\sqrt{\dlat\Delta_{\rm term}+e^{2t}\Delta_t^{2}\,\mathcal E_{t,\mu}(T)}\,\Big]\;\xrightarrow{\ \mathcal E_{t,\mu}\equiv\mathcal E_t\ }\;1-G_{\rm NN}\Big(3\sqrt{\dlat\Delta_{\rm term}+e^{2t}\Delta_t^{2}\,\mathcal E_t(T)}\Big).
\end{equation}
The alternative candidate definition ``condensation at the effective level'', $e^{2t}\Delta_{\rm eff,\mu}(T)<d_{\rm NN}(\mu)^2/(2\log\nsamp)$, is implied by \eqref{eq:mem-criterion} whenever $\dlat>\tfrac29e^{2t}\log\nsamp$ [R: $\dlat>10.6$ at $t=t_c(200)=1.02$, $\nsamp=500$]: the metric threshold, not condensation, binds. [R] The constant $c_t$ depends on the unstated attachment time $t$: $e^{2t}\Delta_t$ spans $0.020$ ($t=t_{\min}$) to $6.7$ ($t=1.02$), so the proposition fixes the functional form $\Delta_{\rm eff}=\Delta_{\rm term}+c\,\mathcal E$ but not $c$ (Open gaps, G2).
\end{proposition}
\begin{proof}[Sketch]
Under (B1) the $\eta=0$ form of \cref{lem:condensation} gives $w_\mu>\tfrac12$ at the centre; the sketch assumes the stronger conclusion $s_{\rm emp}(x_t,t)=-(x_t-e^{-t}x^\mu)/\Delta_t$ up to a negligible $r$, which \cref{lem:condensation} guarantees only under \eqref{eq:cond-rigorous} (Open gaps, G8). Tweedie's identity \citep{efron2011tweedie} for the learned score gives $\hat x=e^{t}\big(x_t+\Delta_t s_T(x_t,t)\big)=x^\mu+e^{t}\Delta_t\,\zeta_T(x_t,t)$, exactly, so $\|\hat x-x^\mu\|^2=e^{2t}\Delta_t^2\|\zeta_T(x_t,t)\|^2$; by (B2) this is $e^{2t}\Delta_t^2\mathcal E_{t,\mu}(T)$ up to concentration, and by (B3) it adds to the terminal variance of \cref{lem:endpoint}, giving \eqref{eq:Deff}. $\chi^2_{\dlat}$ concentration of $\|\delta\|^2$ and the identification $\mathrm{NN}_1=x^\mu$, $\|\mathrm{NN}_1-\mathrm{NN}_2\|=d_{\rm NN}(\mu)$ (valid when $\|\delta\|<d_{\rm NN}(\mu)/2$) give \eqref{eq:mem-criterion}; averaging over the attachment index gives \eqref{eq:Fmem}. Comparison of thresholds: \eqref{eq:mem-criterion} reads $\Delta_{\rm eff}<d_{\rm NN}^2/(9\dlat)$, condensation reads $\Delta_{\rm eff}<e^{-2t}d_{\rm NN}^2/(2\log\nsamp)$.

\emph{When (B2) is justified.} At $T\gg\taugen$ the signal and noise-dim modes are absorbed and $\zeta_T(x)=-(A^\star-A(T))\phi(x)$ restricted to sample modes. The sample-bulk eigendirections are the feature-space images of the per-sample deviations $\xi^\mu$ with the linear (population) part removed; after the random rotation $Q$ their output directions in $\R^{\dlat}$ are exchangeable across coordinates, and the residual at $x_t$ is a sum over the $N_{\rm res}(T)=\#\{i\in\mathcal M:\lambda_iT\lesssim1\}$ unabsorbed modes. When $N_{\rm res}(T)\gg1$ and $\dlat\gg1$ a CLT over these generic directions makes $\|\zeta_T\|^2$ concentrate. The approximation fails (i) early, when the residual is dominated by the $\dint$ signal modes (then $\mathcal E_t$ is far above the threshold and $F_{\rm mem}=0$ regardless), and (ii) at the very end, when only a few sample modes remain: the residual is then low-rank and localized on a few training points, and \eqref{eq:Fmem} must be read with the per-sample $\mathcal E_{t,\mu}$. [R] In the ``blend'' model $\hat x_T=(1-m)\hat x_{\rm pop}+m\,\hat x_{\rm emp}$ the Tweedie identity holds equally (on the null block $c_t\mathcal E_{t,\mu}=(1-m)^2\Delta_t\signoise^2/\beta_t^2$ per coordinate, the same as $\|\delta\|^2/\dlat$ computed directly), so the two readings differ only through the choice of $t$ and the multi-step accumulation (G2, G3), not through the form of the residual.
\end{proof}

\begin{corollary}[Fraction curve and quantile-indexed memorization time; conditional on \cref{prop:bridge}]\label{cor:taumem-fraction}
Under \cref{prop:bridge} with a $\mu$-independent residual $\mathcal E_t(T)$ decreasing in $T$, the fraction first reaches level $q\in(0,F^{\max}_{\rm mem})$ at
\begin{equation}\label{eq:taumem-q}
 \taumem(q)=\inf\Big\{T:\ \mathcal E_t(T)\le \frac{\big(d_{\rm NN}^{(1-q)}\big)^2-9\dlat\Delta_{\rm term}}{9\,e^{2t}\Delta_t^{2}}\Big\}:
\end{equation}
the most isolated training points are memorized first. If $\mathcal E_t(T)=\mathcal E^{\rm samp}_t(0)e^{-2\lambda_{\rm samp}T}$ then
\begin{equation}\label{eq:taumem-spread}
 \taumem(q_2)-\taumem(q_1)=\frac{1}{\lambda_{\rm samp}}\log\frac{d_{\rm NN}^{(1-q_1)}}{d_{\rm NN}^{(1-q_2)}}+O\!\Big(\frac{\dlat\Delta_{\rm term}}{\lambda_{\rm samp}(d^{(1-q)}_{\rm NN})^2}\Big),
\end{equation}
a difference in \emph{linear} $T$. [R] Under this single-exponential, data-spread-only model, $T(5\%)/T(1\%)=1+\log(d^{(0.99)}_{\rm NN}/d^{(0.95)}_{\rm NN})/(\lambda_{\rm samp}T(1\%))\approx1+0.15/(\lambda_{\rm samp}T(1\%))$, so the measured ratios $2.2$--$3.5$ (anchor (6)) already refute a data-level-only width: the within-bulk spread of $\lambda_i$ (\cref{app:theory-predictor}) and, at small $\dlat$, the $\chi^2_{\dlat}$ spread of $\|\delta\|^2$ are required. A single edge with a degenerate $G_{\rm NN}$ would give a step function.
\end{corollary}
\begin{proof}
$1-G_{\rm NN}(u)\ge q$ iff $u\le d_{\rm NN}^{(1-q)}$; substitute $u=3\sqrt{\dlat\Delta_{\rm term}+e^{2t}\Delta_t^2\mathcal E_t(T)}$ and use monotonicity. For \eqref{eq:taumem-spread} take logarithms of $3e^{t}\Delta_t\sqrt{\mathcal E^{\rm samp}_t(0)}\,e^{-\lambda_{\rm samp}T}=\sqrt{(d^{(1-q)}_{\rm NN})^2-9\dlat\Delta_{\rm term}}$ and expand in $9\dlat\Delta_{\rm term}/(d^{(1-q)}_{\rm NN})^2$.
\end{proof}

\begin{corollary}[What $\taumem$ depends on through $d_{\rm NN}$; sketched]\label{cor:dnn-scaling}
For the synthetic data of \cref{sec:design}, with labels uniform at random and the $\nsamp/k$-per-cluster Poisson heuristic for the signal part,
\begin{equation}\label{eq:dnn-scaling}
 d_{\rm NN}(\mu)^2\;\le\;\underbrace{C_{\dint}\,\sigsig^2\,(\nsamp/k)^{-2/\dint}\,\kappa_\mu}_{\text{signal part}}\;+\;\underbrace{2\signoise^2(\dlat-\dint)\big(1-O(\sqrt{\log\nsamp/(\dlat-\dint)})\big)}_{\text{null part}},
\end{equation}
with $C_{\dint}$ depending only on $\dint$ and $\kappa_\mu$ an $O(1)$ random factor with a fixed law. [R] The right side is $\min_\nu S_\nu+\min_\nu N_\nu$, a \emph{lower} bound on $\min_\nu(S_\nu+N_\nu)$ only when the extreme-value correction is dropped; on the repository generator (seed 42, $\signoise=0.5$) the median $d_{\rm NN}^2$ is $6.60/14.72/40.0/85.0$ at $\dlat=20/40/100/200$ against the uncorrected sum $8.61/18.61/48.61/98.61$, i.e.\ $15$--$30\%$ lower, and the null part alone already exceeds the measured median. Hence: (i) if $(\dlat-\dint)\signoise^2\ll\sigsig^2(\nsamp/k)^{-2/\dint}$ (e.g.\ $\signoise=0.01$) the threshold in \eqref{eq:taumem-q} scales as $\nsamp^{-2/\dint}$ and is independent of $\dlat$: all $\dlat$-dependence of $\taumem$ enters through $\mathcal E_t(T)$, and since a larger $\nsamp$ lowers the threshold, $\partial\log\taumem/\partial\log\nsamp$ acquires the additive term $+\frac{1}{\dint\lambda_{\rm samp}\taumem}$ [R: sign corrected; denser data memorizes \emph{later}] beyond the spectral $\nsamp$-dependence, a $\dint$, not $\dlat$, exponent; (ii) if $(\dlat-\dint)\signoise^2\gg\sigsig^2(\nsamp/k)^{-2/\dint}$ (e.g.\ $\signoise=0.5$, $\dlat\gtrsim20$), $d_{\rm NN}^2\simeq2\signoise^2\dlat\times(0.7$--$0.85)$, \eqref{eq:mem-criterion} becomes approximately a per-coordinate condition on $\Delta_{\rm eff}$, the spread $d^{(0.9)}_{\rm NN}/d^{(0.1)}_{\rm NN}\to1$, and the same metric drives the NN-ratio of non-memorized samples toward $1$; (iii) the ceiling \eqref{eq:ceiling} collapses with $\dlat$ in regime (i) and tends to $1$ in regime (ii).
\end{corollary}
\begin{proof}
Within a cluster the signal coordinates of the $\approx\nsamp/k$ points are i.i.d.\ $\mathcal N(0,\sigsig^2I_{\dint})$; the NN distance of $m$ i.i.d.\ points from a smooth density in $\dint$ dimensions scales as $m^{-1/\dint}$ with a limiting law for the rescaled distance \citep{levina2004maximum,facco2017twonn}. The null coordinates of two distinct points are independent $\mathcal N(0,\signoise^2I_{\dlat-\dint})$, so their squared distance is $2\signoise^2\chi^2_{\dlat-\dint}$; the minimum over $\nsamp-1$ such variables is $2\signoise^2(\dlat-\dint)(1-O(\sqrt{\log\nsamp/(\dlat-\dint)}))$ by Gaussian concentration and a union bound. The minimum of a sum is at least the sum of the minima; equality fails because the minimising $\nu$ differ, which is the observed $15$--$30\%$ shortfall.
\end{proof}

\begin{remark}[Real data: latent versus pixel space]\label{rem:pixel}
For VAE latents the sampler runs in $\R^{\dlat}$ but $\varrho$ is computed in pixel space after decoding with $D:\R^{\dlat}\to\R^{P}$. If $\ell\|z-z'\|\le\|D(z)-D(z')\|\le L\|z-z'\|$ on the relevant region, then $(\ell/L)\varrho_{\rm lat}\le\varrho_{\rm pix}\le(L/\ell)\varrho_{\rm lat}$ (assuming the pixel-space $\mathrm{NN}_1,\mathrm{NN}_2$ are the decoded latent ones; otherwise the denominator needs a global minimum over training pairs). [R] This two-sided sandwich imposes no ordering between the two fractions at the common threshold $\tfrac13$: it gives only $F_{\rm lat}(\ell/(3L))\le F_{\rm pix}(\tfrac13)\le F_{\rm lat}(L/(3\ell))$. Two consequences survive: (a) $L/\ell$ is not controlled across the sweep (one VAE per $\dlat$), so pixel-space fractions at different $\dlat$ are compared at effectively different latent thresholds; (b) posterior-collapsed coordinates ($\lambda_{\min}=\Delta_t$ exactly for $\dlat\ge140$ on CelebA, \cref{sec:real-data}) contribute to $\dlat\Delta_{\rm eff}$ in latent space but are annihilated by $D$, so in regime (ii) the relevant dimension is the decoder-visible one ($\approx0.5\dlat$, the measured effective rank), not $\dlat$.
\end{remark}

\paragraph{Empirical anchors (synthetic data, $\nsamp=500$, $k=10$, $\dint=5$, $s=3$, $\sigsig=1$; five replicate training sets).}
(1) Metric ceiling \eqref{eq:ceiling} in the concentrated form with $\Delta_{\rm term}=\Delta_{t_{\min}}$ at $\signoise=0.5$: $0.676,\ 0.790,\ 0.870,\ 0.924,\ 0.950,\ 0.979,\ 0.998$ for $\dlat=5,8,10,12,15,20,25$ and $\approx1$ beyond. [R] With the code's E--M terminal variance and metric, the exact empirical-score sampler run in the code's own loop gives fractions $0.508/0.597/0.667/0.715/0.757/0.808/0.873/0.960$ at $\dlat=5/8/10/12/15/20/25/40$; these, not the concentrated ceilings, are the operative bounds. Measured final fractions of the 5M-step MLP sweep (seed 42): $0.290,0.224,0.168,0.127,0.079,0.044,0.017,0.006,0.004,0.002$ for $\dlat=5,\dots,40$---all below both. At $\signoise=0.01$ the concentrated ceiling is $0.668,0.364,0.250,0.165,0.100,0.034,0.015,0.008,0.002$ for the same $\dlat$, but the exact endpoint-law Monte Carlo gives $0.478,0.349,\dots,0.095$ ($\dlat=8,10,\dots,20$), $0.050$ ($25$), $0.008$ ($40$), and the E--M loop $0.531/0.235/0.041/0.001$ at $\dlat=5/10/20/40$; the available $\signoise=0.01$ runs (older statistic, $300$k steps) report fractions $\le0.006$ and exactly $0$ for $\dlat\ge20$, which is consistent with but not a like-for-like test of the ceiling. (2) $d_{\rm NN}$ deciles $(q_{10}/q_{50}/q_{90})$ at $\signoise=0.5$: $0.72/1.10/1.63$ ($\dlat=5$), $2.13/2.56/3.00$ ($20$), $8.75/9.16/9.56$ ($200$); at $\signoise=0.01$: $0.70/1.07/1.60$, flat in $\dlat$. (3) $\mathrm{median}(d_{\rm NN})\cdot\nsamp^{1/5}=3.66,3.72,3.77,3.76$ for $\nsamp=100,500,1000,5000$ at $\signoise=0.01$. (4) Condensation times of the \emph{exact} empirical score in the $\eta=0$ form (median $d_{\rm NN}$): $t_c=0.046,0.105,0.21,0.39,0.72,1.02$ for $\dlat=5,10,20,40,100,200$ at $\signoise=0.5$, and $0.045$ flat at $\signoise=0.01$. [R] Typical-$\eta$ condensation (median off-attachment mass $1-w_\mu<0.05$) occurs later: $0.053/0.19/0.315/0.757$ at $\dlat=5/20/40/200$; at the $\eta=0$ value $t_c=1.02$ ($\dlat=200$) the median off-attachment mass is $0.83$, i.e.\ not condensed. Either way the exact empirical measure condenses \emph{later in $t$} (earlier in the reverse process) as $\dlat$ grows; the observed delay of memorisation is a property of the learned residual $\mathcal E_t(T)$, not of the data. (5) Mean NN-ratio of generated samples $0.524$ ($\dlat=5$) $\to0.868$ ($40$) follows the training-set self-ratio trend $0.83\to0.96$ (distance concentration, regime (ii)). (6) Crossing times $T(1\%)/T(5\%)/T(10\%)$ in steps: $150$k$/350$k$/550$k ($\dlat=5$), $250$k$/550$k$/1.0$M ($8$), $300$k$/700$k$/1.7$M ($10$), $400$k$/1.0$M$/2.65$M ($12$), $600$k$/2.1$M$/$-- ($15$), $1.05$M ($20$), $2.55$M ($25$).

\paragraph{Predictions and how to test them.}
\begin{enumerate}
\item \emph{One-constant fraction prediction (synthetic):} $F_{\rm mem}(T)=1-G_{\rm NN}(3\sqrt{\dlat\Delta_{\rm term}+c\,\mathcal E_t(T)})$ with $\mathcal E_t(T)$ measured against the closed form \eqref{eq:semp} and $G_{\rm NN}$ from the training set; $c$ is fixed by the attachment time (G2) and is the one constant. Test: requires re-running training with checkpoints (none are saved); curves for different $\dlat$ should collapse; systematic over-prediction indicates (B2)/(B3) fail.
\item \emph{Metric ceiling at $\signoise=0.01$:} the reported fraction is bounded by the E--M endpoint-law ceiling ($0.04$--$0.10$ at $\dlat=20$, $<0.01$ at $\dlat=40$) regardless of training time. Test: (a) train the $\signoise=0.01$ sweep to $5$M steps and compare peak fractions to the ceilings; (b) sharper: add a terminal Tweedie denoising step (or lower $t_{\min}$ to $10^{-3}$); the $\signoise=0.01$ fractions should then jump for $\dlat\in[8,30]$ while the $\signoise=0.5$ ones barely change.
\item \emph{$\dint$-not-$\dlat$ scaling of the data-side threshold:} $d_{\rm NN}\nsamp^{1/\dint}$ is constant in the signal-dominated regime, so at $\signoise\to0$ the threshold scales as $\nsamp^{-2/\dint}$ and is $\dlat$-independent; denser data memorises later by the extra term $+1/(\dint\lambda_{\rm samp}\taumem)$ in $\partial\log\taumem/\partial\log\nsamp$. Test: sweep $\nsamp\in\{100,500,1000,5000\}$ and $\dint\in\{3,5,8\}$ at $\signoise=0.01$.
\item \emph{Fraction-curve width:} the data-level spread $d^{(0.9)}_{\rm NN}/d^{(0.1)}_{\rm NN}$ falls $2.28\to1.09$ over $\dlat=5\to200$ at $\signoise=0.5$ and stays at $2.3$ at $\signoise=0.01$, but the measured $T(5\%)/T(1\%)=2.2$--$3.5$ cannot come from it; the width is model-level (eigenvalue spread) and grows with $\dlat$. Test: extract $\taumem(q)$ for $q\in\{1,2,5\%\}$, normalise by the RFNN sample-bulk $1/\lambda$ and compare the residual spread with $\log(d^{(0.99)}_{\rm NN}/d^{(0.95)}_{\rm NN})$.
\item \emph{The exact empirical score memorises at the ceiling for every $\dlat$}, the opposite of the learned-score trend. Test: run the closed-form $s_{\rm emp}$ sampler for each $\dlat$ ([R] done by a referee: $0.508\to0.960$ over $\dlat=5\to40$ at $\signoise=0.5$).
\end{enumerate}

\paragraph{Open gaps.}
(G1) \emph{Residual model.} (B2) is stated as a concentration hypothesis with its justification and failure windows, not proved. (G2) \emph{Attachment time and multi-step accumulation.} $c_t$ is not determined by the proposition: numerical integration of the offset equation with a persistent relative error $r^2=\mathcal E_s\Delta_s/\dlat$ gives per-coordinate displacement variance (units $\Delta_{t_{\min}}r^2$) $1.02$ for Tweedie at $t_{\min}$, $4.9/26/338$ for Tweedie at $t_c=0.046/0.21/1.02$, $0.61/2.64/8.42$ ODE-integrated, and $1.18/2.58/3.57$ SDE-integrated (the sampler actually used); the constant is $\dlat$-dependent through $t_c$ and the earlier ``accurate to a logarithmic factor'' understates the spread. (G3) The RFNN spectrum is at $t=0.01$ while attachment occurs near $t_c$; relating $\mathcal E_{t_c}(T)$ to the fixed-$t$ spectrum is not done. (G4) $\mathrm{NN}_1$ equals the attachment point only for $\|\delta\|<d_{\rm NN}/2$; the non-memorized tail of $\varrho$ is not modelled. (G5) The projection floor $\mathcal E_t(\infty)>0$ is not computed; the return of memorisation as $\dlat\to\nsamp$ is outside the bridge (\cref{app:theory-scope}). (G6) The per-sample residual $\mathcal E_{t,\mu}$ is not derived from the sample-bulk modes; isolated points may carry the least-supported, smallest-$\lambda_i$ modes, partially cancelling ``most isolated first''. (G7) $L/\ell$ of the decoders is unknown. (G8) \cref{lem:condensation} is rigorous with $32\log(\nsamp/\delta)$; that condition is never met on the sampler interval where memorisation is largest ($e^{-t}d_{\rm NN}/\sqrt{\Delta_t}=7.4$ vs the required $16.5$ at $\dlat=5$, $t=t_{\min}$), so the sketch's step ``(B1) $\Rightarrow$ $r$ negligible'' is not a consequence of the lemma; the $\eta=0$ form is not a typical-$\eta$ criterion ($\E_\eta[w_\nu/w_\mu]=1$ identically), and a typical-$\eta$ criterion needs $e^{-2t}d_{\rm NN}^2/(2\Delta_t)\gtrsim5.8\log\nsamp$. (G9) Both sweeps use the same NN-ratio statistic (\cref{def:bridge-objects}), but the older $\signoise=0.01$ sweep runs on the step clock $T=0.01k$ and the 5M-step sweep on $T=0.01\dlat k$ (\eqref{eq:lin-steps}), so their anchors are not like-for-like in training time. (G10) [R] The E--M terminal variance is $1.36\Delta_{t_{\min}}$, not $\Delta_{t_{\min}}$, and the $\chi^2_{\dlat}$ spread is $O(1)$ at $\dlat\le40$; concentrated ceilings quoted to three digits are not meaningful there. (G11) [R] The zero-parameter test (Prediction~1) is not executable from the repository: no score-model checkpoints are saved. (G12) The label \texttt{prop:bridge} must be issued once; the main text should cite this statement rather than restate it.

% ----------------------------------------------------------------------
\subsection{From the frozen-latent spectrum to a memorisation-fraction predictor}
\label{app:theory-predictor}

Throughout, $\phi(x)=\pwidth^{-1/2}\tanh(Wx/\sqrt{\dlat})$ with $W_{ij}\sim\mathcal N(0,1)$, $x_t=e^{-t}x+\sqrt{\Delta_t}\,\eta$, $M_t=\frac1\nsamp\sum_\mu\E[x_t^\mu x_t^{\mu\top}]$, $q_d=\mathrm{tr}(M_t)/\dlat$, and $\Umat=\frac1\nsamp\sum_\mu\E_\xi[\phi(x_t^\mu)\phi(x_t^\mu)^\top]$; ``code units'' means $\pwidth\times$ theory units. This subsection derives what the released spectral predictor (\cref{sec:spectral-predictor}) actually measures, supplies the sample-bulk term it lacks, and resolves the sample bulk per training point, which is what a fraction-level prediction needs (\cref{def:taumem-p1}(iii)).

\subsubsection{Gaussian-equivalence coefficients}

\begin{lemma}[proved]\label{lem:ge-coeff-pred}
Let $\sigma=\tanh$, $z\sim\mathcal N(0,1)$, $q>0$, $a_1(q)=\E[\sigma'(\sqrt q z)]$, $a_*(q)^2=\E[\sigma(\sqrt q z)^2]-a_1(q)^2q$, and $\sigma(\sqrt qz)=\sum_{k\ \mathrm{odd}}c_k(q)\mathrm{He}_k(z)/k!$ (unnormalised Hermite coefficients; $c_k$ here equals $\sqrt{k!}$ times the $c_k$ of \cref{lem:fourbulk-kernel}). Then (i) $\E[\sigma(\sqrt qz)z]=\sqrt q\,a_1(q)$; (ii) $a_*(q)^2=\E[(\sigma(\sqrt qz)-a_1(q)\sqrt qz)^2]\ge0$; (iii) $a_1^2q=c_1^2$ and $a_*^2=\sum_{k\ \mathrm{odd}\ge3}c_k^2/k!$; (iv) $a_*^2=\tfrac23q^3+O(q^4)$, increasing on $[0.05,3.5]$, with $a_*^2=0.0047,\,0.0274,\,0.0359,\,0.0544,\,0.0784$ and $a_1^2=0.623,\,0.367,\,0.320,\,0.246,\,0.180$ at $q=0.33,\,1,\,1.247,\,1.834,\,2.76$; (v) for jointly Gaussian $(z,z')$ with correlation $\varrho\in[0,1]$,
\[ g(q,\varrho):=\E[\sigma(\sqrt qz)\sigma(\sqrt qz')]-a_1(q)^2q\varrho=\sum_{k\ \mathrm{odd}\ge3}\frac{c_k(q)^2}{k!}\varrho^k\in[0,\varrho^3a_*(q)^2], \]
increasing in $\varrho$ with $g(q,1)=a_*^2$. With $r_t(q):=g(q,1-\Delta_t/q)/a_*(q)^2$ (the $\theta(q)$ of \cref{lem:coeffs-buf}(v)): at $t=0.1$, $r_t=0.086,\,0.505,\,0.574,\,0.675,\,0.755$ at $q=0.33,\,1,\,1.247,\,1.834,\,2.76$; at $t=0.01$, $r_t=0.822,\,0.943,\,0.969$ at $q=0.33,\,1.247,\,2.76$. $g(q,1-\Delta_t/q)$ is increasing in $q$ on this range for both $t$.
\end{lemma}
\begin{proof}
(i) Stein's lemma. (ii) Expand and use (i). (iii) $\tanh$ is odd; orthogonality with $\E\mathrm{He}_k^2=k!$. (iv) $\E\sigma^2=q-2q^2+\tfrac{17}3q^3+O(q^4)$ and $a_1^2q=q-2q^2+5q^3+O(q^4)$; the table is 120-node Gauss--Hermite quadrature and reproduces the $0.180\to0.625$ range of $a_1^2$ of \cref{lem:q-p1}. (v) Mehler's formula $\E[\mathrm{He}_k(z)\mathrm{He}_l(z')]=\delta_{kl}k!\varrho^k$; the subtraction removes $k=1$, oddness the even terms; each term is nonnegative and increasing in $\varrho$, and $\varrho^k\le\varrho^3$. The $r_t$ values agree between two-dimensional quadrature and the Hermite series truncated at $k=39$ to four digits.
\end{proof}

\subsubsection{Mode clocks under gradient flow}

\begin{proposition}[proved]\label{prop:mode-clocks-pred}
Let $L(A)=\frac1{2\nsamp}\sum_\mu\E_\xi\|A\phi(x_t^\mu)-y_\mu(x_t^\mu)\|^2$ for square-integrable targets $y_\mu$, $\Umat=\sum_i\lambda_iu_iu_i^\top$, $V=\frac1\nsamp\sum_\mu\E_\xi[y_\mu\phi^\top]$, and let $\dot A=-\nabla L(A)$, $A(0)=0$, $T=\eta s$ ($s$ = steps). With $a_i(T)=A(T)u_i$ and $a_i^*=Vu_i/\lambda_i$ ($\lambda_i>0$),
\[ a_i(T)=a_i^*\,(1-e^{-\lambda_iT}),\qquad L(T)-L(\infty)=\tfrac12\sum_i\lambda_i\|a_i^*\|^2e^{-2\lambda_iT}. \]
Mode $i$ thus has absorbed amplitude $1-e^{-\eta\lambda_is}$, absorbed loss $1-e^{-2\eta\lambda_is}$ and learned-energy fraction $(1-e^{-\eta\lambda_is})^2$, ordered $(1-e^{-x})^2\le1-e^{-x}\le1-e^{-2x}$. The exponent is normalisation-invariant: in code units $T=\eta_{\rm code}\pwidth s$ and $\eta\lambda s=\eta_{\rm code}\lambda_{\rm code}s$. If $u$ is a unit eigenvector localised on training point $\mu$, $u\propto\bar\rho_\mu=\E_\xi\rho(x_t^\mu)$ with eigenvalue $\Lambda_\mu=\|\bar\rho_\mu\|^2/\nsamp$ (\cref{prop:sample-bulk-persample-pred}), and $\bar y_\mu$ is the sample-specific target, then at leading order $a^*=\bar y_\mu/\|\bar\rho_\mu\|$ and the model output at $x_t^\mu$ along $u$ is $\bar y_\mu(1-e^{-\Lambda_\mu T})$: the amplitude fraction is the fraction of the sample-specific score output produced at the training point.
\end{proposition}
\begin{proof}
$\nabla_AL=A\Umat-V$; project on $u_i$ and integrate (this is \cref{lem:gf-p1}). The loss identity is the quadratic expansion about $A^*=V\Umat^+$. Units: substitute $A_{\rm code}=A/\sqrt{\pwidth}$, $\Umat_{\rm code}=\pwidth\Umat$. Meaning: only the $\nu=\mu$ term of $V$ correlates with $u$ at leading order, so $Vu=\bar y_\mu\bar\rho_\mu^\top\bar\rho_\mu/(\nsamp\|\bar\rho_\mu\|)=\bar y_\mu\|\bar\rho_\mu\|/\nsamp$ and $a(T)\,u^\top\bar\phi_\mu=a^*(1-e^{-\Lambda_\mu T})\|\bar\rho_\mu\|$. For the RFNN, $L$ with fresh noise at every step is exactly the $\E_\xi$-averaged loss; full-batch gradient descent agrees with the flow to $O(\eta\lambda_{\max})$ per step.
\end{proof}

\begin{remark}[The square is not derived]\label{rem:square-pred}
The released clock $P_d(s)=\mathrm{mean}_i(1-e^{-\kappa\lambda_i(M_t)s})^2$ is the learned-energy fraction averaged over the eigenvalues of $M_t$ rather than of $\Umat$. For a single mode the three fractions are monotone reparametrisations of $\eta\lambda_is$, so the exponent is absorbed into the threshold; it matters only through the average over modes, where the square suppresses slow modes at early times quadratically. A derivation would have to show that the NN-ratio memorisation statistic is a quadratic functional of the learned sample-specific displacement at small absorbed amplitude; we do not have one, and use the amplitude clock, which has the output-space meaning above.
\end{remark}

\subsubsection{Data-side rates: what $\mathrm{spec}(M_t)$ labels}

\begin{proposition}[sketched]\label{prop:data-side-rates-pred}
Let $W_{ij}\sim\mathcal N(0,\sigma_w^2)$, $\gamma:=\dlat\sigma_w^2$ ($\gamma=1$ for the RFNN), $\pwidth/\dlat\to\infty$, and assume Gaussian equivalence of the tanh random-feature kernel for the input distribution \citep{pennington2017nonlinear,benigni2021eigenvalue,hu2022universality,goldt2022gaussian}. Then the $\dlat$ data-side eigenvalues of $\Umat$ are $\lambda_i(\Umat)=\gamma a_1(\gamma q_d)^2\lambda_i(M_t)/\dlat\,(1+o(1))$, $i\le\dlat$, with trace $a_1(\gamma q_d)^2\gamma q_d$. The data-side amplitude clock is $1-\exp[-\eta\gamma a_1(\gamma q_d)^2\lambda_i(M_t)s/\dlat]$, so the released $\kappa$ equals $\eta\gamma a_1(\gamma q_d)^2/\dlat$ in theory units; in code units on the $\pwidth=64\dlat$ sweep the rate is $64\,\eta_{\rm code}a_1(q_d)^2\lambda_i(M_t)$, and the $\dlat$-dependence is carried by $a_1(q_d)^2$ alone ($0.180\to0.625$).
\end{proposition}
\begin{proof}[Proof sketch]
Decompose $\tanh(w^\top x/\sqrt{\dlat})=a_1(q_x)w^\top x/\sqrt{\dlat}+r_w(x)$, $q_x=\gamma\|x\|^2/\dlat$; by \cref{lem:ge-coeff-pred}(i) the residual is uncorrelated with the linear part over $w$, so the linear--residual cross block vanishes at leading order. The linear block is $\sigma_w^2W'\tilde M_tW'^\top/(\pwidth\dlat)$ with $\tilde M_t=\frac1\nsamp\sum_\mu a_1(q_{x_t^\mu})^2\E x_t^\mu x_t^{\mu\top}$; replacing $a_1(q_{x_t^\mu})$ by $a_1(\gamma q_d)$ is a heterogeneity correction of relative size $\mathrm{cv}(q_\mu)|a_1'|q/a_1$, and the nonzero spectrum of $W'\tilde M_tW'^\top/\pwidth$ is that of $\tilde M_t(W'^\top W'/\pwidth)\to\tilde M_t$. This is the spectral law of \citet{bonnaire2025} evaluated on the data-side block, and the block structure of \citet{george2025dsm} with $\signoise>0$ (\cref{prop:fourbulk}). Finite-$\pwidth$ corrections are the $O(\sqrt{\dlat/\pwidth})$ Marchenko--Pastur broadening of $W'^\top W'/\pwidth$ \citep{marchenko1967distribution}, consistent with the $<0.5\%$ median motion under a $5.8\times$ change of $\pwidth$. Trace: $\gamma a_1^2\mathrm{tr}M_t/\dlat=a_1^2\gamma q_d$.
\end{proof}

\begin{remark}[What a clock built from $\mathrm{spec}(M_t)$ measures]\label{rem:spec-Mt-misses-mem-pred}
The $\dlat$ eigenvalues of $M_t$ label the signal and noise-dimension bulks of $\Umat$, whose targets are population-score components. The $\nsamp-\dlat$ sample modes, which carry the sample-specific components and hence memorisation, are absent from $\mathrm{spec}(M_t)$. The released $P_d(s)$ therefore measures absorption of the generalising modes and tracks memorisation only insofar as their clock co-varies with the sample-mode clock, which happens through $q_d$ (\cref{cor:shape-from-qmu-pred}(iii)). The two clocks couple to posterior collapse through different channels: the sample-mode clock through $q$-dilution---the channel of the synthetic $\signoise\to0$ buffer---and the released clock through the collapsed eigenvalues $\lambda_i=\Delta_t=0.181$ themselves ($\dlat\ge140$ CelebA, $\ge160$ CIFAR-10; excess $<2\times10^{-3}$), its slowest modes, which dominate the late-time growth of $P_d(s)$ although their score target is $-x_i/\Delta_t$. CelebA's effective rank saturates ($89.5,\,90.3,\,87.9$ at $\dlat=160,180,200$) while memorisation keeps falling; a clock reading $\mathrm{spec}(M_t)$ reproduces this only through the growing number of collapsed coordinates, a channel carrying no data-side information, whereas the sample-mode clock reproduces it through $q_d=1.83\to1.25$.
\end{remark}

\subsubsection{Per-sample resolution of the sample bulk}

\begin{proposition}[sketched]\label{prop:sample-bulk-persample-pred}
In the setting of \cref{prop:data-side-rates-pred} with $\gamma=1$ and $\pwidth>\nsamp$, let $\rho(x)=\phi(x)-\pwidth^{-1/2}a_1(q_x)Wx/\sqrt{\dlat}$, $q_\mu=e^{-2t}\|x^\mu\|^2/\dlat+\Delta_t$, $\varrho_\mu=1-\Delta_t/q_\mu$, $\bar\rho_\mu=\E_\xi\rho(x_t^\mu)$. Distinguish the single-draw $\Umat^{(1)}=\frac1\nsamp\sum_\mu\phi(x_t^\mu)\phi(x_t^\mu)^\top$ from the noise-averaged $\Umat$ (the training loss; the $n_\xi\to\infty$ limit of \texttt{compute\_U}, which averages $n_\xi=50$ draws). Then:
\begin{enumerate}
\item $\E_W\|\rho(x)\|^2=a_*(q_x)^2$ exactly, with relative concentration $\pwidth^{-1/2}$, and $\E_{W,\xi}\|\rho(x_t^\mu)\|^2=a_*(q_\mu)^2(1+O(\dlat^{-1}))$.
\item For $\varrho_{xy}=x^\top y/(\dlat\sqrt{q_xq_y})$: $\E_W[\rho(x)^\top\rho(y)]=\sum_{k\ \mathrm{odd}\ge3}c_k(q_x)c_k(q_y)\varrho_{xy}^k/k!=O(\varrho_{xy}^3)$; two noise draws of the same point have overlap $\varrho_\mu$, so $\E_W\|\bar\rho_\mu\|^2=g(q_\mu,\varrho_\mu)(1+O(\dlat^{-1/2}))=r_t(q_\mu)a_*(q_\mu)^2(1+O(\dlat^{-1/2}))$.
\item $\Umat_*=\frac1\nsamp\sum_\mu\E_\xi[\rho\rho^\top]=\Umat_{\rm loc}+\Umat_{\rm cloud}$, $\Umat_{\rm loc}=\frac1\nsamp\sum_\mu\bar\rho_\mu\bar\rho_\mu^\top$, $\Umat_{\rm cloud}=\frac1\nsamp\sum_\mu\mathrm{Cov}_\xi(\rho_\mu)\succeq0$. To leading order in $\max_{\mu\ne\nu}|\varrho_{\mu\nu}|^3$ and $\sqrt{\nsamp/\pwidth}$, $\Umat_{\rm loc}$ has one eigenvalue per training point, eigenvector $\propto\bar\rho_\mu$,
\begin{equation}\label{eq:Lambda-mu-pred}
 \Lambda_\mu=\frac{r_t(q_\mu)\,a_*(q_\mu)^2}{\nsamp},
\end{equation}
and $\Umat_{\rm cloud}$ has trace $\approx(1-r_t(q_d))a_*(q_d)^2$ over up to $\pwidth$ directions (at finite $n_\xi$, at most $\min\{\pwidth,(n_\xi-1)\nsamp\}$), i.e.\ per-direction mass $O((1-r_t)a_*^2/\pwidth)\ll\Lambda_\mu$ for $\pwidth\gg\nsamp$. For $\Umat^{(1)}$ the same argument with $\varrho_\mu\to1$ gives $\Lambda^{(1)}_\mu=a_*(\hat q_\mu)^2/\nsamp$, no cloud and $\mathrm{rank}\le\nsamp$. Trace check for both objects: $a_1^2q_d+r_ta_*^2+(1-r_t)a_*^2=\E\tanh^2(\sqrt{q_d}z)$.
\item Corrections: Marchenko--Pastur spread of the $\Lambda_\mu$ with relative standard deviation $\varsigma=\sqrt{\nsamp/\pwidth}$ and edges $(1\pm\sqrt{\nsamp/\pwidth})^2$ ($\sqrt{\nsamp/\pwidth}=0.198$ at $\dlat=200$; at $\dlat\le7$ on the $\pwidth=64\dlat$ sweep $\pwidth<\nsamp$ and the proposition does not apply: there are $\pwidth-\dlat$ sample modes with mean $a_*^2/\pwidth$ and no rank-null block). Off-diagonal entries $O(\varrho_{\mu\nu}^3)$, $\varrho_{\mu\nu}=O(\dlat^{-1/2})$ of random sign, have operator norm $O(\sqrt\nsamp\,\dlat^{-3/2})$, so per-eigenvalue attachment requires $\nsamp\ll\dlat^3r_t^2a_*^4$ (synthetic $\dlat=200$: $\dlat^3a_*^4=173<500$, fails; CelebA $\dlat=200$: $3.4\times10^3>10^3$, marginal); otherwise the attachment holds only in the quantile sense of \cref{cor:shape-from-qmu-pred}. The $\dlat$ data-side modes lie above the sample modes because $a_1^2\lambda_i(M_t)/\dlat\gg r_ta_*^2/\nsamp$ (synthetic $\dlat=200$: $8\times10^{-4}$ vs.\ $7.7\times10^{-6}$).
\end{enumerate}
\end{proposition}
\begin{proof}[Proof sketch]
(1) Condition on $x$; row $j$ contributes $\pwidth^{-1}[\tanh(\sqrt{q_x}z_j)-a_1\sqrt{q_x}z_j]^2$, $z_j$ i.i.d.\ standard Gaussian; \cref{lem:ge-coeff-pred}(ii) and the LLN over rows. Over $\xi$, $\|x_t^\mu\|^2/\dlat$ concentrates at $q_\mu$ with $O(\dlat^{-1/2})$ fluctuations. (2) Mehler, \cref{lem:ge-coeff-pred}(v); for two draws of one point $x_t(\xi)^\top x_t(\xi')/\dlat=e^{-2t}\|x^\mu\|^2/\dlat+O(\dlat^{-1/2})$, so $\varrho=\varrho_\mu$. (3) Variance identity $\E[\rho\rho^\top]=\bar\rho\bar\rho^\top+\mathrm{Cov}(\rho)$; the nonzero eigenvalues of $\Umat_{\rm loc}$ are those of the Gram $\bar G/\nsamp$, $\bar G_{\mu\nu}=\bar\rho_\mu^\top\bar\rho_\nu=\mathrm{diag}(r_ta_*^2)+E$ with $|E_{\mu\nu}|=O(|\varrho_{\mu\nu}|^3)$, $\varrho_{\mu\nu}=e^{-2t}x^{\mu\top}x^\nu/(\dlat\sqrt{q_\mu q_\nu})$. (4) Fixed-$W$ entry fluctuations of $\bar G$ are $O(\pwidth^{-1/2})$, the Marchenko--Pastur setting with ratio $\nsamp/\pwidth$ \citep{marchenko1967distribution}; semicircle bound for $E$.

\emph{Numerical check} (repository spectra, $\signoise=0.5$, $\pwidth=64\dlat$, $\nsamp=500$, $t=0.01$, $n_\xi=50$, seed 42, code units). At $\dlat=200$: sample-bulk median $0.093$ vs.\ $\pwidth r_ta_*^2/\nsamp=0.098$ (single-draw formula: $0.119$); sample-bulk mass $28.9$ vs.\ $(\nsamp-\dlat)\pwidth r_ta_*^2/\nsamp=29.4$; mass beyond index $\nsamp$ $11.7$ vs.\ cloud $(1-r_t)\pwidth a_*^2=10.8$; $\lambda_{\nsamp-1}/\lambda_{\nsamp}=1.7$. At $\dlat=150,100$: medians $0.077,0.059$ vs.\ $0.084,0.070$. Below $\dlat\approx60$ the per-sample prediction fails progressively ($0.032$ vs.\ $0.057$ at $\dlat=40$; $0.0018$ vs.\ $0.076$ at $\dlat=5$): the mass migrates into a few collective modes (\cref{rem:collective-pred}).
\end{proof}

\begin{remark}[Single-draw versus noise-averaged objects]\label{rem:single-vs-averaged-pred}
The machine-precision rank cliff at index $\nsamp$ is a property of $\Umat^{(1)}$. The saved spectra use $n_\xi=50$ draws, so $\mathrm{rank}\le\min\{\pwidth,n_\xi\nsamp\}=\pwidth$: all $\pwidth$ eigenvalues exceed $10^{-8}$ at every $\dlat$ and $\lambda_{\nsamp-1}/\lambda_{\nsamp}$ is $1.01$ at $\dlat=10,20$ and $1.7$ at $\dlat=100,200$---a shoulder whose position and mass are those of the cloud. The training loss is the noise-averaged object, so the clock of a training point is $r_t(q_\mu)a_*(q_\mu)^2/\nsamp$; $r_t$ is $0.82$--$0.97$ on the synthetic sweep ($t=0.01$) and $0.09$--$0.76$ at $t=0.1$, the diffusion time of the real-data spectra. Cloud modes carry the part of the residual feature not reproducible across noise draws; they are essentially never absorbed within the training horizon.
\end{remark}

\subsubsection{The two-source fraction predictor}

\begin{definition}\label{def:two-source-predictor-pred}
Given latents $\{z^\mu\}_{\mu\le\nsamp}$, diffusion time $t$, an odd first-layer activation $\sigma$ and gain $\gamma=\dlat\sigma_w^2$ read from the initialisation, set $q_\mu=\gamma(e^{-2t}\|z^\mu\|^2/\dlat+\Delta_t)$, $\varrho_\mu=1-\gamma\Delta_t/q_\mu$, $\Lambda_\mu=g(q_\mu,\varrho_\mu)/\nsamp$ (with $g$ for $\sigma$), $d_{\rm NN}(\mu)=\min_{\nu\ne\mu}\|z^\mu-z^\nu\|$ and $\mathcal D_{t_c}=\{\mu:d_{\rm NN}(\mu)^2\ge2\Delta_{t_c}\log\nsamp\}$ \citep{biroli2024dynamical}. The predicted memorised fraction at step $s$ is
\[ \widehat F(s)=\frac1\nsamp\sum_\mu\mathbf 1[\mu\in\mathcal D_{t_c}]\;\mathbf 1[\Lambda_\mu s\ge C],\qquad C=\log\tfrac1{1-\theta}/\eta, \]
optionally with the second indicator replaced by $\Phi(\log(\Lambda_\mu s/C)/\varsigma)$, $\varsigma=\sqrt{\nsamp/\pwidth}$ (not fitted). $\theta$ is the fraction of the sample-specific score output that must be produced for the sampler to return $z^\mu$ (\cref{prop:mode-clocks-pred}). If the first layer also receives a fixed-norm time embedding $e(t)\in\R^{m_e}$ with fan-in init $\sigma_w^2=c/(\dlat+m_e)$, then $q_\mu=c(e^{-2t}\|z^\mu\|^2+\dlat\Delta_t+\|e(t)\|^2)/(\dlat+m_e)$ and $\varrho_\mu=1-c\Delta_t\dlat/((\dlat+m_e)q_\mu)$. The intersection encodes that a point is returned only if its sample-specific component is learned and wins the softmax competition of the empirical score (\cref{lem:condensation}); under approximate independence $\widehat F(s)\approx(|\mathcal D_{t_c}|/\nsamp)\cdot\frac1\nsamp\#\{\mu:\Lambda_\mu s\ge C\}$. Hypotheses: $\pwidth>\nsamp$; lazy first layer; the decoder approximately preserves nearest-neighbour ordering ($\mathcal D_{t_c}$ is latent-space, the statistic is pixel-space; \cref{rem:pixel}).
\end{definition}

\begin{corollary}[sketched]\label{cor:shape-from-qmu-pred}
With $\mathcal D_{t_c}$ fixed: (i) memorisable points are memorised in decreasing order of $q_\mu$, i.e.\ of latent norm; (ii) for $f\le|\mathcal D_{t_c}|/\nsamp$, $s_f=C\nsamp/g(q_{(f)},\varrho_{(f)})$ with $q_{(f)}$ the upper $(f\nsamp/|\mathcal D_{t_c}|)$-quantile of $\{q_\mu:\mu\in\mathcal D_{t_c}\}$, so the shape of the curve in $\log s$ is the empirical distribution of $\log g(q_\mu,\varrho_\mu)$; (iii) the sample-mode clock depends on $\mathrm{spec}(M_t)$ only through $q_\mu$. Collapsed coordinates ($\lambda_i(M_t)=\Delta_t$) therefore enter only through $q$-dilution---their share of $\Delta_t$ and the $1/\dlat$ fan-in normalisation---the same channel as the synthetic $\signoise\to0$ buffer. At CelebA $\dlat=200$ ($q_d=1.247$, $d_{\rm eff}\approx88$), dropping them at fixed gain raises $q_d$ to $2.60$ and $g$ by $2.7\times$ (memorisation earlier); dropping them with the gain rescaled to keep $q_\mu$ fixed leaves the curve unchanged.
\end{corollary}
\begin{proof}
$g(q,1-\Delta_t/q)$ is increasing in $q$ (\cref{lem:ge-coeff-pred}(v)); invert the threshold within $\mathcal D_{t_c}$. (iii) $\Lambda_\mu$ is a function of $q_\mu$ alone; $q'=e^{-2t}\overline{\|z\|^2}/d_{\rm eff}+\Delta_t$ with $\overline{\|z\|^2}/\dlat=1.302$.
\end{proof}

\begin{remark}[Parameter count and fit sets]\label{rem:param-count-pred}
Released: one global $\kappa$ (121-point $\log_{10}$ grid on $[-12,-4]$, step $1.17\times$) plus six PAVA threshold levels, 7 fitted numbers per dataset, fitted on \emph{all uncensored} rows including pre-threshold dims ($\dlat=20,30$ CelebA; $40$ CIFAR-10) with no hold-out; the main text's $(\kappa,a,b)$ sigmoid is not what the code fits. Derived: $\gamma$ (0 fitted), $\varsigma$ (0), $\mathcal D_{t_c}$ (0 or 1), one scalar $C$ (the single unknown $\theta$). All one-constant fits below use $C=\exp\{\mathrm{median}_{\dlat}\log(\Lambda_{(f)}\hat s_f)\}$ on declared dim sets.
\end{remark}

\paragraph{Practitioner recipe.}
Inputs: frozen encoder, $\nsamp\approx10^3$ training latents $Z\in\R^{\nsamp\times\dlat}$, diffusion time $t$, the diffusion network's first-layer activation, fan-in and initialisation variance (giving $\gamma$, $c$, $m_e$), learning rate $\eta$, sampler condensation time $t_c$.
\begin{enumerate}
\item $q_\mu$ and $\varrho_\mu$ from \cref{def:two-source-predictor-pred} (surrogate: $q_\mu=e^{-2t}\|z^\mu\|^2/\dlat+\Delta_t$).
\item $g(q,\varrho)$ by two-dimensional Gauss--Hermite quadrature ($40$--$80$ nodes per axis): $\sum_{jk}w_jw_k\sigma(\sqrt q x_j)\sigma(\sqrt q(\varrho x_j+\sqrt{1-\varrho^2}x_k))-a_1(q)^2q\varrho$ (for a non-odd $\sigma$ subtract also $a_0^2$, $a_0=\E\sigma(\sqrt qz)$); $\Lambda_\mu=g(q_\mu,\varrho_\mu)/\nsamp$.
\item $d_{\rm NN}(\mu)$ in latent space; $\mathcal D_{t_c}=\{d_{\rm NN}^2\ge2\Delta_{t_c}\log\nsamp\}$.
\item Calibrate $C$ once as the median over training dims of $\log(\Lambda_{(5\%)}\hat s_{5\%})$; report $\hat s_f$ for all levels and all $\dlat$ from the same $C$.
\end{enumerate}

\paragraph{Validation protocol and current status.}
\emph{Leave-one-$\dlat$-out.} Fit $C$ on all dims of a declared set but one; predict the held-out hitting times; score by $|\log\hat s_f/s_f|$. \emph{Cross-dataset.} Fit $C$ on CelebA, predict CIFAR-10. \emph{Collapsed-coordinate test.} Retrain on latents with collapsed coordinates removed, with and without gain rescaling (\cref{cor:shape-from-qmu-pred}(iii)); requires new runs. \emph{Ordering test.} Memorised training points should have larger $\|z^\mu\|$ early in training; requires logging the NN1 index.

\emph{Status on the released tables} (mean $|\log$ residual$|$, in-sample / leave-one-out; released 7-parameter in-sample on the same set). CelebA $5\%$, $\dlat\ge50$ (11 dims): $g$-clock $0.309/0.407$; $a_*^2$-clock $0.347/0.449$; $\tau\propto\dlat$ $0.125/0.125$; released $0.339$. CelebA $25\%$, $\dlat\ge50$: $0.207/0.262$; $0.246/0.308$; $0.083/0.095$; released $0.169$. CelebA $5\%$, $\dlat\ge70$ (9 dims): $0.277/0.384$ vs.\ released $0.311$. CelebA $5\%$, all uncensored $\dlat\ge20$ (14): $0.346/0.417$ vs.\ released $0.365$, with $\tau\propto\dlat$ the worst model ($0.407/0.477$). CIFAR-10 $5\%$, $\dlat\ge60$ (11): $0.130/0.151$ vs.\ released $0.146$ ($\tau\propto\dlat$ $0.155/0.198$); $\dlat\ge100$: $0.120/0.145$ vs.\ $0.116$; $\dlat\ge140$ (7): $0.112/0.136$ vs.\ $0.106$. CIFAR-10 $25\%$, $\dlat\ge60$ (8): $0.168/0.171$ vs.\ $0.143$. The $g$-clock is thus comparable to the released fit on CIFAR-10 and somewhat worse on CelebA, with no fitted response function and one constant. It accounts for $1.78\times$ of the observed $3.9\times$ (CelebA $\dlat=50\to200$), $1.74\times$ of $2.9\times$ ($100\to200$) and $1.85\times$ of $2.1\times$ (CIFAR-10 $100\to240$; $\dlat=260$ has $\tau_5=2520$k with s.d.\ $1392$k and is not used). The fitted $C$ differs between datasets by $2.0\times$ ($5\%$: $14.5$ vs.\ $29.1$) and $1.6\times$ ($25\%$) with identical optimiser settings, so one-constant transfer fails and a dataset-dependent factor ($|\mathcal D_{t_c}|/\nsamp$, decoder distortion of $d_{\rm NN}$) is missing. The width from the learning side alone is too narrow: at CelebA $\dlat=200$ the Gaussian-moment estimate (mean $1.302$, s.d.\ $0.2905$ of $\|z\|^2/\dlat$; per-sample latents are not in the repository) gives a $1.46\times$ span between the levels matching $5\%$ and $25\%$ absolute memorisation ($11\%$ and $56\%$ of the $0.448$ asymptote) against the observed $\tau_{25}/\tau_5=2.14$.

\emph{Surrogate caveat.} All real-data numbers use $\sigma=\tanh$, $\gamma=1$, $m_e=0$. The deployed network's first layer is $\mathrm{Linear}(\dlat+64,1024)$ with GELU on $[x_t;\,32\ \sin,\,32\ \cos]$ and PyTorch default init ($\sigma_w^2=1/(3\,\mathrm{fan\_in})$), for which $q_\mu\approx(e^{-2t}\|z\|^2+\dlat\Delta_t+32)/(3(\dlat+64))=0.36,\,0.43,\,0.36$ at CelebA $\dlat=50,100,200$ and $0.43,\,0.33$ at CIFAR-10 $\dlat=100,260$---nearly $\dlat$-independent---and GELU is not odd (its $k=0,2$ Hermite terms add a rank-one spike and $O(\varrho^2)$ off-diagonals; at $q=1.834$ the proper GELU $a_*^2$ is $0.191$ against $0.380$ from the odd-activation formula). The derived predictor has therefore not been evaluated with the deployed network's own first-layer statistics; either the surrogate stands in for feature learning in deeper layers \citep{ba2022highdimensional}, or the real-data $\dlat$-slope arises mainly on the dynamics side and from the NN-ratio detection drift.

\begin{remark}[Conjecture, numerically supported: cluster-collective modes at small $\dlat$]\label{rem:collective-pred}
Code units, $\signoise=0.5$, $\pwidth=64\dlat$, $\nsamp=500$, $t=0.01$, seed 42. The per-sample picture predicts a sample-bulk median $\pwidth r_ta_*^2/\min\{\nsamp,\pwidth\}$ rising $0.054\to0.098$ ($1.8\times$) and a top edge $\pwidth r_ta_*^2(1+\sqrt{\nsamp/\pwidth})^2/\nsamp$ falling $0.193\to0.141$ ($1.4\times$) over $\dlat=10\to200$ ($\dlat=5$ has $\pwidth<\nsamp$). Observed: median $0.0039\to0.093$ ($24\times$), largest sample-bulk eigenvalue $0.930\to0.173$ ($5.4\times$; $0.957$ at $\dlat=5$); the top ten sample-bulk modes carry $47\%,\,35\%,\,22\%,\,6.3\%,\,5.4\%$ of the sample-bulk mass at $\dlat=5,10,20,100,200$ (flat reference $3.3\%$), with $\lambda_{\rm top}/\lambda_{\rm median}=532,\,236,\,42,\,3.4,\,1.9$. At $\dlat\ge100$ the per-sample+MP picture is quantitatively right; at small $\dlat$ about $k=10$ modes absorb half the mass, and the per-sample clock explains essentially none of the small-$\dlat$ onset behaviour ($1.46\times$ predicted vs $221\times$ observed in code units over $\dlat=5\to200$). Conjectured mechanism: same-cluster points have noise-averaged overlap $\varrho_c=e^{-2t}s^2/(\dlat q_d)=0.64,\,0.58,\,0.50,\,0.23,\,0.135$, so $\bar G$ acquires per cluster a block of $\nsamp/k$ entries $g(q_d,\varrho_c)$ and $k$ collective eigenvalues $\approx\pwidth[r_ta_*^2+(\nsamp/k-1)g(q_d,\varrho_c)]/\nsamp=0.53,\,0.50,\,0.36,\,0.11,\,0.11$, reproducing the observed trend from $\dlat=10$ within a uniform factor $1.6$--$1.9$ (cross-cluster overlaps, also $O(1)$ at $\dlat=5$, and the MP edge are omitted; this block can contribute at most $\sim10\times$ at $\dlat=5$, so the remainder is the $g$ factor of \cref{prop:onset-buf}). Such modes learn the nonlinear shape of a cluster and are generalising; if confirmed, the small-$\dlat$ onset statistic is not a memorisation clock and the buffer claim must be re-based on the localised part of the sample bulk (the $\gamma$-cut of \cref{def:taumem-p1}). Test: training-point-basis IPR of the top 20 sample-bulk eigenvectors ($\approx1$ localised, $\approx k/\nsamp$ collective); requires re-running the sweep with eigenvectors saved.
\end{remark}

\paragraph{Predictions and how to test them.}
\begin{enumerate}
\item \emph{Training points are memorised in decreasing order of latent norm.} Test: log the NN1 index of memorised samples; compare $\|z^\mu\|^2/\dlat$ of those training latents with the full distribution (norm-ranked ROC AUC $>0.5$ early, $\to0.5$ at saturation).
\item \emph{Collapsed coordinates enter the memorisation clock only through $q_\mu$:} dropping them at fixed gain makes CelebA $\dlat=200$ memorise $\sim2.7\times$ faster; dropping them with the gain rescaled leaves the curve unchanged; the released $P_d(s)$ moves in both cases. Test: re-encode and retrain with and without collapsed coordinates and with/without gain rescaling.
\item \emph{The one-constant sample-mode clock $\tau_f(\dlat)=C\nsamp/g(q_d,1-\Delta_t/q_d)$ reproduces hitting times with leave-one-out log-error comparable to the released 7-parameter in-sample fit but under-predicts the $\dlat$-slope} ($1.78\times$ of $3.9\times$, CelebA $50\to200$). Test: executable now from the released tables; then add the detection-drift correction and the deployed-network $q_\mu$.
\item \emph{The learning-side width is too narrow} ($1.46\times$ vs observed $2.14\times$ at CelebA $\dlat=200$); a second source (the $d_{\rm NN}$ quantile, \cref{app:theory-bridge}) is required. Test: requires per-sample latents.
\item \emph{Data-side rate check:} on the synthetic RFNN the data-side eigenvalues scale as $64a_1(q_d)^2\lambda_i(M_t)$ in code units, so a per-$\dlat$ refit of $\kappa$ scales as $a_1(q_d)^2$. Test: regress $\log$ block medians on $\log a_1(q_d)^2$.
\item \emph{Top sample-bulk eigenvectors at small $\dlat$ are cluster-collective}, $\sim k$ of them with eigenvalue $\sim\pwidth[r_ta_*^2+(\nsamp/k-1)g(q_d,\varrho_c)]/\nsamp$ within a factor $2$. Test: IPR of the top 20 sample-bulk eigenvectors.
\item \emph{Cross-dataset transfer at the one-constant level currently fails} ($C$ differs $2.0\times$ at $5\%$, $1.6\times$ at $25\%$). Test: fit $C$ on CelebA, predict CIFAR-10; check whether $|\mathcal D_{t_c}|/\nsamp$ accounts for the offset.
\item \emph{Recomputing $\Umat$ with one noise draw restores the rank cliff at $\nsamp$ and raises the sample-bulk median by $1/r_t(q_d)$} ($1.22\times$ at $\dlat=200$, $1.08\times$ at $\dlat=20$). Test: \texttt{compute\_U} with $n_\xi\in\{1,50,500\}$.
\end{enumerate}

\paragraph{Open gaps.}
(1) The noise-averaged eigenvalue $r_t(q_\mu)a_*(q_\mu)^2/\nsamp$ replaces $a_*(q_\mu)^2/\nsamp$ throughout; at $t=0.01$ (the synthetic spectra) $r_t$ ranges only $0.97\to0.82$, at $t=0.1$ (real data) $0.09\to0.76$. In code units the per-sample clock explains essentially none of the small-$\dlat$ synthetic behaviour (\cref{rem:collective-pred}). (2) Only unit-free eigenvalue ratios are used here; absolute step counts depend on the clock convention of \eqref{eq:lin-steps}. (3) The cliff at index $\nsamp$ is a single-draw statement; the saved $50$-draw spectra have full rank $\pwidth$. (4) Per-eigenvalue attachment needs $\nsamp\ll\dlat^3r_t^2a_*^4$, which fails for the synthetic data at every $\dlat\le200$ and holds only marginally for real latents at $\dlat\ge200$; moreover for the real runs $\nsamp/\pwidth=1000/1024$, so the MP spread $2\sqrt{\nsamp/\pwidth}\approx2$ is $O(1)$ and the sample bulk is a delocalised MP-type density, not the empirical distribution of $\log a_*(q_\mu)^2$: the ordering-by-norm claim has no support there, and the quantile-level statement of \cref{cor:shape-from-qmu-pred} is what the predictor uses and is not proved. (5) The learning-side clock explains $1.7$--$1.9\times$ of the $2.1$--$3.9\times$ observed real-data slowdown; candidate sources of the rest: NN-ratio detection drift, the $\dlat$-dependence of $\mathcal D_{t_c}$, and a $\dlat$-proportional factor. (6) The deployed real-data first layer (GELU, fan-in $\dlat+64$ with a fixed-norm time embedding, PyTorch default init, momentum $0.8$ so $\eta_{\rm eff}=\eta/(1-0.8)$) has $q_\mu\approx0.36$ nearly independent of $\dlat$; applied literally, the recipe predicts almost no $\dlat$-dependence from the first layer. This is the largest open issue for the real-data reading. (7) The threshold $\theta$ has an output-space meaning but is not related to the softmax condensation basin of the empirical score. (8) The square in the released clock is not derived (\cref{rem:square-pred}); the released fit uses 7 parameters with no hold-out (\cref{rem:param-count-pred}).

% ----------------------------------------------------------------------
\subsection{Beyond two atoms: general latent spectra, effective dimension, and posterior collapse}
\label{app:theory-general}

This subsection drops the two-atom assumption. Nothing in the random-feature analysis requires it: \citet{bonnaire2025} state their Theorem~3.1 for an arbitrary spectral measure of the input covariance (their supplement derives the Gaussian-equivalent problem under arbitrary input covariance), and the four-bulk picture is that theorem evaluated at a two-atom measure; \citet{george2025dsm} treat data on a $\dint$-dimensional subspace with $\signoise=0$. Here we record what the same analysis gives for the continuous, partially collapsed spectra of real VAE latents, and confront it with the spectra measured in \cref{sec:spectral-predictor}.

\paragraph{Setting.}
Let $Z\in\R^{\nsamp\times\dlat}$ hold the encoder means $z^\mu$ of the $\nsamp$ training images, $\Sigma_0:=Z^\top Z/\nsamp=\sum_{i=1}^{\dlat}s_iv_iv_i^\top$ with $s_1\ge\dots\ge s_{\dlat}\ge0$, $x_t^\mu=e^{-t}z^\mu+\sqrt{\Delta_t}\eta^\mu$ with one diffused copy per training point, and
\begin{equation}\label{eq:Mt-real}
M_t:=e^{-2t}\Sigma_0+\Delta_tI_{\dlat},\qquad S_t:=\frac1\nsamp\sum_\mu x_t^\mu x_t^{\mu\top},\qquad \E_\eta S_t=M_t,\qquad \lambda_i(M_t)=e^{-2t}s_i+\Delta_t .
\end{equation}
$M_t$ is the matrix used by the predictor of \cref{sec:spectral-predictor}; its eigenvectors are those of $\Sigma_0$, and $\Delta_t$ is the smallest value any of its eigenvalues can take. The feature second-moment matrix $\Umat=\nsamp^{-1}\sum_\mu\phi(x_t^\mu)\phi(x_t^\mu)^\top$ is formed here from the same single copies (so $\operatorname{rank}\Umat\le\nsamp$; this is $\Umat^{(1)}$ of \cref{app:theory-fourbulk}), and its data-side spectrum is controlled by $S_t$, not by $M_t$; the difference is a sampling broadening quantified in \cref{lem:sampling-real}. At $t=0.1$, $\Delta_t=0.1813$, $a_1(\Delta_t)^2=0.746$, $a_\ast(\Delta_t)^2=1.39\times10^{-3}$.

\begin{definition}[Latent spectral measure and effective dimensions]\label{def:deff-real}
Let $\rho_{\dlat}:=\dlat^{-1}\sum_i\delta_{s_i}$ be the empirical spectral measure of $\Sigma_0$ and
\begin{equation}\label{eq:q-real}
q[\rho_{\dlat}]:=\frac{\operatorname{tr}M_t}{\dlat}=e^{-2t}\!\int\! s\,\mathrm d\rho_{\dlat}(s)+\Delta_t,\qquad \hat q:=\frac{\operatorname{tr}S_t}{\dlat}=\frac1\nsamp\sum_\mu\frac{\|x_t^\mu\|^2}{\dlat},\qquad \E\hat q=q .
\end{equation}
For $\delta>0$ define the \emph{floor count}, the \emph{effective dimension}, and the \emph{excess}
\begin{equation}\label{eq:deff-real}
d_{\mathrm{floor}}(\delta):=\#\{i:\lambda_i(M_t)\le(1+\delta)\Delta_t\},\qquad
d_{\mathrm{eff}}:=\frac{\big(\sum_i(\lambda_i(M_t)-\Delta_t)\big)^2}{\sum_i(\lambda_i(M_t)-\Delta_t)^2}=\frac{(\sum_is_i)^2}{\sum_is_i^2},\qquad
d_{\mathrm{exc}}:=\dlat-d_{\mathrm{eff}} .
\end{equation}
\end{definition}
$d_{\mathrm{eff}}$ is the participation ratio of $\Sigma_0$; because $\Delta_t$ is known exactly, it involves no estimated floor. In the two-atom model, $d_{\mathrm{eff}}=(\dint S+(\dlat-\dint)\signoise^2)^2/(\dint S^2+(\dlat-\dint)\signoise^4)$ with $S=s^2/\dint+\sigsig^2$; this tends to $\dint$ as $\signoise\to0$ but equals $76.6$ at $(\dlat,\signoise)=(200,0.5)$: a participation ratio counts a heavy excess block partially, which is why it is the wrong quantity to substitute for $\dint$ when $\signoise$ is not small. The released estimator of \cref{sec:spectral-predictor} (\texttt{effective\_rank\_excess}) takes the participation ratio of $(\lambda_i-\hat\beta)_+$ with $\hat\beta$ the median of the lowest $30\%$ of eigenvalues; on the \emph{population} two-atom spectrum $\hat\beta=\beta_t^2$ and it returns $\dint$ exactly (on a one-copy finite-$\nsamp$ spectrum the noise block is MP-spread, \cref{lem:sampling-real}, and it does not). On the real spectra $\hat\beta$ switches from $0.72$ (CelebA, $\dlat=140$) to $0.189,0.182,0.181$ ($\dlat=160,180,200$), so the estimator coincides with $d_{\mathrm{eff}}$ of \eqref{eq:deff-real} only where $\hat\beta=\Delta_t$, and its values on the two sides of that boundary measure different things. We recommend reporting $d_{\mathrm{eff}}$ and $d_{\mathrm{floor}}$ separately (\cref{rem:xaxis-real}).

\begin{lemma}[Sampling broadening of a collapsed block; (a)--(b) proved, (c) sketched]\label{lem:sampling-real}
Let $F$ be a set of $d_F$ coordinates with $z^\mu_i=0$ for all $i\in F$ and all $\mu$, $L$ its complement ($d_L=\dlat-d_F$), $H_F\in\R^{\nsamp\times d_F}$ the matrix of noises $\eta^\mu_F$ and $X_L\in\R^{\nsamp\times d_L}$ the matrix of diffused live coordinates.
\begin{enumerate}
\item[(a)] $(S_t)_{FF}=\Delta_tH_F^\top H_F/\nsamp$ exactly, with $H_F$ iid $\mathcal N(0,1)$ independent of $X_L$. For $u>0$, with probability $\ge1-2e^{-u^2/2}$,
$\operatorname{spec}(S_t)_{FF}\subset\Delta_t\big[(1-\sqrt{d_F/\nsamp}-u/\sqrt{\nsamp})^2,\,(1+\sqrt{d_F/\nsamp}+u/\sqrt{\nsamp})^2\big]$.
\item[(b)] (Interlacing) $\lambda_{d_L+k}(S_t)\le\lambda_k((S_t)_{FF})$ for $k\le d_F$; the $d_F$ smallest eigenvalues of $S_t$ lie below the upper edge in (a).
\item[(c)] (Schur complement) The floor eigenvalues $\lambda$ of $S_t$ solve $\det\big(\Delta_t\nsamp^{-1}H_F^\top(I-P_L(\lambda))H_F-\lambda\big)=0$ with $P_L(\lambda):=X_L(X_L^\top X_L-\nsamp\lambda)^{-1}X_L^\top$. Replacing $P_L(\lambda)$ by the orthogonal projector $P_L:=P_L(0)$ onto $\operatorname{col}(X_L)$ (relative error $\le\lambda/(\lambda_{\min}((S_t)_{LL})-\lambda)$), the matrix $\Delta_t\nsamp^{-1}H_F^\top(I-P_L)H_F$ is $\Delta_t\frac{\nsamp-d_L}{\nsamp}$ times a normalized Wishart$(d_F,\nsamp-d_L)$; hence, to leading order, the floor block of $S_t$ is a Marchenko--Pastur bulk on
\begin{equation}\label{eq:mp-floor-real}
\Delta_t\Big(1-\frac{d_L}{\nsamp}\Big)\Big(1\pm\sqrt{\frac{d_F}{\nsamp-d_L}}\Big)^2 .
\end{equation}
\end{enumerate}
The same statements hold with $\Delta_t\to\beta_t^2$ for the null block of the two-atom model, whose null coordinates are iid $\mathcal N(0,\beta_t^2)$ independent of the signal block.
\end{lemma}
\begin{proof}
(a) $x^\mu_{t,i}=\sqrt{\Delta_t}\eta^\mu_i$ for $i\in F$, and $\eta_F$ is independent of $(z^\mu,\eta_L)$; the singular-value bound is \citet[Cor.~7.3.3]{vershynin2018hdp} applied to $H_F$. (b) Cauchy interlacing for the principal submatrix. (c) $\det(S_t-\lambda)=\det((S_t)_{LL}-\lambda)\det\big((S_t)_{FF}-\lambda-(S_t)_{FL}((S_t)_{LL}-\lambda)^{-1}(S_t)_{LF}\big)$ and $(S_t)_{LF}=\sqrt{\Delta_t}X_L^\top H_F/\nsamp$, so the correction equals $\Delta_t\nsamp^{-1}H_F^\top X_L(X_L^\top X_L-\nsamp\lambda)^{-1}X_L^\top H_F=\Delta_t\nsamp^{-1}H_F^\top P_L(\lambda)H_F$. $I-P_L$ is a rank-$(\nsamp-d_L)$ orthogonal projector independent of $H_F$, so $(I-P_L)H_F$ is an $(\nsamp-d_L)\times d_F$ Gaussian matrix in a rotated basis, whose normalized Gram matrix has MP edges $(1\pm\sqrt{d_F/(\nsamp-d_L)})^2$.
\end{proof}
\emph{Numbers.} CelebA $\dlat=200$, $\nsamp=1000$, $d_F\ge30$: the block alone gives $[0.68,1.38]\Delta_t$; with the projection factor $1-170/1000=0.83$ and aspect $30/830$, \eqref{eq:mp-floor-real} gives $[0.54,1.18]\Delta_t$. Synthetic null block ($d_F=195$, $d_L=5$, $\nsamp=500$): $[0.14,2.62]\beta_t^2$. Both dwarf the width-$\varepsilon$ term of \cref{prop:pushforward-real} ($\varepsilon=0.27$ at $\pwidth=64\dlat$). The signal block ($\dint=5$, aspect $0.01$) has a sampling width of order $\sqrt{\dint/\nsamp}=0.1$, of the same order as $\varepsilon$, which is why the signal edge statistic is sharp while the noise-dimension block is broad.

\begin{assumption}[Gaussian equivalence on latents; conjectured]\label{assump:gep-real}
As $\dlat,\nsamp,\pwidth\to\infty$ with $\pwidth>\nsamp>\dlat$ and $\varepsilon_0:=\sqrt{\dlat/\pwidth}+u/\sqrt{\pwidth}<1$, the feature second-moment matrix admits the decomposition
\begin{equation}\label{eq:Usplit-real}
\Umat=\Umat^{\mathrm{lin}}+\Umat^{\mathrm{res}},\qquad
\Umat^{\mathrm{lin}}=\frac{a_1(q)^2}{\pwidth\,\dlat}\,WS_tW^\top,\qquad \Umat^{\mathrm{res}}\succeq0,\qquad
\operatorname{tr}\Umat^{\mathrm{res}}=a_\ast(q)^2(1+o(1)),\qquad \|\Umat^{\mathrm{res}}\|=:c_+\frac{a_\ast(q)^2}{\nsamp},
\end{equation}
with $q=q[\rho_{\dlat}]$ and $c_+=O(1)$.
\end{assumption}
This is the Gaussian-equivalence decomposition of \cref{app:theory-fourbulk,app:theory-buffer} \citep{bonnaire2025,goldt2022gaussian,hu2022universality,mei2022generalization}, now asserted for diffused latent inputs and written, as the one-copy $\Umat$ requires, with the sample second moment $S_t$. Three things make it an assumption rather than a theorem. Real latents are not Gaussian. The pre-diffusion $\|z^\mu\|^2/\dlat$ has coefficient of variation $0.22$--$0.23$ across $\mu$ on CelebA for $\dlat\ge160$ (hence $\|x_t^\mu\|^2/\dlat$ about $0.20$), so $q$ is the average of per-sample values $q_\mu$ and every edge below is correspondingly broadened ($\hat q=q$ to relative $0.2/\sqrt{\nsamp}$). And with one copy per point the residual $r_\mu=\phi(x_t^\mu)-a_1Wx_t^\mu/\sqrt{\pwidth\dlat}$ is correlated with $Wx_t^\mu$: since $\operatorname{rank}\Umat\le\nsamp$, $\operatorname{rank}\Umat^{\mathrm{lin}}=\dlat$ and $\Umat^{\mathrm{res}}\succeq0$, one has $\operatorname{range}\Umat^{\mathrm{lin}}\subset\operatorname{range}\Umat^{\mathrm{res}}$, so $\Umat^{\mathrm{res}}$ is not an isotropic Wishart independent of $W$, and the Marchenko--Pastur value $(1+\sqrt{\nsamp/\pwidth})^2$ is not the right reference for $c_+$ (\cref{prop:onset-buf}). We carry $c_+$ as a fitted $O(1)$ constant; on the synthetic sweep it is $\dlat$-dependent (it is $g(\dlat)\theta$ of \cref{app:theory-buffer}).

\begin{proposition}[The data-side spectrum of $\Umat$ is a pushforward of the spectrum of $S_t$; proved given \cref{assump:gep-real}]\label{prop:pushforward-real}
Under \cref{assump:gep-real}, with probability at least $1-2e^{-u^2/2}$ over $W$,
\begin{align}
\lambda_i(\Umat^{\mathrm{lin}})&=\frac{a_1(q)^2}{\dlat}\,\lambda_i(S_t)\,(1+\epsilon_i),\qquad |\epsilon_i|\le\varepsilon:=2\varepsilon_0+\varepsilon_0^2,\qquad i\le\dlat,\label{eq:push-real}\\
\lambda_i(\Umat^{\mathrm{lin}})&\le\lambda_i(\Umat)\le\lambda_i(\Umat^{\mathrm{lin}})+c_+\frac{a_\ast(q)^2}{\nsamp},\qquad i\le\pwidth .\label{eq:weyl-real}
\end{align}
Consequently the $\dlat$ largest eigenvalues of $\Umat$ are, to relative accuracy $\varepsilon+c_+(a_\ast^2/a_1^2)(\dlat/\nsamp)/\lambda_i(S_t)$, the pushforward $\big(\lambda\mapsto a_1(q)^2\lambda/\dlat\big)_\#$ of the spectrum of $S_t$, which is the spectrum of $M_t$ broadened by one-copy sampling (\cref{lem:sampling-real}). The next $\nsamp-\dlat$ eigenvalues form the sample bulk of $\Umat^{\mathrm{res}}$, with mass $a_\ast(q)^2-\operatorname{tr}(P\Umat^{\mathrm{res}}P)\in[a_\ast^2(1-c_+\dlat/\nsamp),a_\ast^2]$ ($P$ the projector on the top $\dlat$ eigenvectors), mean in $[a_\ast^2(1-c_+\dlat/\nsamp),a_\ast^2]/(\nsamp-\dlat)$ (whose $\dlat/\nsamp\to0$ limit is $a_\ast^2/\nsamp$), and upper edge $c_+a_\ast(q)^2/\nsamp$; the remaining $\pwidth-\nsamp$ are zero. The trace identity $\sum_i\lambda_i(\Umat)=a_1(q)^2\hat q+a_\ast(q)^2=\E[\tanh^2(\sqrt qz)]$ holds to $o(1)$.
\end{proposition}
\emph{Special cases.} (a) Two-atom $\rho_{\dlat}$ with atoms $\alpha_t^2$ ($\dint$ of them) and $\beta_t^2$ ($\dlat-\dint$): \eqref{eq:push-real} with \cref{lem:sampling-real} gives the signal block at $a_1(q)^2\alpha_t^2/\dlat$ with sampling width of order $\sqrt{\dint/\nsamp}$ and the noise-dimension bulk on $a_1(q)^2\beta_t^2(1-\dint/\nsamp)(1\pm\sqrt{(\dlat-\dint)/(\nsamp-\dint)})^2/\dlat$, median $\approx a_1(q)^2\beta_t^2/\dlat$; this is the single-copy form of \cref{prop:fourbulk}, whose noise-dimension medians ($5.22\to9.22$ in code units) the formula reproduces. (b) A general $\rho_{\dlat}$ gives a data-side \emph{continuum} on $[a_1^2(e^{-2t}s_{\dlat}+\Delta_t)/\dlat,\;a_1^2(e^{-2t}s_1+\Delta_t)/\dlat]$ together with a \emph{floor bulk} of $d_{\mathrm{floor}}(0)$ eigenvalues with median $\approx a_1(q)^2\Delta_t/\dlat$ and relative half-width $\approx2\sqrt{d_{\mathrm{floor}}/\nsamp}$ by \eqref{eq:mp-floor-real}; the noise-dimension bulk is replaced by the part of the spectrum near the floor. (c) For CelebA at $\dlat=200$, $\nsamp=1000$, the additive term in \eqref{eq:weyl-real} is $0.12\,c_+$ of the floor median; for the synthetic noise-dimension edge at the $t$ where $q=0.33$ ($\beta_t^2=0.268$, $\dlat=200$, $\nsamp=500$) it is $0.011\,c_+$.
\begin{proof}
(i) The nonzero eigenvalues of $WS_tW^\top/(\pwidth\dlat)$ equal those of $\dlat^{-1}S_t^{1/2}GS_t^{1/2}$ with $G:=W^\top W/\pwidth$, since $AB$ and $BA$ share their nonzero spectrum. By \citet[Cor.~7.3.3]{vershynin2018hdp}, $\sqrt{\pwidth}-\sqrt{\dlat}-u\le s_{\min}(W)\le s_{\max}(W)\le\sqrt{\pwidth}+\sqrt{\dlat}+u$ with probability $\ge1-2e^{-u^2/2}$, hence $\operatorname{spec}G\subset[(1-\varepsilon_0)^2,(1+\varepsilon_0)^2]$ and $\|G-I_{\dlat}\|\le\varepsilon$; $\varepsilon_0<1$ makes $G\succ0$. Ostrowski's theorem for the congruence $S_t^{1/2}GS_t^{1/2}$ gives $\lambda_i(S_t^{1/2}GS_t^{1/2})=\theta_i\lambda_i(S_t)$ with $\theta_i\in[\lambda_{\min}(G),\lambda_{\max}(G)]\subset[1-\varepsilon,1+\varepsilon]$. (ii) is Weyl's inequality for the positive-semidefinite perturbation $\Umat^{\mathrm{res}}$. The counts follow from $\operatorname{rank}\Umat^{\mathrm{lin}}=\dlat$ and $\operatorname{rank}\Umat\le\nsamp$. Mass: by \eqref{eq:weyl-real} the residual trace absorbed by the top $\dlat$ eigenvalues is at most $\dlat c_+a_\ast^2/\nsamp$. Trace: $\operatorname{tr}\Umat^{\mathrm{lin}}=a_1^2\operatorname{tr}(S_tG)/\dlat=a_1^2\hat q(1+O(\varepsilon))$ and \eqref{eq:Usplit-real}. The identification of the intermediate $\nsamp-\dlat$ eigenvalues with the residual bulk and of its edge with $c_+a_\ast^2/\nsamp$ rests on the assumption and is the sketched part.
\end{proof}

\begin{definition}[Timescales and buffer ratio for a general spectrum]\label{def:ratio-real}
Per-mode gradient-flow timescales are $\tau_i=1/\lambda_i(\Umat)$ (\cref{lem:gf-p1}). Fix $\theta\in(0,1]$ and let $\lambda_{(\theta)}$ be the largest level such that the eigendirections of $M_t$ at or above it carry a fraction $\theta$ of the latent variance, $\sum_{i:\lambda_i(M_t)\ge\lambda_{(\theta)}}s_i\ge\theta\sum_is_i$. Define
\begin{equation}\label{eq:taus-real}
\taugen^{(\theta)}:=\frac{\dlat}{a_1(q)^2\,\lambda_{(\theta)}},\qquad
\taumem^{\mathrm{onset}}:=\frac1{\lambda_+^{\mathrm{samp}}}=\frac{\nsamp}{c_+\,a_\ast(q)^2},\qquad
R^{(\theta)}[\rho_{\dlat}]:=\frac{\taumem^{\mathrm{onset}}}{\taugen^{(\theta)}}=\frac{\nsamp}{c_+\dlat}\;\frac{a_1(q)^2}{a_\ast(q)^2}\;\lambda_{(\theta)} .
\end{equation}
\end{definition}
For the two-atom model, $\lambda_{(\theta)}=\alpha_t^2$ and $\taugen^{(\theta)}=1/\lambda^{\mathrm{signal}}_{\min}$ (the statistic of \cref{def:abs-times-buf}) only for $\theta$ below the signal variance share $\dint S/(\dint S+(\dlat-\dint)\signoise^2)$, which is $0.22$ at $(\dlat,\signoise)=(200,0.5)$ and $0.05$ at $\signoise=1$; for larger $\theta$, $\lambda_{(\theta)}=\beta_t^2$ and $\taugen^{(\theta)}$ is the noise-dimension timescale. On a continuous spectrum there is no canonical $\theta$; the results below are stated so that $\theta$ drops out. No count of coordinates appears in \eqref{eq:taus-real}: $\dlat$ enters through the fan-in $1/\dlat$, through $q[\rho_{\dlat}]$, and through $c_+$ (\cref{rem:nonsequitur-buf}). \emph{What the synthetic sweep says about $c_+$.} With the Marchenko--Pastur guess $c_+=(1+\sqrt{\nsamp/\pwidth})^2$ and $\pwidth=64\dlat$, the sweep $\dlat=5\to200$ gives an $a_\ast^2$ ratio of $16.5$ and a $c_+$ ratio of $3.5$, i.e.\ a predicted onset growth of $58\times$ against the measured $221\times$: $c_+$ on clustered data carries a further $\approx4\times$ $\dlat$-dependence (the $g$ factor of \cref{prop:onset-buf}). Every real-data ratio below is therefore taken under the explicit hypothesis
\begin{equation}\label{eq:Hc-real}
(\mathrm H_c)\qquad c_+\ \text{is constant across the compared widths.}
\end{equation}

\begin{lemma}[A collapsed coordinate sits exactly on the floor; proved]\label{lem:dead-real}
Suppose coordinate $i$ of the encoder means vanishes on the training set, $z^\mu_i=0$ for all $\mu$ (posterior collapse, \citealp{lucas2019dontblame}). Then:
\begin{enumerate}
\item[(a)] $\Sigma_0e_i=0$ and $M_te_i=\Delta_te_i$: $\lambda=\Delta_t$ exactly, the minimum possible eigenvalue of $M_t$, with eigenvector $e_i$.
\item[(b)] $x^\mu_{t,i}=\sqrt{\Delta_t}\,\eta_i$, and the exact score of the diffused empirical measure \citep{scarvelis2023closedform,gu2023memorization,biroli2024dynamical} has $s^\ast_i(x_t)=-x_{t,i}/\Delta_t$, a linear function of $x_{t,i}$ alone.
\item[(c)] The coordinate contributes exactly $\Delta_t/\dlat$ to $q$.
\item[(d)] Under \cref{assump:gep-real}, the $d_F$ collapsed coordinates contribute to $\Umat^{\mathrm{lin}}$ the floor bulk $a_1(q)^2\lambda(1+\epsilon)/\dlat$ with $\lambda$ in \eqref{eq:mp-floor-real}, the slowest data-side modes; the median timescale is $\tau_{\mathrm{floor}}=\dlat/(a_1(q)^2\Delta_t)$ and the slowest floor mode is slower by the further factor $\big[(1-d_L/\nsamp)(1-\sqrt{d_F/(\nsamp-d_L)})^2\big]^{-1}$.
\item[(e)] For a linear score model $s(x)=Bx$ trained by gradient flow on $\E\|Bx_t-s^\ast(x_t)\|^2$, column $i$ of $B$ evolves independently of all other entries: $B_{ii}(T)=-\Delta_t^{-1}+(B_{ii}(0)+\Delta_t^{-1})e^{-2\Delta_tT}$ and $B_{ji}(T)=B_{ji}(0)e^{-2\Delta_tT}$ for $j\ne i$ (the $\signoise=0$ case of \cref{prop:lin-gf-closed}).
\end{enumerate}
\end{lemma}
\begin{proof}
(a) $(\Sigma_0)_{ij}=\nsamp^{-1}\sum_\mu z^\mu_iz^\mu_j=0$ for every $j$. (b) The diffused empirical measure is $\nsamp^{-1}\sum_\mu\mathcal N(e^{-t}z^\mu,\Delta_tI)$ and every centre has zero $i$-th component, so the softmax weights $w_\mu(x)\propto\exp(-\|x-e^{-t}z^\mu\|^2/2\Delta_t)$ do not depend on $x_{t,i}$, the $i$-th marginal is $\mathcal N(0,\Delta_t)$ independent of the other coordinates, and $s^\ast_i=\sum_\mu w_\mu(x)(0-x_{t,i})/\Delta_t=-x_{t,i}/\Delta_t$. (c) is \eqref{eq:q-real}. (d) is \eqref{eq:push-real} with \cref{lem:sampling-real}(c). (e) The gradient with respect to column $i$ is $2\E[(Bx_t-s^\ast)x_{t,i}]=2\big(BM_te_i-\E[s^\ast x_{t,i}]\big)$. Here $M_te_i=\Delta_te_i$; for $j\ne i$, $s^\ast_j$ depends only on $x_{t,-i}$, so $\E[s^\ast_jx_{t,i}]=\E[s^\ast_j]\E[x_{t,i}]=0$; and $\E[s^\ast_ix_{t,i}]=-\E[x_{t,i}^2]/\Delta_t=-1$. Hence $\frac{d}{dT}(Be_i)=-2\Delta_t\big(Be_i+e_i/\Delta_t\big)$.
\end{proof}
The contrast between (d) and (e) is a normalization: the RFNN rate carries the fan-in factor $1/\dlat$ of $\phi$, the linear model does not. Ratios of timescales are normalization-free. (a) and (c) are exact statements about $M_t$; the one-copy eigenvalue of a collapsed coordinate is not exactly $\Delta_t$ but MP-spread around it, (d).

\begin{proposition}[Effect of a block of $n_{\mathrm{dead}}$ collapsed coordinates; proved given \cref{assump:gep-real} and \eqref{eq:Hc-real}]\label{prop:deadblock-real}
Let a latent representation have $d_0$ live coordinates with spectrum $\rho^{\mathrm{live}}$ and $q_0=q[\rho^{\mathrm{live}}]\le50$, and append $n_{\mathrm{dead}}$ collapsed coordinates, $\dlat=d_0+n_{\mathrm{dead}}$. Then, for any $\theta$:
\begin{enumerate}
\item[(i)] $q(\dlat)=\Delta_t+\dfrac{d_0}{\dlat}\,(q_0-\Delta_t)$, strictly decreasing in $n_{\mathrm{dead}}$ with limit $\Delta_t$; $d_{\mathrm{eff}}$ and $\lambda_{(\theta)}$ are unchanged; $d_{\mathrm{floor}}(0)$ increases by $n_{\mathrm{dead}}$.
\item[(ii)] $a_1(q)$ is strictly decreasing in $q$, so $a_1(q(\dlat))^2$ increases with $n_{\mathrm{dead}}$ toward $a_1(\Delta_t)^2$; $a_\ast(q)^2$ is increasing in $q$ on $[\Delta_t,50]$ (verified numerically), so $a_\ast(q(\dlat))^2$ decreases toward $a_\ast(\Delta_t)^2$.
\item[(iii)] Every live data-side eigenvalue is rescaled by $\dfrac{a_1(q(\dlat))^2}{a_1(q_0)^2}\dfrac{d_0}{\dlat}$, and $\taugen^{(\theta)}(\dlat)=\dfrac{\dlat}{a_1(q(\dlat))^2\lambda_{(\theta)}}\ge\dfrac{\dlat}{a_1(\Delta_t)^2\lambda_{(\theta)}}$ grows at least linearly in $n_{\mathrm{dead}}$; the floor bulk has median timescale $\dlat/(a_1(q)^2\Delta_t)$, slower than $\taugen^{(\theta)}$ by $\lambda_{(\theta)}/\Delta_t$, and its lower edge falls by the further factor $(1-d_0/\nsamp)(1-\sqrt{n_{\mathrm{dead}}/(\nsamp-d_0)})^2$.
\item[(iv)] $\taumem^{\mathrm{onset}}(\dlat)=\nsamp/(c_+a_\ast(q(\dlat))^2)$ increases with $n_{\mathrm{dead}}$ but is bounded: $\taumem^{\mathrm{onset}}(\dlat)/\taumem^{\mathrm{onset}}(d_0)\le a_\ast(q_0)^2/a_\ast(\Delta_t)^2$.
\item[(v)] $\dfrac{R^{(\theta)}(\dlat)}{R^{(\theta)}(d_0)}=\dfrac{d_0}{\dlat}\,\dfrac{a_1(q)^2}{a_1(q_0)^2}\,\dfrac{a_\ast(q_0)^2}{a_\ast(q)^2}$, independent of $\theta$ and $\nsamp$ (and of $c_+$ by \eqref{eq:Hc-real}); and $R^{(\theta)}(\dlat)\le\dfrac{\nsamp\lambda_{(\theta)}}{c_+\dlat}\,\dfrac{a_1(\Delta_t)^2}{a_\ast(\Delta_t)^2}=536\,\dfrac{\nsamp\lambda_{(\theta)}}{c_+\dlat}\to0$ as $n_{\mathrm{dead}}\to\infty$.
\end{enumerate}
Collapsed coordinates therefore change the buffer ratio only through $q$, by a bounded factor, and eventually erode it. The count $n_{\mathrm{dead}}$ never appears except through $q$ and the MP width of the floor bulk.
\end{proposition}
\begin{proof}
(i) \cref{lem:dead-real}(c) applied $n_{\mathrm{dead}}$ times; $\lambda_{(\theta)}$ depends only on the live $s_i$. (ii) $a_1(q)=\E[\operatorname{sech}^2(\sqrt qz)]$ is strictly decreasing (\cref{lem:coeffs-buf}(ii)). Monotonicity of $a_\ast^2$ on $[\Delta_t,50]$ is checked by $200$-node Gauss--Hermite quadrature on a $5000$-point grid (minimum increment $3.3\times10^{-5}$; sample values $0.0014,0.031,0.036,0.043,0.048,0.054,0.066,0.074$ at $q=0.181,1.11,1.25,1.47,1.62,1.83,2.27,2.58$; $a_\ast^2\to1-2/\pi$ as $q\to\infty$); all sweeps in this paper have $q\le2.6$. (iii)--(v) are substitutions into \eqref{eq:push-real}, \eqref{eq:mp-floor-real} and \eqref{eq:taus-real}; the bounds use $a_1(q)^2\le a_1(\Delta_t)^2$ and $a_\ast(q)^2\ge a_\ast(\Delta_t)^2$ for $q\in[\Delta_t,50]$.
\end{proof}

\begin{corollary}[Flat effective dimension: what the mechanism predicts; conjectural]\label{cor:flatdeff-real}
Treat the RFNN surrogate at $t=0.1$ with unit-variance first-layer weights as a proxy for the depth-5, width-1024 MLP actually trained (\cref{rem:blunt-real}), and assume \eqref{eq:Hc-real}. On CelebA ($t=0.1$, $\nsamp=1000$) the measured spectra give $\operatorname{tr}\Sigma_0=252.5,255.2,260.4$ at $\dlat=160,180,200$ (increments $0.35,0.14,0.26$ per added coordinate against a live mean $\operatorname{tr}\Sigma_0/d_{\mathrm{eff}}=2.8$--$3.0$), estimator $d_{\mathrm{eff}}=89.5,90.3,87.9$, $\lambda_{\min}-\Delta_t=6\times10^{-5},2\times10^{-5},1\times10^{-5}$, and the median of the lowest $30\%$ of eigenvalues within $2.4\times10^{-4}$ of $\Delta_t$ for $\dlat\ge180$, so at least $15\%$ of coordinates are on the floor. The live spectrum is flat: the hypotheses of \cref{prop:deadblock-real} hold to this accuracy. With $q=1.473,1.342,1.247$, $a_1(q)^2=0.286,0.305,0.320$, $a_\ast(q)^2=0.0434,0.0391,0.0359$, the proposition predicts the factors in \cref{tab:flatdeff-real}. If the memorization-fraction curve depends on training only through the sample bulk, uniformly rescaled by $a_\ast(q)^2$, then every iso-memorization hitting time scales by $a_\ast(q_{160})^2/a_\ast(q_d)^2$ and the final fraction at budget $T_{\max}$ is $F_d(T_{\max})=F_{160}\big(T_{\max}\,a_\ast(q_d)^2/a_\ast(q_{160})^2\big)$, i.e.\ the $\dlat=160$ curve read at $4.51$M and $4.14$M steps. On CIFAR-10 ($\operatorname{tr}\Sigma_0=281$--$299$ for $\dlat=160$--$260$; floor reached to $1.3\times10^{-3}$ at $\dlat=220$ and to $10^{-4}$ for $\dlat\ge240$; $q\;1.621\to1.108$, $a_\ast^2\;0.0480\to0.0312$) the same construction gives the CIFAR-10 rows.
\end{corollary}

\begin{table}[h]
\centering\small
\caption{Predicted (pure $q$-rescaling, \cref{prop:deadblock-real} with (H$_c$)) versus observed changes across the posterior-collapsed part of the real-data sweeps ($t=0.1$, $\nsamp=1000$, five seeds). Hitting-time ratios are relative to $\dlat=160$; thresholds in parentheses. Final fractions at the $5$M-step budget.}
\label{tab:flatdeff-real}
\begin{tabular}{llcc}
\toprule
 & quantity & predicted & observed\\
\midrule
CelebA $160\to180$ & $\taumem^{\rm onset}$ / hitting-time ratio & $\times1.11$ & $\times1.17$--$1.33$ (thr.\ $0.01$--$0.5$)\\
 & $\taugen$, $R$ & $\times1.06$, $\times1.05$ & --\\
 & final fraction & $0.595$ & $0.509$\\
CelebA $160\to200$ & $\taumem^{\rm onset}$ / hitting-time ratio & $\times1.21$ ($\times1.33$ with fan-in init) & $\times1.33,1.26,1.37,1.41$ (thr.\ $0.01,0.05,0.10,0.25$)\\
 & $\taugen$, $R$ & $\times1.12$, $\times1.08$ & --\\
 & final fraction & $0.559$ & $0.448$\\
CIFAR-10 $160\to260$ & $\taumem^{\rm onset}$ / hitting-time ratio & $\times1.54$ ($\times1.93$ with fan-in init) & $\times1.6$ (thr.\ $0.01$), $\times2.4$ ($0.05$, CI $\pm1.7$M), $\times2.7$ ($0.10$, one seed censored)\\
 & $\taugen$, $R$ & $\times1.26$, $\times1.22$ & --\\
 & final fractions ($\dlat=180,\dots,260$) & $0.557,0.552,0.517,0.505,0.472$ & $0.456,0.346,0.213,0.183,0.125$\\
\bottomrule
\end{tabular}
\end{table}
The absolute $q$ depends on the first-layer weight scale of the surrogate (unit variance here; PyTorch fan-in initialization gives $q/3$). With $q\to q/3$ the CelebA $160\to200$ factors become $\taumem\times1.33$, $R\times1.14$ and the CIFAR-10 $160\to260$ factor $\taumem\times1.93$: the sign of every prediction is scale-free, the magnitude is not. The proposition's hypotheses hold only marginally for the real model ($\pwidth=1024$, $\nsamp=1000$, $\dlat\le260$, $\varepsilon_0\approx0.4$--$0.5$), and its features are learned; hence the corollary is a conjectural prediction, not a consequence.

\begin{remark}[What the mechanism can and cannot account for]\label{rem:blunt-real}
When $d_{\mathrm{eff}}$ is flat, added coordinates change $q$ and the floor bulk, not any buffer count. The released $\texttt{buffer\_proxy}=\dlat-\mathrm{round}(d_{\mathrm{eff}})=71,90,112$ enters no timescale. What the $q$-mechanism does deliver is the sign: hitting times lengthen and the final fraction falls as coordinates collapse. Quantitatively it accounts for $\log1.21/\log\{1.41,1.34,1.26\}=0.55$--$0.82$ of the log-slowdown of CelebA hitting times, for $\log(0.559/0.608)/\log(0.448/0.608)=0.28$ of the log-decline of the CelebA final fraction, and on CIFAR-10 for $\log1.54/\log2.55=0.46$ of the hitting-time log-slowdown and $\log(0.472/0.600)/\log(0.125/0.600)=0.15$ of the final-fraction log-decline. The remainder is outside the surrogate. The real-data score model is a depth-5, width-1024 MLP with sinusoidal time features, trained by SGD (learning rate $10^{-3}$, momentum $0.8$) across $t\in[0.01,3]$ on $\nsamp=1000$ encoder means; $\pwidth=1024>\nsamp=1000$, so an $\nsamp-\dlat$ residual bulk is only marginally accommodated, and its first-layer features are learned, so the RFNN quantities are a surrogate evaluated at one $t$. The residual bulk of clustered real features is not Marchenko--Pastur (and $c_+$ is $\dlat$-dependent even on the synthetic sweep); and the decoder's response to imperfectly learned floor coordinates (the slowest data-side modes) is not modelled, although the FID worsening from $38.9$ ($\dlat=140$) to $47$--$50$ ($\dlat=160$--$200$) suggests it is not inert. Over $\dlat=20$--$120$ on CelebA, $q$ is flat ($1.71$--$1.84$, $a_\ast^2=0.051$--$0.055$) because $\operatorname{tr}\Sigma_0$ grows by $\approx2$ per coordinate; there the surrogate predicts no change in $\taumem^{\mathrm{onset}}$ and $\taugen\propto\dlat$, and the observed rise of memorization belongs to the decoder identity-threshold regime of \cref{sec:real-data}. Finally, memorization is not a function of $q$ alone across datasets: at $q\approx1.25$--$1.29$ the final fractions are $0.45$ (CelebA) and $0.21$ (CIFAR-10). The pixel-space NN-ratio has no ambient-dimension drift across the $\dlat$ sweep (decoded images have fixed dimension), but a decoder-dependent drift of its baseline has not been checked.
\end{remark}

\begin{remark}[Encoder means versus posterior samples]\label{rem:samples-real}
The latents here are encoder means, for which a collapsed coordinate is identically zero. Had the diffusion model been trained on posterior samples $z\sim\mathcal N(\mu(x),\operatorname{diag}\sigma^2(x))$, a collapsed coordinate ($\mu_i\equiv0$, $\sigma_i\equiv1$) would have $s_i=1$ and $\lambda_i(M_t)=e^{-2t}+\Delta_t=1$: it would be an isotropic noise-dimension coordinate with $\signoise=1$, the isotropic control, rather than a floor coordinate, and the floor bulk would move up by $1/\Delta_t\approx5.5$ at $t=0.1$. Then $q=\Delta_t+[d_0(q_0-\Delta_t)+n_{\mathrm{dead}}e^{-2t}]/\dlat$ exceeds the mean-latent value, so $a_\ast(q)^2$ is larger and $\taumem^{\mathrm{onset}}$ shorter, while $a_1(q)^2$ is smaller and $\taugen$ longer; $R$ shrinks by both factors. The theory therefore predicts earlier onset memorization for sampled than for mean latents at large $\dlat$, with the caveat that the training latents are then a stochastic set, so the empirical-score target of \cref{lem:dead-real}(b) and the memorization criterion itself change, not only $q$. Untested.
\end{remark}

\begin{remark}[Reporting]\label{rem:xaxis-real}
For every real-data run we recommend reporting, alongside $\dlat$: $d_{\mathrm{eff}}$ (participation ratio of $\Sigma_0$, floor at the known $\Delta_t$), $d_{\mathrm{floor}}(\delta)$ with $\delta=2\times10^{-3}$, $\operatorname{tr}\Sigma_0$, and $q[\rho_{\dlat}]$; and plotting memorization against $d_{\mathrm{eff}}$ and $q$ as well as against $\dlat$. On the present sweeps the points with $\dlat\ge160$ (CelebA) and $\dlat\ge220$ (CIFAR-10) differ from one another in $q$ and $d_{\mathrm{floor}}$ only, so any effect of $\dlat$ there is a $q$ effect, a $c_+$ effect, or is outside the theory.
\end{remark}

\paragraph{Predictions and how to test them.}
\begin{enumerate}
\item \emph{Pure $q$-rescaling:} for CelebA $\dlat\ge160$ every iso-memorization hitting time scales by $a_*(q_{160})^2/a_*(q_d)^2=1.11$ ($180$), $1.21$ ($200$), threshold-independent, and the curves collapse under the step rescaling. Test: executable now from the released mean curves and hitting-time tables (observed $1.26$--$1.41$, roughly threshold-independent, $5$--$20\%$ above the prediction).
\item \emph{The coordinates added beyond $\dlat\sim140$ (CelebA) / $\sim200$ (CIFAR-10) are posterior-collapsed:} $d_{\mathrm{floor}}(2\times10^{-3})\ge0.15\dlat$ and $\operatorname{tr}\Sigma_0$ is flat. Test: regenerate the spectra from the VAE checkpoints (not in the local repository), count floor eigenvalues, and check the MP width \eqref{eq:mp-floor-real}.
\item \emph{Training on posterior samples instead of means at large $\dlat$ increases onset memorization.} Test: new training at CelebA $\dlat=200$ on sampled latents.
\item \emph{Projecting out the floor coordinates restores memorization toward the $\dlat\sim120$--$140$ level}, with $\taumem$ shortened by $a_*(q_{\rm live})^2/a_*(q_{200})^2\approx0.7$--$0.8$. Test: new training on the projected latents, decode by zero-padding.
\item \emph{At fixed nominal $\dlat$, changing the VAE KL capacity reorders memorization by $q$, not by $\dlat$.} Test: two VAEs at $\dlat=200$ with different capacity targets.
\item \emph{Memorization is not a function of $q$ alone across datasets} (already contradicted: $0.448$ vs $0.213$ at matched $q\approx1.25$--$1.29$); recorded as a falsification of the $q$-only reading.
\item \emph{$c_+$ is $\dlat$-dependent on clustered data} ($\sim13\times$ over $\dlat=5..200$, not the MP $3.5\times$). Test: executable now from the saved synthetic spectra.
\end{enumerate}

\paragraph{Open gaps.}
(1) \cref{assump:gep-real} (Gaussian equivalence on real, non-Gaussian latents) is not proved; with one copy per point $\Umat^{\rm res}$ is correlated with $W$, so $c_+$ is a fitted $O(1)$ constant, and on the synthetic spectra it carries a $\sim4\times$ $\dlat$-dependence beyond MP; all real-data ratios are conditional on (H$_c$) and are sign-and-order-of-magnitude statements. (2) The real-data score model is a depth-5, width-1024 MLP (not depth-3/width-256 as an earlier draft said; the main text must match): $\pwidth=1024>\nsamp=1000$, so the hypotheses $\pwidth>\nsamp>\dlat$ and $\varepsilon_0<1$ hold only marginally, the features are learned, and the RFNN quantities at $t=0.1$ are a surrogate whose absolute $q$ depends on the first-layer scale. (3) The pure time-rescaling hypothesis behind the final-fraction predictions is a modelling choice; it accounts for $0.28$ (CelebA) and $0.15$ (CIFAR-10) of the final-fraction log-decline and $0.46$--$0.82$ of the hitting-time log-slowdown. (4) The quoted $d_{\mathrm{eff}}$ values are the released estimator's and equal \cref{def:deff-real} only where $\hat\beta=\Delta_t$; the spectra, VAE checkpoints and diagnostics are absent from the local repository, so $d_{\mathrm{floor}}(\delta)$ can only be bounded, and Predictions~2--5 require those assets or new training. (5) The $\theta$-quantile $\taugen$ is a convention; only $\theta$-free ratios are used. (6) The decoder's response to floor coordinates and the pixel-space NN-ratio baseline drift are unmodelled.

% ----------------------------------------------------------------------
\subsection{Regime boundaries and scope}
\label{app:theory-scope}

This subsection delimits the four-bulk theory: the block counts and edges are leading-order statements under hypotheses (S1)--(S5) below, and several of the paper's configurations sit at or beyond their boundaries. We first fix which feature second moment is meant, then treat three boundaries ($\dlat\to\nsamp$, $\signoise\to0$, $\pwidth$ close to $\dlat+\nsamp$), the fixed-null-energy control and feature learning; the closing \cref{app:theory-outside} lists the out-of-scope configurations explicitly.

\begin{definition}[Aspect ratios, null-data fraction, the two feature second moments, onset statistic]\label{def:scope-ratios}
Let $D:=\dlat-\dint$, $\psi_n:=\nsamp/\dlat$, $\psi_p:=\pwidth/\dlat$, $\psi:=D/\nsamp$ (the $\gamma_\perp$ of \cref{app:theory-fourbulk}), and $\alpha_t^2,\beta_t^2$ as before.
\begin{enumerate}
\item[(i)] The \emph{null-data fraction} $\rho_\perp(t):=e^{-2t}\signoise^2/\beta_t^2\in[0,1)$ is the share of null-block variance that comes from data rather than from diffusion noise.
\item[(ii)] For the population second moment $M_t$ the pre-activation variance obeys the exact identity
\begin{equation}\label{eq:scope-q}
q=\frac{\operatorname{tr}M_t}{\dlat}=\beta_t^2+\frac{\dint(\alpha_t^2-\beta_t^2)}{\dlat},\qquad q_\infty:=\lim_{\dlat\to\infty}q=\beta_t^2
\end{equation}
(\cref{lem:q-p1}). The empirical trace differs by $O(\nsamp^{-1/2})$, and the per-neuron variance $w_i^\top M_tw_i/\dlat$ fluctuates around $q$ with relative width $O(\dlat^{-1/2})$; by Jensen on the concave $q\mapsto\E\tanh^2(\sqrt qz)$ this makes measured traces slightly low at small $\dlat$ ($3\%$ at $\dlat=10$).
\item[(iii)] \emph{Two feature second moments.} $\Umat:=\frac1\nsamp\sum_\mu\E_\eta[\phi(x_t^\mu)\phi(x_t^\mu)^\top]$ is the $\eta$-averaged object; it governs the gradient-flow ODE and is what the saved spectra estimate (Monte Carlo over $50$ draws). $\Umat^{(1)}:=\Phi\Phi^\top/\nsamp$ is the single-draw Gram. All spectral statements below concern $\Umat$ unless $\Umat^{(1)}$ is named.
\item[(iv)] $\lambda_{\rm onset}$ is the largest eigenvalue of $\Umat$ outside the span of the $\dlat$ Gaussian-equivalent linear modes, i.e.\ the top of the residual (``sample'') bulk; $\taumem^{\rm onset}:=1/\lambda_{\rm onset}$ and $\taugen:=1/\lambda^{\rm sig}_{\min}$. For the RFNN these are spectral proxies: the RFNN runs log no sample-based memorization.
\item[(v)] For the null block, with $\hat C_\perp=\frac1\nsamp\sum_\mu\xi^\mu_\perp\xi^{\mu\top}_\perp$ (rank $\min(D,\nsamp)$), $\hat\Sigma_{t,\perp}=e^{-2t}\hat C_\perp+\Delta_t I_D$ and $\Sigma_{t,\perp}=\beta_t^2I_D$, the \emph{sample-specificity} of the converged linear null score is
\begin{equation}\label{eq:scope-epsdef}
\varepsilon_\perp^2:=\frac{\E\|\hat\Sigma_{t,\perp}^{-1}-\Sigma_{t,\perp}^{-1}\|_F^2}{\|\Sigma_{t,\perp}^{-1}\|_F^2}.
\end{equation}
\end{enumerate}
\end{definition}

\begin{assumption}[Four-bulk regime]\label{asm:scope-standing}
\begin{itemize}
\item[(S1)] \emph{Frozen features.} $W$ is fixed at initialization; only $A$ is trained, from $A(0)=0$, by gradient flow on score matching at fixed $t$, so each eigenmode of $\Umat$ decays independently at rate $\lambda_i$.
\item[(S2)] \emph{Ordering.} $\dint<\dlat<\nsamp<\pwidth$.
\item[(S3)] \emph{Spiked-plus-atom spectrum with data-dominated null block.} $\rho_\Sigma$ consists of $\dint$ finite-rank spikes of mean $\alpha_0^2$ (the $k$ centers give distinct spike locations; measured spread $\approx2.3\times$) and one atom $\signoise^2$ of weight $D/\dlat$; and $\rho_\perp(t)\ge\tfrac12$, i.e.\ $e^{-2t}\signoise^2\ge\Delta_t$.
\item[(S4)] \emph{Gaussian equivalence at leading order} \citep{goldt2022gaussian,hu2022universality}: $\phi(x)=a_1(q)\,Wx/\sqrt{\dlat\pwidth}+a_*(q)\,g(x)$ with $g$ uncorrelated with $x$; the signal--null cross covariance of the empirical $\hat M_t$ is neglected (block-diagonal approximation).
\item[(S5)] \emph{Separation.} The noise-dim block lies at least one decade above $\lambda_{\rm onset}$ (\cref{prop:scope-sigperp}).
\end{itemize}
Under (S1)--(S5) the leading-order block positions (block \emph{means}) are
\begin{equation}\label{eq:scope-edges}
\lambda^{\rm sig}\asymp\frac{a_1(q)^2\alpha_t^2}{\dlat},\qquad
\lambda^{\rm nd}\asymp\frac{a_1(q)^2\beta_t^2}{\dlat},\qquad
\bar\lambda^{\rm samp}:=\frac{a_*(q)^2}{\nsamp}\ \text{(residual mass per training point)},
\end{equation}
with the trace check, valid for block means to leading order in $\dlat/\nsamp$ and $1/\nsamp$,
\begin{equation}\label{eq:scope-trace}
\dint\lambda^{\rm sig}+D\lambda^{\rm nd}+a_*(q)^2\approx\E[\tanh(\sqrt q z)^2]
\end{equation}
(the exact form is \eqref{eq:fourbulk-trace}). The residual mass $a_*(q)^2$ is shared between the residual bulk (measured share $0.5$--$0.7$) and $O(a_*^2/\nsamp)$ per linear mode; the residual bulk is heavy-tailed, so its median is not a location statistic and the check must use means. Measured ($\signoise=0.5$, code units$/\pwidth$): $0.456$ vs $\E\tanh^2=0.470$ ($\dlat=10$), $0.2915$ vs $0.2974$ ($40$), $0.2092$ vs $0.2093$ ($200$).
\end{assumption}

\subsubsection{Which $\Umat$: block counts for the two feature second moments}

\begin{proposition}[Block counts; proved given (S4)]\label{prop:scope-counts}
Assume (S1), (S4) and $\pwidth>\nsamp$.
\begin{itemize}
\item[(a)] \emph{Single-draw Gram.} $\operatorname{rank}\Umat^{(1)}=\nsamp$ a.s., and its non-zero spectrum splits as
\begin{equation}\label{eq:scope-counts}
N_{\rm sig}=\dint,\qquad N_{\rm nd}=\min(\dlat,\nsamp)-\dint,\qquad N_{\rm samp}=(\nsamp-\dlat)_+,\qquad N_{\rm null}=\pwidth-\nsamp ,
\end{equation}
with a hard cliff at index $\nsamp$; the residual bulk is empty for $\psi_n\le1$.
\item[(b)] \emph{$\eta$-averaged $\Umat$.} The linear part is $a_1(q)^2W\hat M_tW^\top/(\dlat\pwidth)$ with $\hat M_t=e^{-2t}\hat C+\Delta_tI_{\dlat}$ of full rank $\dlat$ for every $\dlat$. Hence $N_{\rm sig}=\dint$ and $N_{\rm nd}=D$ for all $\dlat$: for $D\le\nsamp$ all $D$ noise-dim modes lie in the data-broadened band of \eqref{eq:scope-ndmin}, and for $D>\nsamp$ exactly $\nsamp$ of them do while $D-\nsamp$ sit on the diffusion floor $a_1(q)^2\Delta_t(1\pm\psi_p^{-1/2})^2/\dlat$. The residual part $\frac{a_*^2}{\nsamp}\sum_\mu\E_\eta[g(x_t^\mu)g(x_t^\mu)^\top]$ has rank $>\nsamp$: there is no cliff at index $\nsamp$ and no rank-null block, only a soft residual tail (measured $e_{\nsamp-1}/e_{\nsamp}=1.32$ at $\dlat=40$, $1.73$ at $\dlat=200$; $0.3\%$ of the trace beyond index $\nsamp$).
\item[(c)] \emph{Eigenvectors.} In both cases the noise-dim eigenvectors are $W_\perp v/\|W_\perp v\|+O(\psi_p^{-1/2})$ with $v$ ranging over eigenvectors of the empirical noised null second moment; for $\Umat$ that matrix is $\hat\Sigma_{t,\perp}$, whose non-floor eigenvectors are those of $\hat C_\perp$, so they are determined by the training noise vectors $\{\xi^\mu_\perp\}$.
\end{itemize}
\end{proposition}
\begin{proof}
(a) Under (S4), $\Phi=a_1(q)WX_t/\sqrt{\dlat\pwidth}+a_*(q)G$ with $X_t\in\R^{\dlat\times\nsamp}$ a single draw; $\operatorname{rank}(WX_t)=\operatorname{rank}X_t=\min(\dlat,\nsamp)$ a.s.\ (the $\sqrt{\Delta_t}\eta$ term makes $X_t$ generic; $W$ has full column rank since $\pwidth>\dlat$), of which $\dint$ directions carry the signal spikes. $\Phi$ has $\nsamp$ columns, so $\operatorname{rank}\Umat^{(1)}\le\nsamp$, with equality a.s.\ since $G$ is generic; the residual therefore adds $\nsamp-\min(\dlat,\nsamp)$ directions and $\pwidth-\nsamp$ are annihilated (\cref{lem:fourbulk-counts}). (b) $\E_\eta[x_tx_t^\top]=e^{-2t}xx^\top+\Delta_tI$, so the $\eta$-average of the linear Gram is $a_1^2W\hat M_tW^\top/(\dlat\pwidth)$ with $\hat M_t\succeq\Delta_tI$ of full rank; its null block is $\hat\Sigma_{t,\perp}=e^{-2t}\signoise^2\Xi\Xi^\top/\nsamp+\Delta_tI_D$ with $\Xi\in\R^{D\times\nsamp}$ standard Gaussian, whose $\min(D,\nsamp)$ non-trivial eigenvalues are $e^{-2t}\signoise^2m_i+\Delta_t$ ($m_i$ Wishart eigenvalues) and whose remaining $(D-\nsamp)_+$ eigenvalues equal $\Delta_t$ exactly. The residual $\eta$-average is a sum of $\nsamp$ PSD matrices each of rank $>1$, so no cliff occurs. (c) With $W_\perp^\top W_\perp/\pwidth=I_D+O(\psi_p^{-1/2})$ the eigenpairs of $W_\perp\hat\Sigma_{t,\perp}W_\perp^\top$ are $(\lambda_i\pwidth,W_\perp v_i)$ to that order; the isotropic $\Delta_t$ shift does not change the eigenvectors of $\hat\Sigma_{t,\perp}$ relative to $\hat C_\perp$.
\end{proof}

\subsubsection{The boundary $\dlat\to\nsamp$}

\begin{proposition}[Sample-specificity of the learned null score; proved in the block-diagonal approximation]\label{prop:scope-nullspec}
Model the null block of the score linearly, $s_\perp(x)=B_\perp x_\perp$, trained by gradient flow on denoising score matching over the $\nsamp$ training points at fixed $t$, in the block-diagonal approximation (\cref{prop:lin-emp}). Then $B_\perp(T)\to-\hat\Sigma_{t,\perp}^{-1}$ (population target $-\Sigma_{t,\perp}^{-1}=-I_D/\beta_t^2$), and in the proportional limit $D,\nsamp\to\infty$, $D/\nsamp\to\psi$,
\begin{equation}\label{eq:scope-eps}
\varepsilon_\perp^2\;\to\;\int\frac{\rho_\perp^2\,(m-1)^2}{\big(1+\rho_\perp(m-1)\big)^2}\,d\mu_{\psi}(m),
\end{equation}
$\mu_\psi$ the Marchenko--Pastur law of ratio $\psi$ with unit mean \citep{marchenko1967distribution}. Consequently:
\begin{itemize}
\item[(i)] $\varepsilon_\perp^2=\rho_\perp^2\psi\,(1-2\rho_\perp\psi+O(\rho_\perp^2\psi))$ for $\rho_\perp\psi\to0$; $\varepsilon_\perp^2$ depends on $(\dlat,\nsamp)$ only through $\psi=D/\nsamp$ and on $\signoise$ only through $\rho_\perp$. (The expansion is poor for $\rho_\perp\approx1$, $\psi\gtrsim0.1$; use \eqref{eq:scope-eps} numerically there.)
\item[(ii)] For $\psi>1$ the atom of $\mu_\psi$ at $0$ gives $\varepsilon_\perp^2\ge(1-\psi^{-1})\rho_\perp^2/(1-\rho_\perp)^2$: on the $D-\nsamp$ directions orthogonal to $\operatorname{span}\{\xi^\mu_\perp\}$ the learned score is $-x_\perp/\Delta_t$ instead of $-x_\perp/\beta_t^2$, a contraction stronger by $(1-\rho_\perp)^{-1}$; the reverse dynamics confine generated null components to the span of the training noise vectors.
\item[(iii)] $\varepsilon_\perp^2\to0$ as $\rho_\perp\to0$ at fixed $\psi$: at $\signoise=0.01$, $t=0.01$, $\rho_\perp=4.9\times10^{-3}$ and $\varepsilon_\perp^2<10^{-4}$ for all $\dlat\le1000$.
\end{itemize}
The neglected signal--null cross covariance $\frac1\nsamp\sum_\mu(m_{c(\mu)}+\xi^\mu_{\rm sig})\xi^{\mu\top}_\perp$ has entries of order $\alpha_t\beta_t/\sqrt\nsamp$ and contributes a sample-specific error of the same order as $\varepsilon_\perp^2$ when $\psi$ is small.
\end{proposition}
\begin{proof}
With DSM target $y=-\eta/\sqrt{\Delta_t}$ and $x_t=e^{-t}x+\sqrt{\Delta_t}\eta$, the null-block loss over the empirical measure is $\E\|B_\perp x_{t,\perp}-y_\perp\|^2$; its gradient flow is $\dot B_\perp=-2(B_\perp\hat\Sigma_{t,\perp}+I_D)$ because $\E[x_{t,\perp}x_{t,\perp}^\top]=\hat\Sigma_{t,\perp}$ over the training points (cluster means have zero null component) and $\E[y_\perp x_{t,\perp}^\top]=-I_D$. Hence $B_\perp(T)=-\hat\Sigma_{t,\perp}^{-1}+(B_\perp(0)+\hat\Sigma_{t,\perp}^{-1})e^{-2\hat\Sigma_{t,\perp}T}$. $\hat C_\perp=\signoise^2\,\Xi\Xi^\top/\nsamp$, so its eigenvalues are $\signoise^2m_i$ with $\{m_i\}$ Wishart eigenvalues of ratio $\psi$ whose empirical law converges weakly to $\mu_\psi$. Since $e^{-2t}\signoise^2m+\Delta_t=\beta_t^2(1+\rho_\perp(m-1))$,
$$\|\hat\Sigma_{t,\perp}^{-1}-\Sigma_{t,\perp}^{-1}\|_F^2=\frac1{\beta_t^4}\sum_i\frac{\rho_\perp^2(m_i-1)^2}{(1+\rho_\perp(m_i-1))^2}.$$
Dividing by $\|\Sigma_{t,\perp}^{-1}\|_F^2=D/\beta_t^4$ gives the normalized trace of a bounded continuous function of $\hat C_\perp$ (denominator $\ge1-\rho_\perp>0$ on $m\ge0$), so weak convergence yields \eqref{eq:scope-eps}. (i): expand $(1+\rho_\perp(m-1))^{-2}$ and use the MP central moments $\int(m-1)^2d\mu_\psi=\psi$, $\int(m-1)^3d\mu_\psi=\psi^2$, $\int(m-1)^4d\mu_\psi=2\psi^2+\psi^3$. (ii): $\mu_\psi=(1-\psi^{-1})\delta_0+\psi^{-1}\tilde\mu$ for $\psi>1$, and the integrand at $m=0$ equals $\rho_\perp^2/(1-\rho_\perp)^2$; the $m_i=0$ eigenvectors are the complement of $\operatorname{span}\{\xi^\mu_\perp\}$, on which $\hat\Sigma_{t,\perp}=\Delta_tI$. (iii): the integrand is $O(\rho_\perp^2)$ uniformly on the support once $\rho_\perp\le\tfrac12$.
\end{proof}
Numerically, at $\signoise=0.5$, $t=0.01$ ($\rho_\perp=0.925$, $\nsamp=500$), \eqref{eq:scope-eps} gives $\varepsilon_\perp^2=0.028,\,0.103,\,0.302,\,0.638,\,1.19$ at $\dlat=20,50,100,150,200$; the first-order form $\rho_\perp^2\psi$ gives $0.026,0.077,0.163,0.248,0.334$ and under-estimates from $\dlat\approx100$ on because small eigenvalues of $\hat C_\perp$ are amplified by $(1+\rho_\perp(m-1))^{-2}$. At $\signoise=1$ ($\rho_\perp=0.98$): $0.032,0.122,0.379,0.867,1.80$.

\begin{proposition}[Onset statistics across the $\dlat\to\nsamp$ boundary; sketched]\label{prop:scope-rebound}
Assume (S1), (S3), (S4) and $\pwidth>\nsamp$; $\Umat$ is the $\eta$-averaged second moment.
\begin{itemize}
\item[(a)] \emph{Residual-bulk onset.} In the isotropic-residual approximation (columns $g(x^\mu_t)$ with uncorrelated entries of variance $1/\pwidth$),
\begin{equation}\label{eq:scope-onset-iso}
\lambda_{\rm onset}^{\rm iso}=\frac{a_*(q)^2}{\nsamp}\Big(1+\sqrt{\nsamp/\pwidth}\Big)^2,\qquad \taumem^{\rm onset,iso}=\frac{\nsamp}{a_*(q)^2(1+\sqrt{\nsamp/\pwidth})^2},
\end{equation}
which depends on $\dlat$ only through $q\downarrow\beta_t^2$ in $a_*(q)^2$. On the repo spectra ($\signoise=0.5$, $\pwidth=64\dlat$) the measured onset exceeds \eqref{eq:scope-onset-iso} by $4.6,4.2,2.8,2.2,1.7,1.5,1.1,1.0$ at $\dlat=10,20,40,60,80,100,150,200$: the excess is a small-$\dlat$ effect of cluster-level residual spikes (the $g(\dlat)$ of \cref{prop:onset-buf}, the collective modes of \cref{rem:collective-pred}) and disappears at $\dlat\ge150$, where \eqref{eq:scope-onset-iso} is quantitatively right ($0.173$ vs $0.171$ code units at $\dlat=200$). Because the excess itself decays, \eqref{eq:scope-onset-iso} under-predicts the $\dlat$-trend: the measured onset falls $107\times$ from $\dlat=10$ to $200$ ($221\times$ from $\dlat=5$), the formula predicts $24\times$ ($\approx60\times$).
\item[(b)] \emph{Noise-dim top for $D>\nsamp$.} The residual bulk does not empty at $\psi_n=1$ (\cref{prop:scope-counts}(b)); its top stays at \eqref{eq:scope-onset-iso}. The noise-dim block keeps $\nsamp$ data-broadened modes with top
\begin{equation}\label{eq:scope-onset-rd}
\lambda^{\rm nd}_{\max}\asymp\frac{a_1(q)^2}{\dlat}\Big[e^{-2t}\signoise^2\big(1+\sqrt{\psi}\big)^2+\Delta_t\Big]\big(1+\psi_p^{-1/2}\big)^2\;\xrightarrow{\psi_n\to0}\;\frac{a_1(\beta_t^2)^2\,e^{-2t}\signoise^2\,(1+\psi_p^{-1/2})^2}{\nsamp},
\end{equation}
so $1/\lambda^{\rm nd}_{\max}\to\nsamp/(a_1(\beta_t^2)^2e^{-2t}\signoise^2(1+\psi_p^{-1/2})^2)$: $\dlat$-independent and linear in $\nsamp$. Its ratio to the large-$\dlat$ limit of $\taumem^{\rm onset,iso}$ tends to $a_*(\beta_t^2)^2/(a_1(\beta_t^2)^2e^{-2t}\signoise^2)=0.019$ at $\signoise=0.5$, $t=0.01$ as $\pwidth/\dlat,\pwidth/\nsamp\to\infty$ ($0.015$ at $\psi_p=64$). These modes are sample-specific in the sense of \eqref{eq:scope-epsdef}: their eigenvectors are fixed by $\{\xi^\mu_\perp\}$ and their learned content deviates from the population score by $\varepsilon_\perp^2\ge\rho_\perp^2(1-\psi^{-1})/(1-\rho_\perp)^2$ (\cref{prop:scope-nullspec}(ii)). In this regime the earliest sample-specific modes are the noise-dim modes, learned $\approx50\times$ faster than the residual bulk: the population/sample distinction of the noise-dim ``buffer'' dissolves.
\end{itemize}
\end{proposition}
\begin{proof}[Proof sketch]
(a) Under (S4) the residual part of $\Umat$ is $\frac{a_*^2}{\nsamp}\sum_\mu\E_\eta[gg^\top]$; with uncorrelated columns of variance $1/\pwidth$ it is $a_*^2/\nsamp$ times a Wishart matrix of ratio $\nsamp/\pwidth$ with top edge $(1+\sqrt{\nsamp/\pwidth})^2$; the residual bulk is this spectrum projected onto the complement of the linear span and by interlacing its top is at most the unprojected edge. The measured/iso ratios use the sample-bulk tops $0.93,0.374,0.198,0.173$ (code units, $\dlat=10,40,100,200$) against $\pwidth\lambda^{\rm iso}_{\rm onset}=0.203,0.133,0.135,0.171$ with $a_*(q)^2=0.0447,0.0125,0.00644,0.00467$. (b) The null-block linear part $a_1^2W_\perp\hat\Sigma_{t,\perp}W_\perp^\top/(\pwidth\dlat)$ has non-zero eigenvalues equal to those of $a_1^2\hat\Sigma_{t,\perp}^{1/2}(W_\perp^\top W_\perp/\pwidth)\hat\Sigma_{t,\perp}^{1/2}/\dlat$; the top edge of $\hat\Sigma_{t,\perp}$ is $e^{-2t}\signoise^2(1+\sqrt\psi)^2+\Delta_t$ (for $\psi>1$ the non-zero Wishart eigenvalues lie in $[(\sqrt\psi-1)^2,(\sqrt\psi+1)^2]$) and that of $W_\perp^\top W_\perp/\pwidth$ is $(1+\sqrt{D/\pwidth})^2$; the product of top edges is the top edge of the free multiplicative convolution to leading order in the broadening. As $\psi_n\to0$, $(1+\sqrt\psi)^2/\dlat\to1/\nsamp$ and $q\to\beta_t^2$. Numerics: $a_*(\beta_t^2)^2=3.09\times10^{-3}$, $a_1(\beta_t^2)^2e^{-2t}\signoise^2=0.671\times0.245=0.164$.
\end{proof}

\begin{remark}[Conjecture: return of memorization through the null block at large $\dlat$]\label{rem:scope-rebound-conj}
\cref{prop:scope-rebound}(b) shows that for $\psi_n\lesssim1$ the fast noise-dim modes are already sample-specific. We \emph{conjecture} that this is the channel through which memorization returns at large $\dlat$, and state what the conjecture does and does not contain.
\begin{itemize}
\item[(i)] \emph{What a linear null score cannot do.} The converged null score $-\hat\Sigma_{t,\perp}^{-1}x_\perp$ is linear; under the reverse dynamics it generates null components with the empirical covariance $\hat C_\perp$ (confined to $\operatorname{span}\{\xi^\mu_\perp\}$ when $D>\nsamp$), i.e.\ $x_\perp\approx\Xi c/\sqrt\nsamp$ with $c\sim N(0,I_\nsamp)$, whose squared distance to the nearest training point is $\approx2D\signoise^2\,(1-O(\sqrt{\log\nsamp/D}))$ and whose NN-ratio is $\approx1$. A linear null score therefore \emph{cannot} by itself produce NN-ratio memorization (ratio $<1/3$); that requires the softmax concentration of the empirical score \citep{biroli2024dynamical,scarvelis2023closedform,gu2023memorization}, which lives in the nonlinear (residual) part of the features (\cref{rem:lin-nomem}). $\varepsilon_\perp$ is proposed only as a \emph{necessary-condition proxy}: it measures how much sample-specific null-block information the fast modes carry.
\item[(ii)] \emph{Conditional functional form.} If a metric-dependent threshold $\varepsilon_c$ exists such that memorization appears once $\varepsilon_\perp\ge\varepsilon_c$, then, since $\varepsilon_\perp^2$ depends on $(\dlat,\nsamp)$ only through $\psi$ and on $\signoise$ only through $\rho_\perp$,
\begin{equation}\label{eq:scope-rebound}
\dlat^{\rm reb}=\dint+\nsamp\,\psi^{\rm reb}(\rho_\perp),\qquad \varepsilon_\perp^2\big(\rho_\perp,\psi^{\rm reb}\big)=\varepsilon_c^2 ,
\end{equation}
linear in $\nsamp$ at fixed $(\signoise,t)$, absent as $\rho_\perp\to0$, and earlier for larger $\rho_\perp$.
\item[(iii)] \emph{Evidence and confounds.} In the scaled-width MLP sweep (hidden $8\dlat=800,1200,1600$; $\signoise=0.5$, $\nsamp=500$, $3\times10^5$ steps) the memorization fraction is $0.0000,0.0022,0.0144$ at $\dlat=100,150,200$ and still rising ($0.004\to0.0144$ over the final $5\times10^4$ steps at $\dlat=200$), with the mean NN-ratio falling $0.755\to0.713$; the matching $\signoise=0.01$ scaled sweep gives $0$ at all three with the NN-ratio rising $1.98\to2.59$ (distance concentration). In the paper's \emph{fixed-width} sweep (hidden $256$) at $\signoise=0.5$ the fraction is $0.0000$ at every evaluation for $\dlat=100,150,200$. Hence the rebound requires \emph{both} $\rho_\perp\approx1$ and capacity growing with $\dlat$: the $\signoise$ contrast at equal capacity rules out capacity as the sole cause, the fixed-width null rules out $\rho_\perp$ as the sole cause. Calibrating on the scaled sweep gives $\varepsilon_c^2\in(0.30,0.64)$ (the values at $\dlat=100$ and $150$) and $\psi^{\rm reb}(0.925)\in(0.19,0.29)$; this is a fit to those two points, not a test. Under \eqref{eq:scope-rebound} the discriminating predictions are $\dlat^{\rm reb}\approx55$--$80$ at $\nsamp=250$ and $200$--$300$ at $\nsamp=1000$ (same capacity schedule), and $\dlat^{\rm reb}\in[88,130]$ at $\signoise=1$, $\nsamp=500$.
\end{itemize}
\end{remark}

\subsubsection{The boundary $\signoise\to0$ and the merge condition}

\begin{proposition}[$\signoise\to0$; separation threshold; sketched]\label{prop:scope-sigperp}
Assume (S1), (S2), (S4); $\Umat$ is the $\eta$-averaged second moment.
\begin{itemize}
\item[(i)] As $\signoise\to0$: $\beta_t^2\to\Delta_t$, $\rho_\perp\to0$, $q\to\Delta_t+\dint(\alpha_t^2-\Delta_t)/\dlat$, all eigenvalues of $\hat\Sigma_{t,\perp}$ tend to $\Delta_t$ (no data-side broadening), and the noise-dim block sits on the diffusion floor
\begin{equation}\label{eq:scope-floor}
\lambda^{\rm nd}\to\frac{a_1(q)^2\Delta_t}{\dlat}\big(1\pm\psi_p^{-1/2}\big)^2,
\end{equation}
broadened only by $W$ (factor $1.65$ at $\psi_p=64$; measured $1.5$--$1.6$ at $\signoise=0.01$), with $N_{\rm nd}=D$. These modes carry no data variance, the learned null score equals the population one ($\varepsilon_\perp\to0$), and the data lie in the $\dint$-subspace, the setting of the linear pencil of \citet{george2025dsm} (\cref{rem:fourbulk-sigperp0}).
\item[(ii)] Including the data-side (ratio $\psi$) and width-side (ratio $D/\pwidth\approx\psi_p^{-1}$) Marchenko--Pastur broadenings, the lower edge of the noise-dim block is
\begin{equation}\label{eq:scope-ndmin}
\lambda^{\rm nd}_{\min}\asymp\frac{a_1(q)^2}{\dlat}\Big[e^{-2t}\signoise^2(1-\sqrt\psi)^2+\Delta_t\Big](1-\psi_p^{-1/2})^2\qquad(\psi<1),
\end{equation}
and with $G:=\log_{10}(\lambda^{\rm nd}_{\min}/\lambda_{\rm onset})$ the one-decade criterion $G\ge1$ of (S5) holds iff
\begin{equation}\label{eq:scope-sigthr}
\signoise^2\;\ge\;\frac{e^{2t}}{(1-\sqrt\psi)^2}\Big[\frac{10\,\dlat\,\lambda_{\rm onset}}{a_1(q)^2(1-\psi_p^{-1/2})^2}-\Delta_t\Big]=:(\signoise^*)^2 .
\end{equation}
\item[(iii)] With the measured onsets of the repo spectra ($t=0.01$, $\pwidth=64\dlat$, $\nsamp=500$, $\signoise=0.5$): $\signoise^*=0.91,0.70,0.50,0.45,0.41,0.40,0.42,0.52$ at $\dlat=10,20,40,60,80,100,150,200$. Hence $\signoise=0.5$ is below threshold for $\dlat\le20$ and at or above it for $\dlat\ge40$, matching the measured $G<1$ for $\dlat\le20$ and $G\ge1$ for $\dlat\ge30$ (measured $G=0.64,0.89,1.12,1.19,1.27,1.26,1.26,1.09$; predicted from \eqref{eq:scope-ndmin}: $0.51,0.73,0.99,1.08,1.14,1.15,1.11,0.98$, the product-of-edges lower edge under-estimating $\lambda^{\rm nd}_{\min}$ by $\approx25\%$). $\signoise=0.01$ lies below the floor ($\signoise^2=5\times10^{-3}\Delta_t$) and the measured separation there is $G=0.004$--$0.056$ decades: a three-bulk spectrum. The four-bulk picture applies only for $\signoise\gtrsim\signoise^*$.
\end{itemize}
\end{proposition}
\begin{proof}[Proof sketch]
(i) Substitute in \cref{def:scope-ratios} and \eqref{eq:scope-edges}; $\varepsilon_\perp^2\to0$ by \cref{prop:scope-nullspec}(iii); the null block of $\hat M_t$ is $e^{-2t}\hat C_\perp+\Delta_tI_D\to\Delta_tI_D$, so the only remaining broadening is that of $W_\perp^\top W_\perp/\pwidth$, with edges $(1\pm\sqrt{D/\pwidth})^2$. (ii) The non-zero eigenvalues of $a_1^2W_\perp\hat\Sigma_{t,\perp}W_\perp^\top/(\pwidth\dlat)$ equal those of $a_1^2\hat\Sigma_{t,\perp}^{1/2}(W_\perp^\top W_\perp/\pwidth)\hat\Sigma_{t,\perp}^{1/2}/\dlat$; the lower edges of the two factors are $e^{-2t}\signoise^2(1-\sqrt\psi)^2+\Delta_t$ and $(1-\sqrt{D/\pwidth})^2$, and their product is the lower edge of the free multiplicative convolution to leading order in the broadening (exact edges from the Silverstein equation, cf.\ \citealp{ELKaroui2010spectrum}). Solve $\lambda^{\rm nd}_{\min}\ge10\lambda_{\rm onset}$ for $\signoise^2$. (iii) Code units are $\pwidth\times$ theory units. $\dlat=10$: $\lambda_{\rm onset}=0.93/640=1.45\times10^{-3}$, $a_1(q)^2=0.281$, $\psi=0.01$, giving $(\signoise^*)^2=e^{0.02}(0.9)^{-2}[100\cdot1.45\times10^{-3}/(0.281\cdot0.766)-0.0198]=0.83$. $\dlat=40$: $\lambda_{\rm onset}=0.374/2560=1.46\times10^{-4}$, $a_1(q)^2=0.494$, $\psi=0.07$, giving $(\signoise^*)^2=1.0202\cdot1.849\cdot[400\cdot1.46\times10^{-4}/(0.494\cdot0.766)-0.0198]=0.256$, $\signoise^*=0.50$. The remaining entries are the same evaluation at the other $\dlat$. Measured $G$ uses only the noise-dim minimum and the sample-bulk top (index $\dlat+1$) of the released block summaries (whose block partition otherwise follows the superseded cumulative count): $\signoise=0.5$: $0.669,0.644,0.784,0.888,0.996,1.12,1.19,1.27,1.25,1.26,1.09$ at $\dlat=8,10,15,20,30,40,60,80,100,150,200$; $\signoise=0.01$: $0.045,0.032,0.019,0.006,0.004,0.056$ at $\dlat=8,\dots,40$.
\end{proof}

\subsubsection{$\pwidth=\dlat+\nsamp+300$: $\psi_p$ decreasing toward $1$}

\begin{remark}[$\pwidth=\dlat+\nsamp+300$]\label{rem:scope-psip}
This configuration was run only at $\dlat=20$ (the $\dint$ sweep); the statements for $\dlat\ge40$ are predictions. One has $\pwidth>\nsamp$ for all $\dlat$ and $\psi_p=1+800/\dlat\ge4.3$ for $\dlat\le240$; the margin $\psi_p-1-\psi_n=300/\dlat$ of the proportional asymptotics of \citet{bonnaire2025} stays positive but is not uniform in $\dlat$. The block \emph{means} \eqref{eq:scope-edges} are $\pwidth$-independent (a $5.8\times$ change of $\pwidth$ moves them by $<0.5\%$). What changes is the block \emph{width}: the factor $(1\pm\sqrt{D/\pwidth})^2$ in \eqref{eq:scope-ndmin} has $\sqrt{D/\pwidth}=0.22$ at $\dlat=40$ (noise-dim block spanning a factor $2.4$) and $0.44$ at $\dlat=200$ (factor $6.8$, i.e.\ $0.83$ decades from $W$ alone), against $0.125$ (factor $1.65$) at $\psi_p=64$. (S5) then requires $G\ge1$ \emph{after} this broadening, which the product-of-edges estimate predicts to fail at $\signoise=0.5$ for $\dlat\gtrsim100$. On the residual side $\nsamp/\pwidth\to0.5$ at $\dlat=200$ and the MP lower edge $(1-\sqrt{\nsamp/\pwidth})^2=0.086$ pushes the slowest residual modes toward zero (the random-features interpolation peak of \citealp{mei2022generalization}); this stretches the time to \emph{complete} memorization but leaves $\lambda_{\rm onset}$ within $(1+\sqrt{\nsamp/\pwidth})^2\le2.9$ of the residual mean.
\end{remark}

\subsubsection{The fixed-null-energy control}

\begin{proposition}[$\signoise^2=c/D$; sketched]\label{prop:scope-control}
Assume (S1), (S2), (S4) and $\signoise^2=c/D$ with $c>0$ fixed. Then
\begin{equation}\label{eq:scope-ctrl}
\beta_t^2=\Delta_t+\frac{e^{-2t}c}{D},\qquad \rho_\perp=\frac{e^{-2t}c}{e^{-2t}c+\Delta_tD},\qquad q=\Delta_t+\frac{\dint(\alpha_t^2-\Delta_t)+e^{-2t}c}{\dlat},
\end{equation}
\begin{equation}\label{eq:scope-ctrl-edges}
\lambda^{\rm sig}\asymp\frac{a_1(q)^2\alpha_t^2}{\dlat},\qquad
\lambda^{\rm nd}\asymp\frac{a_1(q)^2}{\dlat}\Big(\Delta_t+\frac{e^{-2t}c}{D}\Big),\qquad
\bar\lambda^{\rm samp}=\frac{a_*(q)^2}{\nsamp},
\end{equation}
with crossover $D_c:=e^{-2t}c/\Delta_t\approx49.5\,c$ at $t=0.01$. For $D\gg D_c$ the control is spectrally the $\signoise\to0$ limit of \cref{prop:scope-sigperp}: (i) the noise-dim block collapses onto the diffusion floor and the four-bulk separation is lost; (ii) $\taugen=\dlat/(a_1(q)^2\alpha_t^2)$ keeps its linear growth, with $a_1(q)^2\to a_1(\Delta_t)^2=0.96$ (vs.\ $0.63$ at $\signoise=0.5$, $\dlat=200$); (iii) $a_*(q)^2=\tfrac23q^3+O(q^4)$ collapses faster than in the $\signoise=0.5$ sweep, so the frozen-feature delay of the residual bulk persists (cf.\ \cref{cor:fixedenergy-buf}); (iv) $\varepsilon_\perp^2\approx\rho_\perp^2\psi\approx c^2e^{-4t}/(\Delta_t^2D\nsamp)\to0$, so no null-block rebound (\cref{rem:scope-rebound-conj}) is predicted. The control thus separates the \emph{count} $D$ acting through $q$-dilution (retained) from the \emph{null data variance} acting through $\rho_\perp$ (removed): the frozen-feature theory predicts that the memorization delay survives the control while the four-bulk structure and the rebound do not.
\end{proposition}
\begin{proof}[Proof sketch]
Substitute $\signoise^2=c/D$ into \cref{def:scope-ratios} and \eqref{eq:scope-edges}; the trace check \eqref{eq:scope-trace} holds by \eqref{eq:scope-q}. (i) is \cref{prop:scope-sigperp}(i)--(ii). (ii) $\alpha_t^2$ is $\signoise$-independent and $a_1$ is decreasing in $q$; $a_1(0.0198)^2=0.962$. (iii) is \cref{lem:q-p1}. (iv) is \cref{prop:scope-nullspec}(i). The prediction that the delay persists is a frozen-feature statement, not a statement about the trainable MLP.
\end{proof}

\subsubsection{Feature learning}

\begin{remark}[Conjecture: one-step feature learning]\label{rem:scope-feature}
Everything above is a frozen-feature (S1) result. For the trainable MLP we conjecture the following, modelled on the one-gradient-step analysis of \citet{ba2022highdimensional}. Let $A$ first be trained for a time $T_0$ (so that $A\neq0$; at $A(0)=0$ the $W$-gradient vanishes identically), then take one step $W_1=W_0-\eta\nabla_WL$ with
$$\nabla_WL=\frac{1}{\nsamp\sqrt{\dlat\pwidth}}\sum_\mu\big[(A^\top r_\mu)\odot\tanh'(W_0x^\mu_t/\sqrt{\dlat})\big]\,x_t^{\mu\top},\qquad r_\mu=s_A(x^\mu_t)-y_\mu .$$
The update is a sum of rank-one terms with row space in $\operatorname{span}\{x^\mu_t\}$ whose signal- and null-block energies scale as $\alpha_t^2$ and $\beta_t^2$ per dimension, so the spike it adds to the conjugate kernel $\Umat_1=\frac1\nsamp\sum_\mu\E_\eta\phi_1\phi_1^\top$ loads on the signal subspace with per-dimension null leakage $O(\beta_t^2/\alpha_t^2)$ ($0.096$ at $\signoise=0.5$, $0.007$ at $0.01$). Because $\E[rx_t^\top]=A\E[\phi x_t^\top]+I$ contains an isotropic term from the DSM target, the \emph{total} null energy of the spike, $D\beta_t^2$ against $\dint\alpha_t^2$, can be comparable at large $\dlat$; whether this shifts the noise-dim block is what (P-i) asks to be measured.
\textbf{(P-i)} After the step the signal-block eigenvalues increase by $O(\eta)$ (the cross term of $\Phi_1\Phi_1^\top$ is first order in the step; a rank-$\le\dint$ spike, BBP-type, \citealp{baik2005phase}) while the noise-dim edges move by at most $O(\eta\beta_t^2/\alpha_t^2)$: $\taugen$ decreases, the noise-dim timescale does not.
\textbf{(P-ii)} Under continued training the first layer acquires components aligned with individual training points (the nonlinear part of the score target is the softmax assignment of the empirical score, \citealp{biroli2024dynamical,scarvelis2023closedform}), raising the top of the residual bulk and lowering $\taumem^{\rm onset}$; the buffer in absolute steps, $\taumem^{\rm onset}-\taugen$, is then \emph{shorter} than the frozen-feature value even though (P-i) alone lengthens it in ratio.
Status: conjecture. The saved RFNN pre/post spectra cannot test it: $W$ is frozen there and the pre/post block medians agree to $<0.5\%$, which only confirms (S1). The test requires logging the first-layer conjugate-kernel spectrum of the MLP at initialization and during training (\texttt{compute\_U} applies verbatim to $W_1$); the MLP runs saved only configurations and metrics.
\end{remark}

\paragraph{Predictions and how to test them.}
\begin{enumerate}
\item \emph{The large-$\dlat$ return of memorization is a null-block effect that additionally requires capacity:} it appears at $\rho_\perp\approx1$ ($\signoise=0.5$) with hidden $=8\dlat$, is absent at $\rho_\perp\approx0$ ($\signoise=0.01$) at the same scaled capacity, and is absent at $\rho_\perp\approx1$ with fixed hidden $=256$. Test: rerun $\dlat\in\{100,150,200\}$ at hidden $=1600$ fixed for $\signoise\in\{0.01,0.5,1\}$, beyond $300$k steps; the isotropic run should rebound earlier ($\dlat^{\rm reb}\in[88,130]$) and more strongly.
\item \emph{Conditional on \cref{rem:scope-rebound-conj}, the rebound point scales linearly with $\nsamp$:} $\dlat^{\rm reb}-\dint=\nsamp\psi^{\rm reb}(\rho_\perp)$, $\psi^{\rm reb}(0.925)\in(0.19,0.29)$. Test: rerun the scaled-width sweep at $\nsamp=250$ and $1000$ (predicted onset $\dlat\approx55$--$80$ and $200$--$300$); the null-block signature needs generated-sample dumps (none exist).
\item \emph{Four-bulk separation threshold:} the noise-dim/residual gap crosses one decade at $\signoise^*=0.91$ ($\dlat=10$), $0.70$ ($20$), $0.50$ ($40$), $0.40$--$0.45$ ($60$--$100$), $0.52$ ($200$); below $\sqrt{\Delta_te^{2t}}=0.14$ the block is diffusion-noise dominated and sits on the $W$-broadened floor. Test: RFNN spectra at $\dlat\in\{10,40,100\}$ over $\signoise\in\{0.05,\dots,1.0\}$; expect $G=1$ near $0.9$ at $\dlat=10$ and near $0.5$ at $\dlat=40$.
\item \emph{Fixed-null-energy control $\signoise^2=c/D$:} for $D\gg50c$ the spectrum coincides with the $\signoise\to0$ spectrum, $\taugen$ keeps growing linearly with $a_1(q)^2\to0.96$, the RFNN residual-bulk delay persists, and no rebound appears. Test: RFNN spectral sweep and scaled-width MLP sweep with $\signoise^2=1/D$.
\item \emph{$\pwidth=\dlat+\nsamp+300$:} block means unchanged, noise-dim width grows as $((1+\sqrt{D/\pwidth})/(1-\sqrt{D/\pwidth}))^2$ ($2.4\times$ at $\dlat=40$, $6.8\times$ at $200$), and four-bulk separation fails at $\signoise=0.5$ for $\dlat\gtrsim100$. Test: spectra at $\dlat\in\{40,100,200\}$ for $\pwidth\in\{\dlat+\nsamp+300,8\dlat,64\dlat\}$.
\item \emph{Feature-learning conjecture:} after one $W$-step following readout pre-training, signal-block eigenvalues rise by $O(\eta)$ while noise-dim edges move by $O(\eta\beta_t^2/\alpha_t^2)$; with continued training the residual-bulk top rises and $\taumem^{\rm onset}-\taugen$ shrinks. Test: log the MLP first-layer conjugate-kernel spectrum during training.
\end{enumerate}

\paragraph{Open gaps.}
(1) The measured residual-bulk top exceeds \eqref{eq:scope-onset-iso} by $4.6\times$ at $\dlat=10$, decaying to $1.0\times$ at $\dlat=200$; attributed to cluster-level residual spikes but not derived (\cref{prop:onset-buf}, \cref{rem:collective-pred}); whether the small-$\dlat$ top modes are cluster-level or sample-specific is unresolved and affects what ``onset'' means. This also covers the isotropic control's reported $10.7\times$ gap growth, which the frozen-feature iso-residual model does not reproduce. (2) The rebound mechanism is a conjecture: a linear null score cannot produce NN-ratio $<1/3$; $\varepsilon_\perp$ is a necessary-condition proxy and $\varepsilon_c$ is calibrated on the same two points previously claimed as a match; Prediction~2 is the only non-circular test. (3) Capacity confound: the fixed-width $\signoise=0.5$ runs show $0.0000$ throughout; the rebound exists only with hidden $=8\dlat$ and is still rising at $300$k steps; the main text must not say the theory ``matches the MLP sweep''. (4) The cliff at index $\nsamp$ holds for $\Umat^{(1)}$; the saved $50$-draw spectra show a soft shoulder ($1.32$--$1.73$) and have $N_{\rm nd}=D$ for all $\dlat$; the released block-summary script still partitions with the superseded cumulative count and only its noise-dim minimum and sample top are used here. (5) Block widths in \eqref{eq:scope-ndmin} use a product-of-edges approximation that under-estimates the noise-dim lower edge by $\sim25\%$; exact edges need the Silverstein fixed point. (6) \cref{prop:scope-nullspec} neglects the signal--null cross covariance, of the same order as $\varepsilon_\perp^2$ at small $\psi$. (7) Feature learning is entirely conjectural; no MLP spectra exist; the one-step scaling is $O(\eta)$ (corrected from $\Theta(\eta^2)$) and whether the first-order term survives depends on alignment. (8) The isotropic control and $\dlat=240$ do not exist as runs. (9) The real-data regime (continuous $\rho_\Sigma$, collapsed coordinates) is placed outside (S3) and handled in \cref{app:theory-general}. (10) Spectra are computed at the single time $t=0.01$ while generation covers $t\in[0.01,3]$; $\rho_\perp(t)\ge\tfrac12$ only for $t\le\tfrac12\log(1+\signoise^2)$ ($0.11$ at $\signoise=0.5$), so the null block is data-dominated only on a short final segment of the reverse trajectory.

% ----------------------------------------------------------------------
\subsection{Configurations outside the hypotheses}
\label{app:theory-outside}

\begin{remark}[Out-of-scope configurations]\label{rem:scope-outside}
\begin{enumerate}
\item $\dlat=\dint=5$ (both sweeps): $N_{\rm nd}=0$ and $\pwidth=320<\nsamp$, so (S2) fails twice; the residual bulk has $\le\pwidth-\dlat=315$ modes. A baseline, not a four-bulk spectrum.
\item $\signoise=0.01$ at $t=0.01$ (RFNN and MLP): $\rho_\perp=4.9\times10^{-3}$, $\signoise^2=5\times10^{-3}\Delta_t$, measured noise-dim/residual gap $0.004$--$0.056$ decades. (S3), (S5) fail: three-bulk regime. Counts and $\taugen$ still apply; \cref{prop:scope-sigperp}(i) is the operative statement and predicts no rebound (observed in the scaled sweep: memorization $0$ for $20\le\dlat\le200$). Floor medians match $\pwidth a_1(q)^2\Delta_t/\dlat$ for $\dlat\ge30$ ($0.67,0.77$ predicted vs $0.74,0.81$ measured at $\dlat=30,40$); at $\dlat=10$ the floor block is merged with the residual bulk ($G=0.03$) and the comparison is not meaningful.
\item $\dlat\ge150$ at $\nsamp=500$, $\signoise=0.5$: $\psi_n\le3.3$, $\varepsilon_\perp^2\ge0.64$; regime of \cref{rem:scope-rebound-conj}. (S5) also degrades (gap $1.26\to1.09$ decades from $\dlat=150$ to $200$). Memorization returns only in the scaled-width MLP sweep ($0.0022,0.0144$ at $\dlat=150,200$, hidden $8\dlat$) and not in the fixed-width sweep (hidden $256$: $0.0000$).
\item $\dint$ sweep at $\dlat=20$: $\dint\in\{16,20\}$ gives $N_{\rm nd}\in\{4,0\}$; the signal block is a finite-rank spike for every $\dint$, so ``signal edges'' are spike locations, not bulk edges.
\item RFNN with $\pwidth=\dlat+\nsamp+300$: run only at $\dlat=20$; \cref{rem:scope-psip} predicts that width-side broadening breaks (S5) for $\dlat\gtrsim100$ at $\signoise=0.5$.
\item Isotropic control $\signoise=\sigsig=1$ (no released training run at $\signoise=1$ exists; the spectral statements are checked by the simulation in \cref{rem:isotropic-buf}): (S3) holds ($\rho_\perp=0.98$) and the $D$ noise-dim modes exist at $a_1(q)^2\beta_t^2/\dlat$ between signal and residual modes, but the signal/noise-dim ratio is $\alpha_t^2/\beta_t^2=2.76$ ($0.44$ decades) for all $\dlat$: the two-atom \emph{separation} of the signal block is absent, not the block (\cref{rem:fourbulk-isotropic,rem:isotropic-buf}). The reported $10.7\times$ growth of the spectral gap under isotropy must then be a residual-bulk effect (fall of $a_*(q)^2$ and of the cluster-spike excess), not a noise-dim one. \cref{rem:scope-rebound-conj} with $\varepsilon_c^2\in(0.30,0.64)$ gives $\psi^{\rm reb}(0.98)\in(0.166,0.251)$, i.e.\ $\dlat^{\rm reb}\in[88,130]$ at $\nsamp=500$, earlier than at $\signoise=0.5$.
\item All trainable-MLP runs (synthetic, hidden $256$ fixed or $8\dlat$ scaled; real, depth 5, width 1024): (S1) fails; only \cref{rem:scope-feature} applies, and the clock statements of \cref{rem:clock-p1,rem:lin-gd} do not cover Adam or momentum SGD.
\item Real VAE latents ($t=0.1$, $\nsamp=1000$): $\rho_\Sigma$ is continuous, effective rank $\approx0.5\,\dlat$, and for $\dlat\ge140$ (CelebA) / $160$ (CIFAR-10) a set of coordinates has $\lambda=\Delta_t$ exactly (posterior collapse). (S3) fails as stated; the collapsed coordinates are the $\signoise=0$ case of \cref{prop:scope-sigperp}(i): free modes at the floor with $\varepsilon_\perp=0$, counted in $N_{\rm nd}$ but not contributing to $\rho_\perp$. The general-$\rho_\Sigma$ form of \citet[Thm.~3.1]{bonnaire2025} applies (\cref{app:theory-general}); the two-atom edges do not.
\item Time mismatch: spectra are computed at the single time $t=0.01$, while training and generation cover $t\in[0.01,3]$ (the reverse SDE is integrated from $t_{\max}=3$ to $t_{\min}=0.01$; $t_{\rm eval}=0.1$ enters only the score-error metric). $\rho_\perp(t)\ge\tfrac12$ only for $t\le\tfrac12\log(1+\signoise^2)$ ($0.11$ at $\signoise=0.5$, $5\times10^{-5}$ at $0.01$), so the null block is data-dominated only on a short final segment of the reverse trajectory; the bridge of \cref{app:theory-bridge} places the attachment time $t_c$ between $0.05$ and $1.0$ across the sweep, outside that segment for $\dlat\gtrsim20$.
\item Synthetic score-error curves: all saved values were computed with a transposed whitening matrix and are artifacts; the predictions of \cref{app:theory-linear} are for the corrected metric and are currently unanchored.
\end{enumerate}
\end{remark}
